\chapter{Definite unitary groups and iterated theta lifts}\label{chapter-unitary-theta}

Over the host field of Chapter~\ref{chapter-host-fields} the argument fixes
totally definite Hermitian spaces $X_0,\ldots,X_m$ of ranks $k+j(t+1)$,
indefinite companions $Y_j$, and Harish-Chandra parameters that interlace in
the sense of Gan, Gross and Prasad. An algebraic modular form on $U(X_0)$ is
lifted $m$ times by kernels attached to pairs of lattice classes, and the
periods of the lifted forms make up the auxiliary lattice measured in
Chapters~\ref{chapter-period-divisibility} to~\ref{chapter-covolumes}. This
chapter sets up these data, states the inner product formula of Rallis and the
nonvanishing of the first theta lifts in the forms used, and proves that the
nonvanishing of a period is carried up a tower of theta lifts to a Bessel
period.

\section{Ranks and Harish-Chandra parameters}\label{section-th-ranks}

\begin{definition}\label{th-ranks}
A \emph{rank datum} consists of:
\begin{enumerate}
\item a Galois datum (Definition~\ref{hf-frame}) with $2$-power exponent
  $e\geq0$ whose group has order $n$;
\item integers $k,m,t\geq0$ with $k$ and $m$ even, $t$ odd, and
  \[
  2\leq m,\quad m+3\leq k,\quad k+m+4\leq t,\quad 3m<t,\quad
  m\leq 2n\leq k,\quad m^2\leq k+m(t+1).
  \]
\end{enumerate}
For a rank datum put $f=2n$, $r=k+m-1$, $r_S=m/2-1$ and
\[
N(j)=k+j(t+1),\qquad \operatorname{rk}Y(j)=r+jt\qquad(j\geq0),
\]
and $\operatorname{sig}W(j)=(k+j-1,\,m-j)$ for $1\leq j\leq m$. Then
$N(m)=r+mt+1$, every $N(j)$ is even and $r$ is odd.
\end{definition}
\leanref{Formalism/LocalRingCardinality/Params.lean}{Formalism.BasicNotions.Ranks}
\leanref{Formalism/SharpRestriction/RamificationFrobenius/TotientGaloisClasses/ZsmulGeneratorEqZeroOfProjection.lean}{Formalism.HC.Frame}

In the argument the ranks are adapted to a scale $x$ as in
Definition~\ref{hf-adapted-ranks}: $n$ is of size $x^{2+8\varepsilon}$, the
number $m$ of lifts is of size $x^{1/2-\varepsilon/16}$, and $k$ is of size
$x^{9/2+10\varepsilon}$. The integer $N(j)$ is the rank of the $j$-th definite
space, and each lift adds a block of rank $t+1$.

\begin{definition}\label{th-hc-parameter}
Let a rank datum be given (Definition~\ref{th-ranks}). A
\emph{Harish-Chandra parameter} consists of rational numbers
$u_1,\ldots,u_k$, $v_1,\ldots,v_{k-1}$, $b_1,\ldots,b_m$ and
$e_1,\ldots,e_m$ such that:
\begin{enumerate}
\item $u_i\in\tfrac12+\mathbf Z$ and $v_i,b_j,e_j\in\mathbf Z$;
\item $u_1>v_1>u_2>v_2>\cdots>u_{k-1}>v_{k-1}>u_k$;
\item $-b_1<u_k$;
\item $e_j-b_j\geq(t+3)/2$ for $1\leq j\leq m$, and
  $b_{j+1}-e_j\geq(t+3)/2$ for $1\leq j<m$.
\end{enumerate}
\end{definition}
\leanref{Formalism/LocalRingCardinality/Params.lean}{Formalism.BasicNotions.HCParam}
\leanref{Formalism/LocalRingCardinality/Params.lean}{Formalism.BasicNotions.factHCMono}

Condition (2) says that $(u_i)$ and $(v_i)$ interlace, the pattern of the
branching from $U(k)$ to $U(k-1)$ that governs the Gan-Gross-Prasad periods
\cite{GGP12}. By (3) and (4) the numbers
$u_k>-b_1>-e_1>-b_2>\cdots>-e_m$ decrease, with gaps at least $(t+3)/2$ in
the $(b,e)$-part. The weights of the forms on $U(X_j)$ are given in
Definition~\ref{th-definite-datum}(7): the $u$-part followed by $j$ constant
blocks of length $t+1$, the $b$-th of which is determined by $e_b$.

\section{Spaces, column chains and the definite-side datum}\label{section-th-datum}

\begin{definition}\label{th-spaces}
Let $E$ be a CM field with totally real subfield $F$, and write
$x\mapsto\bar x$ for complex conjugation on $E$. A \emph{Hermitian space} over
$E$ is a finite-dimensional $E$-vector space $V$ with a nondegenerate Hermitian
form $h$, conjugate-linear in the first variable, together with a basis of
$V$, a \emph{signature} $\tau\mapsto(p_\tau,q_\tau)$ on the complex embeddings
$\tau$ of $E$, and a unimodular $\mathcal O_E$-lattice $L_0\subset V$. Its
Gram matrix $G_V$ is the matrix of $h$ in the basis. The signature is
\emph{correct} if for every $\tau$ the Hermitian matrix $\tau(G_V)$ has $p_\tau$
positive and $q_\tau$ negative eigenvalues. The space is \emph{definite} if
$(p_\tau,q_\tau)$ is $(\operatorname{rk}V,0)$ or $(0,\operatorname{rk}V)$ for
every $\tau$. For a rank datum (Definition~\ref{th-ranks}), a \emph{system of
spaces} over $E/F$ consists of:
\begin{enumerate}
\item a prime $T$;
\item Hermitian spaces $X_j$ ($j\geq0$) with correct signature, of rank $N(j)$
  and definite for $j\leq m$;
\item a Hermitian space $X'$ with correct signature, of rank $r+mt$, an
  isometric embedding $\iota'\colon X'\to X_m$, and a vector $z\in X_m$ with
  $h(z,z)\neq0$ and $z\perp\iota'(X')$;
\item Hermitian spaces $Y_j$ ($j\geq0$) with correct signature, of rank
  $r+jt$ for $j\leq m$, isometric embeddings $X_j\to Y_j$ for $j<m$, and an
  isometry $X'\cong Y_m$;
\item in $Y_0$: a vector $z_0$ with $h(z_0,z_0)=-T$, and vectors
  $e_1,\ldots,e_{r_S}$, $f_1,\ldots,f_{r_S}$ with $h(e_a,e_b)=h(f_a,f_b)=0$
  and $h(e_a,f_b)=\delta_{ab}$, all orthogonal to the image of $X_0$, with
  $z_0$ orthogonal to the $e_a$ and $f_a$; and
  $(p_\tau,q_\tau)(Y_0)=(k+r_S,\,r_S+1)$ for every $\tau$;
\item Hermitian spaces $W_j$ ($j\geq0$) with correct signature such that for
  $1\leq j\leq m$, $W_j$ has rank $r$ and signature $\operatorname{sig}W(j)$
  at every $\tau$, and there is $g\in GL_r(E\otimes_F\mathbf A_{F,f})$ with
  $\bar g^{\mathsf T}G_{Y_0}g=G_{W_j}$ after reordering the basis of $Y_0$;
\item Hermitian spaces $V_j$ ($j\geq0$) with correct signature and signs
  $\epsilon_j=\pm1$ such that for $1\leq j\leq m$, $V_j$ has rank $r+t$ and
  $L_0$ has an orthonormal basis for $\epsilon_jh$.
\end{enumerate}
\end{definition}
\leanref{Formalism/LocalRingCardinality/DataSpaces.lean}{Formalism.BasicNotions.HSpace}
\leanref{Formalism/LocalRingCardinality/DataSpaces.lean}{Formalism.BasicNotions.Spaces}

The rank of $Y_0$ is $k+1+2r_S=r$, so (5) decomposes $Y_0$ as the image of
$X_0$, the line spanned by $z_0$, and $r_S$ hyperbolic planes. Since $z_0$ has
negative norm at every real place and the hyperbolic planes contribute
$(r_S,r_S)$, the signature condition forces $X_0$ to be positive definite at
every place.

\begin{definition}\label{th-column-chain}
Let a rank datum be given (Definition~\ref{th-ranks}), let $\mathcal T$ be a
finite set, and let $\lambda$ assign to each $\tau\in\mathcal T$ a
Harish-Chandra parameter $\lambda(\tau)$ (Definition~\ref{th-hc-parameter}),
with $e$-part $e_1(\tau),\ldots,e_m(\tau)$. A \emph{column chain} for
$\lambda$ consists of:
\begin{enumerate}
\item finite sets $\mathrm{cls}_i$ for $i\geq0$, the \emph{classes};
\item for each $c\in\mathrm{cls}_i$ a finite group $\Gamma_c$, and an integer
  $D_\Gamma>0$ divisible by $|\Gamma_c|$ for every $c\in\mathrm{cls}_i$ with
  $i\leq m$;
\item sets $Y_i$ and functions $\delta_i\colon Y_i\times\mathcal T\to\mathbf
  C^\times$;
\item complex numbers $\rho_i(c)$ for $c\in\mathrm{cls}_i$ and $\rho'_i(c')$
  for $c'\in\mathrm{cls}_{i+1}$, of absolute value $1$ when $i<m$;
\item for $c\in\mathrm{cls}_i$ and $c'\in\mathrm{cls}_{i+1}$ a linear map
  $\vartheta_i(c,c')$ from functions on $Y_i$ to functions on $Y_{i+1}$, the
  \emph{kernel} of the pair;
\item integers $s_i(\tau)$ and $s'_i(\tau)$ with
  \[
  s_i(\tau)=e_{i+1}(\tau)-\tfrac{N(i+1)}2,\qquad
  s'_i(\tau)=e_{i+1}(\tau)-\tfrac{N(i)}2\qquad(0\leq i<m).
  \]
\end{enumerate}
\end{definition}
\leanref{Formalism/LocalRingCardinality/DataKernels.lean}{Formalism.BasicNotions.ColumnChain}

Since $N(i+1)-N(i)=t+1$, the exponents satisfy $s'_i-s_i=(t+1)/2$. The
names reflect the intended reading: $\mathrm{cls}_i$ stands for a set of
classes of lattices in $X_i$, $\Gamma_c$ for the stabilizer of a lattice of
class $c$, and a function on $\mathrm{cls}_i\times Y_i$ for a vector-valued
algebraic modular form on $U(X_i)$, written as a function of a class and of a
point of $Y_i$. The datum below attaches a lattice to each class
(Definition~\ref{th-definite-datum}(7)); the other identifications are
conditions imposed where they are used (Remark~\ref{th-lift-kernels}).

\begin{definition}\label{th-iterated-lift}
Let a column chain be given (Definition~\ref{th-column-chain}). For
$0\leq i$ and a function $\varphi$ on $\mathrm{cls}_i\times Y_i$, written
$c\mapsto\varphi_c$ with $\varphi_c$ a function on $Y_i$, define the
function $K_i\varphi$ on $\mathrm{cls}_{i+1}\times Y_{i+1}$ by
\[
(K_i\varphi)_{c'}(y')=\sum_{c\in\mathrm{cls}_i}\frac{\rho_i(c)\rho'_i(c')}{|\Gamma_c|}
\prod_{\tau\in\mathcal T}\delta_{i+1}(y',\tau)^{-s'_i(\tau)}\;
\vartheta_i(c,c')\Bigl(y\mapsto\prod_{\tau\in\mathcal T}\delta_i(y,\tau)^{s_i(\tau)}\varphi_c(y)\Bigr)(y').
\]
The \emph{$i$-th lift} is $\varphi\mapsto K_i\varphi$, and the
\emph{iterated lift} is $\operatorname{col}=K_{m-1}\circ\cdots\circ K_0$,
from functions on $\mathrm{cls}_0\times Y_0$ to functions on
$\mathrm{cls}_m\times Y_m$.
\end{definition}
\leanref{Formalism/LocalRingCardinality/DataKernels.lean}{Formalism.BasicNotions.ColumnChain.K}
\leanref{Formalism/LocalRingCardinality/DataKernels.lean}{Formalism.BasicNotions.ColumnChain.col}
\leanref{Formalism/ThetaLiftIndependence/Statement.lean}{Formalism.ThetaLiftIndependence.chainOpM}

\begin{definition}\label{th-definite-datum}
A \emph{definite-side datum} consists of the following data, subject to the
conditions stated with them.
\begin{enumerate}
\item \emph{Field and spaces.} A rank datum (Definition~\ref{th-ranks}); a
  host for it, consisting of a CM field $E$ with totally real subfield $F$,
  a CM embedding $\varphi_0\colon E\to\mathbf C$, auxiliary primes $q\neq3$ and
  $d$ with $4tr+1<q,d$, primes $g_0\neq g_1$, a finite set of \emph{level
  primes} containing $g_0$ and $g_1$, two further level primes
  $p_{de}\neq p'_{de}$ congruent to $1$ modulo $15$ and different from $g_0,g_1$,
  an element $\xi_0=-1/(\mathrm N\mathfrak c\cdot\sqrt{-d})\in E$ for an
  integral ideal $\mathfrak c$, and the further host data of
  Chapter~\ref{chapter-host-fields}; and a system of spaces over $E/F$
  (Definition~\ref{th-spaces}) with signature
  $(r+jt-\lfloor(m-j)/2\rfloor,\ \lfloor(m-j)/2\rfloor)$ for $Y_j$, $j\leq m$,
  at every embedding.
\item \emph{Parameters.} A set $\mathcal T$ of $f$ elements, identified with a
  subset of the group of Chapter~\ref{chapter-host-fields} that indexes the
  complex embeddings of $E$, so that each $\tau\in\mathcal T$ is a complex
  embedding of $E$; a Harish-Chandra parameter
  $\lambda(\tau)$ for each $\tau\in\mathcal T$
  (Definition~\ref{th-hc-parameter}); a finite set of \emph{words} with a base
  word $u_0$; and for each word $u$ parameters $\lambda_u(\tau)$ with
  $\lambda_{u_0}=\lambda$, such that $\lambda_u(\tau)$ and $\lambda(\tau)$
  have the same $v$, $b$ and $e$ and the same $u_1$ and $u_k$.
\item \emph{Members.} For each word $u$: a column chain $C_u$ for $\lambda_u$
  (Definition~\ref{th-column-chain}); data at an auxiliary prime, namely a
  prime $T_u$, a set of $f$ finite fields $\mathbf F_v$ with $T_u$ elements,
  nontrivial additive characters of the $\mathbf F_v$, actions $\rho_v$ of
  $\mathbf F_v^\times$ on functions on $\mathrm{cls}_0\times Y_0$, a
  finite-dimensional subspace $M_u$ stable under them, and numbers
  $d^*_u\geq1$, $D^*>0$ and $\Gamma_P>1$; and a period map, consisting of
  complex vector spaces $\mathcal W_u$ and $\mathcal M^{\mathrm t}_u$, a
  linear map $\operatorname{ev}_u$ from functions on
  $\mathrm{cls}_m\times Y_m$ to $\mathcal W_u$, and a bilinear pairing
  $\mathcal P_u$ on $\mathcal W_u\times\mathcal M^{\mathrm t}_u$. The
  \emph{period of the iterated lift} of $\varphi$ against $F$ is
  $\mathcal P_u(\operatorname{ev}_u(\operatorname{col}\varphi),F)$. Put
  $M_{u,\mathrm{ram}}=M_u\cap\bigcap_v\ker\sum_{a\in\mathbf F_v^\times}\rho_v(a)$,
  $X_u=(D^*/d^*_u)^{1/(fk)}$ and $T^*_u=X_u\Gamma_P^{-2/(fk)}$.
\item \emph{Conditions at the auxiliary primes.} For every word $u$: $T_u$ is
  a level prime different from $p_{de}$ and $p'_{de}$;
  $T_u=\pi\bar\pi$ for some $\pi\in\mathcal O_E$ lying in the subfield $K$
  of the host, with $\pi\equiv1$ modulo $Q_Sqd(g_0g_1)^2$, where $Q_S$ is the
  integer of the host; $T^*_u\leq T_u\leq(1+1/(fk))T^*_u$; $\Gamma_P$ does not depend on $u$;
  $T_{u_0}=T$; and there is $c_u\in E^\times$ with $c_u\bar c_uT_u=T_{u_0}$.
\item \emph{Coefficient field and periods.} A totally complex number field
  $L$, Galois over $\mathbf Q$, with an involution and an embedding
  $L\to\mathbf C$ under which the involution is complex conjugation, and an
  embedding $E\to L$ that carries complex conjugation of $E$ to the involution
  of $L$ and whose composite with $L\to\mathbf C$ is $\varphi_0$;
  finite-dimensional Hilbert spaces
  $\mathcal E$ (\emph{top forms}), $\mathcal M_u$ for the words $u$, and
  $\mathcal A_\mu$ for the tuples $\mu$ of Hecke characters, with a base
  tuple $\mu_0$; periods $\operatorname{per}_{u,\mu}$ on
  $\mathcal M_u\times\mathcal A_\mu$; and \emph{Bessel pairings} $B_u$ on
  $\mathcal M_u\times\mathcal E$, linear in each variable. For each word $u$
  there are a linear isomorphism $\mathcal M_u\cong M_{u,\mathrm{ram}}$ and an
  identification of $\mathcal M^{\mathrm t}_u$ with $\mathcal A_{\mu_0}$ under
  which $\operatorname{per}_{u,\mu_0}$ is the period of the iterated lift of
  $C_u$, and the inner product of $\mathcal M_u$ is
  $\langle\varphi,\psi\rangle=\sum_c|\Gamma_c|^{-1}\langle\varphi_c,\psi_c\rangle$
  for the forms of (7); for every tuple $\mu$ there is a twisted
  column chain $C_{u,\mu}$, with $C_{u,\mu_0}=C_u$ and with a linear
  identification of the functions in degree $0$, whose iterated lift gives
  $\operatorname{per}_{u,\mu}$ in the same way.
\item \emph{Determinants and twists.} An abelian group $\Delta$, a class
  $d_i(c)\in\Delta$ for every lattice class $c\in\mathrm{cls}_i$ of every
  member, characters $\tilde\mu_i$ ($i\geq0$) and $\chi_\Delta$ of $\Delta$,
  and a homomorphism from the norm-one elements of $E$ to $\Delta$ through
  which $\chi_\Delta(z)=\prod_{\tau\in\mathcal T}\tau(z)^{-1}$. The twists of
  the member $C_u$ are
  \[
  \rho_i(c)=\tilde\mu_{i+1}(d_i(c))^{-1}\chi_\Delta(d_i(c))^{-N(i+1)/2},\qquad
  \rho'_i(c')=\tilde\mu_{i+1}(d_{i+1}(c'))\,\chi_\Delta(d_{i+1}(c'))^{N(i)/2}.
  \]
  The datum also contains relative determinants of pairs of self-dual lattices,
  multiplicative in the pair and given by $\det\gamma$ on translates by a
  unitary matrix $\gamma$.
\item \emph{Lattices of forms.} For each word $u$, each $i$ and each
  $c\in\mathrm{cls}_i$ an $\mathcal O_E$-lattice $L_u(i,c)\subset X_i$; and for
  each word $u$ and each $j$ dominant weights
  $\lambda_{u,j}(\tau)\in\mathbf Z^{N(j)}$ with
  \[
  \lambda_{u,j}(\tau)_i=
  \begin{cases}
  u_i(\tau)-\tfrac{k+1}2+i, & 1\leq i\leq k,\\[2pt]
  -e_b(\tau)+\tfrac k2+\tfrac{(2b-1)(t+1)}2, & i=k+(b-1)(t+1)+s,\ 1\leq s\leq t+1,
  \end{cases}
  \]
  for $j\leq m$, where $u_i(\tau)$ and $e_b(\tau)$ are the entries of
  $\lambda_u(\tau)$; together with Hermitian lattices over $L$ whose images
  span subspaces $F_{u,j}$ of the functions on $\mathrm{cls}_j\times Y_j$, with
  $F_{u,0}=M_u$, and sesquilinear forms $\langle\ ,\ \rangle$ on the
  functions on $Y_j$, one for each class, such that the norm of a lattice
  vector at the embedding $L\to\mathbf C$ is
  $\sum_c|\Gamma_c|^{-1}\langle\varphi_c,\varphi_c\rangle$.
\item \emph{Bessel-period data.} A rational $\Delta_S>0$ with
  $v_l(\tau)-v_{l+1}(\tau)=4\Delta_S$ for $1\leq l\leq k-2$; an injective map
  $u\mapsto(\varepsilon_1(u),\ldots,\varepsilon_{k-2}(u))\in\{0,1\}^{k-2}$ on
  words with
  $u_{l+1}(\tau)=v_l(\tau)-\tfrac12-\varepsilon_l(u)\Delta_S$ for the
  parameter $\lambda_u(\tau)$ and $1\leq l\leq k-2$; a finite-dimensional
  Hilbert space $\mathcal H$ with a conjugate-linear isometry
  $s\colon\mathcal E\to\mathcal H$; and for each word a finite-dimensional
  Hilbert space $\mathcal C_u$, a linear isometry $\mathcal C_u\cong M_u$ onto
  the space of (3) with the inner product of (7), and a continuous linear map
  $P_u\colon\mathcal C_u\to\mathcal H$. The \emph{standard Bessel period} of
  $\varphi\in\mathcal M_u$ against $F\in\mathcal E$ is
  $\mathcal B_u(\varphi,F)=\langle s(F),P_u(\tilde\varphi)\rangle$, where
  $\tilde\varphi\in\mathcal C_u$ corresponds to the image of $\varphi$ in
  $M_{u,\mathrm{ram}}\subset M_u$.
\item \emph{Further data used in later chapters:} division data (lattices
  $N_u$ spanning $\mathcal M_u$ and a lattice $R$ spanning $\mathcal E$, of
  Chapter~\ref{chapter-integral-structure}), local data at the primes and at
  the places of $L$ with their compatibilities, the twisting data of
  Definition~\ref{th-twisting-data}, theta-jet data, and the measures and
  constants that enter the size conditions.
\end{enumerate}
\end{definition}
\leanref{Formalism/LocalRingCardinality/RhoDetPrimeEqMuTildeMulChiPow.lean}{Formalism.BasicNotions.D7}
\leanref{Formalism/LocalRingCardinality/RhoDetPrimeEqMuTildeMulChiPow.lean}{Formalism.BasicNotions.Member}
\leanref{Formalism/LocalRingCardinality/DataKernels.lean}{Formalism.BasicNotions.ColumnsAtT}
\leanref{Formalism/LocalRingCardinality/DataKernels.lean}{Formalism.BasicNotions.PeriodMap}
\leanref{Formalism/LocalRingCardinality/DataRows.lean}{Formalism.BasicNotions.Rows.Data}
\leanref{Formalism/LocalRingCardinality/DataForms.lean}{Formalism.BasicNotions.MemberForms}
\leanref{Formalism/SharpRestriction/CuspSpanNormBounds/DataS.lean}{Formalism.MainNotions.S.Draft.SData}
\leanref{Formalism/LocalRingCardinality/RhoDetPrimeEqMuTildeMulChiPow.lean}{Formalism.BasicNotions.D7.besS}

By (7) the weight of the forms on $U(X_j)$ is the weight attached to the
$u$-part of the parameter, followed by $j$ constant blocks of length $t+1$.
A constant block is the weight of a one-dimensional representation of
$U(t+1)$; this is the archimedean trace of the shape
$\pi\mapsto\pi\boxplus\mu_1[t+1]\boxplus\cdots\boxplus\mu_j[t+1]$ of the
iterated lift. By (8) distinct words have distinct binary labels of length
$k-2$. Each letter moves one interior entry $u_{l+1}$ down by $\Delta_S$
inside the interval between $v_{l+1}$ and $v_l$, which has length
$4\Delta_S$, so all parameters $\lambda_u(\tau)$ interlace with the same
$v(\tau)$.

\begin{definition}\label{th-twisting-data}
Let $D$ be a definite-side datum (Definition~\ref{th-definite-datum}), with host
field $E$, ranks $r$, $m$ and $t$, level primes $g_0,g_1$ and coefficient field
$L$, and let $L^0\subset L$ be the subfield fixed by the datum. Let $\Omega$ be
an algebraic closure of $E$, put $G_E=\operatorname{Gal}(\Omega/E)$, and let
$G_0=\prod_{k\in\{0,1\}}\prod_vGL_r(\mathbf F_{g_k})$, the product over the
places $v$ of the datum over $g_k$. The \emph{twisting data} of $D$ consist of
the following.
\begin{enumerate}
\item \emph{Indices and coefficients.} A finite set of \emph{twisting indices}
  $i$, each with: a prime $\ell_i$; a complete discrete valuation ring
  $\mathcal O^0_i$, finite over $\mathbf Z_{\ell_i}$, with two embeddings of
  $L^0$ in its fraction field that are exchanged by complex conjugation; a
  complete discrete valuation ring $\mathcal O_i\supset\mathcal O^0_i$ over
  $\mathbf Z_{\ell_i}$ with algebraically closed residue field and fraction
  field $K_i$; for $k\geq0$, a finite extension $\mathcal O_{i,k}$ of
  $\mathcal O_i$ that is a discrete valuation ring, containing a primitive
  $\ell_i^k$-th root of unity $\zeta_{i,k}$ that generates its fraction field
  over $K_i$; the set $P_D(i)$ of places of $E$ over $p_{de}$, or over
  $p'_{de}$ if $\ell_i=p_{de}$; and a finite set $S_0(i)$ of places of $E$
  containing $P_D(i)$ and the places over $\ell_i$, $g_0$ and $g_1$.
\item \emph{Ray class groups.} For each $i$ and each finite set $Q$ of places
  of $E$, a finite abelian $\ell_i$-group $\mathrm{RCG}_i(Q)$ and a continuous
  surjective homomorphism $\operatorname{art}\colon G_E\to\mathrm{RCG}_i(Q)$,
  unramified outside $Q$ and inverted by conjugation with complex conjugation;
  when no place of $Q$ lies over $\ell_i$, every continuous homomorphism from
  $G_E$ to a finite abelian $\ell_i$-group with these two properties factors
  through $\operatorname{art}$; and the maps are compatible with restriction
  from $Q$ to $Q'\subset Q$.
\item \emph{Signatures.} For each $i$, a character
  $\delta_i\colon G_E\to(\mathcal O^0_i)^\times$ with open kernel, inverted by
  conjugation with complex conjugation, nontrivial on the inertia groups at a
  place $v$ of $E$ if and only if $v\in P_D(i)$, and mapping each inertia group
  at a place of $P_D(i)$ onto a group of order $15$; every $v\in P_D(i)$ has
  $Nv\equiv1\pmod{15}$. Continuous characters
  $\chi_{i,j,b}\colon G_E\to(\mathcal O^0_i)^\times$ for $j\geq0$ and signs
  $b$ such that: $\chi_{i,j,b}$ is unramified at every place of $E$ that is not
  in $P_D(i)$ and lies over none of $\ell_i$, $g_0$ and $g_1$; on the inertia
  groups at the places of $P_D(i)$ it equals $\delta_i$ or $\delta_i^{-1}$,
  according to $b$; and $\chi_{i,j,b}(\sigma)^{g_k-1}=1$ for $\sigma$ in an
  inertia group at a place over $g_k$ not over $\ell_i$. A finite set of
  \emph{signatures} $s$, each with a
  lift $1\leq j_s\leq m$, a sign $b_s$, a continuous representation
  $\rho_s\colon G_E\to GL_r(\mathcal O^0_i)$ and \emph{Hecke eigenvalues}
  $\lambda_s(k,P)\in\mathcal O_i$ for $k\geq1$ and places $P$ of $E$. Put
  $\chi_s=\chi_{i,j_s,b_s}$.
\item \emph{Theta values.} For each $i$, a finite set $\Lambda_i$, each
  $\lambda\in\Lambda_i$ with a lift, a sign $b_\lambda$, an integer and a
  continuous character $G_E\to\mathcal O_i^\times$ unramified outside $S_0(i)$;
  and for every finite set $Q$ of places, every $\lambda\in\Lambda_i$ and every
  $\delta\in\mathrm{RCG}_i(Q)$, a \emph{theta value}
  $\vartheta_i(Q,\lambda,\delta)\in K_i$.
\item \emph{Spaces of forms.} For each $i$, each lift $j$ and each sign $b$, a
  finite-dimensional $K_i$-vector space $F_{i,j,b}$ with a linear action of
  $G_0$ and $K_i$-linear \emph{Hecke operators} $T_k(P)$ for $k\geq0$ and places
  $P$; a subgroup $U\subset G_0$ with characters
  $\psi_{U,i,b}\colon U\to K_i^\times$; a central idempotent $e_{i,b}$ of the
  group algebra $K_i[G_0]$ whose coefficients are the images of elements of
  $L^0$ under the embedding of (1) attached to $b$; and an
  $\mathcal O_i$-submodule of $F_{i,j,b}$.
\item \emph{Theta data.} For each $i$, $j$ and $b$, a set of \emph{theta data}
  $a$, each with a \emph{theta element} $\theta(a)\in F_{i,j,b}$ and a finite
  set of places of $E$, its \emph{ramification set}; for $k\geq0$, a finite set $Q$ of
  places and a character $\psi\colon G_E\to\mathcal O_{i,k}^\times$, a set of
  \emph{$\psi$-twisted theta elements}
  $\theta_{\psi,Q}(a)\subset\mathcal O_{i,k}\otimes_{\mathcal O_i}F_{i,j,b}$,
  with $\theta_{1,\emptyset}(a)=\{1\otimes\theta(a)\}$; and a finitely generated
  $\mathcal O_{i,k}$-module $M_{i,j,b}(k,Q,\psi)$ spanning a vector space over
  the fraction field of $\mathcal O_{i,k}$, with an $\mathcal O_{i,k}$-linear map
  $\theta_{\psi,Q}\colon M_{i,j,b}(k,Q,\psi)\to\mathcal O_{i,k}\otimes_{\mathcal O_i}F_{i,j,b}$.
\end{enumerate}
The data also contain: for each $i$, each finite set $Q$ and each $M\geq0$, a
homomorphism from $G_E$ to $\Delta_i(Q,M)=(\mathbf Z/\ell_i^M\mathbf Z)^Q$,
unramified outside $Q$ and trivial on the decomposition groups over $\ell_i$,
and for each theta datum an element of
$F_{i,j,b}\otimes_{\mathcal O_i}\mathcal O_i[\Delta_i(Q,M)]$; the
identification of the characters of the tuples at the twisted places with
products of the characters $\chi_{i,j,b}$ and characters of $\ell_i$-power
order; and, in each $F_{i,j,b}$, the image of a free $\mathcal O_i$-module of
finite rank containing a lattice of theta elements.

With these data, define for each $i$, $j$ and $b$:
\begin{enumerate}\setcounter{enumi}{6}
\item the \emph{space of the twisting data} $V_0(i,j,b)=e_{i,b}F^0_{i,j,b}$,
  where $F^0_{i,j,b}$ is the subspace of $x\in F_{i,j,b}$ with
  $gx=\psi_{U,i,b}(g)x$ for $g\in U$;
\item for a set $\Sigma$ of theta data, the \emph{theta lattice}
  $\Theta(\Sigma)$, the $\mathcal O_i$-span of the $\theta(a)$ with
  $a\in\Sigma$; its image $e_{i,b}\Theta(\Sigma)$; the \emph{$\psi$-twisted
  theta lattice} $e_{i,b}\Theta_{\psi,Q}(\Sigma)$, the image under $e_{i,b}$ of
  the $\mathcal O_{i,k}$-span of the $\theta_{\psi,Q}(a)$ with $a\in\Sigma$; and
  the \emph{twisted lattice of the twisting data}
  $e_{i,b}\,\theta_{\psi,Q}(M_{i,j,b}(k,Q,\psi))$;
\item a signature $s$ is \emph{good} if some nonzero $x\in F_{i,j_s,b_s}$ has
  $e_{i,b_s}x=x$ and $T_k(P)x=\lambda_s(k,P)x$ for $1\leq k\leq r$ and every
  place $P\notin S_0(i)$ split in $E$ over $F$;
\item for a signature $s$, $k\geq0$, a finite set $Q$ of places and a character
  $\psi\colon G_E\to\mathcal O_{i,k}^\times$, the \emph{Bloch-Kato length} of
  $s$ twisted by $\psi$ is the length over $\mathcal O_{i,k}$, divided by the
  valuation of $\ell_i$, of the Selmer group of the representation
  $\rho_s\chi_s\psi$ of $G_E$ on $(K_{i,k}/\mathcal O_{i,k})^r$: the classes in
  $H^1(G_E,-)$ that lie in the image of the Bloch-Kato finite classes at the
  places over $\ell_i$, are arbitrary at the places of $Q$ and at the places
  over $g_0$, $g_1$ and in $P_D(i)$ that are not over $\ell_i$, and are
  unramified elsewhere; here $K_{i,k}$ is the fraction field of
  $\mathcal O_{i,k}$, and the finite classes at a place over $\ell_i$ are those
  of $\rho_s\chi_s$, carried over by the comparison data of $D$ at that place;
\item the \emph{bound} $\nu_i$ is the least $\nu\geq0$ such that
  $\varepsilon(g)^{tr}\not\equiv1\pmod{\ell_i^{\nu+1}}$ for some $g\in G_E$,
  where $\varepsilon$ is the $\ell_i$-adic cyclotomic character;
\item for $\lambda\in\Lambda_i$ and a character $\psi$ of $\mathrm{RCG}_i(Q)$
  with values in $\mathcal O_{i,k}^\times$, the \emph{theta value of $\lambda$
  twisted by $\psi$} is
  $\vartheta_i(Q,\lambda,\psi)=\sum_{\delta\in\mathrm{RCG}_i(Q)}\vartheta_i(Q,\lambda,\delta)\,\psi(\delta)$,
  and the \emph{bound} for $\lambda$ is
  $c_0(\lambda)=v\bigl(\sum_{\delta\in\mathrm{RCG}_i(\emptyset)}\vartheta_i(\emptyset,\lambda,\delta)\bigr)$,
  where $v$ is the valuation normalized by $v(\ell_i)=1$.
\end{enumerate}
\end{definition}
\leanref{Formalism/SharpRestriction/LevelLetterCocycles/DataN.lean}{Formalism.BasicNotions.NTw}
\leanref{Formalism/SharpRestriction/LevelLetterCocycles/DataN.lean}{Formalism.BasicNotions.NExp}
\leanref{Formalism/SharpRestriction/LevelLetterCocycles/IntegralHeckeForms.lean}{Formalism.MainNotions.IntegralHeckeForms.V0}
\leanref{Formalism/SharpRestriction/LevelLetterCocycles/IntegralHeckeForms.lean}{Formalism.MainNotions.IntegralHeckeForms.eThPsi}
\leanref{Formalism/SharpRestriction/LevelLetterCocycles/IntegralHeckeForms.lean}{Formalism.MainNotions.IntegralHeckeForms.eThPMt}
\leanref{Formalism/SharpRestriction/LevelLetterCocycles/BlochKatoPhases.lean}{Formalism.MainNotions.BlochKatoReduction.Gen.goodE}
\leanref{Formalism/SharpRestriction/LevelLetterCocycles/BlochKatoPhases.lean}{Formalism.MainNotions.BlochKatoReduction.Gen.sLen}
\leanref{Formalism/SharpRestriction/LevelLetterCocycles/BlochKatoPhases.lean}{Formalism.MainNotions.BlochKatoReduction.Gen.nu}
\leanref{Formalism/SharpRestriction/LevelLetterCocycles/BlochKatoPhases.lean}{Formalism.MainNotions.BlochKatoReduction.Gen.thEval}
\leanref{Formalism/SharpRestriction/LevelLetterCocycles/BlochKatoPhases.lean}{Formalism.MainNotions.BlochKatoReduction.Gen.c0}

\begin{remark}\label{th-lift-kernels}
Let a definite-side datum be given (Definition~\ref{th-definite-datum}). The
kernels $\vartheta_i(c,c')$ of its members are part of the data. At the cores
of the construction of the families, whose definite-side data are built from
true columns and have integral periods
(Section~\ref{section-pd-period-integrality}),
the lifts are required to come with \emph{lift objects}: lattices
$L_i(c)\subset E^{N(i)}$ attached to the classes and self-dual for the standard
Hermitian form,
the identification of $Y_i$ with the tuples $(g_\tau)_{\tau\in\mathcal T}$ of
invertible complex matrices of size $N(i)$ with $\delta_i(y,\tau)=\det g_\tau$,
and for each pair $(c,c')$ a function $\Theta_i(c,c')$ on tuples of complex
$N(i)\times N(i+1)$ matrices. The function $\Theta_i(c,c')$ is a finite
combination, with algebraic coefficients integral away from a finite set $S$
of primes, of the translates of a fixed function $e(L_i(c),L_{i+1}(c'))$ by
points that become points of the lattice attached to the pair
$(L_i(c),L_{i+1}(c'))$ after multiplication by an integer with prime factors
in $S$; and $\vartheta_i(c,c')$ sends the polynomial
function $Q$ of the entries of the matrices $g_\tau$ to the function whose value at $y'$
is the constant-coefficient differential operator determined by $Q$ and $y'$
applied to $\Theta_i(c,c')$ at the origin. A lift object for a member of the
datum is moreover required to use the determinant classes and characters of
the datum and an isometry of $X_i$ with $E^{N(i)}$, for the standard form,
under which $L_i(c)=L+\bar\pi L^\vee$ for $L=L_u(i,c)$ and $\pi$ as in
Definition~\ref{th-definite-datum}(4).
\end{remark}
\leanref{Formalism/SharpRestriction/GlobalSchurBounds/GL3Theta.lean}{Formalism.MainNotions.GL3Theta.liftTie}
\leanref{Formalism/SharpRestriction/GlobalSchurBounds/GL3Theta.lean}{Formalism.MainNotions.GL3Theta.lift_formula}
\leanref{Formalism/SharpRestriction/GlobalSchurBounds/GL3Theta.lean}{Formalism.MainNotions.GL3Theta.LiftObj.IsMem}

\section{Weil data and the first theta lifts}\label{section-th-weil}

\subsection{Weil data, theta grids and bottom triples}\label{subsection-th-weil-grids}

\begin{definition}\label{th-weil-datum}
Let $E/F$ be the host field of a definite-side datum
(Definition~\ref{th-definite-datum}), with element $\xi_0\in E$. Write
$\mathbf A_F$ for the adeles of $F$ and $\mathbf A_E=\mathbf A_F\otimes_FE$.
For a Hermitian space $S$ over $E$ (Definition~\ref{th-spaces}) let
$U(S)(\mathbf A)$ be the group of $g\in GL_{\operatorname{rk}S}(\mathbf A_E)$
with $\bar g^{\mathsf T}G_Sg=G_S$, and $U(S)(F)$ its subgroup of matrices
with entries in $E$. Let $S,T$ be Hermitian spaces, put
$\mathbb W=M_{\operatorname{rk}T\times\operatorname{rk}S}(\mathbf A_E)$ and
\[
\langle x,y\rangle=\operatorname{Tr}_{\mathbf A_E/\mathbf A_F}\bigl(\xi_0\operatorname{tr}(G_S^{-1}\bar y^{\mathsf T}G_Tx)\bigr)\qquad(x,y\in\mathbb W),
\]
and let $(g,h)\in U(S)(\mathbf A)\times U(T)(\mathbf A)$ act on $\mathbb W$ by
$x\mapsto hxg^{-1}$. Let $\psi$ be a character of $\mathbf A_F$, and for
$w_0\in\mathbb W$ and a function $f$ on $\mathbb W$ put
$(\rho(w_0)f)(w)=\psi(\tfrac12\langle w,w_0\rangle)f(w+w_0)$.
\begin{enumerate}
\item The character $\psi$ is \emph{standard} if it is continuous, unitary,
  nontrivial and trivial on $F$, and
  $\psi(x)=\exp\bigl(-2\pi i\sum_{\tau\in\mathcal T}\tau(x)\bigr)$ for
  $x\in F\otimes\mathbf R$, where each $\tau\in\mathcal T$
  (Definition~\ref{th-definite-datum}(2)) is extended to $F\otimes\mathbf R$.
\item A \emph{theta function} for $(S,T)$ is a continuous function $f$ on
  $\mathbb W$ such that
  $f(\gamma+w)=\psi(-\tfrac12\langle\gamma,w\rangle)f(w)$ for every
  $\gamma\in M_{\operatorname{rk}T\times\operatorname{rk}S}(E)$; there is an
  open subgroup $U\subset\mathbb W$ with $\rho(w_0)f=f$ for every $w_0\in U$
  with archimedean component $0$; and $t\mapsto f(w+\ell(t))$ is smooth for
  every $w$ and every continuous additive map $\ell\colon\mathbf R^a\to\mathbb W$
  with finite component $0$.
\item Let $\chi_S,\chi_T$ be continuous unitary characters of
  $\mathbf A_E^\times$. A \emph{Weil datum} for $(S,T)$ with characters
  $(\chi_S,\chi_T)$ is a homomorphism $\omega$ from
  $U(S)(\mathbf A)\times U(T)(\mathbf A)$ to the linear endomorphisms of the
  space of functions on $\mathbb W$ such that: $\psi$ is standard; $\chi_S$ and
  $\chi_T$ are trivial on $E^\times$ and their restrictions to
  $\mathbf A_F^\times$ are $\eta^{\operatorname{rk}S}$ and
  $\eta^{\operatorname{rk}T}$, where $\eta$ is the quadratic character of $E/F$;
  $\omega(g,h)$ preserves theta functions; on theta functions
  $\omega(g,h)\rho(w_0)=\rho(hw_0g^{-1})\omega(g,h)$;
  $(\omega(\gamma,\gamma')f)(0)=f(0)$ for $\gamma\in U(S)(F)$,
  $\gamma'\in U(T)(F)$ and theta functions $f$; $(g,h)\mapsto(\omega(g,h)f)(w)$
  is continuous; $\omega(g,h)$ preserves the $L^2$-norm of theta functions on a
  fundamental domain for $M_{\operatorname{rk}T\times\operatorname{rk}S}(E)$
  in $\mathbb W$; and at the places of $F$ split in $E$, $\omega(g,h)$ acts on
  linear functionals invariant under the corresponding half of $\mathbb W$
  through the scalar $r\,\chi_T(x)^{-1}\chi_S(y)$, with $r>0$ and $x,y$ the
  determinants of $g$ and $h$ there.
\item The \emph{theta kernel} of a Weil datum is
  $\theta_\omega(g,h;f)=(\omega(g,h)f)(0)$. For Haar measures and fundamental
  domains $\mathcal F_S$ and $\mathcal F_T$ for $U(S)(F)$ and $U(T)(F)$, the
  \emph{theta lifts} of functions $\varphi$ on $U(S)(\mathbf A)$ and $G$ on
  $U(T)(\mathbf A)$ are
  \[
  \theta_f(\varphi)(h)=\int_{\mathcal F_S}\theta_\omega(g,h;f)\varphi(g)\,dg,\qquad
  \theta_f(G)(g)=\int_{\mathcal F_T}\theta_\omega(g,h;f)G(h)\,dh.
  \]
\end{enumerate}
\end{definition}
\leanref{Formalism/LocalToGlobal/IndefiniteReal/LocalNewformFunction.lean}{Formalism.IndefR.Datum.IsThetaFunO}
\leanref{Formalism/LocalToGlobal/IndefiniteReal/LocalNewformFunction.lean}{Formalism.IndefR.Datum.IsWeilDatumO}
\leanref{Formalism/LocalToGlobal/IndefiniteReal/LocalNewformFunction.lean}{Formalism.IndefR.Datum.thetaKO}
\leanref{Formalism/LocalToGlobal/BesselHermitianTheta/TiedWeil.lean}{Formalism.L2.Theta.IsStdChar}
\leanref{Development/AbelianVarietyTheory/Definitions/TwistedInnerProducts/WittSmallGroupEmbeddingMemPointwiseStabilizerIso.lean}{GlobalTheory.MixPer.liftS}
\leanref{Development/AbelianVarietyTheory/Definitions/TwistedInnerProducts/WittSmallGroupEmbeddingMemPointwiseStabilizerIso.lean}{GlobalTheory.MixPer.liftT}

This is the Weil representation in the model of functions on the whole
symplectic space that are quasi-periodic under its rational points, a model
in which the theta distribution is evaluation at $0$. The conditions in (3)
are the defining properties of the Weil representation of \cite{Wei64}, with
the splitting of \cite{Kud94}. At the archimedean places the Fock model is
used, with the annihilation operators of Definition~\ref{th-grid-fock}(3); a
polynomial $P$ in the creation operators gives the vector $f\otimes P$.

\begin{definition}\label{th-grid-fock}
Let $D$ be a definite-side datum (Definition~\ref{th-definite-datum}) with host
field $E/F$, set $\mathcal T$ of embeddings and subfield
$K=\mathbf Q(\sqrt{-d})$ of $E$, and let $\mathfrak c$ and $\mathfrak b$ be the
ideals of $\mathcal O_K$ and of $\mathcal O_E$ of its host
(Definition~\ref{hf-host}). The \emph{ideal of the host} is the fractional
ideal $\mathfrak A=\mathfrak c\mathcal O_E\cdot\mathfrak b^{-1}$ of $E$. Let
$\psi$ be a standard character, let $S$ and $T$ be Hermitian spaces over $E$,
and let $\mathbb W=M_{\operatorname{rk}T\times\operatorname{rk}S}(\mathbf A_E)$,
its pairing $\langle\ ,\ \rangle$ and the translations $\rho$ be those of the
pair $(S,T)$ (Definition~\ref{th-weil-datum}). For $e\in\mathbf A_E$ let
$e\mathbb W$ be the set of $x\in\mathbb W$ with $ex=x$, the \emph{$e$-half} of
$\mathbb W$.
\begin{enumerate}
\item An \emph{archimedean family} is a family
  $\theta=(\theta_\tau)_{\tau\in\mathcal T}$ of ring homomorphisms
  $\theta_\tau\colon(F\otimes_{\mathbf Q}\mathbf R)\otimes_FE\to\mathbf C$ such
  that $\theta_\tau(1\otimes x)=\tau(x)$ for $x\in E$ and
  $a\mapsto\theta_\tau(a\otimes1)$ is continuous. For such a family, for
  $\tau\in\mathcal T$ and for $z\in M_{\operatorname{rk}T\times\operatorname{rk}S}(\mathbf C)$,
  let $z_\tau\in\mathbb W$ be the element with finite component $0$ whose
  archimedean component $\theta_\tau$ reads, entry by entry, as $z$, and
  $\theta_{\tau'}$ as $0$ for $\tau'\neq\tau$.
\item Let $\tau\in\mathcal T$. A \emph{Cartan involution} of $S$ at $\tau$ is a
  matrix $s\in M_{\operatorname{rk}S}(\mathbf C)$ with $s^2=1$ and
  $\bar s^{\mathsf T}\tau(G_S)s=\tau(G_S)$ such that $\tau(G_S)s$ is Hermitian
  and positive definite. An \emph{orthonormal basis} for $s$ is a matrix
  $B\in M_{\operatorname{rk}S}(\mathbf C)$ with $\bar B^{\mathsf T}\tau(G_S)B=J$
  and $sB=BJ$ for some diagonal matrix $J$ with diagonal entries $\pm1$.
\item Let $\theta$ be an archimedean family, $\tau\in\mathcal T$, and $s_S$ and
  $s_T$ Cartan involutions of $S$ and $T$ at $\tau$. For a function $f$ on
  $\mathbb W$ and $z\in M_{\operatorname{rk}T\times\operatorname{rk}S}(\mathbf C)$
  let $\partial_zf$ be the function
  $x\mapsto\frac{d}{dt}(\rho(tz_\tau)f)(x)|_{t=0}$. The \emph{annihilation
  operator} of $z$ at $\tau$ is $a_\tau(z)=\partial_z+i\,\partial_{J_\tau(z)}$,
  where $J_\tau(z)=i\,s_Tzs_S$. A function $f$ is a \emph{Fock vacuum vector at
  $\tau$} for $s_S$ and $s_T$ if $a_\tau(z)f=0$ for every $z$.
\item Let $\Lambda_S\subset E^{\operatorname{rk}S}$ and
  $\Lambda_T\subset E^{\operatorname{rk}T}$ be $\mathcal O_E$-lattices. The
  \emph{lattice of the pair} for $\Lambda_S$ and $\Lambda_T$ is the subgroup of
  $\mathbb W$ generated by the elements with archimedean component $0$ and
  finite component $\alpha a\,y\bar x^{\mathsf T}G_S$, for
  $\alpha\in\hat{\mathcal O}_F$, $a\in\mathfrak A$, $x\in\Lambda_S$ and
  $y\in\Lambda_T$.
\item A \emph{split idempotent} is an $e\in\mathbf A_E$ with $e^2=e$ and
  $e\bar e=0$ such that $e+\bar e$ is the idempotent of $\mathbf A_F$ at a finite
  place $v$ of $F$, with component $1$ at $v$ and $0$ at the other places. For a
  subgroup $\Gamma\subset\mathbb W$, a \emph{lattice vacuum} for $\Gamma$ is a
  theta function $f$ for $(S,T)$ (Definition~\ref{th-weil-datum}(2)) with
  $\rho(2u)f=f$ for every $u\in\Gamma$, and $\rho(u)f=f$ for every split
  idempotent $e$ and every $u\in\Gamma\cap e\mathbb W$.
\item Let $\theta$ be an archimedean family, let $s_S=(s_{S,\tau})_\tau$ and
  $s_T=(s_{T,\tau})_\tau$ be Cartan involutions of $S$ and $T$ at the
  $\tau\in\mathcal T$, and let $\Lambda$ be the lattice of the pair for lattices
  $\Lambda_S$ and $\Lambda_T$. A \emph{Fock vacuum vector} for $s_S$, $s_T$ and
  $\Lambda$ is a lattice vacuum $f$ for $\Lambda$ that is a Fock vacuum vector
  at every $\tau\in\mathcal T$ for $s_{S,\tau}$ and $s_{T,\tau}$, such that there
  are a Haar measure $\mu$ on $\mathbb W$ and a fundamental domain $\mathcal F$
  for $M_{\operatorname{rk}T\times\operatorname{rk}S}(E)$ in $\mathbb W$ with
  $\mu(\mathcal F)=1$, and $\int_{\mathcal F}|f|^2\,d\mu=1$ for every such $\mu$
  and $\mathcal F$.
\item Let $\Lambda$ be the lattice of the pair for lattices $\Lambda_S$ and
  $\Lambda_T$, let $e\in\mathbf A_E$, and let $f_0$ be a function on
  $\mathbb W$. A function $\Phi$ on $\mathbb W$ is \emph{Schwartz-Bruhat on
  $e\mathbb W$} if it vanishes off $e\mathbb W$ and there is an integer
  $k\geq1$ with $\Phi(x+ku)=\Phi(x)$ for all $x\in\mathbb W$ and
  $u\in\Lambda\cap e\mathbb W$, and with $kx\in\Lambda$ whenever
  $\Phi(x)\neq0$. A \emph{place model} of $f_0$ at $e$ is a linear map $\pi$
  on the functions on $\mathbb W$ such that, for every $\Phi$ that is
  Schwartz-Bruhat on $e\mathbb W$, $\pi$ carries $x'\mapsto\Phi(x'+x)$ to
  $\rho(x)\pi(\Phi)$ for $x\in e\mathbb W$, and
  $x'\mapsto\psi(\langle x',y\rangle)\Phi(x')$ to $\rho(y)\pi(\Phi)$ for
  $y\in\bar e\mathbb W$; and $\pi$ carries the indicator function of
  $\Lambda\cap e\mathbb W$ to $f_0$.
\end{enumerate}
% Lean: the ideal of the host is Host.AA = ccE * AA0 (AA0 = b^-1). (1) is IsInfFam theta (IsInfEmb (tauE D
% tau) (theta tau)) on Ainf F ⊗ E, and z_tau is lineAt theta tau z 1 (IsLineAt). (2) is IsCartanInv
% (gramAt D S tau) s; the orthonormal basis is the clause onb_is of GridData.Weil. (3) is annih and
% IsFockVacAt, with J = cplxO sS sT. (4) is latO S T LS LT (fun _ => True). (5) is IsSplitIdem (with
% IsPlaceIdem) and IsLatVac, together with IsThetaFunO. (6) is IsVacO psi theta S T sS sT LS LT, that is
% IsStdChar psi (-1) and IsVacuum (inf, theta, lat, fock, norm: AllSome IsCovolOne, nrm2 = 1). (7) is
% IsSBHalf and IsPlaceModel psi (symplO S T) (heisO psi S T) e Lambda 1 f0.
\end{definition}
\leanref{Formalism/LocalRingCardinality/DataHost.lean}{Formalism.BasicNotions.Host.AA}
\leanref{Formalism/LocalToGlobal/BesselHermitianTheta/IsAutFunDef.lean}{Formalism.L2.Theta.IsInfFam}
\leanref{Formalism/LocalToGlobal/IndefiniteReal/LocalNewformFunction.lean}{Formalism.IndefR.Datum.lineAt}
\leanref{Formalism/LocalToGlobal/IndefiniteReal/LocalNewformFunction.lean}{Formalism.IndefR.Datum.IsCartanInv}
\leanref{Formalism/LocalToGlobal/IndefiniteReal/LocalNewformFunction.lean}{Formalism.IndefR.Datum.annih}
\leanref{Formalism/LocalToGlobal/IndefiniteReal/LocalNewformFunction.lean}{Formalism.IndefR.Datum.IsFockVacAt}
\leanref{Formalism/LocalToGlobal/IndefiniteReal/LocalNewformFunction.lean}{Formalism.IndefR.Datum.latO}
\leanref{Formalism/LocalToGlobal/BesselHermitianTheta/TiedWeil.lean}{Formalism.L2.Theta.IsSplitIdem}
\leanref{Formalism/LocalToGlobal/IndefiniteReal/LocalNewformFunction.lean}{Formalism.IndefR.Datum.IsLatVac}
\leanref{Formalism/LocalToGlobal/IndefiniteReal/LocalNewformFunction.lean}{Formalism.IndefR.Datum.IsVacuum}
\leanref{Formalism/LocalToGlobal/IndefiniteReal/LocalNewformFunction.lean}{Formalism.IndefR.Datum.IsVacO}
\leanref{Formalism/LocalToGlobal/IndefiniteReal/LocalNewformFunction.lean}{Formalism.IndefR.Datum.IsSBHalf}
\leanref{Formalism/LocalToGlobal/IndefiniteReal/LocalNewformFunction.lean}{Formalism.IndefR.Datum.IsPlaceModel}

\begin{definition}\label{th-grid-places}
Let $D$ be a definite-side datum (Definition~\ref{th-definite-datum}) with
host field $E/F$, the subfield $K=\mathbf Q(\sqrt{-d})$ of $E$, the integers
$Q_S$, $q$ and $d$ and the ideal $\mathfrak c$ of $\mathcal O_K$ of its host,
and the level primes $g_0,g_1,p_{de},p'_{de}$. Let $r_\ast$ be the least prime
that does not divide $Q_Sqd(p_{de}-1)(p'_{de}-1)\,\mathrm N\mathfrak c$, and
let the \emph{auxiliary prime} $q_0$ be the least prime with
$q_0\equiv1\pmod{r_\ast Q_Sqd}$ that does not divide $\mathrm N\mathfrak c$. Let
$S_1$ be the set of the finite places of $F$ that contain one of $g_0$, $g_1$,
$p_{de}$ and $p'_{de}$. Let $P_g$ and $Q_t$, for $t$ in a set
$\mathcal T_\mu$, be finite sets of finite places of $E$. \emph{Places of a
theta grid} for $P_g$ and the sets $Q_t$ consist of a finite place $v_0$ of $F$
that contains $q_0$, the \emph{base place}, and a finite place $w(v)$ of $E$
for every finite place $v$ of $F$, such that:
\begin{enumerate}
\item for $v=v_0$ and for $v\in S_1$, $w(v)$ lies over $v$ and
  $w(v)\neq\overline{w(v)}$;
\item $w(v)\in P_g$ for every $v$ that contains $g_0$ or $g_1$;
\item for each $p\in\{p_{de},p'_{de}\}$ there is a prime $\mathfrak p$ of $K$
  such that $w(v)$ lies over $\mathfrak p$ for every $v$ that contains $p$;
\item every place $P$ of a set $Q_t$ that lies over a place $v\in S_1$ is
  $w(v)$;
\item no place of a set $Q_t$ lies over $v_0$.
\end{enumerate}
% Lean: ChiGrid.Places L uses only L.Pg and L.Q of the level system L. (1) is w_over and w_split
% (a split idempotent at w v, IsIdemAt, which exists exactly when w v is not conjugate to itself:
% MovedIdem.exists_isIdemAt and ModulusNorm.conjPl_ne_of_isIdemAt), (2) is w_g, (3) w_de, (4) w_Q
% and (5) levelPlace_ne_basePlace; v0 contains q0 D by basePlace_over_qZero. S1 D is primesOf of
% g0 g1 pde pde' (S1_mem_OmF); q_0 is q0 D = q0Av D and r_* is rPrAv D
% (leastSplitAuxiliaryPrimeThetaChain_spec, rPr_spec).
% Names (STYLE rule 13): q_0 and r_* name these primes throughout the book (Definition
% asm-twist-condition restates them). The helper prime r_* is not r = k+m-1, the rank of D
% (Definition th-ranks), which the book writes r and not r_0.
\end{definition}
\leanref{Development/AbelianVarietyTheory/Definitions/Divisors/BasePlaceNotMemChiGridBadPlaces.lean}{GlobalTheory.ChiGrid.Places}
\leanref{Development/AbelianVarietyTheory/Definitions/Divisors/LeastSplitAuxiliaryPrimeThetaChainSpec.lean}{GlobalTheory.ChiGrid.S1}
\leanref{Development/AbelianVarietyTheory/Definitions/Divisors/LeastSplitAuxiliaryPrimeThetaChainSpec.lean}{GlobalTheory.ChiGrid.q0}
\leanref{Development/AbelianVarietyTheory/Definitions/Divisors/LeastSplitAuxiliaryPrimeThetaChainSpec.lean}{GlobalTheory.ChiGrid.leastSplitAuxiliaryPrimeThetaChain_spec}
\leanref{Formalism/SharpRestriction/LevelLetterCocycles/TwGood.lean}{Formalism.MainNotions.rPrAv}
\leanref{Formalism/SharpRestriction/LevelLetterCocycles/TwGood.lean}{Formalism.MainNotions.q0Av}

\begin{definition}\label{th-grid-centres}
Let $D$ be a definite-side datum (Definition~\ref{th-definite-datum}) with
rank datum $(k,m,t)$, $r=k+m-1$ (Definition~\ref{th-ranks}) and set
$\mathcal T$ of embeddings, and for $\tau\in\mathcal T$ let
$v_1(\tau),\ldots,v_{k-1}(\tau)$ be the $v$-part of its Harish-Chandra
parameter $\lambda(\tau)$ (Definition~\ref{th-hc-parameter}). Put $n_0=r+1$
and $N=k+m(t+r+7)$, and let $\ell$ be the largest of the integers
$v_1(\tau)-v_{k-1}(\tau)$, $\tau\in\mathcal T$. For a real number $c$ let
$\gamma(c)$ be the least integer $g$ with
\[
cN^3+cn_0^7+cn_0^3\sqrt\ell\leq g,\qquad 2N^3+2r+14<g,\qquad
48n_0^2\sqrt m\leq g.
\]
For $\tau\in\mathcal T$ and $0\leq l\leq m$ put
$a_l(\tau)=v_1(\tau)+\tfrac12+(m+1-l)\gamma(c)$, and for $1\leq l\leq m$ let
the \emph{centre} $c_l(\tau)$ be the integer
$\lfloor(a_{l-1}(\tau)+a_l(\tau))/2\rfloor$.
% Lean: gamma(c) is ChiGrid.gap D c, the least g with GapBound D c g (isLeast_gap), with nRk D = n_0,
% bigN D = N (sLen D = r + 6) and cloudLen D = ell. For G = GridData.ofHCParam (D.hc tau) (sLen D)
% (gap D c) 0, whose v 2 is v_1(tau), a_l(tau) is G.a l and c_l(tau) = ChiGrid.cen D c tau l is the
% floor of G.mid l.
\end{definition}
\leanref{Development/AbelianVarietyTheory/Definitions/Divisors/LeastSplitAuxiliaryPrimeThetaChainSpec.lean}{GlobalTheory.ChiGrid.cloudLen}
\leanref{Development/AbelianVarietyTheory/Definitions/Divisors/LeastSplitAuxiliaryPrimeThetaChainSpec.lean}{GlobalTheory.ChiGrid.gap}
\leanref{Development/AbelianVarietyTheory/Definitions/Divisors/LeastSplitAuxiliaryPrimeThetaChainSpec.lean}{GlobalTheory.ChiGrid.isLeast_gap}
\leanref{Development/AbelianVarietyTheory/Definitions/Divisors/LeastSplitAuxiliaryPrimeThetaChainSpec.lean}{GlobalTheory.ChiGrid.cen}
\leanref{Development/AbelianVarietyTheory/Definitions/Foundations/StrictInterlaceIffDescEnum.lean}{GlobalTheory.SignedHC.GridData.a}
\leanref{Development/AbelianVarietyTheory/Definitions/Foundations/StrictInterlaceIffDescEnum.lean}{GlobalTheory.SignedHC.GridData.mid}

\begin{definition}\label{th-grid}
Let $D$ be a definite-side datum (Definition~\ref{th-definite-datum}) with
host field $E/F$, set $\mathcal T$ of embeddings and Harish-Chandra parameters
$\lambda(\tau)$ with $e$-parts $e_1(\tau),\ldots,e_m(\tau)$. Let
$K=\mathbf Q(\sqrt{-d})$ be the subfield of $E$, $Q_S$, $q$ and $d$ the
integers of the host and $g_0,g_1,p_{de},p'_{de}$ the level primes; let $\eta$
be the quadratic character of $E/F$, and $k_P$ the residue field of a finite
place $P$ of $E$. A Hecke character of $E$ is \emph{tamely ramified} at $P$ if
it is trivial on the units of $E_P$ congruent to $1$ modulo $P$. A
\emph{level system} for $D$ consists of a set $\mathcal T_\mu$ of
\emph{tuples} with a \emph{base tuple} $\mu_0$ and the following data, subject
to the conditions stated with them.
\begin{enumerate}
\item[(L1)] \emph{Primes over $g_0$ and $g_1$.} Primes $\mathfrak p_0$ and
  $\mathfrak p_1$ of $K$ with $g_l\in\mathfrak p_l$, and injective characters
  $\omega_l$ of $(\mathbf Z/g_l)^\times$, for $l=0,1$. The set $P_g$ is the set
  of the places of $E$ over $\mathfrak p_0$ or $\mathfrak p_1$.
  % Lean: C.P (PiBar: pb, pb_over) and C.W (Omg: om, gen) of the CharFrame C = LevCh.C; P_g is
  % LevCh.Pg with LevCh.Is.Pg (PP in Pg iff C.P.Div l PP, that is PP.under K = pb l, for some l).
\item[(L2)] \emph{The characters $\mu^0_j$.} Unitary Hecke characters
  $\mu^0_1,\ldots,\mu^0_m$ of $E$. Each $\mu^0_j$ is trivial on
  $\mathbf A_F^\times$; its archimedean component is
  $k\mapsto\prod_{\tau\in\mathcal T}(\tau(k)/|\tau(k)|)^{-2e_j(\tau)}$ on
  $E^\times$; it is unramified at every finite place that contains neither
  $g_0$ nor $g_1$, and tamely ramified at every finite place; and at every
  place $w$ of $E$ over $\mathfrak p_l$ its component at $w$ takes the value
  $\omega_l(a)^{-2j}$ at every positive integer $a$ prime to $g_l$.
  % Lean: C.mu0 with C.isMu0 j : IsMu0 C.P C.W (eExp D j) j (C.mu0 j), 0-based (C.mu0 j is mu^0_{j+1}):
  % unitary, triv (TrivOn F), inf (HasInfQ embPsi (-eExp D j), eExp D j = (D.hc tau).e j, with embPsi
  % at eTau tau equal to embTau tau), unr, tame, atg (locK w a = (om l a ^ (2 (j+1)))^-1). The level
  % system's LevCh.mu0 (j+1), a unitary character of the ideles, is tied to C.mu0 j by LevCh.Is.mu0
  % (ChrEq); the book writes mu^0_{j+1} for both.
\item[(L3)] \emph{The characters $\eta_L$, $\chi_\Psi$, $\delta$ and
  $\delta'$.} Unitary Hecke characters $\eta_L$ and $\chi_\Psi$ of $E$,
  unramified at every finite place, with archimedean components
  $k\mapsto\prod_{\tau\in\mathcal T}\tau(k)/|\tau(k)|$ and
  $k\mapsto\prod_{\tau\in\mathcal T}(\tau(k)/|\tau(k)|)^2$ on $E^\times$; the
  restriction of $\eta_L$ to $\mathbf A_F^\times$ is $\eta$, and $\chi_\Psi$ is
  trivial on $\mathbf A_F^\times$. Hecke characters $\delta$ and $\delta'$ of
  $E$ of finite order, trivial on $\mathbf A_F^\times$ and tamely ramified at
  every finite place, such that $\delta$ is unramified at every finite place
  that does not contain $p_{de}$, and at every place $P$ that contains
  $p_{de}$ its restriction to the units of $E_P$ has order $15$; and $\delta'$
  likewise with $p'_{de}$.
  % Lean: C.eta with C.isEta : IsEtaT (unitary, unr, res = ResEta F, inf = HasInf embPsi 1); the level
  % system's LevCh.eta is tied to it by LevCh.Is.eta (ChrEq), and the book writes eta_L for both.
  % C.chiPsi with C.isChi : IsChiPsi (unitary, unr, triv, inf = HasInf embPsi 2). C.del and C.del' with
  % C.isDel : IsDelta pde and C.isDel' : IsDelta pde' (fin, triv, tame, unr, and ine: at a place over
  % p, for every natural number n, del(x)^n = 1 for all units x if and only if 15 divides n).
\item[(L4)] \emph{Tuple characters.} For each tuple $t$ a Hecke character
  $\xi_t$ of $E$ of finite order, trivial on $\mathbf A_F^\times$ and tamely
  ramified at every finite place, with $t\mapsto\xi_t$ injective and
  $\xi_{\mu_0}=1$. The tuple $t$ is $\mu_{t,j}=\mu^0_j\xi_t$ ($1\leq j\leq m$).
  % Lean: LevCh.xiH t, tied to LevCh.xi t by LevCh.Is.xi (ChrEq), with xi_fin, xi_triv, xi_tame,
  % xi_inj and xi_base; mu_{t,j} is LevCh.mu t j = Lift.tup LevCh.mu0 (LevCh.xi t) j.
\item[(L5)] \emph{Ramified sets.} For each tuple $t$ a finite set $Q_t$ of
  finite places of $E$ such that $\xi_t$ is ramified at a finite place $P$ if
  and only if $P$ or $\bar P$ lies in $Q_t$. Each $P\in Q_t$ has a residue
  field of prime order $q_P\neq g_0,g_1$, with $q_P\notin P^2$, and contains
  none of $Q_S$, $q$, $d$, the prime of the system of spaces
  (Definition~\ref{th-spaces}(1)) and the primes $T_u$ of the words $u$
  (Definition~\ref{th-definite-datum}(3)); no place of $Q_t$ is the complex
  conjugate of a place of $Q_t$, itself included; and for each
  $p\in\{p_{de},p'_{de}\}$ there is a prime $\mathfrak p$ of $K$ under every
  place that contains $p$ and lies in some $Q_t$.
  % Lean: LevCh.Q t with LevCh.Is.Q_ram, Q_pl (IsTwistPlaces D.host.g: card k_P = q prime, q not in
  % P^2, q != g l, and not IsConjPl PP PP' for PP, PP' in Q t), Q_off ((condN * D.sp.T : O_E) not in PP,
  % condN = QS q d, and (D.mem u).A.T not in PP for every word u) and Q_de.
\item[(L6)] \emph{Residue characters.} Characters $\chi_{j,P}$ of $k_P^\times$
  for $1\leq j\leq m$ and $P\in P_g$, and $\xi_{t,P}$ of $k_P^\times$ for each
  tuple $t$ and $P\in Q_t$, such that the components at $P$ of $\mu^0_j$ and
  of $\xi_t$ take the values $\chi_{j,P}(a\bmod P)^{-1}$ and
  $\xi_{t,P}(a\bmod P)$ at every $a\in\mathcal O_E\setminus P$.
  % Lean: LevCh.chi (j+1) P with LevCh.Is.chi (IsRes (C.mu0 j) P (LevCh.chi (j+1) P)^-1), and
  % LevCh.xiR t P with LevCh.Is.xiR (IsRes (xiH t) P (xiR t P)); IsRes chi P c says c (a mod P) =
  % chi.locK P a for a in O_E \ P. LevCh.chi and xiR are families over all places, constrained only
  % at P_g and at Q t.
\item[(L7)] \emph{Lattices.} Self-dual $\mathcal O_E$-lattices
  $\Lambda_{Y_j}\subset Y_j$ for $j<m$, $\Lambda_{X'}\subset X'$, and
  $\Lambda_{V_i}\subset V_i$ and $\Lambda_{W_i}\subset W_i$ for $1\leq i\leq m$
  (Definition~\ref{th-spaces}), each with an \emph{adelic matrix}: a matrix
  $\beta\in GL_n(E\otimes_F\mathbf A_{F,f})$, $n$ the rank of the space, that
  carries the integral vectors of $(E\otimes_F\mathbf A_{F,f})^n$ onto the
  closure of the lattice. The lattices $\Lambda_{X'}$, $\Lambda_{V_i}$ and
  $\Lambda_{W_i}$ are the lattices $L_0$ of their spaces, and for each $j<m$
  there is an integer $h\geq0$ with $(g_0g_1)^hL_0\subset\Lambda_{Y_j}$ and
  $(g_0g_1)^h\Lambda_{Y_j}\subset L_0$.
  % Lean: LevCh.LambdaY j, LambdaX, LambdaV i, LambdaW i (the fields named with a capital lambda) and
  % betaY, betaX, betaV, betaW, with LevCh.Is.betaY (j < m), betaX, betaV, betaW: IsSDFrame S beta Lam N,
  % that is IsBaseFrame S beta Lam (N^c) for some c (IsFrame beta Lam, dualL Lam = Lam, N^c L0c in Lam and
  % N^c Lam in L0c), with N = g0 g1 for Y_j and N = 1 otherwise, which makes Lam = L0c. Lattices are in the
  % coordinates of the basis of the space; V_i and W_i are Vc D (i-1) = D.sp.V i and Wc D (i-1) = D.sp.W i.
\end{enumerate}
A \emph{theta grid} for a level system consists of Hermitian spaces with
correct signatures (Definition~\ref{th-spaces}) and self-dual lattices:
\begin{enumerate}
\item $V_{0,j}=Y_j$ with the lattice $\Lambda_{Y_j}$ for $j<m$, and
  $V_{0,m}=X'$ with $\Lambda_{X'}$ (Definition~\ref{th-spaces} and (L7));
\item $V_{i,j}$ for $1\leq i\leq m$ and $0\leq j\leq m$, of rank
  $r+jt+i(r+6)$ and signature
  $(\operatorname{rk}V_{i,j}-\lfloor(m-j)/2\rfloor,\ \lfloor(m-j)/2\rfloor)$
  at every embedding;
\item \emph{core spaces} $P_i$ for $0\leq i\leq m$, of rank $r+1+i(r+5)$ and
  signature $(\operatorname{rk}P_i-\lfloor(m-i+1)/2\rfloor,\ \lfloor(m-i+1)/2\rfloor)$;
\item \emph{top spaces} $Q_0=P_m$ and $Q_j$ for $1\leq j\leq m$, of rank
  $r+1+m(r+5)+j(t+1)$ and signature $(\operatorname{rk}Q_j,0)$ at every
  embedding;
\end{enumerate}
with isometric embeddings $Y_0\subset P_0$, $P_i\subset V_{i,0}$ for
$1\leq i\leq m$, $Q_j\subset V_{m,j}$ for $j<m$ (for $j=0$ this is
$P_m\subset V_{m,0}$), and $V_{m,m}\subset Q_m$, whose orthogonal complements
have ranks $c'$ equal to $1$, $i-1$, $m-1-j$ and $1$ and bases
$a_1,\ldots,a_{c'}$ with antidiagonal Gram matrices:
$\langle a_k,a_l\rangle=1$ if $k+l=c'+1$ and $k\neq l$, and
$\langle a_k,a_l\rangle=0$ if $k+l\neq c'+1$; if $c'$ is odd, the middle
vector has $\langle a_k,a_k\rangle=1$ for the first and the last embedding and
$\langle a_k,a_k\rangle=-1$ for the others.
% Lean: Vertices: V_is, P_is and X_is (LSp.IsV: rank, SigOK, signature, IsSelfDualC), with rankV,
% rankP and rankX; the embeddings are dec0, colDec i, rowDecS j (1 <= j < m; for j = 0 the embedding
% is colDec m) and topDec, each a HermitianOrthogonalExtensionData A B c eps (an isometric embedding,
% a matrix a whose columns are orthogonal to A with Gram matrix antiDiagonalSignGramMatrix c eps) with
% eps = 1 for dec0 and topDec and eps = -1 for colDec and rowDecS.
A theta grid also has a triple $(c,C_F,b)$ of
real numbers, its \emph{constants}; places of a theta grid for the set $P_g$
and the sets $Q_t$ of the level system (Definition~\ref{th-grid-places}),
with base place $v_0$ over the auxiliary prime $q_0$, places $w(v)$ of $E$ and
the set $S_1$ of places of $F$ given there; and unitary Hecke characters
$\chi_1,\ldots,\chi_m$ of $E$. Put $S_0=\{v_0\}\cup S_1$, the set of
\emph{distinguished places}. For $1\leq i\leq m$:
\begin{enumerate}
\item[(a)] $\chi_i$ is trivial on $\mathbf A_F^\times$;
\item[(b)] the archimedean component of $\chi_i$ is
  $k\mapsto\prod_{\tau\in\mathcal T}(\tau(k)/|\tau(k)|)^{2c_i(\tau)}$ on
  $E^\times$, where $c_i(\tau)$ is the centre of
  Definition~\ref{th-grid-centres} for the constant $c$;
\item[(c)] $\chi_i$ is unramified at every finite place of $E$ that lies over
  no place of $S_0$;
\item[(d)] at every place $P$ of $E$ over $v_0$, $\chi_i$ is ramified and is
  trivial on the units of $E_P$ congruent to $1$ modulo $P$;
\item[(e)] at every place $P$ of $E$ over a place of $S_1$, $\chi_i$ is trivial
  on the units of $E_P$ congruent to $1$ modulo $P^2$ and nontrivial on those
  congruent to $1$ modulo $P$;
\item[(f)] at every place $P$ of $E$ that contains $p_{de}$ or $p'_{de}$,
  $\chi_i$ is trivial on the roots of unity in $E_P$ of order dividing
  $\mathrm NP-1$.
\end{enumerate}
The \emph{depth} of $v\in S_0$ is $d(v)=2$ for $v\neq v_0$, and
$d(v_0)=\lfloor\max(s,0)/\log q_0\rfloor+1$, where
$s=C_Ff^2(r+1)^6\log\bigl(b(r+1)g_0g_1p_{de}p'_{de}|d_E|\bigr)$ and $d_E$ is
the discriminant of $E$. The \emph{modulus} of the grid is
$\mathfrak m=\prod_{v\in S_0}w(v)^{d(v)}$.
% Lean: the characters are ChiGrid.Chars pl c with c = cGap, for 1 <= i <= m: chi_eq ties chi i
% (unitary, on the ideles) to the Hecke character chiH i; (a) triv (HChar.TrivOn F), (b) inf
% (HChar.HasInf embTau with exponents 2 * cen D c tau i), (c) unr (off Places.S0), (d) lev0
% (LevelLe 1 and not UnrAt), (e) lev1 (LevelLe 2 and not LevelLe 1), (f) wild. The depth is
% Places.depth (depth0 D cDep b at v0, from sizeP, and 2 elsewhere), the modulus Places.modul.

The lattices of the spaces are tied by the embeddings. Let
$\mathfrak A=\mathfrak c\mathcal O_E\cdot\mathfrak b^{-1}$ be the ideal of the
host (Definition~\ref{th-grid-fock}), for its ideals $\mathfrak c$ of
$\mathcal O_K$ and $\mathfrak b$ of $\mathcal O_E$ (Definition~\ref{hf-host}).
For each of the embeddings above,
of a space with lattice $\Lambda$, with complement basis $a_1,\ldots,a_{c'}$,
the lattice of the larger space is $\Lambda+\sum_{k=1}^{c'}I_ka_k$ for
fractional ideals $I_k$ of $E$: $I_k=\mathcal O_E$ for $Y_0\subset P_0$ and
$V_{m,m}\subset Q_m$, and $I_k=\mathfrak a^{c'+1-k}\mathfrak a'^{-k}$
with
$(\mathfrak a,\mathfrak a')=(\mathfrak A\bar{\mathfrak m}^{-1},\bar{\mathfrak A}\mathfrak m^{-1})$
for $P_i\subset V_{i,0}$, and with
$(\mathfrak a,\mathfrak a')=(\mathfrak A,\bar{\mathfrak A})$ for
$Q_j\subset V_{m,j}$, $1\leq j<m$.
% Lean: Vertices L (pl.stA sc) (pl.stAbar sc): lat0 and topLat (the ideals 1), colLat (stairId with
% Places.stA = AA * modulBar^-1 and Places.stAbar = AAbar * modul^-1) and rowLat (stairId with AA and
% AAbar); HermitianOrthogonalExtensionData.lat is sumLat, and stairId a a' c p = a^(c-p) (a'^(p+1))^-1
% for 0-based p, the book's I_{p+1}. AA is Host.AA = ccE * b^-1 and AAbar its conjugate.

The \emph{grid pairs} are the following pairs $p=(S,T)$ of spaces of the grid,
each with a character $c_p$:
\begin{enumerate}
\item[(G1)] $(V_{i+1,j},V_{i,j})$ for $0\leq i<m$ and $0\leq j\leq m$, with
  $c_p=\chi_{m-i}$;
\item[(G2)] for each tuple $t$, the pair $(V_{i,j},V_{i,j+1})$ at $t$, for
  $0\leq i\leq m$ and $0\leq j<m$, with $c_p=\mu_{t,j+1}$;
\item[(G3)] $(P_{i+1},P_i)$ for $0\leq i<m$, with $c_p=\chi_{m-i}$;
\item[(G4)] for each tuple $t$, the pair $(Q_j,Q_{j+1})$ at $t$, for
  $0\leq j<m$, with $c_p=\mu_{t,j+1}$.
\end{enumerate}
The \emph{characters} of a grid pair $p=(S,T)$ are
$\chi_{S,p}=c_p\eta_L^{-\epsilon(S)}$ and $\chi_{T,p}=c_p\eta_L^{-\epsilon(T)}$,
where $\epsilon(X)\in\{0,1\}$ is the parity of $\operatorname{rk}X$.
% Lean: the grid pairs are ChiGrid.Pair T: (G1) chi i j, (G2) mu t i j, (G3) core i, (G4) top t j, with
% Pair.InRange m for the ranges; the spaces are GridData.pS and pT (pS of chi i j and core i is negSp);
% c_p is GridData.chr p (G.ch.chi (m - i) or LevCh.mu t (j + 1)), and chi_{S,p}, chi_{T,p} are
% GridData.chrS p and chrT p = Lift.attChr (chr p) LevCh.eta n (chr p * eta^-1 for n odd, chr p for n
% even) at the ranks n of pS p and pT p.

Write $\Lambda_X$ for the lattice of a space $X$ of the grid. A \emph{grid Weil
datum} consists of the following data, subject to the conditions stated with
them (Definition~\ref{th-grid-fock}).
\begin{enumerate}
\item[(W1)] A standard character $\psi$ (Definition~\ref{th-weil-datum}(1)).
  % Lean: GridData.Weil.psi with psi_std : IsStdChar D psi (-1).
\item[(W2)] An archimedean family $\theta$.
  % Lean: theta with theta_is : IsInfFam theta.
\item[(W3)] \emph{Cartan involutions.} For each space $X$ of the grid, that is
  $V_{i,j}$, $P_i$ and $Q_j$ for $0\leq i,j\leq m$, with $Q_0=P_m$ taken twice,
  once as $Q_0$ and once as $P_m$, and for each $\tau\in\mathcal T$, a Cartan
  involution $s_X(\tau)$ of $X$ at $\tau$.
  % Lean: cart X tau with cart_is X : IsCartanInv (gramAt D (G.V X) tau) (cart X tau) for the nodes X with
  % Node.InRange m (grid i j, core i, top j); top 0 and core m are distinct nodes with the same space
  % (node_top_zero), so they carry separate data.
\item[(W4)] \emph{Orthonormal bases.} For each such $X$ and $\tau$, an
  orthonormal basis $B_X(\tau)$ for $s_X(\tau)$.
  % Lean: onb X tau with onb_is X : there is eps : Fin rk -> Bool with (onb X tau)^H (gramAt D (G.V X) tau)
  % (onb X tau) = Jm eps and cart X tau * onb X tau = onb X tau * Jm eps; Jm eps is the diagonal sign matrix.
\item[(W5)] \emph{Weil data.} For each grid pair $p=(S,T)$, a Weil datum
  $\omega_p$ for $(S,T)$ with the character $\psi$ and the characters
  $(\chi_{S,p},\chi_{T,p})$ (Definition~\ref{th-weil-datum}(3)), the first space
  of the pairs (G1) and (G3) taken with the negated form.
  % Lean: om p with om_is p : IsWeilDatumO psi (pS p) (pT p) (chrS p) (chrT p) (om p) for Pair.InRange m.
\item[(W6)] \emph{Vacuum vectors.} For each $0\leq i<m$, a Fock vacuum vector
  $f_i$ on $\mathbb W(-V_{i+1,0},V_{i,0})$ for $\theta$, the Cartan involutions
  $-s_{V_{i+1,0}}(\tau)$ and $s_{V_{i,0}}(\tau)$, and the lattice $\Lambda_i$ of
  the pair for $\Lambda_{V_{i+1,0}}$ and $\Lambda_{V_{i,0}}$, the \emph{lattice
  of the row pair}.
  % Lean: vac i with vac_is i : IsVacO psi theta (pS (chi i 0)) (pT (chi i 0)) (cartS cart (chi i 0))
  % (onT cart (chi i 0)) (latS (chi i 0)) (latT (chi i 0)) (vac i) for i < m, with pS (chi i 0) =
  % negSp (V (grid (i+1) 0)), cartS cart (chi i 0) = -cart (grid (i+1) 0) and onT cart (chi i 0) =
  % cart (grid i 0).
\item[(W7)] \emph{Coset vectors.} Let $e=\sum_{v\in S_0}1_{w(v)}$, where, for a
  place $P$ of $E$ that is not conjugate to itself, $1_P\in\mathbf A_E$ is the
  idempotent with component $1$ at $P$ and $0$ at the other places; then
  $e\bar e=0$, since the places $w(v)$ lie over distinct places of $F$ and are
  not conjugate to themselves (Definition~\ref{th-grid-places}(1)). For
  $0\leq i<m$ let $M_i=\operatorname{Hom}(\Lambda_{V_{i+1,0}},\Lambda_{V_{i,0}})$
  be the module of the matrices
  $x\in M_{\operatorname{rk}V_{i,0}\times\operatorname{rk}V_{i+1,0}}(E)$ with
  $x\Lambda_{V_{i+1,0}}\subset\Lambda_{V_{i,0}}$, and let the \emph{closure} of
  $\mathfrak mM_i$ be the subgroup of $\mathbb W(-V_{i+1,0},V_{i,0})$ generated
  by the elements with archimedean component $0$ and finite component
  $\alpha y$, for $\alpha\in\hat{\mathcal O}_F$ and $y\in\mathfrak mM_i$. For
  each $0\leq i<m$ and each class $\kappa\in M_i/\mathfrak mM_i$, a function
  $f_{i,\kappa}$ on $\mathbb W(-V_{i+1,0},V_{i,0})$, the \emph{coset vector} of
  $\kappa$, such that $f_i$ has a place model at $e$ for the lattice
  $\Lambda_i$, and $f_{i,\kappa}=\pi(\Phi_\kappa)$ for every such place model
  $\pi$. Here $\Phi_\kappa$ is the indicator function of the set of
  $x\in e\mathbb W(-V_{i+1,0},V_{i,0})$ such that $x-ex_0$ lies in the closure
  of $\mathfrak mM_i$ for some $x_0\in M_i$ in the class $\kappa$.
  % Lean: cosVec i kappa with cosVec_is : there is e with IsIdemSum e G.pl.W0 and, for i < m and kappa :
  % G.CosAt i = Cos (Λ (grid (i+1) 0)) (Λ (grid i 0)) modul, IsThetaAt psi (symplO _ _) (heisO psi _ _) e
  % (latO _ _ (latS (chi i 0)) (latT (chi i 0)) fun _ => True) 1 (vac i) (cosInd e kappa) (cosVec i kappa).
  % W0 is the image of S0 under w; IsIdemSum e W0 makes e a sum of elements IsIdemAt (es P) P, each of which
  % is 1_P (a split idempotent at P), with e ebar = 0. M_i is homLat, cosInd e kappa is the indicator of
  % cosSet e kappa, and the closure is hatMx.
\end{enumerate}
\end{definition}
\leanref{Development/AbelianVarietyTheory/Definitions/Divisors/BasePlaceNotMemChiGridBadPlaces.lean}{GlobalTheory.ChiGrid.GridData}
\leanref{Development/AbelianVarietyTheory/Definitions/Divisors/BasePlaceNotMemChiGridBadPlaces.lean}{GlobalTheory.ChiGrid.Vertices}
\leanref{Development/AbelianVarietyTheory/Definitions/Divisors/BasePlaceNotMemChiGridBadPlaces.lean}{GlobalTheory.ChiGrid.Chars}
\leanref{Development/AbelianVarietyTheory/Definitions/Divisors/BasePlaceNotMemChiGridBadPlaces.lean}{GlobalTheory.ChiGrid.Consts}
\leanref{Development/AbelianVarietyTheory/Definitions/Divisors/BasePlaceNotMemChiGridBadPlaces.lean}{GlobalTheory.ChiGrid.Places.depth}
\leanref{Development/AbelianVarietyTheory/Definitions/Divisors/BasePlaceNotMemChiGridBadPlaces.lean}{GlobalTheory.ChiGrid.Places.modul}
\leanref{Development/AbelianVarietyTheory/Definitions/Divisors/LeastSplitAuxiliaryPrimeThetaChainSpec.lean}{GlobalTheory.ChiGrid.depth0}
\leanref{Development/AbelianVarietyTheory/Definitions/Divisors/BasePlaceNotMemChiGridBadPlaces.lean}{GlobalTheory.ChiGrid.GridData.Weil}
\leanref{Formalism/LocalToGlobal/IndefiniteReal/LinksDefs.lean}{Formalism.IndefR.Links.LevCh}
\leanref{Formalism/LocalToGlobal/IndefiniteReal/LinksDefs.lean}{Formalism.IndefR.Links.LevCh.Is}
\leanref{Formalism/LocalToGlobal/IndefiniteReal/LinksDefs.lean}{Formalism.IndefR.Links.LevCh.mu}
\leanref{Formalism/LocalToGlobal/UnitaryNorms/CharTied.lean}{Formalism.L2.Host.CharFrame}
\leanref{Formalism/LocalToGlobal/UnitaryNorms/CharHost.lean}{Formalism.L2.Host.PiBar}
\leanref{Formalism/LocalToGlobal/UnitaryNorms/CharHost.lean}{Formalism.L2.Host.Omg}
\leanref{Formalism/LocalToGlobal/UnitaryNorms/CharHost.lean}{Formalism.L2.Host.IsMu0}
\leanref{Formalism/LocalToGlobal/UnitaryNorms/CharHost.lean}{Formalism.L2.Host.IsEtaT}
\leanref{Formalism/LocalToGlobal/UnitaryNorms/CharHost.lean}{Formalism.L2.Host.IsChiPsi}
\leanref{Formalism/LocalToGlobal/UnitaryNorms/CharHost.lean}{Formalism.L2.Host.IsDelta}
\leanref{Formalism/LocalToGlobal/IndefiniteReal/SubgroupCharExtensionOneApply.lean}{Formalism.IndefR.Lev.IsTwistPlaces}
\leanref{Formalism/LocalToGlobal/IndefiniteReal/LinksDefs.lean}{Formalism.IndefR.Links.IsRes}
\leanref{Formalism/LocalToGlobal/IndefiniteReal/LinksDefs.lean}{Formalism.IndefR.Links.IsSDFrame}
\leanref{Development/AbelianVarietyTheory/Definitions/Divisors/CmHostImaginaryGeneratorNeZero.lean}{GlobalTheory.ChiGrid.HermitianOrthogonalExtensionData}
\leanref{Development/AbelianVarietyTheory/Definitions/Divisors/CmHostImaginaryGeneratorNeZero.lean}{GlobalTheory.ChiGrid.stairId}
\leanref{Development/AbelianVarietyTheory/Definitions/Divisors/BasePlaceNotMemChiGridBadPlaces.lean}{GlobalTheory.ChiGrid.Places.stA}
\leanref{Formalism/LocalRingCardinality/DataHost.lean}{Formalism.BasicNotions.Host.AA}
\leanref{Development/AbelianVarietyTheory/Definitions/Divisors/BasePlaceNotMemChiGridBadPlaces.lean}{GlobalTheory.ChiGrid.Pair}
\leanref{Development/AbelianVarietyTheory/Definitions/Divisors/BasePlaceNotMemChiGridBadPlaces.lean}{GlobalTheory.ChiGrid.GridData.chr}
\leanref{Development/AbelianVarietyTheory/Definitions/Divisors/BasePlaceNotMemChiGridBadPlaces.lean}{GlobalTheory.ChiGrid.GridData.chrS}
\leanref{Development/AbelianVarietyTheory/Definitions/Divisors/BasePlaceNotMemChiGridBadPlaces.lean}{GlobalTheory.ChiGrid.GridData.chrT}
\leanref{Formalism/LocalToGlobal/IndefiniteReal/OrthogonalComplementCharacterisation.lean}{Formalism.IndefR.Lift.attChr}
\leanref{Formalism/LocalToGlobal/IndefiniteReal/LocalNewformFunction.lean}{Formalism.IndefR.Datum.IsThetaAt}
\leanref{Formalism/LocalToGlobal/IndefiniteReal/DatumVec.lean}{Formalism.IndefR.Datum.IsIdemSum}
\leanref{Formalism/LocalToGlobal/HeckeThetaElements/IsHeckeElt.lean}{Formalism.L2.NTheta.IsIdemAt}
\leanref{Development/AbelianVarietyTheory/Definitions/Divisors/BasePlaceNotMemChiGridBadPlaces.lean}{GlobalTheory.ChiGrid.Places.W0}
\leanref{Development/AbelianVarietyTheory/Definitions/Divisors/BasePlaceNotMemChiGridBadPlaces.lean}{GlobalTheory.ChiGrid.GridData.CosAt}
\leanref{Development/AbelianVarietyTheory/Definitions/Divisors/CmHostImaginaryGeneratorNeZero.lean}{GlobalTheory.ChiGrid.homLat}
\leanref{Development/AbelianVarietyTheory/Definitions/Divisors/CmHostImaginaryGeneratorNeZero.lean}{GlobalTheory.ChiGrid.Cos}
\leanref{Development/AbelianVarietyTheory/Definitions/Divisors/CmHostImaginaryGeneratorNeZero.lean}{GlobalTheory.ChiGrid.cosSet}
\leanref{Development/AbelianVarietyTheory/Definitions/Divisors/CmHostImaginaryGeneratorNeZero.lean}{GlobalTheory.ChiGrid.cosInd}

The spaces $V_{i,j}$ form a grid in which a step in $j$ adds rank $t$, as in
the passage from $Y_j$ to $Y_{j+1}$, and a step in $i$ adds rank $r+6$. For
the square with corner $(i,j)$, $0\leq i,j<m$, put
\[
a(i,j)=\operatorname{rk}V_{i+1,j+1}-2\lfloor\tfrac{m-j-1}2\rfloor,\quad
a'(i,j)=\operatorname{rk}V_{i,j}-2\lfloor\tfrac{m-j}2\rfloor,\quad
e_0=\tfrac{a-a'}2,\quad h_0=\tfrac{a+a'}2,
\]
and $h_1(i,j)=(\operatorname{rk}V_{i+1,j+1}+\operatorname{rk}V_{i,j})/2$. The
square carries the characters
\[
\alpha_{t}=\mu_{t,j+1}\eta_L^{\epsilon},\qquad
\beta=\chi_{m-i}\eta_L^{\epsilon},\qquad
\nu_{t,i,j}=\alpha_t/\beta=\mu^0_{j+1}\xi_t/\chi_{m-i},
\]
where $\epsilon\in\{0,1\}$ is the parity of $\operatorname{rk}V_{i,j}$.

\begin{definition}\label{th-bottom-triple}
Let $D$ be a definite-side datum with a level system, a theta grid and a grid
Weil datum with character $\psi$ (Definition~\ref{th-grid}), and let
$0\leq i,j<m$. A \emph{bottom triple} for $(i,j)$ consists of Hermitian spaces
$A$, $A'$ and $W$ over $E$ with identity Gram matrices, standard lattices
$\mathcal O_E^n$ and trivial archimedean data, of ranks
$\operatorname{rk}A=a(i,j)$, $\operatorname{rk}A'=a'(i,j)$ and
$\operatorname{rk}W=a+a'$; and, for each tuple $t$:
\begin{enumerate}
\item a Weil datum $\omega_t$ for $(-A,W)$ with characters
  $(\beta,\beta/\alpha_t)$ and a Weil datum $\omega'_t$ for $(A',W)$ with
  characters $(\alpha_t,\alpha_t/\beta)$ (Definition~\ref{th-weil-datum});
\item an idempotent $\varepsilon_t$ of $\mathbf A_E$ supported at the finite
  set $S_t$ of places of the grid attached to $t$, and for $w\in S_t$ the
  exponent $c_t(w)$ of the conductor of $\nu_{t,i,j}$ at $w$;
\item Fock vacuum vectors for $(-A,W)$ and $(A',W)$, and theta functions
  $f_t$ and $f'_t$ obtained from them by replacing the component at
  $\varepsilon_t$ by a standard function: the function on
  $\varepsilon_t\mathbb W(-A,W)$ supported on the integral matrices whose last
  $a'$ rows are divisible by $\mathfrak f_t=\prod_{w\in S_t}w^{c_t(w)}$, with
  value the inverse of $\alpha_t/\beta$ at the determinant of the top
  $a\times a$ block; and the analogous function on
  $\bar\varepsilon_t\mathbb W(A',W)$, in dual coordinates, supported on the
  matrices whose first $a$ columns are divisible by $\mathfrak f_t$, with value
  given by the last $a'\times a'$ block;
\item Haar measures and fundamental domains for $U(A)(F)$ and $U(A')(F)$.
\end{enumerate}
For a polynomial $P$ in the creation operators at the archimedean places of the
pair, write $f_t\otimes P$ for the corresponding vector, and put
\[
I_t(P,h)=\int_{\mathcal F_A}\theta_{\omega_t}(g,h;f_t\otimes P)\,dg,\qquad
I'_t(P,h)=\int_{\mathcal F_{A'}}\theta_{\omega'_t}(g,h;f'_t\otimes P)\,dg
\]
for $h\in U(W)(\mathbf A)$: the theta lifts of the constant function $1$ from
$U(A)$ and $U(A')$ to $U(W)$.
\end{definition}
\leanref{Development/AbelianVarietyTheory/BlockTheory/TwistedDiscreteSeries/UnitCharTriangle.lean}{GlobalTheory.IndefiniteRowSideFullHecke.UnitCharTriangle.Tri}
\leanref{Development/AbelianVarietyTheory/BlockTheory/TwistedDiscreteSeries/UnitCharTriangle.lean}{GlobalTheory.IndefiniteRowSideFullHecke.UnitCharTriangle.Tri.Is}
\leanref{Development/AbelianVarietyTheory/Definitions/TwistedInnerProducts/MemHalfCosetSelf.lean}{GlobalTheory.RowTest.IsBottomTri}
\leanref{Development/AbelianVarietyTheory/Definitions/Divisors/FinitePrimesUnder.lean}{GlobalTheory.FunObj.SqNum.botBig}
\leanref{Development/AbelianVarietyTheory/Definitions/Divisors/FinitePrimesUnder.lean}{GlobalTheory.FunObj.SqNum.botSmall}

\begin{definition}\label{th-large-data}
For a Hecke character $\nu$ of $E$, an \emph{$L$-function} of $\nu$ is an
entire function $\Lambda$ with
$\Lambda(s)=\prod_{w}(1-\nu(w)\mathrm Nw^{-s})^{-1}$ for $\operatorname{Re}s>1$,
the product over the finite places $w$ of $E$ at which $\nu$ is unramified. Let
$P$ be a property of a theta grid with its level system and base tuple and of
numbers $C,x,b$. It holds \emph{for all sufficiently large data} if there is
$C_F^0$ such that for all $\varepsilon$ and $C_F\geq C_F^0$ there are $C_0$
and $x_0$ with the following property: for every definite-side datum $D$
(Definition~\ref{th-definite-datum}) and size data for $D$ at a scale
$x\geq x_0$ with exponent $\varepsilon$ and constant $b$, satisfying the size conditions on the
rows, the host hypotheses, $F$ the maximal real subfield of $E$, and
$|\tfrac{L'}{L}(1)|\leq x^\varepsilon$ for the logarithmic derivative recorded
in the size data, for every level system with base tuple and every theta
grid with constants $(1,C_F,b)$ (Definition~\ref{th-grid}), and for every
$C\geq C_0$, the property $P$ holds at $(C,x,b)$.
\end{definition}
\leanref{Development/AbelianVarietyTheory/Statements/PlaceShiftStatements/NormSmall.lean}{GlobalTheory.SqTop.GBig}
\leanref{Development/AbelianVarietyTheory/Statements/PlaceShiftStatements/NormSmall.lean}{GlobalTheory.SqTop.Rec}
\leanref{Development/AbelianVarietyTheory/Definitions/Divisors/EntireEqOfHalfPlane.lean}{GlobalTheory.CMVal.IsHeckeL}

\begin{theorem}[Rallis inner product formula]\label{th-rallis-inner-product}
Assume the published local inputs of the theta correspondence
(Theorem~\ref{lit-bundle-theta}). Then for all sufficiently large data
(Definition~\ref{th-large-data}) the following holds. Let a grid Weil datum be
given, let $0\leq i,j<m$, let $(A,A',W)$ be a bottom triple for $(i,j)$
(Definition~\ref{th-bottom-triple}) and let $t$ be a tuple. Then there are a
polynomial $P$, a Haar measure $\nu$ on $U(W)(\mathbf A)$, a fundamental
domain $\mathcal F$ for $U(W)(F)$ and $c\in\mathbf C$ with $c\neq0$ such that
\[
\int_{\mathcal F}|I_t(P,h)|^2\,d\nu(h)=c\prod_{s=1-e_0(i,j)}^{h_0(i,j)}\Lambda(s)
\]
for every $L$-function $\Lambda$ of $\nu_{t,i,j}$.
\end{theorem}
\leanref{Development/AutomorphicArithmetic/ThetaCorrespondence/ThetaPairLifts/LocalTypeDictionary/BigThetaNormEqRallisInnerProduct.lean}{GlobalTheory.RallisBottom.big_theta_norm_eq_rallis_inner_product}

The product runs over the integers $s$ from $1-e_0$ to $h_0$. All of $A$,
$A'$ and $W$ are totally definite, so the integrals are over compact quotients.
The formula is the Rallis inner product formula \cite{Ral84,Ral87} for the
theta lift of the trivial representation of $U(A)$, in the doubling form of
\cite{HKS96,Yam14,GQT14}.

\begin{definition}\label{th-theta-hypotheses}
A \emph{triple} over a CM extension $E/F$ consists of Hermitian spaces $V$,
$V'$ and $W$ with $\operatorname{rk}V+\operatorname{rk}V'=\operatorname{rk}W$,
Fock-model data at the archimedean places, and Weil data $\omega$ for
$(V,W)$ and $\omega'$ for $(V',W)$ (Definition~\ref{th-weil-datum}), with theta
integrals $I(\Phi)$ and $I'(\Phi')$ over $U(V)$ and $U(V')$. The
\emph{agreement of global and local theta integrals} holds if:
\begin{enumerate}
\item the theta integrals of every triple converge, with moderate growth;
\item whenever local maps from the Schwartz functions to spaces of
  coinvariants are matched to a triple, there is a unique constant $c$ with
  $I(\Phi)=c\,I'(\Phi')$ for all pairs $(\Phi,\Phi')$ with matching images,
  $c\neq0$, some $I'(\Phi')$ is nonzero, and at the places of $F$ that split
  completely in $E$ and lie outside the finite set of places of the triple,
  the matched maps describe the coinvariants;
\item when a hyperbolic plane is added to $V'$ and $W$, or to $V$ and $W$, the
  constant changes by an explicit nonzero factor, built from a partial Hecke
  $L$-value of the character of the triple, its values on ideals, and local
  epsilon factors.
\end{enumerate}
The \emph{local matching} holds if every triple has matched local maps.
\end{definition}
\leanref{Development/AbelianVarietyTheory/BlockTheory/DivisorBase/WeilTensorModels.lean}{GlobalTheory.IndefiniteRowSideFullHecke.WeilTensorModels.ThetaIntegralsAgree}
\leanref{Development/AbelianVarietyTheory/BlockTheory/DivisorBase/WeilTensorModels.lean}{GlobalTheory.IndefiniteRowSideFullHecke.WeilTensorModels.LocalMatching}
\leanref{Development/AutomorphicArithmetic/ThetaCorrespondence/ThetaPairLifts/ThirdPlaceSeries/EveryCharacterOfLiteratureOfVWin.lean}{UpperTheory.LiteratureTheorems.theta_integrals_agree_and_local_matching_of_literature}
\leanref{Development/AbelianVarietyTheory/Proofs/DivisorProofs/SquareBottomHalfRankDiffLeHalfRank.lean}{GlobalTheory.SqTop.lNe_of_functionalEquation}

In Alpöge's proof both hypotheses are derived from published results on
Siegel Eisenstein series and the Siegel-Weil formula
(Chapter~\ref{chapter-literature}). The nonvanishing of $L$-values in the next
theorem is derived from the functional equation of Hecke $L$-functions and the
boundedness of completed Hecke $L$-functions in vertical strips.

\subsection{Matchings of square triples}\label{subsection-th-square-matchings}

\begin{definition}\label{th-first-lift-matching}
Let $D$ be a definite-side datum (Definition~\ref{th-definite-datum}) with CM
field $E$, totally real subfield $F$ and set $\mathcal T$ of complex
embeddings, let $\psi$ be a standard character
(Definition~\ref{th-weil-datum}), and fix an archimedean family $\theta$
(Definition~\ref{th-grid-fock}(1)). For a Hermitian space $S$
(Definition~\ref{th-spaces}) let $U(S)(\mathbf A_f)$ be the subgroup of
elements of $U(S)(\mathbf A)$ with archimedean component $1$; for
$\tau\in\mathcal T$ let $\mathfrak u_\tau(S)$ be the space of complex matrices
$Y$ with $\bar Y^{\mathsf T}\tau(G_S)+\tau(G_S)Y=0$, and for
$Y\in\mathfrak u_\tau(S)$ and $s\in\mathbf R$ let $\exp_\tau(sY)$ be the element
of $U(S)(\mathbf A)$ with component $\exp(sY)$ at the place of $\tau$ and $1$
elsewhere. \emph{Archimedean data} on $S$ consist of, at each $\tau$, a
diagonal matrix $J_\tau$ of signs and an invertible matrix $B_\tau$ with
$\bar B_\tau^{\mathsf T}J_\tau B_\tau=\tau(G_S)$. They give the Cartan
involution $s_\tau=B_\tau^{-1}J_\tau B_\tau$ of $S$ at $\tau$, with orthonormal
basis $B_\tau^{-1}$ (Definition~\ref{th-grid-fock}(2)); conversely a Cartan
involution $s$ with an orthonormal basis $B$, $sB=BJ$, is given in this way by
the data $(J,B^{-1})$. The space $-S$ carries the same data with $J_\tau$
replaced by $-J_\tau$.
\begin{enumerate}
\item Let $S$ and $T$ be Hermitian spaces with archimedean data. An
  \emph{archimedean vacuum} for $(S,T)$ is a theta function for $(S,T)$
  (Definition~\ref{th-weil-datum}) that is a Fock vacuum vector at every
  $\tau\in\mathcal T$ for the Cartan involutions of the data
  (Definition~\ref{th-grid-fock}(3)). Let $\mathcal P(S,T)$ be the space of polynomials
  in the creation operators of the pair, and $\mathcal A(S,T)$ the span of the
  vectors $f\otimes P$ for archimedean vacua $f$ and $P\in\mathcal P(S,T)$.
\item A Hecke character $\chi$ of $E$ has \emph{archimedean exponent}
  $x\in\frac12\mathbf Z$ at $\tau\in\mathcal T$ if $\chi(y)=(z/|z|)^{2x}$ for
  every idele $y$ whose components away from the place of $\tau$ are $1$, where
  $z\in\mathbf C^\times$ is the component of $y$ at that place read through
  $\tau$. It has archimedean exponents $a\colon\mathcal T\to\mathbf Q$ if it
  has archimedean exponent $a(\tau)$ at every $\tau$.
\item Let $A$, $A'$ and $W$ be Hermitian spaces, let $\alpha$ and $\beta$ be
  Hecke characters of $E$, trivial on $E^\times$, and let $\eta$ be the
  quadratic character of $E/F$. Then $(A,A',W;\alpha,\beta)$ is a
  \emph{square triple} if:
  \begin{enumerate}
  \item[(a)] $\alpha$ and $\beta$ restrict to $\eta^{\operatorname{rk}A'}$ and
    $\eta^{\operatorname{rk}A}$ on $\mathbf A_F^\times$;
  \item[(b)] $\operatorname{rk}W=\operatorname{rk}A+\operatorname{rk}A'$, and
    this rank is even;
  \item[(c)] there are $s,s'\geq0$ such that at every $\tau$ the signatures of
    $A$, $A'$ and $W$ are $(\operatorname{rk}A-s,s)$,
    $(\operatorname{rk}A'-s',s')$ and $(\operatorname{rk}W-s-s',s+s')$, and
    $\operatorname{rk}A'+2(2s+2s'+1)<\operatorname{rk}A$ and
    $2s+2s'+1<\operatorname{rk}A'$;
  \item[(d)] at every $\tau$, $\alpha$ and $\beta$ have archimedean exponents
    $x$ and $x'$ with $2(x-x')<-\operatorname{rk}W$;
  \item[(e)] for every finite place $v$ of $F$, either $\alpha$ and $\beta$ are
    unramified at the places over $v$, or $v$ splits in $E$, $\alpha/\beta$ is
    ramified at a place over $v$, and $v$ does not divide the different of $F$.
  \end{enumerate}
\item Let $A$, $A'$ and $W$ carry archimedean data, and let $\alpha,\beta$ be
  Hecke characters of $E$. A \emph{matching} of $(A,A',W;\alpha,\beta)$ is a
  tuple $(\omega,\omega',\Pi,\varrho,p,p')$ of: a Weil datum $\omega$ for
  $(-A,W)$ with characters $(\beta,\beta/\alpha)$ and a Weil datum $\omega'$
  for $(A',W)$ with characters $(\alpha,\alpha/\beta)$
  (Definition~\ref{th-weil-datum}); a complex vector space $\Pi\neq0$ with
  linear operators $\varrho(h)$ for $h\in U(W)(\mathbf A_f)$ and
  $\varrho_\tau(Y)$ for $\tau\in\mathcal T$ and $Y\in\mathfrak u_\tau(W)$, such
  that $0$ and $\Pi$ are the only subspaces stable under all of them; and
  linear maps $p\colon\mathcal A(-A,W)\to\Pi$ and
  $p'\colon\mathcal A(A',W)\to\Pi$. These are subject to the following
  conditions, for $(\omega_0,\mathcal A,p_0,S)$ equal to
  $(\omega,\mathcal A(-A,W),p,A)$ and to $(\omega',\mathcal A(A',W),p',A')$:
  \begin{enumerate}
  \item[(a)] $\mathcal A$ is stable under $\omega_0(g,1)$ and $\omega_0(1,h)$
    for $g\in U(S)(\mathbf A_f)$ and $h\in U(W)(\mathbf A_f)$; and for
    $\Phi\in\mathcal A$, $\tau\in\mathcal T$ and $Y$ in $\mathfrak u_\tau(S)$,
    respectively $\mathfrak u_\tau(W)$, the derivative at $s=0$ of
    $\omega_0(\exp_\tau(sY),1)\Phi$, respectively of
    $\omega_0(1,\exp_\tau(sY))\Phi$, exists at every point and lies in
    $\mathcal A$;
  \item[(b)] $p_0(\mathcal A)=\Pi$;
  \item[(c)] $p_0(\omega_0(g,1)\Phi)=p_0(\Phi)$ for $g\in U(S)(\mathbf A_f)$,
    and $p_0$ vanishes on the derivatives along $\mathfrak u_\tau(S)$ in (a);
  \item[(d)] $p_0(\omega_0(1,h)\Phi)=\varrho(h)p_0(\Phi)$ for
    $h\in U(W)(\mathbf A_f)$, and $p_0$ maps the derivative of $\Phi$ along
    $Y\in\mathfrak u_\tau(W)$ in (a) to $\varrho_\tau(Y)p_0(\Phi)$.
  \end{enumerate}
\item Let $M$ and $W$ be Hermitian spaces with archimedean data and $\omega_0$
  a Weil datum for $(M,W)$. Let $\mathcal R(\omega_0)$ be the span of the
  functions $\omega_0(u,1)\Psi-\Psi$ with $u\in U(M)(\mathbf A_f)$ and
  $\Psi\in\mathcal A(M,W)$. For a Haar measure $\nu$ on $U(M)(\mathbf A)$ and a
  fundamental domain $\mathcal F$ for $U(M)(F)$, the \emph{theta integral} of a
  function $\Phi$ on $\mathbb W(M,W)$ is the function
  $I_{\omega_0}(\Phi)(h)=\int_{\mathcal F}\theta_{\omega_0}(g,h;\Phi)\,d\nu(g)$
  on $U(W)(\mathbf A)$, the theta lift of the constant function $1$.
\end{enumerate}
\end{definition}
\leanref{Development/AbelianVarietyTheory/Proofs/TwistedOperators/BlkGLMem.lean}{GlobalTheory.SquareComm.IsFinF}
\leanref{Development/AbelianVarietyTheory/Proofs/TwistedOperators/BlkGLMem.lean}{GlobalTheory.SquareComm.algF}
\leanref{Formalism/LocalToGlobal/BesselHermitianTheta/DefsReal.lean}{Formalism.L2.Theta.HasInfExp}
\leanref{Development/AbelianVarietyTheory/Proofs/TwistedOperators/BlkGLMem.lean}{GlobalTheory.SquareComm.IsTriple}
\leanref{Development/AbelianVarietyTheory/Proofs/TwistedOperators/BlkGLMem.lean}{GlobalTheory.SquareComm.IsMatch}
\leanref{Development/AbelianVarietyTheory/BlockTheory/TwistedDiscreteSeries/MatchingAdjoints.lean}{GlobalTheory.IndefiniteRowSideFullHecke.MatchingAdjoints.IsIrred}
\leanref{Development/AbelianVarietyTheory/BlockTheory/TwistedDiscreteSeries/MatchingAdjoints.lean}{GlobalTheory.IndefiniteRowSideFullHecke.MatchingAdjoints.IsCoinvMap}
\leanref{Development/AbelianVarietyTheory/Proofs/TwistedOperators/BlkGLMem.lean}{GlobalTheory.SquareComm.finRel}
\leanref{Development/AbelianVarietyTheory/Proofs/TwistedOperators/BlkGLMem.lean}{GlobalTheory.SquareComm.thetaInt}

\begin{definition}\label{th-square-fock-operators}
Let $D$, $\psi$ and the archimedean embeddings be as in
Definition~\ref{th-first-lift-matching}, and let $S$ and $T$ be Hermitian
spaces with archimedean data $(J_\tau,B_\tau)$ and $(J'_\tau,B'_\tau)$. The
space $\mathcal P(S,T)$ is the polynomial ring over $\mathbf C$ in variables
$x_{\tau,i,j}$ for $\tau\in\mathcal T$, $1\leq i\leq\operatorname{rk}T$ and
$1\leq j\leq\operatorname{rk}S$. Let $\epsilon_{S,\tau}(j)$ and
$\epsilon_{T,\tau}(i)$ be the diagonal entries of $J_\tau$ and $J'_\tau$.
\begin{enumerate}
\item The \emph{raising} and \emph{lowering operators} on $\mathcal P(S,T)$ are
  $x^+_{\tau,i,j}=\epsilon_{S,\tau}(j)\,x_{\tau,i,j}$ and
  $x^-_{\tau,i,j}=\epsilon_{S,\tau}(j)\,\partial/\partial x_{\tau,i,j}$ if
  $\epsilon_{T,\tau}(i)=\epsilon_{S,\tau}(j)$, and
  $x^+_{\tau,i,j}=-\epsilon_{S,\tau}(j)\,\partial/\partial x_{\tau,i,j}$ and
  $x^-_{\tau,i,j}=\epsilon_{S,\tau}(j)\,x_{\tau,i,j}$ otherwise, where
  $x_{\tau,i,j}$ acts by multiplication.
\item Let $a\colon\mathcal T\to\mathbf Q$ and $\tau\in\mathcal T$. For
  $Y\in\mathfrak u_\tau(T)$ put $Y'=B'_\tau Y{B'_\tau}^{-1}$ and
  \[
  \partial^a_{T,\tau}(Y)=\sum_{i,l}Y'_{il}\sum_j\bigl(x^+_{\tau,i,j}x^-_{\tau,l,j}+\tfrac12\delta_{il}\bigr)+a(\tau)\operatorname{tr}(Y).
  \]
  For $Z\in\mathfrak u_\tau(S)$ put $Z'=\overline{B_\tau ZB_\tau^{-1}}$, the
  entrywise complex conjugate, and
  \[
  \partial^a_{S,\tau}(Z)=\sum_{j,k}Z'_{jk}\sum_i\bigl(x^+_{\tau,i,j}x^-_{\tau,i,k}+\tfrac12\delta_{jk}\bigr)-a(\tau)\operatorname{tr}(Z).
  \]
  The \emph{coinvariant space} $\mathcal C_a(S,T)$ is the span of the
  polynomials $\partial^a_{S,\tau}(Z)P$ with $\tau\in\mathcal T$,
  $Z\in\mathfrak u_\tau(S)$ and $P\in\mathcal P(S,T)$.
\item Let $Q$ be a complex vector space, $\rho=(\rho_i)_{i\in I}$ a family of
  linear operators on $Q$ and $C\subset Q$ a subspace. A subspace of $Q$ is
  \emph{admissible} for $(\rho,C)$ if it is stable under every $\rho_i$ and
  contains $C$. The pair $(\rho,C)$ satisfies the \emph{Howe condition} if
  there is $L$ such that every strictly increasing chain of admissible
  subspaces has at most $L$ members, and the sum of two admissible subspaces
  different from $Q$ is different from $Q$. The \emph{radical}
  $\operatorname{rad}(\rho,C)$ is the sum of all admissible subspaces
  different from $Q$.
\end{enumerate}
\end{definition}
\leanref{Development/AbelianVarietyTheory/BlockTheory/TwistedDiscreteSeries/RaisingLoweringOperators.lean}{GlobalTheory.IndefiniteRowSideFullHecke.RaisingLoweringOperators.raiO}
\leanref{Development/AbelianVarietyTheory/BlockTheory/TwistedDiscreteSeries/RaisingLoweringOperators.lean}{GlobalTheory.IndefiniteRowSideFullHecke.RaisingLoweringOperators.lowO}
\leanref{Development/AbelianVarietyTheory/BlockTheory/TwistedDiscreteSeries/RaisingLoweringOperators.lean}{GlobalTheory.IndefiniteRowSideFullHecke.RaisingLoweringOperators.fockT}
\leanref{Development/AbelianVarietyTheory/BlockTheory/TwistedDiscreteSeries/RaisingLoweringOperators.lean}{GlobalTheory.IndefiniteRowSideFullHecke.RaisingLoweringOperators.fockS}
\leanref{Development/AbelianVarietyTheory/Proofs/TwistedOperators/HasInfExpDiv.lean}{GlobalTheory.RealQuot.actT}
\leanref{Development/AbelianVarietyTheory/Proofs/TwistedOperators/HasInfExpDiv.lean}{GlobalTheory.RealQuot.actS}
\leanref{Development/AbelianVarietyTheory/Proofs/TwistedOperators/HasInfExpDiv.lean}{GlobalTheory.RealQuot.coinv}
\leanref{Development/AbelianVarietyTheory/Proofs/TwistedOperators/HasInfExpDiv.lean}{GlobalTheory.RealQuot.HoweCond}
\leanref{Development/AbelianVarietyTheory/Proofs/TwistedOperators/HasInfExpDiv.lean}{GlobalTheory.RealQuot.rad}

\begin{lemma}\label{th-square-fock-derivatives}
Let $D$, $\psi$ and the archimedean embeddings be as in
Definition~\ref{th-first-lift-matching}, let $S$ and $T$ be Hermitian spaces
with archimedean data, and let $\omega$ be a Weil datum for $(S,T)$
(Definition~\ref{th-weil-datum}) with characters $(\chi_S,\chi_T)$. Let $f$ be
an archimedean vacuum for $(S,T)$, let $P\in\mathcal P(S,T)$ and
$w\in\mathbb W$.
\begin{enumerate}
\item Let $\tau\in\mathcal T$ and $a_S,a_T\colon\mathcal T\to\mathbf Q$, and
  assume that $\chi_S$ and $\chi_T$ have archimedean exponents $a_S(\tau)$ and
  $a_T(\tau)$ at $\tau$. For $Z\in\mathfrak u_\tau(S)$ the function
  $s\mapsto(\omega(\exp_\tau(sZ),1)(f\otimes P))(w)$ on $\mathbf R$ is
  differentiable at $0$ with derivative
  $(f\otimes\partial^{a_T}_{S,\tau}(Z)P)(w)$, and for $Y\in\mathfrak u_\tau(T)$
  the function $s\mapsto(\omega(1,\exp_\tau(sY))(f\otimes P))(w)$ is
  differentiable at $0$ with derivative
  $(f\otimes\partial^{a_S}_{T,\tau}(Y)P)(w)$
  (Definition~\ref{th-square-fock-operators}).
\item For $h\in U(T)(\mathbf A_f)$ the function $\omega(1,h)f$ is an
  archimedean vacuum for $(S,T)$, and
  $\omega(1,h)(f\otimes P)=(\omega(1,h)f)\otimes P$.
\end{enumerate}
\end{lemma}
\leanref{Development/AbelianVarietyTheory/Proofs/TwistedOperators/LieActionFockOperatorsScalarsTrace.lean}{GlobalTheory.LieScalar.lieFockAll}
\leanref{Development/AbelianVarietyTheory/Definitions/TwistedInnerProducts/IsVacFOfIsVacVec.lean}{GlobalTheory.ChainFr.LieFockWith}
\leanref{Development/AbelianVarietyTheory/Proofs/AdelicLatticeProofs/EqZeroOfLintegralEqZero.lean}{GlobalTheory.OrientedHeis.om_dRho}
\leanref{Development/AbelianVarietyTheory/Proofs/TwistedOperators/CompAtMtOfMemGfC.lean}{GlobalTheory.OrientedHeis.om_fin_pureTensorO}
\begin{proof}
Part (1) is a computation in the Fock model, and its result is stated. For
(2), let $w_0\in\mathbb W$ have finite component $0$. By the conditions on a
Weil datum (Definition~\ref{th-weil-datum}(3)), $\omega(1,h)$ preserves theta
functions and $\omega(1,h)\rho(sw_0)=\rho(shw_0)\omega(1,h)$ on them, and
$hw_0=w_0$ since $h$ has archimedean component $1$. The difference quotients
$s^{-1}(\rho(sw_0)f-f)$ are theta functions that converge to the derivative of
$f$ along $w_0$, uniformly on $\mathbb W$. Since $\omega(1,h)$ preserves the
$L^2$-norm of theta functions on a fundamental domain of measure $1$ for the
rational points of $\mathbb W$, it carries the derivative of $f$ along $w_0$
to the derivative of $\omega(1,h)f$ along $w_0$; the two continuous theta
functions agree almost everywhere, hence everywhere. The annihilation and
creation operators at the archimedean places are linear combinations of these
derivatives. So $\omega(1,h)f$ is annihilated by the annihilation operators,
and $\omega(1,h)(f\otimes P)=(\omega(1,h)f)\otimes P$, since $f\otimes P$ is
obtained from $f$ by applying the polynomial $P$ in the creation operators.
\end{proof}

\begin{lemma}\label{th-square-fock-maps}
Let $D$, $\psi$ and the archimedean embeddings be as in
Definition~\ref{th-first-lift-matching}, let $S$ and $T$ be Hermitian spaces
with archimedean data, and let $\omega$ be a Weil datum for $(S,T)$ whose
characters have archimedean exponents $a_S$ and $a_T$. Let $\Pi$ be a complex
vector space with linear operators $\varrho(h)$ for $h\in U(T)(\mathbf A_f)$
and $\varrho_\tau(Y)$ for $\tau\in\mathcal T$ and $Y\in\mathfrak u_\tau(T)$,
and let $p_0\colon\mathcal A(S,T)\to\Pi$ be a linear map satisfying conditions
(a), (c) and (d) of a matching (Definition~\ref{th-first-lift-matching}(4)),
read with $\omega$, $\mathcal A(S,T)$, $p_0$, $U(S)$ and $U(T)$ in place of
$\omega_0$, $\mathcal A$, $p_0$, $U(S)$ and $U(W)$. Let $f$ be an archimedean
vacuum for $(S,T)$. Then:
\begin{enumerate}
\item $p_0(f\otimes\partial^{a_S}_{T,\tau}(Y)P)=\varrho_\tau(Y)\,p_0(f\otimes P)$
  for $\tau\in\mathcal T$, $Y\in\mathfrak u_\tau(T)$ and
  $P\in\mathcal P(S,T)$;
\item $p_0(f\otimes P)=0$ for $P\in\mathcal C_{a_T}(S,T)$;
\item for $h\in U(T)(\mathbf A_f)$ the function $f'=\omega(1,h)f$ is an
  archimedean vacuum for $(S,T)$, and
  $\varrho(h)\,p_0(f\otimes P)=p_0(f'\otimes P)$ for every $P$.
\end{enumerate}
\end{lemma}
\leanref{Development/AbelianVarietyTheory/Proofs/TwistedOperators/OrientedPairFockMapsLieEquivariant.lean}{GlobalTheory.RealQuot.oriented_maps}
\leanref{Development/AbelianVarietyTheory/Proofs/TwistedOperators/OrientedPairFockMapsLieEquivariantStatement.lean}{GlobalTheory.T.orientedPair_fockMaps_lie_equivariant}
\begin{proof}
The vectors $f\otimes P$ lie in $\mathcal A(S,T)$. For (1), by
Lemma~\ref{th-square-fock-derivatives}(1) the derivative of
$\omega(1,\exp_\tau(sY))(f\otimes P)$ at $s=0$ is
$f\otimes\partial^{a_S}_{T,\tau}(Y)P$, which lies in $\mathcal A(S,T)$. By
condition (d) of a matching, $p_0$ maps it to
$\varrho_\tau(Y)p_0(f\otimes P)$.

For (2), let $\tau\in\mathcal T$ and $Z\in\mathfrak u_\tau(S)$. By
Lemma~\ref{th-square-fock-derivatives}(1) the derivative of
$\omega(\exp_\tau(sZ),1)(f\otimes P)$ at $s=0$ is
$f\otimes\partial^{a_T}_{S,\tau}(Z)P$, and $p_0$ vanishes on it by condition
(c). The map $P\mapsto p_0(f\otimes P)$ is linear, and
$\mathcal C_{a_T}(S,T)$ is spanned by the polynomials
$\partial^{a_T}_{S,\tau}(Z)P$, so it vanishes on $\mathcal C_{a_T}(S,T)$.

For (3), $f'$ is an archimedean vacuum and
$\omega(1,h)(f\otimes P)=f'\otimes P$ by
Lemma~\ref{th-square-fock-derivatives}(2). By condition (d),
$\varrho(h)p_0(f\otimes P)=p_0(\omega(1,h)(f\otimes P))=p_0(f'\otimes P)$.
\end{proof}

\begin{lemma}\label{th-square-howe-model}
Assume the published local inputs of the theta correspondence
(Theorem~\ref{lit-bundle-theta}). Let $D$, $\psi$ and the archimedean
embeddings be as in Definition~\ref{th-first-lift-matching}, let $S$ and $T$
be Hermitian spaces with archimedean data, and let $\omega$ be a Weil datum for
$(S,T)$ whose characters have archimedean exponents $a_S$ and $a_T$. Then there
are data as in Theorem~\ref{lit-howe-duality-real} for $E/F$: a Hermitian
space $X$ of dimension $\operatorname{rk}S$, a skew-Hermitian space $W_0$ of
dimension $\operatorname{rk}T$, characters, a standard character, a Weil
representation $\omega_0$ and Cartan involutions, with Fock space
$\mathcal P_0$; and an injective linear map
$\Phi\colon\mathcal P(S,T)\to\mathcal P_0$ with image $\mathcal P_0$, such that:
\begin{enumerate}
\item if a subspace $N\subset\mathcal P(S,T)$ is stable under every
  $\partial^{a_S}_{T,\tau}(Y)$ (Definition~\ref{th-square-fock-operators}),
  then $\Phi(N)$ is stable under $d\omega_0(\mathfrak u(W_0))$ and
  $\omega_0(K_{W_0})$;
\item $d\omega_0(Y)f$ and $\omega_0(k)f-f$ lie in
  $\Phi(\mathcal C_{a_T}(S,T))$ for $f\in\mathcal P_0$, $Y\in\mathfrak u(X)$
  and $k\in K_X$.
\end{enumerate}
\end{lemma}
\leanref{Development/AutomorphicArithmetic/ThetaCorrespondence/ThetaPairLifts/LocalTypeDictionary/TwistRealPlaceCorrespondence.lean}{GlobalTheory.T.twist_real_place_correspondence}
\leanref{Development/AutomorphicArithmetic/ThetaCorrespondence/ThetaPairLifts/LocalTypeDictionary/TwistRealPlaceCorrespondence.lean}{GlobalTheory.RealDict.real_place_correspondence}

The lemma is a computation comparing two models of the archimedean Weil
representation. The space $X$ is $S$ with the negated form, the space $W_0$ is
$\xi_0T$ with its Gram matrix transposed, and the Cartan involutions come from
the archimedean data.
For a fixed nonzero archimedean vacuum $f_0$ and a fixed function $r$ at the
finite places, $\Phi(P)$ is the archimedean component of the Schwartz-Bruhat
function whose theta series, with finite component $r$, is $f_0\otimes P$.

\begin{lemma}\label{th-square-howe-condition}
Assume the published local inputs of the theta correspondence
(Theorem~\ref{lit-bundle-theta}), which include Howe duality at the real places
(Theorem~\ref{lit-howe-duality-real}). Let $D$, $\psi$ and the archimedean
embeddings be as in Definition~\ref{th-first-lift-matching}, let $S$ and $T$ be
Hermitian spaces with archimedean data, and let $\omega$ be a Weil datum for
$(S,T)$ whose characters have archimedean exponents $a_S$ and $a_T$. Then the
family of operators $\partial^{a_S}_{T,\tau}(Y)$ on $\mathcal P(S,T)$, indexed
by the pairs $(\tau,Y)$ with $\tau\in\mathcal T$ and
$Y\in\mathfrak u_\tau(T)$, and the subspace $\mathcal C_{a_T}(S,T)$ satisfy
the Howe condition (Definition~\ref{th-square-fock-operators}).
\end{lemma}
\leanref{Development/AutomorphicArithmetic/ThetaCorrespondence/ThetaPairLifts/LocalTypeDictionary/HoweConditionActionCoinvariantsFock.lean}{GlobalTheory.HoweFockReal.howe_condition_action_coinvariants_fock}
\leanref{Development/AutomorphicArithmetic/ThetaCorrespondence/ThetaPairLifts/LocalTypeDictionary/HoweConditionActionCoinvariantsFock.lean}{GlobalTheory.T.howe_condition_of_fock_coinvariants}
\begin{proof}
Let $X$, $W_0$, $\omega_0$, $\mathcal P_0$ and $\Phi$ be given by
Lemma~\ref{th-square-howe-model}, and write $C=\mathcal C_{a_T}(S,T)$. Let $N$
be admissible: stable under every $\partial^{a_S}_{T,\tau}(Y)$ and containing
$C$. Then $\Phi(N)$ is a theta subspace of $\mathcal P_0$ in the sense of
Theorem~\ref{lit-howe-duality-real}. Indeed $\Phi(N)\subset\mathcal P_0$; it
is stable under $d\omega_0(\mathfrak u(W_0))$ and $\omega_0(K_{W_0})$ by
Lemma~\ref{th-square-howe-model}(1); and it contains $d\omega_0(Y)f$ and
$\omega_0(k)f-f$ for $f\in\mathcal P_0$, $Y\in\mathfrak u(X)$ and $k\in K_X$,
since these lie in $\Phi(C)\subset\Phi(N)$ by
Lemma~\ref{th-square-howe-model}(2).

Let $L$ be as in Theorem~\ref{lit-howe-duality-real}(1). Since $\Phi$ is
injective, a strictly increasing chain of admissible subspaces maps to a
strictly increasing chain of theta subspaces, so it has at most $L$ members.
Let $N$ and $N'$ be admissible and different from $\mathcal P(S,T)$. Since
$\Phi$ is injective with image $\mathcal P_0$, the theta subspaces $\Phi(N)$
and $\Phi(N')$ are different from $\mathcal P_0$, so
$\Phi(N+N')=\Phi(N)+\Phi(N')\neq\mathcal P_0$ by
Theorem~\ref{lit-howe-duality-real}(2). Hence
$N+N'\neq\mathcal P(S,T)$.
\end{proof}

\begin{lemma}\label{th-square-radical}
Let $Q$ and $\Pi$ be complex vector spaces, let $\rho=(\rho_i)_{i\in I}$ be
linear operators on $Q$ and $\sigma=(\sigma_i)_{i\in I}$ linear operators on
$\Pi$, and let $C\subset Q$ be a subspace such that $(\rho,C)$ satisfies the
Howe condition (Definition~\ref{th-square-fock-operators}). Put
$R=\operatorname{rad}(\rho,C)$. A \emph{$\sigma$-subspace} of $\Pi$ is a
subspace stable under every $\sigma_i$, and an \emph{intertwiner} from
$(Q,\rho,C)$ is a linear map $\beta\colon Q\to\Pi$ with
$\beta\rho_i=\sigma_i\beta$ for every $i$ and $C\subset\ker\beta$.
\begin{enumerate}
\item If some admissible subspace is different from $Q$, then $R\neq Q$, $R$
  is admissible, and every admissible subspace containing $R$ is $R$ or $Q$.
\item Let $(\beta_f)_{f\in\Sigma}$ be intertwiners from $(Q,\rho,C)$ whose
  images span $\Pi$, and let $(\kappa_j)_{j\in J}$ be linear operators on $\Pi$
  such that for all $j$ and $f$ there is $f'$ with
  $\kappa_j\beta_f=\beta_{f'}$. If $0$ and $\Pi$ are the only subspaces of
  $\Pi$ stable under every $\sigma_i$ and every $\kappa_j$, then every
  $\sigma$-subspace of $\Pi$ has a complement that is a $\sigma$-subspace.
\item Assume that $R\neq Q$ and that every $\sigma$-subspace of $\Pi$ has a
  complement that is a $\sigma$-subspace. Then $R\subset\ker\beta$ for every
  intertwiner $\beta$ from $(Q,\rho,C)$.
\item Let $Q'$ be a complex vector space with linear operators
  $\rho'=(\rho'_i)_{i\in I}$ and a subspace $C'$ such that $(\rho',C')$
  satisfies the Howe condition, and put $R'=\operatorname{rad}(\rho',C')$.
  Assume that some admissible subspace for $(\rho,C)$ is different from $Q$,
  that some admissible subspace for $(\rho',C')$ is different from $Q'$, and
  that every $\sigma$-subspace of $\Pi$ has a complement that is a
  $\sigma$-subspace. Let $(\beta_f)$ be intertwiners from $(Q,\rho,C)$ whose
  images span $\Pi$, and let $\beta_1\neq0$ be an intertwiner from
  $(Q',\rho',C')$. Then there is a linear isomorphism $e\colon Q/R\to Q'/R'$
  with $e\bar\rho_i=\bar\rho'_ie$ for every $i$, where $\bar\rho_i$ and
  $\bar\rho'_i$ are the operators induced by $\rho_i$ and $\rho'_i$.
\end{enumerate}
\end{lemma}
\leanref{Development/AbelianVarietyTheory/Proofs/TwistedOperators/HasInfExpDiv.lean}{GlobalTheory.RealQuot.rad_ne_top}
\leanref{Development/AbelianVarietyTheory/Proofs/TwistedOperators/HasInfExpDiv.lean}{GlobalTheory.RealQuot.rad_le_ker}
\leanref{Development/AbelianVarietyTheory/Proofs/TwistedOperators/ExistsStableComplementOfIrreducibleSpanIntertwin.lean}{GlobalTheory.RealQuot.SemisimpleSpan.exists_stab_compl}
\leanref{Development/AbelianVarietyTheory/Proofs/TwistedOperators/ExistsStableComplementOfIrreducibleSpanIntertwinStatement.lean}{GlobalTheory.T.exists_stable_complement_of_irreducible_span_intertwiners}
\leanref{Development/AbelianVarietyTheory/Proofs/TwistedOperators/SimpleQuotientsIsoOfIntertwinedSemisimple.lean}{GlobalTheory.RealQuot.simpleQuot_iso_of_intertwined}
\leanref{Development/AbelianVarietyTheory/Proofs/TwistedOperators/SimpleQuotientsIsoOfIntertwinedSemisimpleStatement.lean}{GlobalTheory.T.simple_quotients_iso_of_intertwined_semisimple}
\begin{proof}
For (1), let $\mathcal S$ be the set of admissible subspaces different from
$Q$, which is not empty. It has a maximal element $M$, since otherwise it would
contain a strictly increasing chain with more than $L$ members. Let
$N\in\mathcal S$. The sum $M+N$ is admissible, and it is different from $Q$ by
the second half of the Howe condition, so $M+N\in\mathcal S$ and $M+N=M$ by
maximality. Thus every $N\in\mathcal S$ lies in $M$, and $R=M$. So
$R\neq Q$, $R$ is admissible, and an admissible subspace containing $R$ and
different from $Q$ lies in $R$.

For (2), let $K=\bigcap_f\ker\beta_f$. It is stable under every $\rho_i$,
since each $\beta_f$ is an intertwiner, and it contains $C$; so $K$ is
admissible. For a subspace $U\subset Q$ stable under every $\rho_i$, the
subspace $\sum_f\beta_f(U)$ of $\Pi$ is stable under every $\sigma_i$, since
$\beta_f\rho_i=\sigma_i\beta_f$, and under every $\kappa_j$, since
$\kappa_j\beta_f=\beta_{f'}$; so it is $0$ or $\Pi$. If $K=Q$, every $\beta_f$
is $0$, so $\Pi=0$ and there is nothing to prove. Otherwise $Q$ is an
admissible subspace strictly containing $K$. Choose such a subspace $M$ that
is minimal; it exists, since a strictly decreasing chain of admissible
subspaces has at most $L$ members. Then $\sum_f\beta_f(M)$ is not $0$, since
$M\not\subset K$, so it is $\Pi$. Let $V$ be a $\sigma$-subspace contained in
$\beta_f(M)$, and put $U=M\cap\beta_f^{-1}(V)$. Then $U$ is admissible and
contains $K$. If $U\neq K$, then $U=M$ by minimality, and $V=\beta_f(M)$. If
$U=K$, then every element $\beta_f(x)$ of $V$ with $x\in M$ has $x\in K$, so
$V=0$. Thus $\Pi$ is a sum of $\sigma$-subspaces $\beta_f(M)$ whose only
$\sigma$-subspaces are $0$ and themselves. As a module over the free algebra
on the symbols $\sigma_i$, $\Pi$ is a sum of simple submodules, so it is
semisimple and every submodule has a complement \cite[\S2]{Lam01}.

For (3), suppose that $R\not\subset\ker\beta$. The image $\beta(R)$ is a
$\sigma$-subspace, since $R$, a sum of subspaces stable under every $\rho_i$,
is stable under every $\rho_i$; let $N'$ be a complement of it that is a
$\sigma$-subspace, and put $U=\beta^{-1}(N')$. Then $U$ is admissible. Every
$x\in Q$ has $\beta(x)=\beta(r)+n$ with $r\in R$ and $n\in N'$, so $x-r\in U$,
and $R+U=Q$. If $U=Q$, then $\beta(R)\subset\beta(R)\cap N'=0$, which is
false. Otherwise $U\subset R$ by the definition of $R$, and $Q=R+U=R$, which
is also false.

For (4), apply (1) and (3) to $(Q,\rho,C)$ and to $(Q',\rho',C')$. By (1),
$R\neq Q$ and $R'\neq Q'$, every admissible subspace containing $R$ is $R$ or
$Q$, and every admissible subspace containing $R'$ is $R'$ or $Q'$. By (3),
$R\subset\ker\beta_f$ for every $f$ and $R'\subset\ker\beta_1$. Put
$U=\beta_1(Q')$, a $\sigma$-subspace. Every nonzero $\sigma$-subspace $V$ of
$U$ is $U$: the subspace $\beta_1^{-1}(V)$ is admissible and contains $R'$, and
it is not $R'$, since $V\neq0$ and $\beta_1(R')=0$; so it is $Q'$. Let $N'$ be
a complement of $U$ that is a $\sigma$-subspace, and let $\pi$ be the
projection onto $U$ along $N'$, which commutes with every $\sigma_i$. Some
$\pi\beta_f$ is not $0$: otherwise every $\beta_f(Q)$ lies in $N'$, so
$N'=\Pi$, $U=0$ and $\beta_1=0$. Put $\gamma=\pi\beta_f$. It is an
intertwiner, and $\ker\gamma$ is admissible, contains $R$ and is not $Q$, so
$\ker\gamma=R$. In the same way $\ker\beta_1=R'$. The image of $\gamma$ is a
nonzero $\sigma$-subspace of $U$, so it is $U$. Thus $\gamma$ and $\beta_1$
induce isomorphisms $Q/R\to U$ and $Q'/R'\to U$ that carry $\bar\rho_i$ and
$\bar\rho'_i$ to $\sigma_i$, and $e$ is the first followed by the inverse of
the second.
\end{proof}

\begin{definition}\label{th-darboux-model}
Let $E/F$ be a CM extension with complex conjugation $x\mapsto\bar x$, let
$Y$ be a Hermitian space and $Z$ a skew-Hermitian space over $E$, with Gram
matrices $J_Y$ and $J_Z$ and ranks $n$ and $d$, and let $\psi$ be a standard
character of $\mathbf A_F$. Put $\Lambda=\{1,\ldots,n\}\times\{1,\ldots,d\}$,
$\mathbb W=M_{n\times d}(\mathbf A_E)$ with the symplectic form defined by
$J_Y$ and $J_Z$, $F_\infty=F\otimes\mathbf R$, and write $\mathcal S(X)$ for
the Schwartz-Bruhat functions on $X$.
\begin{enumerate}
\item A \emph{Schrödinger model} of $(Y,Z)$ consists of a symplectic
  isomorphism $\mathbb W\cong\mathbf A_F^\Lambda\times\mathbf A_F^\Lambda$
  (\emph{Darboux coordinates}), a function $\phi^0_v$ on $F_v^\Lambda$ for each
  finite place $v$ of $F$ (the \emph{vacuum} at $v$), a measure $\mu$ on
  $U(Y)(\mathbf A)$ and a subset $\mathcal F$ of $U(Y)(\mathbf A)$. The
  coordinates give the Schrödinger action $\rho(u)$, $u\in\mathbb W$, of the
  Heisenberg group on functions on $\mathbf A_F^\Lambda$. The model is
  \emph{good} if the coordinates carry $M_{n\times d}(E)$ into
  $F^\Lambda\times F^\Lambda$; each $\phi^0_v$ is a nonzero element of
  $\mathcal S(F_v^\Lambda)$ on which every element $\gamma$ of the lattice of
  $(Y,Z)$ at $v$ acts by $\rho(\gamma)$ through the scalar
  $\psi(\tfrac12\langle\vartheta\gamma,\gamma\rangle)$, for an integer
  $\vartheta$ of $E$ with $\vartheta+\bar\vartheta=1$; $\phi^0_v$ is the
  characteristic function of the integral points of $F_v^\Lambda$ for almost
  all $v$; $\mu$ is a Haar measure; and $\mathcal F$ is a fundamental domain
  for $U(Y)(F)$.
\item The \emph{theta sum} of $\Phi\in\mathcal S(\mathbf A_F^\Lambda)$ is
  $\sum_{x\in F^\Lambda}\Phi(x)$. Let $c$ be an additive bijection from
  $M_{n\times d}(\mathbf A_E)$ to the space $\mathbb W(S,T)$ of a pair of
  Hermitian spaces (Definition~\ref{th-weil-datum}), and $f$ a function on
  $\mathbb W(S,T)$. Write $\Theta(\Phi)=f$ if $\Phi$ is Schwartz-Bruhat and
  the theta sum of $\rho(u)\Phi$ is $f(c(u))$ for every
  $u\in M_{n\times d}(\mathbf A_E)$.
\item Let the model be good, and let Cartan involutions of $Y$ and $Z$ at the
  archimedean places be given. The \emph{archimedean Fock space}
  $\mathcal P_D\subset\mathcal S(F_\infty^\Lambda)$ is the smallest subspace
  that contains the Gaussian vectors, those annihilated by the annihilation
  operators that the Cartan involutions define, and is stable under the
  derivatives of the Schrödinger action along the archimedean elements of
  $\mathbb W$. The \emph{algebraic functions} $\mathcal A_D$ of the model are
  the span of the functions
  $f\otimes\bigotimes_v\phi_v\colon x\mapsto f(x_\infty)\prod_v\phi_v(x_v)$ with
  $f\in\mathcal P_D$, $\phi_v\in\mathcal S(F_v^\Lambda)$, and $\phi_v=\phi^0_v$
  for almost all $v$; these are \emph{pure tensors}.
\item Let $\varpi$ be a Weil representation of $U(Y)(\mathbf A)\times
  U(Z)(\mathbf A)$ on functions on $\mathbf A_F^\Lambda$ in the Schrödinger
  model, as in a triple over $E/F$ (Definition~\ref{th-theta-hypotheses}).
  Elements of $U(Y)(\mathbf A)$ whose components away from a finite place $v$
  are $1$ form $U(Y)(F_v)$, and they act on $\mathcal S(F_v^\Lambda)$ through a
  representation $\varpi_v$; the same holds for $Z$. The \emph{local
  augmentation} $\mathcal R_v\subset\mathcal S(F_v^\Lambda)$ is the span of
  the functions $\varpi_v(g,1)\phi-\phi$ with $g\in U(Y)(F_v)$ and
  $\phi\in\mathcal S(F_v^\Lambda)$. Let $K_\infty(Y)$ be the elements of
  $U(Y)(\mathbf A)$ with finite components $1$ that commute with the Cartan
  involutions, and call a continuous homomorphism $\gamma\colon\mathbf R\to
  U(Y)(\mathbf A)$ with values in the elements with finite components $1$ a
  \emph{one-parameter subgroup}; the same for $Z$.
\end{enumerate}
\end{definition}
\leanref{Development/AbelianVarietyTheory/BlockTheory/DivisorBase/WeilTensorModels.lean}{GlobalTheory.IndefiniteRowSideFullHecke.WeilTensorModels.Model}
\leanref{Development/AbelianVarietyTheory/BlockTheory/DivisorBase/WeilTensorModels.lean}{GlobalTheory.IndefiniteRowSideFullHecke.WeilTensorModels.Model.IsGood}
\leanref{Development/AbelianVarietyTheory/BlockTheory/DivisorBase/WeilTensorModels.lean}{GlobalTheory.IndefiniteRowSideFullHecke.WeilTensorModels.Model.Pol}
\leanref{Development/AbelianVarietyTheory/BlockTheory/DivisorBase/WeilTensorModels.lean}{GlobalTheory.IndefiniteRowSideFullHecke.WeilTensorModels.augAt}
\leanref{Development/AbelianVarietyTheory/Proofs/TwistedOperators/ThetaEqB.lean}{GlobalTheory.SquareAgree.IsDatum}
\leanref{Development/AbelianVarietyTheory/Proofs/TwistedOperators/ThetaEqB.lean}{GlobalTheory.SquareAgree.algD}

\begin{lemma}\label{th-darboux-reading}
Assume the published local inputs of the theta correspondence
(Theorem~\ref{lit-bundle-theta}). Let $D$, $\psi$ and the archimedean
embeddings be as in Definition~\ref{th-first-lift-matching}, let
$(A,A',W;\alpha,\beta)$ be a square triple of Hermitian spaces with
archimedean data, and let $\omega$ be a Weil datum for $(-A,W)$ with
characters $(\beta,\beta/\alpha)$ and $\omega'$ a Weil datum for $(A',W)$ with
characters $(\alpha,\alpha/\beta)$. Put $a=\operatorname{rk}A$,
$a'=\operatorname{rk}A'$, and let $p$ and $q$ be the numbers $s$ and $s'$ of
condition (c) of a square triple.
\begin{enumerate}
\item Let $P$ be the set of finite places of $F$ below the places of $E$ at
  which $\alpha$ or $\beta$ is ramified. Every $v\in P$ splits in $E$, does
  not divide the different of $F$, and $\alpha/\beta$ is ramified at the
  places above it; choose embeddings $E\to F_v$ over $F$ for $v\in P$. There
  are a finite place $v_1\notin P$ that splits completely in $E$ and an
  embedding $E\to F_{v_1}$. For $\tau\in\mathcal T$ let $x(\tau)$ and
  $x'(\tau)$ be archimedean exponents of $\alpha$ and $\beta$ as in condition
  (d) of a square triple; for the archimedean place $w$ of $\tau$ the integers
  $m_\alpha(w)=-2x(\tau)$ and $m_\beta(w)=-2x'(\tau)$ satisfy
  $\alpha(y)=\prod_w(\bar z_w/|z_w|)^{m_\alpha(w)}$ for ideles $y$ with
  finite components $1$, where $z_w$ is the component of $y$ at $w$ read
  through $\tau$, the same for $\beta$, and $m_\alpha-m_\beta>a+a'$. The
  extension $E/F$ is unramified at the finite places, the places inert in $E$
  do not divide $2$ and have norm greater than $3$, and $\psi$ is unramified
  at the inert places and at the places of $P$.
\item There are a Hermitian space $V$ with Gram matrix $G_A$, a Hermitian
  space $V'$ with Gram matrix $-G_{A'}$, and a skew-Hermitian space $W_\xi$ with
  Gram matrix $\xi_0G_W^{\mathsf T}$, with self-dual lattices, with the Cartan
  involutions of the archimedean data of $A$, their negatives for $A'$ and the
  complex conjugates of those of $W$, and with bases at the places of $P$
  adapted to the splitting. Their signatures are $(a-p,p)$, $(q,a'-q)$ and
  $(a+a'-p-q,p+q)$.
\item For $S=A$ with $c=-1$, and for $S=A'$ with $c=1$, the map
  $u\mapsto u^{\mathsf T}\,(cG_S)$ is an additive bijection $c_S$ from
  $M_{\operatorname{rk}S\times\operatorname{rk}W}(\mathbf A_E)$ to
  $M_{\operatorname{rk}W\times\operatorname{rk}S}(\mathbf A_E)$, the space
  $\mathbb W$ of the pairs $(-A,W)$ and $(A',W)$. The maps $\iota\colon g\mapsto\bar g$
  from $U(A)(\mathbf A)$ to $U(V)(\mathbf A)$ and from $U(A')(\mathbf A)$ to
  $U(V')(\mathbf A)$, and $\kappa\colon h\mapsto h$ from $U(W)(\mathbf A)$ to
  $U(W_\xi)(\mathbf A)$, are isomorphisms; $\iota$ carries the elements with
  components $1$ away from a finite place $v$ onto those of $U(V)(\mathbf A)$,
  and $\iota^{-1}$ carries them into $U(A)(\mathbf A_f)$; the same holds for
  $\kappa$.
\item The map $\iota$ carries a Haar measure on $U(A)(\mathbf A)$ and a
  fundamental domain for $U(A)(F)$ to a Haar measure on $U(V)(\mathbf A)$ and a
  fundamental domain for $U(V)(F)$, and the same holds for $A'$ and $V'$.
\item Let base data with the properties of the end of (1) be given, for any
  finite set of places that split in $E$ and do not divide the different of
  $F$ in place of $P$, and let $Y$ and $Z$ be a Hermitian and a skew-Hermitian
  space with the Gram matrices of $V$ (or $V'$) and $W_\xi$ in (2), with
  self-dual lattices, Cartan involutions at the archimedean places and bases at
  the places of that set adapted to the splitting; the spaces of (2) are an
  example. There is a good Schrödinger model of $(Y,Z)$
  (Definition~\ref{th-darboux-model}), and there is one with any given Haar
  measure $\mu$ on $U(Y)(\mathbf A)$ and fundamental domain $\mathcal F$ for
  $U(Y)(F)$. Its vacuum at $v$ is the characteristic function of the integral
  points when that function is a vacuum, which holds at almost all $v$. For a
  good model and the bijection $c_A$ (or $c_{A'}$) of (3), $\Theta$ is linear,
  and $\Theta(\Phi_D)=\Theta(\Psi_D)$ implies $\Phi_D=\Psi_D$.
\item For a good Schrödinger model of $(Y,Z)$ as in (5), with the Gram matrix
  of $V$, there is a Weil representation $\varpi$ of the model with characters
  $(\beta,\beta/\alpha)$ such that
  $\Theta(\varpi(\iota(g),\kappa(h))\Phi_D)=\omega(g,h)\Theta(\Phi_D)$ for
  every Schwartz-Bruhat function $\Phi_D$, $g\in U(A)(\mathbf A)$ and
  $h\in U(W)(\mathbf A)$, where $\Theta$ is formed with $c_A$. In particular
  $\theta_\varpi(\iota(g),\kappa(h);\Phi_D)=\theta_\omega(g,h;\Theta(\Phi_D))$.
  The same holds with the Gram matrix of $V'$, $(\alpha,\alpha/\beta)$,
  $c_{A'}$ and $\omega'$.
\item Let a good Schrödinger model be as in (5), for spaces with the Gram
  matrix of $V$ and the Cartan involutions of (2), whose archimedean Fock space
  is defined by these Cartan involutions. Every $\Phi\in\mathcal A(-A,W)$ is
  $\Theta(\Phi_D)$ for an algebraic function $\Phi_D$ of the model, and
  $\Theta(\Phi_D)\in\mathcal A(-A,W)$ for every algebraic function $\Phi_D$.
  The same holds for $A'$.
\item Let $\mu$ and $\mathcal F$ be the images under $\iota$ of a Haar measure
  $\nu$ on $U(A)(\mathbf A)$ and a fundamental domain $\mathcal F_A$ for
  $U(A)(F)$, and let the model of (5) and $\varpi$ of (6) be formed with them.
  If $\Theta(\Phi_D)=\Phi$, then the theta integral $I_\omega(\Phi)$ for
  $\nu$ and $\mathcal F_A$ (Definition~\ref{th-first-lift-matching}(5)) is
  $h\mapsto\int_{\mathcal F}\theta_\varpi(g,\kappa(h);\Phi_D)\,d\mu(g)$. The
  same holds for $A'$.
\item The algebraic functions of a good Schrödinger model have countable
  dimension: so do its archimedean Fock space and each
  $\mathcal S(F_v^\Lambda)$.
\item There are a triple over $E/F$ with $\psi$, the characters $\alpha$
  and $\beta$, the Gram matrices and Cartan involutions of (2), a finite set
  $T$ of places of $F$ that split in $E$ and do not divide the different of
  $F$, such that $\alpha/\beta$ is ramified at the places above $T$ and
  $\alpha$ and $\beta$ are unramified at the places above the places outside
  $T$, good Schrödinger models, and Weil representations $\varpi$ and
  $\varpi'$ with characters $(\beta,\beta/\alpha)$ and $(\alpha,\alpha/\beta)$,
  such that
  $\Theta(\varpi(\iota(g),\kappa(h))\Phi_D)=\omega(g,h)\Theta(\Phi_D)$ for
  every $\Phi_D\in\mathcal S(\mathbf A_F^\Lambda)$, $g\in U(A)(\mathbf A)$ and
  $h\in U(W)(\mathbf A)$, and the same for $\varpi'$ and $\omega'$.
\end{enumerate}
\end{lemma}
\leanref{Development/AbelianVarietyTheory/Proofs/TwistedOperators/ExistsThetaTripleBaseDataCharactersStatement.lean}{GlobalTheory.T.exists_thetaTriple_baseData_characters}
\leanref{Development/AbelianVarietyTheory/Proofs/TwistedOperators/ExistsPlacewiseTripleSpacesOfSeesawStatement.lean}{GlobalTheory.T.exists_placewise_triple_spaces_of_seesaw}
\leanref{Development/AbelianVarietyTheory/Proofs/TwistedOperators/ExistsAddEquivSymplecticCoordsOrientedModelStatement.lean}{GlobalTheory.T.exists_addEquiv_symplecticCoords_orientedModel}
\leanref{Development/AbelianVarietyTheory/Proofs/TwistedOperators/ConjUInvMemGfC.lean}{GlobalTheory.SquareAgree.CoinvPass.conjUInv_mem_GfC}
\leanref{Development/AbelianVarietyTheory/Proofs/TwistedOperators/TransposeGramUnitarySurjectiveFinitePlaceSupportStatement.lean}{GlobalTheory.T.transposeGramUnitary_surjective_finitePlaceSupported}
\leanref{Development/AbelianVarietyTheory/Proofs/TwistedOperators/ConjugateUnitaryAdelicIsoMapHaarFundamentalDomaiStatement.lean}{GlobalTheory.T.conjugateUnitaryAdelicIso_map_haar_fundamentalDomain}
\leanref{Development/AbelianVarietyTheory/Proofs/TwistedOperators/ExistsGoodRationalSchroedingerModelStatement.lean}{GlobalTheory.T.exists_good_rational_schroedingerModel}
\leanref{Development/AbelianVarietyTheory/Proofs/TwistedOperators/EqZeroOfThetaDistributionHeisenbergTranslatesStatement.lean}{GlobalTheory.T.eq_zero_of_thetaDistribution_heisenbergTranslates}
\leanref{Development/AbelianVarietyTheory/Proofs/TwistedOperators/SchwartzDatumEqOfThetaFunctionEq.lean}{GlobalTheory.ThetaInj.theta_inj}
\leanref{Development/AbelianVarietyTheory/Proofs/TwistedOperators/ExistsWeilRepOrientedPairRationalSchrodingerStatement.lean}{GlobalTheory.T.exists_weilRep_orientedPair_rationalSchrodinger}
\leanref{Development/AbelianVarietyTheory/Proofs/TwistedOperators/ThetaEqB.lean}{GlobalTheory.SquareAgree.IsOver.theta_eq_A}
\leanref{Development/AbelianVarietyTheory/Proofs/TwistedOperators/AlgebraicVectorsCorrespondSchrodingerThetaMapStatement.lean}{GlobalTheory.T.algebraicVectors_correspond_schrodinger_thetaMap}
\leanref{Development/AbelianVarietyTheory/Proofs/TwistedOperators/ConjUMeasurableEmbedding.lean}{GlobalTheory.SquareAgree.thetaInt_eq_modelI}
\leanref{Development/AbelianVarietyTheory/Proofs/TwistedOperators/RankAlgebraicDataSchrodingerModelLeAleph0Statement.lean}{GlobalTheory.T.rank_algebraicData_schrodingerModel_le_aleph0}
\leanref{Development/AbelianVarietyTheory/Proofs/TwistedOperators/RankPolynomialFockSchrodingerModelLeAleph0Statement.lean}{GlobalTheory.T.rank_polynomialFock_schrodingerModel_le_aleph0}
\leanref{Development/AbelianVarietyTheory/Proofs/TwistedOperators/RankSchwartzBruhatFinitePlaceCompletionLeAleph0Statement.lean}{GlobalTheory.T.rank_schwartzBruhat_finitePlaceCompletion_le_aleph0}
\leanref{Development/AutomorphicArithmetic/ThetaCorrespondence/ThetaPairLifts/LocalTypeDictionary/ExistsDarbouxCoordinatesTriple.lean}{GlobalTheory.T.exists_darboux_coordinates_triple}
\leanref{Development/AbelianVarietyTheory/Proofs/TwistedOperators/TripleFrameBaseExists.lean}{GlobalTheory.TripleDarboux.exists_model}
\leanref{Development/AbelianVarietyTheory/Proofs/TwistedOperators/TripleFrameBaseExists.lean}{GlobalTheory.TripleDarboux.exists_weil}
\begin{proof}
For (1), $\alpha$ and $\beta$ are Hecke characters, so they are ramified at
finitely many places and $P$ is finite. Let $v\in P$. Since $\alpha$ or
$\beta$ is ramified above $v$, the first alternative of condition (e) of a
square triple fails at $v$, so $v$ splits in $E$, $\alpha/\beta$ is ramified
at a place above $v$, and $v$ does not divide the different of $F$. The
character $\alpha/\beta$ is trivial on $\mathbf A_F^\times$, since by
conditions (a) and (b) it restricts to
$\eta^{\operatorname{rk}A'-\operatorname{rk}A}=1$; so it is ramified at both
places above $v$. For $v_1$, take a rational prime $\ell$ larger than the sum
of the norms of the places of $P$ and congruent to $1$ modulo the integer
$N=Q_Sqd$ of the host field (Chapter~\ref{chapter-host-fields}); by
Dirichlet's theorem such $\ell$ exists, it splits completely in $E$, and a
place of $F$ above $\ell$ is not in $P$. By the definition of archimedean exponents, $\alpha(y)=(z/|z|)^{2x(\tau)}
=(\bar z/|z|)^{-2x(\tau)}$ at the place of $\tau$, which gives $m_\alpha$ and
$m_\beta$, and $m_\alpha-m_\beta=-2(x-x')>\operatorname{rk}W=a+a'$ by
conditions (b) and (d). The extension $E/F$ of the host field is unramified at
the finite places, its inert places satisfy the conditions on $2$ and the
norm, and a standard character is unramified at the places that do not
divide the different, which include the inert places and the places of $P$.

Parts (2) to (10) are constructions and computations in the models, and their
results are stated. In (2) the lattices are the conjugates of those of the
frames of $A$ and $A'$ and the lattice of $W$ dual to its frame, transposed,
and the bases at the places of $P$ split the lattices along the embeddings of
(1). In (5) the Darboux coordinates are rational, and at the places where the
characteristic function of the integral points is not a vacuum, a vacuum for
the lattice is chosen; when no measure is given, a Haar measure and a
fundamental domain for $U(Y)(F)$ are chosen, which exist since
$U(Y)(\mathbf A)$ is locally compact and second countable and $U(Y)(F)$ is a
discrete subgroup. The injectivity of $\Theta$ is the statement that a
Schwartz-Bruhat function whose Heisenberg translates all have theta sum $0$
is $0$. In (6) the Weil representation $\varpi$ is obtained from $\omega$ by
transport of structure along $\Theta$, which intertwines the two Heisenberg
actions; the theta kernels agree since $\theta_\varpi(g,h;\Phi_D)$ is the
theta sum of $\varpi(g,h)\Phi_D$, which is
$(\omega(g,h)\Phi)(c_A(0))=\theta_\omega(g,h;\Phi)$. In (7), $\Theta$ carries
a pure tensor with archimedean component in the Fock space to a vector
$f\otimes P$ of Definition~\ref{th-first-lift-matching}(1), and conversely.
Part (8) is the change of variables $g\mapsto\iota(g)$ in the integral, with
the agreement of the theta kernels of (6). In (9), a basis of the algebraic
functions is formed from pure tensors of basis vectors of the Fock space and of
the spaces $\mathcal S(F_v^\Lambda)$ at finitely many places, with the vacua at
the others; there are countably many finite sets of places. Part (10) repeats
the constructions of (1), (2), (5) and (6), with a Haar measure and a
fundamental domain chosen in (5), and with $T$ the set of places of $F$ below
the places at which $\alpha$ or $\beta$ is ramified.
\end{proof}

\begin{definition}\label{th-kernel-theta-subspace}
Let a triple over $E/F$ (Definition~\ref{th-theta-hypotheses}) be given, with
spaces $V$, $V'$ and $W$, good Schrödinger models of $(V,W)$ and $(V',W)$
(Definition~\ref{th-darboux-model}) with Fock spaces $\mathcal P_D$ and
$\mathcal P'_D$, algebraic functions $\mathcal A_D$ and $\mathcal A'_D$ and
vacua $\phi^0_v$ and $\phi'^0_v$, Weil representations $\varpi$ and $\varpi'$,
and finite set $T_X$ of places.
\begin{enumerate}
\item A \emph{theta subspace} is a subspace $N\subset\mathcal A_D$ such that:
  \begin{enumerate}
  \item[(a)] $\varpi(1,h)N\subset N$ for $h\in U(W)(F_v)$, $v$ finite, and for
    $h\in K_\infty(W)$; and $\Phi_1\in N$ whenever $\Phi\in N$,
    $\Phi_1\in\mathcal A_D$ and $\Phi_1(x)$ is the derivative at $t=0$ of
    $(\varpi(1,\gamma(t))\Phi)(x)$ for every $x$, for a one-parameter subgroup
    $\gamma$ of $U(W)$;
  \item[(b)] $\varpi(g,1)\Phi-\Phi\in N$ for $\Phi\in\mathcal A_D$ and
    $g\in U(V)(F_v)$, $v$ finite, or $g\in K_\infty(V)$; and $\Phi_1\in N$
    whenever $\Phi,\Phi_1\in\mathcal A_D$ and $\Phi_1$ is the derivative at
    $t=0$ of $\varpi(\gamma(t),1)\Phi$ for a one-parameter subgroup $\gamma$ of
    $U(V)$.
  \end{enumerate}
\item A \emph{theta subspace of the second kind} is a subspace
  $N\subset\mathcal A'_D$ that satisfies (a) for $\varpi'$ and contains: the
  pure tensors $f\otimes\bigotimes_w\phi_w$ in $\mathcal A'_D$ with
  $\phi_v\in\mathcal R'_v$ for some finite $v$; the pure tensors
  $f_1\otimes\bigotimes_w\phi_w$ with $f,f_1\in\mathcal P'_D$ such that
  $f_1\otimes r$ is the derivative at $0$ of $\varpi'(\gamma(t),1)(f\otimes r)$
  for every function $r$ of the finite components, $\gamma$ a one-parameter
  subgroup of $U(V')$; and the differences
  $f_1\otimes\bigotimes_w\phi_w-f\otimes\bigotimes_w\phi_w$ with
  $\varpi'(k,1)(f\otimes r)=f_1\otimes r$ for every $r$, $k\in K_\infty(V')$.
\item \emph{Matched local maps} for the triple consist of complex vector spaces
  $\Pi_{\mathbf R}$ and $\Pi_v$, $v$ finite, and linear maps $r$ and $r'$
  from the Schwartz functions on $F_\infty^\Lambda$ and $F_\infty^{\Lambda'}$
  to $\Pi_{\mathbf R}$, and $r_v\colon\mathcal S(F_v^\Lambda)\to\Pi_v$ and
  $r'_v\colon\mathcal S(F_v^{\Lambda'})\to\Pi_v$, with the conditions of the
  local matching, in which $f$ and $f'$ range over $\mathcal P_D$ and
  $\mathcal P'_D$: $r(\mathcal P_D)=r'(\mathcal P'_D)=\Pi_{\mathbf R}$; $r$
  vanishes on the derivatives of $f$ along one-parameter subgroups of $U(V)$
  and is invariant under $K_\infty(V)$, and the same for $r'$ and $V'$;
  whenever $r(f)=r'(f')$, $r$ and $r'$ agree on the
  derivatives of $f$ and $f'$ along a one-parameter subgroup of $U(W)$ and on
  their translates by an element of $K_\infty(W)$; $\Pi_{\mathbf R}\neq0$, and
  its only subspaces stable under the operators so induced by $r$ are $0$ and
  $\Pi_{\mathbf R}$; $r_v$ and $r'_v$ are onto and vanish on $\mathcal R_v$ and
  $\mathcal R'_v$; whenever $r_v(\phi)=r'_v(\phi')$,
  $r_v(\varpi_v(1,h)\phi)=r'_v(\varpi'_v(1,h)\phi')$ for $h\in U(W)(F_v)$;
  $\Pi_v\neq0$, and its only subspaces stable under the operators so induced
  by $r_v$ are $0$ and $\Pi_v$; $r_v(\phi^0_v)=r'_v(\phi'^0_v)\neq0$ for
  $v\notin T_X$; and a condition on local zeta integrals at the places of
  $T_X$.
\item Let matched local maps be given. For a finite set $T\supset T_X$ of
  finite places, an \emph{expression over $T$} of $\Phi\in\mathcal A_D$ is a
  sum $\Phi=\sum_if_i\otimes\bigotimes_v\phi_{i,v}$ of pure tensors with
  $\phi_{i,v}=\phi^0_v$ for $v\notin T$; its \emph{image} is
  $\sum_ir(f_i)\otimes\bigotimes_{v\in T}r_v(\phi_{i,v})$ in
  $\Pi_{\mathbf R}\otimes\bigotimes_{v\in T}\Pi_v$. The same applies to
  $\mathcal A'_D$ with $r'$ and $r'_v$. Two functions $\Phi\in\mathcal A_D$
  and $\Phi'\in\mathcal A'_D$ have \emph{agreeing images} if for some $T$ they
  have expressions over $T$ with equal images. The \emph{image} of $\Phi$
  \emph{is $0$} if $\Phi$ and $0$ have agreeing images, and the same for
  $\Phi'$.
\end{enumerate}
\end{definition}
\leanref{Development/AbelianVarietyTheory/Proofs/TwistedOperators/IsThetaSubSelf.lean}{GlobalTheory.CoinvKer.IsThetaSub}
\leanref{Development/AbelianVarietyTheory/Proofs/TwistedOperators/IsThetaSubSelf.lean}{GlobalTheory.CoinvKer.IsStabW}
\leanref{Development/AbelianVarietyTheory/Proofs/TwistedOperators/IsThetaSubSelf.lean}{GlobalTheory.CoinvKer.HasRelY}
\leanref{Development/AbelianVarietyTheory/Proofs/TwistedOperators/IsThetaSubGAlgDOne.lean}{GlobalTheory.SquareAgree.IsThetaSubG}
\leanref{Development/AbelianVarietyTheory/BlockTheory/DivisorBase/WeilTensorModels.lean}{GlobalTheory.IndefiniteRowSideFullHecke.WeilTensorModels.IsMatched}
\leanref{Development/AbelianVarietyTheory/BlockTheory/DivisorBase/WeilTensorModels.lean}{GlobalTheory.IndefiniteRowSideFullHecke.WeilTensorModels.IsRealMaps}
\leanref{Development/AbelianVarietyTheory/BlockTheory/DivisorBase/WeilTensorModels.lean}{GlobalTheory.IndefiniteRowSideFullHecke.WeilTensorModels.IsFinMaps}
\leanref{Development/AbelianVarietyTheory/BlockTheory/DivisorBase/WeilTensorModels.lean}{GlobalTheory.IndefiniteRowSideFullHecke.WeilTensorModels.PolMatch}
\leanref{Development/AbelianVarietyTheory/Proofs/TwistedOperators/SumPureMemAlgD.lean}{GlobalTheory.PolKer.PolZero}

\begin{lemma}\label{th-kernel-dictionary}
Assume the published local inputs of the theta correspondence
(Theorem~\ref{lit-bundle-theta}). Let $(A,A',W;\alpha,\beta)$ be a square
triple, $\omega$ and $\omega'$ Weil data as in a matching of it, and Haar
measures and fundamental domains for $U(A)$ and $U(A')$ given. Let the triple
over $E/F$ of Lemma~\ref{th-square-model}(2) be formed with them, with spaces
$V$, $V'$, $W_\xi$, Weil representations $\varpi$ and $\varpi'$ and the maps
$\iota$, $\kappa$ and $\Theta$ of Lemma~\ref{th-darboux-reading}, and let
matched local maps for it and a matching $(\omega,\omega',\Pi,\varrho,p,p')$ of
$(A,A',W;\alpha,\beta)$ be given. For $\Phi_D\in\mathcal A_D$ put
$p_D(\Phi_D)=p(\Theta(\Phi_D))$, and for $\Phi'_D\in\mathcal A'_D$ put
$p'_D(\Phi'_D)=p'(\Theta(\Phi'_D))$
(Lemma~\ref{th-darboux-reading}(7)).
\begin{enumerate}
\item Let $\Phi_D\in\mathcal A_D$.
  \begin{enumerate}
  \item[(a)] For $k\in K_\infty(W_\xi)$, $\varpi(1,k)\Phi_D\in\mathcal A_D$,
    and $p_D(\Phi_D)=0$ implies $p_D(\varpi(1,k)\Phi_D)=0$.
  \item[(b)] If $\Phi_1\in\mathcal A_D$ is the derivative at $0$ of
    $\varpi(1,\gamma(t))\Phi_D$ for a one-parameter subgroup $\gamma$ of
    $U(W_\xi)$, then $p_D(\Phi_D)=0$ implies $p_D(\Phi_1)=0$.
  \item[(c)] For $k\in K_\infty(V)$, $\varpi(k,1)\Phi_D\in\mathcal A_D$ and
    $p_D(\varpi(k,1)\Phi_D)=p_D(\Phi_D)$.
  \item[(d)] If $\Phi_1\in\mathcal A_D$ is the derivative at $0$ of
    $\varpi(\gamma(t),1)\Phi_D$ for a one-parameter subgroup $\gamma$ of
    $U(V)$, then $p_D(\Phi_1)=0$.
  \item[(e)] If a subspace $N\subset\mathcal A_D$ satisfies condition (a) of a
    theta subspace, then $\varpi(1,\kappa(h))N\subset N$ for
    $h\in U(W)(\mathbf A_f)$; and for $\tau\in\mathcal T$,
    $Y\in\mathfrak u_\tau(W)$ and $\Phi_D\in N$ there is $\Phi_{1,D}\in N$
    such that $\Theta(\Phi_{1,D})\in\mathcal A(-A,W)$ is the derivative at
    $s=0$ of $\omega(1,\exp_\tau(sY))\Theta(\Phi_D)$.
  \end{enumerate}
\item The set of $\Phi'_D\in\mathcal A'_D$ with $p'_D(\Phi'_D)=0$ satisfies
  condition (a) of a theta subspace for $\varpi'$ and contains the elements
  listed in Definition~\ref{th-kernel-theta-subspace}(2). Part (1)(e) holds for
  $\varpi'$ and $\omega'$.
\item Two pairs with agreeing images give agreeing images for
  $(\Phi+c\Psi,\Phi'+c\Psi')$, $c\in\mathbf C$, and $0$ and $0$ have agreeing
  images. The functions in $\mathcal A_D$ with image $0$ form a theta subspace
  different from $\mathcal A_D$, and those in $\mathcal A'_D$ with image $0$
  form a theta subspace of the second kind different from $\mathcal A'_D$.
\item For every $\Phi_D\in\mathcal A_D$ there is $\Phi'_D\in\mathcal A'_D$
  such that they have agreeing images. If $\Phi_D$ and $\Phi'_D$ have agreeing
  images, then so do $\varpi(1,\kappa(h))\Phi_D$ and
  $\varpi'(1,\kappa(h))\Phi'_D$ for $h\in U(W)(\mathbf A_f)$, and so do
  $\Phi_{1,D}$ and $\Phi'_{1,D}$ whenever $\Theta(\Phi_{1,D})$ and
  $\Theta(\Phi'_{1,D})$ are the derivatives at $s=0$ of
  $\omega(1,\exp_\tau(sY))\Theta(\Phi_D)$ and
  $\omega'(1,\exp_\tau(sY))\Theta(\Phi'_D)$, for $\tau\in\mathcal T$ and
  $Y\in\mathfrak u_\tau(W)$, and all four lie in $\mathcal A_D$ and
  $\mathcal A'_D$.
\item If $\Phi_D$ and $\Phi'_D$ have agreeing images and $\Phi'_D$ lies in the
  space $\mathcal S'$ of Lemma~\ref{th-square-model}(3), then
  $(\Phi_D,\Phi'_D)$ is a pair with matching images in the sense of
  Definition~\ref{th-theta-hypotheses}(2).
\end{enumerate}
\end{lemma}
\leanref{Development/AutomorphicArithmetic/ThetaCorrespondence/ThetaPairLifts/LocalTypeDictionary/SquareAgreeStabilizerCompactWeyl.lean}{GlobalTheory.T.square_agree_stabilizer_compact_weyl}
\leanref{Development/AutomorphicArithmetic/ThetaCorrespondence/ThetaPairLifts/LocalTypeDictionary/KerProjectionMaximalCompactStable.lean}{GlobalTheory.T.square_agree_kernel_compact_stable}
\leanref{Development/AutomorphicArithmetic/ThetaCorrespondence/ThetaPairLifts/LocalTypeDictionary/LineKernelProjectionWeyl.lean}{GlobalTheory.T.line_kernel_projection_weyl}
\leanref{Development/AutomorphicArithmetic/ThetaCorrespondence/ThetaPairLifts/LocalTypeDictionary/WeilAlgebraicSubspaceKInftyStable.lean}{GlobalTheory.T.weil_algebraic_subspace_kInfty_stable}
\leanref{Development/AutomorphicArithmetic/ThetaCorrespondence/ThetaPairLifts/LocalTypeDictionary/WeilProjectionKInftyInvariant.lean}{GlobalTheory.T.weil_projection_kInfty_invariant}
\leanref{Development/AutomorphicArithmetic/ThetaCorrespondence/ThetaPairLifts/LocalTypeDictionary/ThetaLineVectorAnnihilation.lean}{GlobalTheory.T.theta_line_vector_annihilation}
\leanref{Development/AutomorphicArithmetic/ThetaCorrespondence/ThetaPairLifts/LocalTypeDictionary/FiniteAdelicStabilizerFixesKappa.lean}{GlobalTheory.T.finite_adelic_stabilizer_fixes_kappa}
\leanref{Development/AutomorphicArithmetic/ThetaCorrespondence/ThetaPairLifts/LocalTypeDictionary/WeilLieAlgebraActionByLineDerivatives.lean}{GlobalTheory.T.weil_lie_algebra_action_by_line_derivatives}
\leanref{Development/AutomorphicArithmetic/ThetaCorrespondence/ThetaPairLifts/LocalTypeDictionary/KerPolarizationProjectionFiniteDimensional.lean}{GlobalTheory.T.polarization_kernel_finiteDimensional}
\leanref{Development/AutomorphicArithmetic/ThetaCorrespondence/ThetaPairLifts/LocalTypeDictionary/LieActionMemKerPolarizationProjection.lean}{GlobalTheory.T.ker_polarizationProjection_lie_stable}
\leanref{Development/AutomorphicArithmetic/ThetaCorrespondence/ThetaPairLifts/LocalTypeDictionary/CompactActionMemKerPolarizationProjection.lean}{GlobalTheory.T.ker_polarizationProjection_compact_stable}
\leanref{Development/AutomorphicArithmetic/ThetaCorrespondence/ThetaPairLifts/LocalTypeDictionary/KerPolarizationProjectionStabilizerFinite.lean}{GlobalTheory.T.ker_polarizationProjection_stabilizer_finite}
\leanref{Development/AutomorphicArithmetic/ThetaCorrespondence/ThetaPairLifts/LocalTypeDictionary/KerPolarizationProjectionStableUnderStabilizerLi.lean}{GlobalTheory.T.ker_polarizationProjection_stable_under_stabilizer_lie}
\leanref{Development/AutomorphicArithmetic/ThetaCorrespondence/ThetaPairLifts/LocalTypeDictionary/KerPolarizationProjectionStableUnderStabilizerCo.lean}{GlobalTheory.T.ker_polarizationProjection_stable_under_stabilizer_compact}
\leanref{Development/AutomorphicArithmetic/ThetaCorrespondence/ThetaPairLifts/LocalTypeDictionary/AdelicStabilizerFiniteOfThetaSubgroup.lean}{GlobalTheory.T.adelic_stabilizer_finite_of_theta_subgroup}
\leanref{Development/AutomorphicArithmetic/ThetaCorrespondence/ThetaPairLifts/LocalTypeDictionary/AdelicStabilizerLieKerPolarization.lean}{GlobalTheory.T.adelic_stabilizer_lie_ker_polarization}
\leanref{Development/AbelianVarietyTheory/Proofs/TwistedOperators/PolMatchLinear.lean}{GlobalTheory.PolMatchLin.polMatch_linear}
\leanref{Development/AutomorphicArithmetic/ThetaCorrespondence/ThetaPairLifts/LocalTypeDictionary/PolKernelEqThetaSubspaceOfMatch.lean}{GlobalTheory.T.thetaSubspace_contains_polKernel}
\leanref{Development/AutomorphicArithmetic/ThetaCorrespondence/ThetaPairLifts/LocalTypeDictionary/ThetaSubgroupOfKerDerivation.lean}{GlobalTheory.T.thetaSubgroup_of_ker_derivation}
\leanref{Development/AbelianVarietyTheory/Proofs/TwistedOperators/TmulSumEqZeroOfPureTensorFunctionZeroStatement.lean}{GlobalTheory.T.tmul_sum_eq_zero_of_pureTensor_function_zero}
\leanref{Development/AbelianVarietyTheory/Proofs/TwistedOperators/PolMatchOnto.lean}{GlobalTheory.SquareAgree.polMatch_onto}
\leanref{Development/AutomorphicArithmetic/ThetaCorrespondence/ThetaPairLifts/LocalTypeDictionary/PolarizationMatchSquareKappa.lean}{GlobalTheory.SquareAgree.polarization_match_square_kappa}
\leanref{Development/AutomorphicArithmetic/ThetaCorrespondence/ThetaPairLifts/LocalTypeDictionary/PolarizationMatchLieDerivative.lean}{GlobalTheory.SquareAgree.LieEq.polarization_match_lie_derivative}
\leanref{Development/AbelianVarietyTheory/Proofs/TwistedOperators/MatchOfPolMatch.lean}{GlobalTheory.SquareAgree.PolMatchConst.match_of_polMatch}
\begin{proof}
Each part reads a condition of the matching or of the matched local maps
through the theta-sum map $\Theta$, which carries $\varpi$ to $\omega$ and the
algebraic functions onto $\mathcal A(-A,W)$
(Lemma~\ref{th-darboux-reading}(6) and (7)); the results are stated. In
(1)(b) and (1)(d), a one-parameter subgroup of $U(W_\xi)$ or $U(V)$ is a
product of one-parameter subgroups at the complex places, the derivative of
$\Theta(\Phi_D)$ along each is a derivative along a subgroup
$\exp_\tau(sY)$ as in Definition~\ref{th-first-lift-matching}(4), and
conditions (c) and (d) of a matching give $p_D(\Phi_1)=\sum\varrho_\tau(Y)
p_D(\Phi_D)=0$ in (b) and $p_D(\Phi_1)=0$ in (d). In (1)(a) and (1)(c), the
groups $K_\infty$ preserve the Fock space, an algebraic function lies in a
finite-dimensional space of algebraic functions stable under $K_\infty$ and
these derivatives, and $K_\infty$ is a product of exponentials of
one-parameter subgroups; on such a space the conditions on derivatives
integrate to the conditions on $K_\infty$. In (1)(e), $\kappa$ carries
$U(W)(\mathbf A_f)$ to elements that are products of elements with components
$1$ away from one finite place (Lemma~\ref{th-darboux-reading}(3)), and a
derivative along $\exp_\tau(sY)$ in $U(W)$ is one along a one-parameter
subgroup of $U(W_\xi)$. Part (2) is the same reading for the second pair.

In (3) and (4) the image of an expression over $T$ depends only on the
function: if a sum of pure tensors over $T$ is $0$, so is its image. The maps
$r$ and $r_v$ vanish on the relations of a theta subspace, and the
equivariance conditions on matched local maps, applied with $0$ on one side,
show that the functions with image $0$ satisfy condition (a); the same holds
for the second pair, and agreeing images are carried to agreeing images in
(4) in the same way. The pure tensor $f\otimes\bigotimes_v\phi_v$ with
$r(f)\neq0$, $r_v(\phi_v)\neq0$ at $v\in T_X$ and $\phi_v=\phi^0_v$ elsewhere
has image $r(f)\otimes\bigotimes_vr_v(\phi_v)\neq0$, since
$r_v(\phi^0_v)\neq0$ for $v\notin T_X$; so the image $0$ does not occur for
all algebraic functions. Since $r'$ and the $r'_v$ are onto, every image of an
expression of $\Phi_D$ is the image of an expression of some $\Phi'_D$, which
gives the first statement of (4). In (5), the pair with matching images of
Definition~\ref{th-theta-hypotheses}(2) is formed with the constant sequences
of the expressions of $\Phi_D$ and $\Phi'_D$.
\end{proof}

\begin{lemma}\label{th-kernel-matching}
In the situation of Lemma~\ref{th-kernel-dictionary}:
\begin{enumerate}
\item the set $N_p$ of $\Phi_D\in\mathcal A_D$ with $p_D(\Phi_D)=0$ is a theta
  subspace, different from $\mathcal A_D$, and every theta subspace that
  contains $N_p$ is $N_p$ or $\mathcal A_D$;
\item the set $N'_p$ of $\Phi'_D\in\mathcal A'_D$ with $p'_D(\Phi'_D)=0$ is a
  theta subspace of the second kind, different from $\mathcal A'_D$, and every
  theta subspace of the second kind that contains $N'_p$ is $N'_p$ or
  $\mathcal A'_D$.
\end{enumerate}
\end{lemma}
\leanref{Development/AutomorphicArithmetic/ThetaCorrespondence/ThetaPairLifts/LocalTypeDictionary/ThetaSubspaceKernelMaximal.lean}{GlobalTheory.T.theta_subspace_kernel_maximal}
\leanref{Development/AutomorphicArithmetic/ThetaCorrespondence/ThetaPairLifts/LocalTypeDictionary/ThetaSubspaceEqKerProjDerivationOfSquare.lean}{GlobalTheory.T.thetaSubspace_eq_ker_projDerivation_of_square}
\leanref{Development/AutomorphicArithmetic/ThetaCorrespondence/ThetaPairLifts/LocalTypeDictionary/KerProjectionMaximalThetaSubspace.lean}{GlobalTheory.T.thetaKernel_projection_isMax}
\begin{proof}
For (1), $N_p$ is a subspace, since $\Theta$ and $p$ are linear. It lies in
$\mathcal A_D$ and satisfies condition (a) of a theta subspace. Let $v$ be a
finite place, $h\in U(W_\xi)(F_v)$ and $\Phi_D\in N_p$. By
Lemma~\ref{th-darboux-reading}(3), $h=\kappa(h')$ with
$h'\in U(W)(\mathbf A_f)$. By Lemma~\ref{th-darboux-reading}(6),
$\Theta(\varpi(1,h)\Phi_D)=\omega(1,h')\Theta(\Phi_D)$, which lies in
$\mathcal A(-A,W)$ by condition (a) of a matching; so
$\varpi(1,h)\Phi_D\in\mathcal A_D$ by Lemma~\ref{th-darboux-reading}(5) and
(7). By condition (d) of a matching,
$p_D(\varpi(1,h)\Phi_D)=\varrho(h')p_D(\Phi_D)=0$. The elements of
$K_\infty(W_\xi)$ and the derivatives along one-parameter subgroups of
$U(W_\xi)$ preserve $N_p$ by Lemma~\ref{th-kernel-dictionary}(1)(a) and (b).

It satisfies condition (b). Let $g\in U(V)(F_v)$ and $\Phi_D\in\mathcal A_D$.
Then $g=\iota(g')$ with $g'\in U(A)(\mathbf A_f)$, and as above
$\varpi(g,1)\Phi_D\in\mathcal A_D$ and $p_D(\varpi(g,1)\Phi_D)=
p(\omega(g',1)\Theta(\Phi_D))=p_D(\Phi_D)$ by condition (c) of a matching; so
$\varpi(g,1)\Phi_D-\Phi_D\in N_p$. For $k\in K_\infty(V)$ the same holds by
Lemma~\ref{th-kernel-dictionary}(1)(c), and the derivatives along
one-parameter subgroups of $U(V)$ lie in $N_p$ by
Lemma~\ref{th-kernel-dictionary}(1)(d).

It is different from $\mathcal A_D$. Since $\Pi\neq0$ and
$p(\mathcal A(-A,W))=\Pi$, some $\Phi\in\mathcal A(-A,W)$ has $p(\Phi)\neq0$,
and $\Phi=\Theta(\Phi_D)$ with $\Phi_D\in\mathcal A_D$
(Lemma~\ref{th-darboux-reading}(7)); then $\Phi_D\notin N_p$.

Let $N\supset N_p$ be a theta subspace, and let
$\Pi_N=p_D(N)\subset\Pi$. It is a subspace. It is stable under $\varrho(h)$
for $h\in U(W)(\mathbf A_f)$: for $\Phi_D\in N$,
$\varpi(1,\kappa(h))\Phi_D\in N$ by Lemma~\ref{th-kernel-dictionary}(1)(e),
and $p_D(\varpi(1,\kappa(h))\Phi_D)=\varrho(h)p_D(\Phi_D)$ as above. It is
stable under $\varrho_\tau(Y)$: for $\Phi_D\in N$, the function
$\Phi_{1,D}\in N$ of Lemma~\ref{th-kernel-dictionary}(1)(e) has
$p_D(\Phi_{1,D})=\varrho_\tau(Y)p_D(\Phi_D)$ by condition (d) of a matching.
Since $0$ and $\Pi$ are the only subspaces stable under all these operators,
$\Pi_N=0$ or $\Pi_N=\Pi$. If $\Pi_N=0$, then $N\subset N_p$, so $N=N_p$. If
$\Pi_N=\Pi$, let $\Phi_D\in\mathcal A_D$; there is $\Psi_D\in N$ with
$p_D(\Psi_D)=p_D(\Phi_D)$, so $\Phi_D-\Psi_D\in N_p\subset N$ and
$\Phi_D\in N$; thus $N=\mathcal A_D$.

For (2), the set $N'_p$ is a theta subspace of the second kind by
Lemma~\ref{th-kernel-dictionary}(2), and it is different from
$\mathcal A'_D$ for the same reason as in (1). If $N'\supset N'_p$ is a theta
subspace of the second kind, then $p'_D(N')$ is stable under the $\varrho(h)$
and $\varrho_\tau(Y)$, by Lemma~\ref{th-kernel-dictionary}(2) and conditions
(d) of a matching for $p'$, and the argument of (1) gives $N'=N'_p$ or
$N'=\mathcal A'_D$.
\end{proof}

\begin{theorem}[Howe duality at the finite places \cite{How79,How89,Wal90,GT16}]\label{th-howe-finite}
Let a triple over $E/F$ with good Schrödinger models be given
(Definitions~\ref{th-theta-hypotheses} and~\ref{th-darboux-model}), let $v$
be a finite place of $F$, let $\varpi_v$ be the representation of
$U(V)(F_v)\times U(W)(F_v)$ on $\mathcal S_v=\mathcal S(F_v^\Lambda)$ defined
by the Weil representation of the first pair, and let $\mathcal R_v$ be its
local augmentation. Then one of the following holds:
\begin{enumerate}
\item the only subspaces of $\mathcal S_v$ that contain $\mathcal R_v$ and are
  stable under every $\varpi_v(1,h)$, $h\in U(W)(F_v)$, are $\mathcal R_v$ and
  $\mathcal S_v$;
\item there is $s\in\mathcal S_v$ such that $\mathcal S_v$ is spanned by
  $\mathcal R_v$ and the vectors $\varpi_v(1,h)s$, $h\in U(W)(F_v)$, and for
  every $\phi\in\mathcal S_v$ there are $h_0,\ldots,h_m\in U(W)(F_v)$ and
  $w_0,\ldots,w_m,c\in\mathbf C$ with $w_i\varpi_v(1,h_i)s=s$ for every $i$ and
  $\sum_iw_i\varpi_v(1,h_i)\phi-c\,s\in\mathcal R_v$.
\end{enumerate}
The same holds for the second pair.
\end{theorem}
\leanref{Development/AutomorphicArithmetic/ThetaCorrespondence/HoweDuality/HoweDualityUnramifiedAllPlaces.lean}{UpperTheory.Lit.HoweDuality.sphTri}
\leanref{Development/WeilBound/ThetaLifts/SphericalLattices/HoweInert.lean}{PositivityTheory.Sph.SphTri}
\leanref{Development/WeilBound/ThetaLifts/SphericalLattices/HoweDuality.lean}{PositivityTheory.Sph.howeInert}
\leanref{Development/WeilBound/ThetaLifts/SphericalSplitting/SphSplitGeo.lean}{PositivityTheory.Sph.sphSplit_geo}

\begin{lemma}\label{th-howe-radical}
Assume the published local inputs of the theta correspondence
(Theorem~\ref{lit-bundle-theta}). Let a triple over $E/F$ with good
Schrödinger models and matched local maps be given, and let $v$ be a finite
place of $F$. Let $\mathcal O_v$ be the set of operators $\varpi_v(1,h)$,
$h\in U(W)(F_v)$, on $\mathcal S_v$, and $K_v=\ker r_v$. Then:
\begin{enumerate}
\item $\mathcal R_v$ and $K_v$ are stable under $\mathcal O_v$,
  $\mathcal R_v\subset K_v\neq\mathcal S_v$, and every subspace stable under
  $\mathcal O_v$ that contains $K_v$ is $K_v$ or $\mathcal S_v$;
\item there is $s\in\mathcal S_v$ such that every subspace stable under
  $\mathcal O_v$ that contains $\mathcal R_v$ and $s$ is $\mathcal S_v$, and
  for all $k_1,\ldots,k_l\in K_v$ there is $e$ in the algebra generated by
  $\mathcal O_v$ with $e(s)=s$ and $e(k_j)\in\mathcal R_v$ for every $j$.
\end{enumerate}
The same holds for the second pair, with $r'_v$.
\end{lemma}
\leanref{Development/AutomorphicArithmetic/ThetaCorrespondence/HoweDuality/KerV.lean}{UpperTheory.Lit.HoweDuality.kerV}
\leanref{Development/AutomorphicArithmetic/ThetaCorrespondence/HoweDuality/KerVp.lean}{UpperTheory.Lit.HoweDuality.kerVp}
\leanref{Development/AutomorphicArithmetic/ThetaCorrespondence/HoweDuality/HoweDualityUnramifiedAllPlaces.lean}{UpperTheory.Lit.HoweDuality.locV}
\leanref{Development/AutomorphicArithmetic/ThetaCorrespondence/HoweDuality/HoweDualityUnramifiedAllPlaces.lean}{UpperTheory.Lit.HoweDuality.locVp}
\leanref{Development/AutomorphicArithmetic/ThetaCorrespondence/HoweDuality/RadSph.lean}{UpperTheory.Lit.HoweDuality.radSph}
\leanref{Development/AutomorphicArithmetic/ThetaCorrespondence/HoweDuality/RadIrr.lean}{UpperTheory.Lit.HoweDuality.radIrr}
\leanref{Development/AutomorphicArithmetic/ThetaCorrespondence/HoweDuality/RelStable.lean}{UpperTheory.Lit.HoweDuality.relStable}
\leanref{Development/AutomorphicArithmetic/ThetaCorrespondence/HoweDuality/RadFin.lean}{UpperTheory.Lit.HoweDuality.radFin}
\begin{proof}
For (1), $\mathcal R_v$ is stable since $\varpi_v(1,h)$ commutes with
$\varpi_v(g,1)$. The map $r_v$ vanishes on $\mathcal R_v$. If $r_v(\phi)=0$,
then $r_v(\phi)=r'_v(0)$, so
$r_v(\varpi_v(1,h)\phi)=r'_v(\varpi'_v(1,h)0)=0$ by the equivariance of
matched local maps; so $K_v$ is stable. Since $r_v$ is onto and
$\Pi_v\neq0$, $K_v\neq\mathcal S_v$. If $N\supset K_v$ is stable, then
$r_v(N)$ is stable under the operators induced on $\Pi_v$, so it is $0$ or
$\Pi_v$; in the first case $N=K_v$, and in the second, for $\phi\in\mathcal
S_v$ there is $n\in N$ with $\phi-n\in K_v\subset N$, so $N=\mathcal S_v$.

For (2), first let $k\in K_v$ be a single element. By
Theorem~\ref{th-howe-finite} one of its two alternatives holds. In the first,
$K_v=\mathcal R_v$ by (1): $K_v$ is stable, contains $\mathcal R_v$ and is
not $\mathcal S_v$. Then any $s\notin\mathcal R_v$ works, with $e=1$; such
$s$ exists since $K_v\neq\mathcal S_v$. In the second, let $s$ be as there;
the first condition on $s$ holds by assumption. For $k$, take $h_i$, $w_i$
and $c$ with $\sum_iw_i\varpi_v(1,h_i)k-cs\in\mathcal R_v$, and put
$e=(m+1)^{-1}\sum_iw_i\varpi_v(1,h_i)$. Then $e(s)=s$, $e(k)\in K_v$ since
$K_v$ is stable, and $e(k)-(m+1)^{-1}cs\in\mathcal R_v\subset K_v$. If
$c\neq0$, then $s\in K_v$, and $K_v=\mathcal S_v$ by the first condition on
$s$, which is false; so $c=0$ and $e(k)\in\mathcal R_v$. For finitely many
$k_1,\ldots,k_l$, choose $e_1$ for $k_1$, then $e_2$ for $e_1(k_2)\in K_v$,
and so on, and put $e=e_l\cdots e_1$. Then $e(s)=s$, and $e(k_j)\in\mathcal
R_v$ for every $j$, since $\mathcal R_v$ is stable under the algebra
generated by $\mathcal O_v$.
\end{proof}

\begin{lemma}\label{th-howe-archimedean}
Assume the published local inputs of the theta correspondence
(Theorem~\ref{lit-bundle-theta}), which include Howe duality at the real places
(Theorem~\ref{lit-howe-duality-real}). Let a triple over $E/F$ with good
Schrödinger models and matched local maps be given. Let $\mathcal O_\infty$ be
the set of operators on $\mathcal P_D$ given by the elements of $K_\infty(W)$
and by the derivatives along one-parameter subgroups of $U(W)$, and let
$\mathcal R_\infty\subset\mathcal P_D$ be the span of the derivatives along
one-parameter subgroups of $U(V)$ and of the differences $f_1-f$ with
$\varpi(k,1)(f\otimes r)=f_1\otimes r$ for every $r$, $k\in K_\infty(V)$.
Call a subspace of $\mathcal P_D$ \emph{admissible} if it is stable under
$\mathcal O_\infty$ and contains $\mathcal R_\infty$. Then $\mathcal P_D$ has
countable dimension; there is $L$ such that every strictly increasing chain of
admissible subspaces has at most $L$ members; $\mathcal R_\infty\subset\ker r$
and $\ker r$ is stable under $\mathcal O_\infty$; and every admissible
subspace other than $\mathcal P_D$ lies in $\ker r$. The same holds for the
second pair, with $r'$.
\end{lemma}
\leanref{Development/AbelianVarietyTheory/Proofs/TwistedOperators/HoweDualityRealPlacesTrivialCharFockLiftStatement.lean}{GlobalTheory.T.howeDuality_realPlaces_trivialChar_fockLift}
\leanref{Development/AbelianVarietyTheory/Proofs/TwistedOperators/HoweDualityRealPlacesTrivialCharFockLiftSmallPaiStatement.lean}{GlobalTheory.T.howeDuality_realPlaces_trivialChar_fockLift_smallPair}
\leanref{Development/AutomorphicArithmetic/ThetaCorrespondence/ThetaPairLifts/LocalTypeDictionary/GlobalHoweTransferRealPlace.lean}{GlobalTheory.T.global_howe_transfer_real_place}
\leanref{Development/AbelianVarietyTheory/Proofs/TwistedOperators/HoweDualityRealPlacesPolynomialVectorsSmallPairStatement.lean}{GlobalTheory.T.howeDuality_realPlaces_polynomialVectors_smallPair}
\begin{proof}
The operators of $\mathcal O_\infty$ and the operators defining
$\mathcal R_\infty$ preserve $\mathcal P_D$, and $\mathcal P_D$ has countable
dimension; these are computations in the Fock model. An admissible subspace is
then a theta subspace of $\mathcal P_D$ in the sense of
Theorem~\ref{lit-howe-duality-real}, whose Fock space is $\mathcal P_D$, so
the bound on chains is Theorem~\ref{lit-howe-duality-real}(1). The map $r$
vanishes on $\mathcal R_\infty$ by the conditions on matched local maps, and
$\ker r$ is stable since $r(f)=r'(0)$ implies that $r$ vanishes on the
derivatives and translates of $f$.

Let $N$ be admissible and different from $\mathcal P_D$, and put
$K=\ker r$. Then $K$ is admissible, and $K\neq\mathcal P_D$ since $r$ is onto
and $\Pi_{\mathbf R}\neq0$. The subspace $r(N)$ of $\Pi_{\mathbf R}$ is stable
under the operators induced by $r$. Indeed, let $f\in\mathcal P_D$ with
$r(f)=r(n_0)$, $n_0\in N$, and let $f_1$ and $n_1$ be the derivatives of $f$
and $n_0$ along a one-parameter subgroup of $U(W)$. Since $r'$ is onto, there
is $f'\in\mathcal P'_D$ with $r'(f')=r(f)$; let $f'_1$ be its derivative. By
the conditions on matched local maps, $r(f_1)=r'(f'_1)=r(n_1)$, and
$n_1\in N$. The same argument applies to the elements of $K_\infty(W)$. So
$r(N)$ is $0$ or $\Pi_{\mathbf R}$. If $r(N)=\Pi_{\mathbf R}$, then every
$f\in\mathcal P_D$ is $n+(f-n)$ with $n\in N$ and $f-n\in K$, so
$N+K=\mathcal P_D$, which contradicts
Theorem~\ref{lit-howe-duality-real}(2). So $r(N)=0$ and $N\subset K$. The
second pair is treated in the same way.
\end{proof}

\begin{lemma}\label{th-howe-tensor}
Let $M_0,M_1,\ldots,M_k$ be complex vector spaces of countable dimension. For
$0\leq i\leq k$ let $\mathcal O_i$ be a set of linear operators on $M_i$,
$R_i\subset M_i$ a subspace and $q_i\colon M_i\to P_i$ a linear map, and call
a subspace of $M_i$ \emph{admissible} if it is stable under $\mathcal O_i$ and
contains $R_i$. Assume:
\begin{enumerate}
\item[(a)] there is $L$ such that every strictly increasing chain of
  admissible subspaces of $M_0$ has at most $L$ members; $R_0\subset\ker q_0$;
  $\ker q_0$ is stable under $\mathcal O_0$; and every admissible subspace of
  $M_0$ other than $M_0$ lies in $\ker q_0$;
\item[(b)] for $1\leq i\leq k$, $R_i$ and $\ker q_i$ are stable under
  $\mathcal O_i$; every subspace stable under $\mathcal O_i$ that contains
  $\ker q_i$ is $\ker q_i$ or $M_i$; and there is $s_i\in M_i$ such that every
  admissible subspace containing $s_i$ is $M_i$, and for every
  $x\in\ker q_i$ there is $e$ in the algebra generated by $\mathcal O_i$ with
  $e(s_i)=s_i$ and $e(x)\in R_i$.
\end{enumerate}
Put $M=M_0\otimes M_1\otimes\cdots\otimes M_k$, and let $N\subset M$ be a
subspace that contains every $x_0\otimes\cdots\otimes x_k$ with some
$x_i\in R_i$, and is stable under the operators that apply an element of
$\mathcal O_i$ to the $i$-th factor. Then $N=M$ or
$N\subset\ker(q_0\otimes q_1\otimes\cdots\otimes q_k)$.
\end{lemma}
\leanref{Development/AutomorphicArithmetic/ThetaCorrespondence/HoweDuality/Imports.lean}{UpperTheory.Lit.HoweDuality.tens}
\leanref{Development/AbelianVarietyTheory/Proofs/TwistedOperators/TensorEqTopOrLeKerFiniteLengthModulesStatement.lean}{GlobalTheory.T.tensor_eq_top_or_le_ker_finiteLength_modules}
\leanref{Development/AutomorphicArithmetic/ThetaCorrespondence/HoweDuality/NakAll.lean}{UpperTheory.Lit.HoweDuality.nakAll}
\leanref{Development/AutomorphicArithmetic/ThetaCorrespondence/HoweDuality/Nak.lean}{UpperTheory.Lit.HoweDuality.nak}
\leanref{Development/AbelianVarietyTheory/Proofs/TwistedOperators/SimpleOverStableSubmoduleTensorPureKernelSpanStatement.lean}{GlobalTheory.T.simpleOverStableSubmodule_tensorPureKernelSpan}
\leanref{Development/AbelianVarietyTheory/Proofs/TwistedOperators/SimpleOverStableSubmodulePiTensorStatement.lean}{GlobalTheory.T.simpleOverStableSubmodule_piTensor}
\begin{proof}
For $i\geq0$ let $J_i\subset M$ be the span of the pure tensors whose $i$-th
factor lies in $\ker q_i$, and put $N_1=N+J_1+\cdots+J_k$. Each $J_i$ is
stable under the operators, since $\ker q_i$ is. Consider the data with $R_i$
replaced by $\ker q_i$ for $i\geq1$. For these, every strictly increasing
chain of admissible subspaces of $M_i$ has at most two members, by the
condition on subspaces containing $\ker q_i$, and $N_1$ satisfies the
hypotheses on $N$. Since $N\subset N_1$, it suffices that $N_1=M$ or
$N_1\subset\ker(q_0\otimes\cdots\otimes q_k)$, and that $N_1=M$ implies
$N=M$.

If some $q_i$ is $0$ there is nothing to prove, so let every $q_i$ be
nonzero, and let $N_1\neq M$. For every $i$, every subspace of $M_i$ stable
under $\mathcal O_i$ and containing $\ker q_i$ is $\ker q_i$ or $M_i$; for
$i=0$ because a proper one lies in $\ker q_0$ by (a). By the bound on chains
there are finite sets $F_i\subset M_i$ such that every admissible subspace
containing $F_i$ is $M_i$, and then every subspace of $M$ that satisfies the
hypotheses on $N$ and contains the pure tensors of elements of the $F_i$ is
$M$. So, by Zorn's lemma, $N_1$ lies in a subspace $N_2\neq M$ that is stable
under the operators and maximal among such subspaces, and the quotient
$Q=M/N_2$ has no stable subspaces other than $0$ and $Q$.

The subspace $N_2$ contains $J_0$. Fix pure tensors in the factors
$1,\ldots,k$, and let $\varphi\colon M_0\to Q$ send $x$ to the class of $x$
tensored with them. It vanishes on $R_0$, intertwines $\mathcal O_0$ with the
operators it induces on $Q$, and these commute with the operators of the other
factors. If $\varphi=0$ there is nothing to prove. Otherwise let
$T\subset\varphi(M_0)$ be a nonzero subspace stable under
$\mathcal O_0$ and minimal among such subspaces; it exists by the bound on
chains, applied to inverse images. The images of such minimal subspaces under
the operators of the other factors are again minimal, so the sum of all
minimal $\mathcal O_0$-stable subspaces of $Q$ is stable under all operators;
it is not $0$, so it is $Q$. Hence every $\mathcal O_0$-stable subspace of $Q$
has an $\mathcal O_0$-stable complement. Let $C$ be one for
$\varphi(\ker q_0)$. Then $\varphi^{-1}(C)+\ker q_0=M_0$, and
$\varphi^{-1}(C)$ is admissible; if it were not $M_0$, it would lie in
$\ker q_0$ by (a), and then $\ker q_0=M_0$, which is false. So
$\varphi(\ker q_0)\subset C\cap\varphi(\ker q_0)=0$. The subspace $N_2$
contains $J_1,\ldots,J_k$, since $N_1$ does.

The quotient $M/(J_0+\cdots+J_k)$ is the tensor product of the spaces
$M_i/\ker q_i$, each of which has countable dimension and no subspaces stable
under $\mathcal O_i$ other than $0$ and itself. Such a tensor product has no
stable subspaces other than $0$ and itself: by
Lemma~\ref{th-first-lift-dixmier} every operator commuting with the
$\mathcal O_i$ on $M_i/\ker q_i$ is a scalar, so by the Jacobson density
theorem \cite[\S11]{Lam01} the algebra generated by $\mathcal O_i$ agrees with
any given operator on any finite set, and a nonzero stable subspace of the
tensor product contains a nonzero pure tensor and then all of them. So
$N_2=J_0+\cdots+J_k$, which lies in $\ker(q_0\otimes\cdots\otimes q_k)$.

Now let $N_1=M$. Remove the $J_i$ one at a time: let $S\subset\{1,\ldots,k\}$
and $i\notin S$ with $N+\sum_{l\in S\cup\{i\}}J_l=M$, and put
$N'=N+\sum_{l\in S}J_l$. It contains the pure tensors with a factor in $R_i$
in the $i$-th place and is stable under the operators of the $i$-th factor.
Let $y$ be a pure tensor with $s_i$ in the $i$-th place, and write $y=n+u$
with $n\in N'$ and $u\in J_i$, a sum of pure tensors whose $i$-th factors are
$x_1,\ldots,x_l\in\ker q_i$. By (b), applied successively as in the proof of
Lemma~\ref{th-howe-radical}(2), there is $e$ in the algebra generated by
$\mathcal O_i$ with $e(s_i)=s_i$ and $e(x_j)\in R_i$ for every $j$. Applying
$e$ in the $i$-th factor fixes $y$, maps $n$ into $N'$, and maps $u$ into the
span of the pure tensors with a factor in $R_i$ in the $i$-th place, which
lies in $N'$. So $y\in N'$. For fixed factors in the other places, the set of
$x\in M_i$ such that the pure tensor with $x$ in the $i$-th place lies in $N'$
is admissible and contains $s_i$, so it is $M_i$. Hence $N'=M$, and after $k$
steps $N=M$.
\end{proof}

\begin{lemma}\label{th-kernel-howe}
Assume the published local inputs of the theta correspondence
(Theorem~\ref{lit-bundle-theta}), which include Howe duality at the real places
(Theorem~\ref{lit-howe-duality-real}). Let a triple over $E/F$ with good
Schrödinger models and matched local maps be given. If $N$ is a theta subspace
different from $\mathcal A_D$, then the image of every element of $N$ under the
matched local maps is $0$. If $N'$ is a theta subspace of the second kind
different from $\mathcal A'_D$, the same holds for the elements of $N'$.
\end{lemma}
\leanref{Development/AutomorphicArithmetic/ThetaCorrespondence/ThetaPairLifts/LocalTypeDictionary/HoweSwap.lean}{UpperTheory.Lit.HoweDuality.twisted_theta_kernel_condition}
\leanref{Development/AutomorphicArithmetic/ThetaCorrespondence/ThetaPairLifts/LocalTypeDictionary/HoweSwap.lean}{UpperTheory.Lit.HoweDuality.twisted_theta_jet_condition}
\leanref{Development/AutomorphicArithmetic/ThetaCorrespondence/ThetaPairLifts/LocalTypeDictionary/HoweSwap.lean}{UpperTheory.Lit.HoweDuality.allV_alt}
\leanref{Development/AutomorphicArithmetic/ThetaCorrespondence/ThetaPairLifts/LocalTypeDictionary/HoweSwap.lean}{UpperTheory.Lit.HoweDuality.stageV_alt}
\leanref{Development/AutomorphicArithmetic/ThetaCorrespondence/HoweDuality/HoweDualityUnramifiedAllPlaces.lean}{UpperTheory.Lit.HoweDuality.allVp}
\leanref{Development/AutomorphicArithmetic/ThetaCorrespondence/HoweDuality/HoweDualityUnramifiedAllPlaces.lean}{UpperTheory.Lit.HoweDuality.stageVp}
\leanref{Development/AbelianVarietyTheory/Proofs/TwistedOperators/ExistsLinearTensorPolyVacuumPureTensorsStatement.lean}{GlobalTheory.T.exists_linear_tensor_polyVacuum_pureTensors}
\leanref{Development/AutomorphicArithmetic/ThetaCorrespondence/ThetaPairLifts/LocalTypeDictionary/GlobalHoweTensorProductTransfer.lean}{GlobalTheory.T.global_howe_tensor_product_transfer}
\leanref{Development/AbelianVarietyTheory/Proofs/TwistedOperators/UnitaryDualPairGlobalToLocalTensorTransferStatement.lean}{GlobalTheory.T.unitaryDualPair_global_to_local_tensor_transfer}
\begin{proof}
For a finite set $T\supset T_X$ of finite places let $\mathcal A_D(T)$ be the
span of the pure tensors in $\mathcal A_D$ whose components outside $T$ are the
vacua. The linear map
\[
\mathcal L_T\colon\mathcal P_D\otimes\bigotimes_{v\in T}\mathcal S_v\to\mathcal A_D(T),\qquad
f\otimes\bigotimes_{v\in T}\phi_v\mapsto f\otimes\bigotimes_{v\in T}\phi_v\otimes\bigotimes_{v\notin T}\phi^0_v,
\]
is onto, and an element of the source has image under
$r\otimes\bigotimes_{v\in T}r_v$ equal to the image of the expression over $T$
that it defines (Definition~\ref{th-kernel-theta-subspace}(4)).

First, for every such $T$, either $\mathcal A_D(T)\subset N$, or every
element of $N\cap\mathcal A_D(T)$ has image $0$. Apply
Lemma~\ref{th-howe-tensor} with $M_0=\mathcal P_D$, $\mathcal O_0$,
$R_0=\mathcal R_\infty$ and $q_0=r$ as in Lemma~\ref{th-howe-archimedean},
and, for $v\in T$, $M_v=\mathcal S_v$, $\mathcal O_v$, $R_v=\mathcal R_v$
and $q_v=r_v$ as in Lemma~\ref{th-howe-radical}; the spaces $\mathcal S_v$
have countable dimension (Lemma~\ref{th-darboux-reading}(9)), and the
hypotheses (a) and (b) are Lemmas~\ref{th-howe-archimedean}
and~\ref{th-howe-radical}. The subspace $\mathcal L_T^{-1}(N)$ satisfies the
hypotheses on $N$ in Lemma~\ref{th-howe-tensor}: since $N$ is a theta
subspace, it contains the images under $\mathcal L_T$ of the pure tensors
with a factor in $\mathcal R_\infty$ or in some $\mathcal R_v$, which are
relations of the kind in condition (b) of a theta subspace, and it is stable
under the operators of $K_\infty(W)$, of the derivatives along one-parameter
subgroups of $U(W)$, and of the $\varpi_v(1,h)$, $h\in U(W)(F_v)$, which act
on the images of $\mathcal L_T$ factorwise. So $\mathcal L_T^{-1}(N)$ is the
whole space, and then $\mathcal A_D(T)\subset N$, or it lies in the kernel of
$r\otimes\bigotimes_{v\in T}r_v$, and then every element of
$N\cap\mathcal A_D(T)$, written as $\mathcal L_T(t)$, has an expression over
$T$ with image $0$.

Now let $\Phi_D\in N$, and suppose that its image is not $0$. Choose
$T\supset T_X$ with $\Phi_D\in\mathcal A_D(T)$. Let
$\Psi_D\in\mathcal A_D$, and choose $T'\supset T$ with
$\Psi_D\in\mathcal A_D(T')$. Since $\Phi_D\in N\cap\mathcal A_D(T')$ and its
image is not $0$, the first step for $T'$ gives
$\mathcal A_D(T')\subset N$, so $\Psi_D\in N$. Thus $N=\mathcal A_D$, which is
false.

For theta subspaces of the second kind the argument is the same, with the
second pair in Lemmas~\ref{th-howe-archimedean} and~\ref{th-howe-radical}:
the pure tensors with a factor in $\mathcal R'_v$, and those defined by the
relations at the archimedean places, lie in $N'$ by
Definition~\ref{th-kernel-theta-subspace}(2).
\end{proof}

\begin{lemma}\label{th-kernel-graph}
Let $\Pi\neq0$ be a complex vector space of countable dimension with a family
$(\sigma_i)_{i\in I}$ of linear operators such that $0$ and $\Pi$ are the only
subspaces stable under every $\sigma_i$. Let $S\subset\Pi\times\Pi$ be a
subspace stable under every $(\sigma_i,\sigma_i)$ such that $(x,0)\in S$
implies $x=0$, $(0,y)\in S$ implies $y=0$, and for every $x\in\Pi$ there is
$y$ with $(x,y)\in S$. Then there is $c\neq0$ such that $y=cx$ for every
$(x,y)\in S$.
\end{lemma}
\leanref{Development/AbelianVarietyTheory/Proofs/TwistedOperators/EqSmulOfStableGraph.lean}{GlobalTheory.SquareAgree.eq_smul_of_stable_graph}
\begin{proof}
For $x\in\Pi$ let $t(x)$ be an element with $(x,t(x))\in S$. It is unique: if
$(x,y)\in S$, then $(0,t(x)-y)\in S$, so $y=t(x)$. So $S$ is the graph of the
map $t$, which is linear since $S$ is a subspace, and $t\sigma_i=\sigma_it$
since $S$ is stable. Since $\Pi$ has countable dimension, there is
$c\in\mathbf C$ with $t-c$ not invertible, by Dixmier's argument in
Lemma~\ref{th-first-lift-dixmier}. The kernel and the image of $t-c$ are stable
under every $\sigma_i$; if the kernel were $0$ and the image $\Pi$, $t-c$ would
be invertible. So the kernel is $\Pi$ or the image is $0$, and in both cases
$t=c$. If $c=0$, then $(x,0)\in S$ for every $x$, so $\Pi=0$, which is false.
\end{proof}

\begin{lemma}\label{th-assemble-matching}
Assume the published local inputs of the theta correspondence
(Theorem~\ref{lit-bundle-theta}), which include Howe duality at the real places
(Theorem~\ref{lit-howe-duality-real}). Let $D$, $\psi$ and the archimedean
embeddings be as in Definition~\ref{th-first-lift-matching}, let
$(A,A',W;\alpha,\beta)$ be a square triple with archimedean data, and let
$\omega$ and $\omega'$ be Weil data for $(-A,W)$ and $(A',W)$ with characters
$(\beta,\beta/\alpha)$ and $(\alpha,\alpha/\beta)$. Let a triple over $E/F$ be
given with $\psi$, the characters $\alpha$ and $\beta$, and the Gram matrices
and Cartan involutions of Lemma~\ref{th-darboux-reading}(2), with good
Schrödinger models, and with Weil representations $\varpi$ and $\varpi'$ such
that $\Theta(\varpi(\iota(g),\kappa(h))\Phi_D)=\omega(g,h)\Theta(\Phi_D)$ for
every Schwartz-Bruhat function $\Phi_D$, and the same for $\varpi'$ and
$\omega'$. Let matched local maps $(\Pi_{\mathbf R},\Pi_v,r,r',r_v,r'_v)$ for
it be given (Definition~\ref{th-kernel-theta-subspace}(3)), and let
$\Pi=\Pi_{\mathbf R}\otimes\bigotimes'_v\Pi_v$, the restricted tensor product
being taken with respect to the vectors $r_v(\phi^0_v)$. Then:
\begin{enumerate}
\item there are linear maps $p$ and $p'$ from the functions on
  $\mathbb W(-A,W)$ and on $\mathbb W(A',W)$ to $\Pi$ such that
  \[
  p\Bigl(\Theta\Bigl(\sum_if_i\otimes\bigotimes_v\phi_{i,v}\Bigr)\Bigr)
  =\sum_ir(f_i)\otimes\bigotimes_vr_v(\phi_{i,v})
  \]
  whenever the sum lies in $\mathcal A_D$, the $\phi_{i,v}$ are
  Schwartz-Bruhat functions and $\phi_{i,v}=\phi^0_v$ for almost all $v$, and
  the same formula holds for $p'$ with $r'$, $r'_v$ and $\mathcal A'_D$;
\item there are linear operators $\varrho(h)$, $h\in U(W)(\mathbf A_f)$, and
  $\varrho_\tau(Y)$, $\tau\in\mathcal T$, $Y\in\mathfrak u_\tau(W)$, on $\Pi$
  such that $(\omega,\omega',\Pi,\varrho,p,p')$, with $p$ and $p'$ restricted to
  $\mathcal A(-A,W)$ and $\mathcal A(A',W)$, is a matching of
  $(A,A',W;\alpha,\beta)$ (Definition~\ref{th-first-lift-matching}(4)).
\end{enumerate}
\end{lemma}
\leanref{Development/AutomorphicArithmetic/ThetaCorrespondence/ThetaPairLifts/LocalTypeDictionary/ExistsMatchingTripleForThetaLift.lean}{GlobalTheory.T.exists_matching_triple_for_theta_lift}
\leanref{Development/AbelianVarietyTheory/Proofs/TwistedOperators/IsReading.lean}{GlobalTheory.TripleRead.IsReading}
\leanref{Development/AbelianVarietyTheory/Proofs/TwistedOperators/IsReading.lean}{GlobalTheory.TripleRead.IsTransp}
\leanref{Development/AbelianVarietyTheory/Proofs/TwistedOperators/ExistsRtensOp.lean}{GlobalTheory.TripleAssembly.PiT}
\leanref{Development/AbelianVarietyTheory/Proofs/TwistedOperators/ExistsRtensOp.lean}{GlobalTheory.TripleAssembly.mapA}
\leanref{Development/AbelianVarietyTheory/Proofs/TwistedOperators/ExistsRtensOp.lean}{GlobalTheory.TripleAssembly.mapB}
\leanref{Development/AutomorphicArithmetic/ThetaCorrespondence/ThetaPairLifts/LocalTypeDictionary/FirstFactorMapThetaSumPureTensor.lean}{GlobalTheory.T.tripleProduct_image_thetaSum_pure}
\leanref{Development/AutomorphicArithmetic/ThetaCorrespondence/ThetaPairLifts/LocalTypeDictionary/TripleProductSecondMapThetaPure.lean}{GlobalTheory.T.tripleProduct_secondMap_theta_pure}
\leanref{Development/AbelianVarietyTheory/Proofs/TwistedOperators/FockAlgebraicThetaVectorsStableFiniteAdelesLieStatement.lean}{GlobalTheory.T.fockAlgebraicThetaVectors_stable_finiteAdeles_lie}
\leanref{Development/AutomorphicArithmetic/ThetaCorrespondence/ThetaPairLifts/LocalTypeDictionary/TripleProductAlgebraWeilRepRangeEqTop.lean}{GlobalTheory.T.triple_algebra_map_surjective}
\leanref{Development/AutomorphicArithmetic/ThetaCorrespondence/ThetaPairLifts/LocalTypeDictionary/TripleMapCompInverseFinite.lean}{GlobalTheory.T.triple_assembly_inverse_finite}
\leanref{Development/AutomorphicArithmetic/ThetaCorrespondence/ThetaPairLifts/LocalTypeDictionary/ThetaTripleMapsLieActionEqZero.lean}{GlobalTheory.T.theta_triple_lift_lie_invariant}
\leanref{Development/AutomorphicArithmetic/ThetaCorrespondence/ThetaPairLifts/LocalTypeDictionary/ThetaTripleMapsEquivariantOnFiniteVectors.lean}{GlobalTheory.T.theta_triple_maps_equivariant}
\leanref{Development/AutomorphicArithmetic/ThetaCorrespondence/ThetaPairLifts/LocalTypeDictionary/ThetaTripleMapsLieEquivariant.lean}{GlobalTheory.T.theta_triple_lift_lie_equivariant}
\leanref{Development/AutomorphicArithmetic/ThetaCorrespondence/ThetaPairLifts/LocalTypeDictionary/ThetaTripleIrreducibleAtRealPlace.lean}{GlobalTheory.T.theta_triple_lift_irreducible}
\leanref{Development/AbelianVarietyTheory/Proofs/TwistedOperators/RestrictedTensorProductIrreducibleCountableDimenStatement.lean}{GlobalTheory.T.restrictedTensorProduct_irreducible_countable_dimension}
\leanref{Development/AbelianVarietyTheory/Proofs/TwistedOperators/TensorIrreducibleOfCountableDimensionIrreducibleStatement.lean}{GlobalTheory.T.tensor_irreducible_of_countable_dimension_irreducible}
\leanref{Development/AutomorphicArithmetic/ThetaCorrespondence/ThetaPairLifts/LocalTypeDictionary/PlacewisePureTensorFinitePlaces.lean}{GlobalTheory.T.theta_pure_tensor_placewise_finite}
\leanref{Development/AbelianVarietyTheory/Proofs/TwistedOperators/ExistsLocalOperatorThetaQuotientFinitePlaceStatement.lean}{GlobalTheory.T.exists_localOperator_thetaQuotient_finitePlace}
\leanref{Development/AutomorphicArithmetic/ThetaCorrespondence/ThetaPairLifts/LocalTypeDictionary/KFiniteSubspaceKStableIffLieStable.lean}{GlobalTheory.T.kFinite_subspace_kStable_iff_lie_stable}
\begin{proof}
The vectors $r_v(\phi^0_v)$ and $r'_v(\phi'^0_v)$ are equal for
$v\notin T_X$, so the restricted tensor products formed with either are the
same. For (1), $\Theta$ is linear and injective on the Schwartz-Bruhat
functions (Lemma~\ref{th-darboux-reading}(5)), and the algebraic functions are
Schwartz-Bruhat functions. The multilinear map
$(f,(\phi_v)_v)\mapsto r(f)\otimes\bigotimes_vr_v(\phi_v)$ gives a linear map
$q$ on the span of the pure tensors $f\otimes\bigotimes_v\phi_v$ with
$\phi_v=\phi^0_v$ for almost all $v$, and $p$ is a linear map with
$p(\Theta(\Phi_D))=q(\Phi_D)$ for $\Phi_D\in\mathcal A_D$; the formula follows
by linearity. The same construction gives $p'$.

For (2), let $\varrho(h)$ be the operator on $\Pi$ with
$\varrho(h)p(\Phi)=p(\omega(1,h)\Phi)$ for $\Phi\in\mathcal A(-A,W)$, and
$\varrho_\tau(Y)$ the operator with $\varrho_\tau(Y)p(\Phi)=p(\Phi_1)$ for
$\Phi_1$ the derivative at $s=0$ of $\omega(1,\exp_\tau(sY))\Phi$; they are
described below and exist by that description. The conditions of a matching
hold.

Condition (a): $\mathcal A(-A,W)$ is stable under $\omega(g,1)$ and
$\omega(1,h)$ for finite adelic $g$ and $h$, and contains the derivatives
along the subgroups $\exp_\tau(sY)$ of both groups; this is a computation in
the Fock model, and the same holds for $\mathcal A(A',W)$.

Condition (b): $p(\mathcal A(-A,W))=\Pi$. By
Lemma~\ref{th-darboux-reading}(7), $\mathcal A(-A,W)=\Theta(\mathcal A_D)$, so
by (1) the image of $p$ is spanned by the vectors
$r(f)\otimes\bigotimes_vr_v(\phi_v)$ with $f\in\mathcal P_D$ and
$\phi_v=\phi^0_v$ for almost all $v$. Since $r$ maps $\mathcal P_D$ onto
$\Pi_{\mathbf R}$ and each $r_v$ is onto, these vectors span $\Pi$. The same
holds for $p'$.

Condition (c): let $g\in U(A)(\mathbf A_f)$ and $\Phi=\Theta(\Phi_D)$ with
$\Phi_D$ a pure tensor. Then
$\omega(g,1)\Phi=\Theta(\varpi(\iota(g),1)\Phi_D)$, and
$\varpi(\iota(g),1)\Phi_D$ is the pure tensor whose component at each finite
place $v$ is $\varpi_v(g_v,1)\phi_v$, which is $\phi^0_v$ at almost all $v$.
Since $r_v$ vanishes on the local augmentation,
$r_v(\varpi_v(g_v,1)\phi_v)=r_v(\phi_v)$, so
$p(\omega(g,1)\Phi)=p(\Phi)$. A derivative along $\exp_\tau(sY)$ in $U(A)$
is, under $\Theta$, a derivative of the archimedean component along a
one-parameter subgroup of $U(V)$, which $r$ maps to $0$; so $p$ vanishes on
these derivatives. The same holds for $p'$.

Condition (d): let $h\in U(W)(\mathbf A_f)$. On pure tensors,
$\varpi(1,\kappa(h))$ acts placewise at the finitely many finite places where
$h$ does not fix the vacuum, and for each finite $v$ there is an operator
$a_v$ on $\Pi_v$ with $r_v(\varpi_v(1,h_v)\phi)=a_v(r_v(\phi))$ and
$r'_v(\varpi'_v(1,h_v)\phi')=a_v(r'_v(\phi'))$, by the equivariance of matched
local maps and since $r_v$ and $r'_v$ are onto; $a_v$ fixes $r_v(\phi^0_v)$
at almost all $v$. Let $\varrho(h)$ act on $\Pi$ by $1$ on $\Pi_{\mathbf R}$
and by $a_v$ on $\Pi_v$. Then $p(\omega(1,h)\Phi)=\varrho(h)p(\Phi)$ and
$p'(\omega'(1,h)\Phi')=\varrho(h)p'(\Phi')$, by (1). In the same way, a
derivative along $\exp_\tau(sY)$ in $U(W)$ is a derivative of the archimedean
component along a one-parameter subgroup of $U(W_\xi)$, and the operator $b$
on $\Pi_{\mathbf R}$ that it induces through $r$, which is also the one
induced through $r'$, gives $\varrho_\tau(Y)=b\otimes1$.

Irreducibility: $\Pi\neq0$, since $\Pi_{\mathbf R}$ and the $\Pi_v$ are
nonzero and $r_v(\phi^0_v)\neq0$ for $v\notin T_X$. The spaces
$\Pi_{\mathbf R}$ and $\Pi_v$ have countable dimension, as images of
$\mathcal P_D$ and of $\mathcal S_v$. Each $\Pi_v$ has no subspaces other than
$0$ and $\Pi_v$ stable under the operators $a_v$, by the conditions on matched
local maps, so the restricted tensor product $\bigotimes'_v\Pi_v$ has none
stable under the operators $\varrho(h)$. A subspace of $\Pi_{\mathbf R}$ stable
under the operators $b$ for all $\tau$ and $Y$ is stable under the operators
induced by all one-parameter subgroups of $U(W_\xi)$, which are products of the
subgroups $\exp_\tau(sY)$, and then under those induced by $K_\infty(W_\xi)$;
so it is $0$ or $\Pi_{\mathbf R}$ by the conditions on matched local maps. By
the same density argument as in Lemma~\ref{th-howe-tensor}, the tensor product
$\Pi$ has no subspaces other than $0$ and $\Pi$ stable under all $\varrho(h)$
and $\varrho_\tau(Y)$.
\end{proof}

\begin{lemma}\label{th-square-model}
Assume the published local inputs of the theta correspondence
(Theorem~\ref{lit-bundle-theta}). Let $D$, $\psi$ and the archimedean
embeddings be as in Definition~\ref{th-first-lift-matching}, and let
$(A,A',W;\alpha,\beta)$ be a square triple of Hermitian spaces with
archimedean data.
\begin{enumerate}
\item There are a Weil datum $\omega$ for $(-A,W)$ with characters
  $(\beta,\beta/\alpha)$ and a Weil datum $\omega'$ for $(A',W)$ with
  characters $(\alpha,\alpha/\beta)$ (Definition~\ref{th-weil-datum}).
\item Let $\omega$ and $\omega'$ be as in (1), and let Haar measures and
  fundamental domains for $U(A)(F)$ and $U(A')(F)$ be given. There is a triple
  over $E/F$ (Definition~\ref{th-theta-hypotheses}) whose spaces are $A$, the
  space $A'$ with the negated form, and the skew-Hermitian space $\xi_0W$ with
  its Gram matrix transposed, with the characters $\alpha$ and $\beta$, the
  character $\psi$ and the corresponding measures and fundamental domains,
  with the following properties. Let $\Theta$ be the map that sends a
  Schwartz-Bruhat function $\Phi$ on the Schrödinger model of the triple to
  the theta function $u\mapsto\sum_x(\rho(u)\Phi)(x)$, the sum over the
  rational points of the model. The \emph{algebraic functions} of the triple
  are the span of the functions $f_\infty\otimes\bigotimes_v\phi_v$ with
  $f_\infty$ in the archimedean Fock space of the triple and $\phi_v$
  Schwartz-Bruhat functions on the completions, equal to the vacuum of the
  triple at almost all places $v$. Then $\Theta$ carries the Weil
  representations of the triple to $\omega$ and $\omega'$; the algebraic
  functions form a space of countable dimension, and $\Theta$ maps them onto
  $\mathcal A(-A,W)$ and onto $\mathcal A(A',W)$; and the theta integrals $I$
  and $I'$ of the triple satisfy $I_\omega(\Theta(\Phi))(h)=I(\Phi)(\kappa(h))$
  and $I_{\omega'}(\Theta(\Phi'))(h)=I'(\Phi')(\kappa(h))$, where
  $I_\omega$ and $I_{\omega'}$ are the theta integrals for the given measures
  and domains (Definition~\ref{th-first-lift-matching}(5)) and $\kappa$
  identifies $U(W)(\mathbf A)$ with the unitary group of $\xi_0W$.
\item Let the triple be as in (2), let local maps be matched to it, and let
  $(\omega,\omega',\Pi,\varrho,p,p')$ be a matching of
  $(A,A',W;\alpha,\beta)$. For $\Phi\in\mathcal A(-A,W)$ and
  $\Phi'\in\mathcal A(A',W)$ let $\Phi_D$ and $\Phi'_D$ be algebraic functions
  with $\Theta(\Phi_D)=\Phi$ and $\Theta(\Phi'_D)=\Phi'$. Let $\mathcal S'$ be
  the span of the functions $f_\infty\otimes\bigotimes_v\phi_v$ as in (2) for
  $A'$ whose component at the auxiliary split place of the triple, read in the
  linear model there, vanishes on the matrices of rank less than
  $\operatorname{rk}A'$.
  \begin{enumerate}
  \item[(a)] $p(\Phi)=0$ if and only if the image of $\Phi_D$ under the
    matched local maps is $0$, and $p'(\Phi')=0$ if and only if the image of
    $\Phi'_D$ is $0$.
  \item[(b)] There is $s\neq0$ such that $(\Phi_D,s\Phi'_D)$ is a pair with
    matching images, in the sense of Definition~\ref{th-theta-hypotheses}(2),
    whenever $p(\Phi)=p'(\Phi')$ and $\Phi'_D\in\mathcal S'$.
  \item[(c)] If the matched local maps describe the coinvariants at the places
    of $F$ that split completely in $E$ and lie outside the finite set of the
    triple, then for every $\Psi\in\mathcal A(A',W)$ there is
    $\Phi_0\in\mathcal A(A',W)$ with $\Psi-\Phi_0\in\mathcal R(\omega')$ and
    $\Phi_0=\Theta(\Phi_{0,D})$ for some algebraic function
    $\Phi_{0,D}\in\mathcal S'$.
  \item[(d)] If $\Phi'_D$ is a Schwartz-Bruhat function whose theta integral
    $I'(\Phi'_D)$ converges and whose theta kernel is bounded by a constant
    times a power of the height of $h$, then the integral defining
    $I_{\omega'}(\Theta(\Phi'_D))(h)$ converges absolutely for every $h$ and
    is bounded by a constant times a power of an adelic height of $h$.
  \end{enumerate}
\item Let $\omega$ and $\omega'$ be as in (1). There is a triple over $E/F$
  with the Gram matrices of the spaces of (2), the characters $\alpha$ and
  $\beta$, the character $\psi$, the Cartan involutions of the archimedean
  data and good Schrödinger models, whose Weil representations $\Theta$
  carries to $\omega$ and $\omega'$ on all Schwartz-Bruhat functions. For every
  such triple, every family of local maps matched to it assembles into a
  matching $(\omega,\omega',\Pi,\varrho,p,p')$ of $(A,A',W;\alpha,\beta)$, with
  $\Pi$, $p$ and $p'$ as in Lemma~\ref{th-assemble-matching}.
\end{enumerate}
\end{lemma}
\leanref{Development/AutomorphicArithmetic/ThetaCorrespondence/ThetaPairLifts/LocalTypeDictionary/WeilRepresentationSquareComm.lean}{GlobalTheory.T.weilRepresentation_square_comm}
\leanref{Development/AbelianVarietyTheory/Statements/HeckeThetaIntegrals/ExistsNeZeroOfIsMatch.lean}{GlobalTheory.SquareSt.Model.exists_weil_of_isTriple}
\leanref{Development/AutomorphicArithmetic/ThetaCorrespondence/ThetaPairLifts/LocalTypeDictionary/ExistsDarbouxCoordinatesTriple.lean}{GlobalTheory.T.exists_darboux_coordinates_triple}
\leanref{Development/AbelianVarietyTheory/Proofs/TwistedOperators/ExistsPlacewiseSeesawTripleOverHermitianStatement.lean}{GlobalTheory.T.exists_placewise_seesaw_triple_over_hermitian}
\leanref{Development/AbelianVarietyTheory/Proofs/TwistedOperators/ThetaEqB.lean}{GlobalTheory.SquareAgree.IsOver}
\leanref{Development/AbelianVarietyTheory/Proofs/TwistedOperators/ThetaEqB.lean}{GlobalTheory.SquareAgree.IsDatum}
\leanref{Development/AbelianVarietyTheory/Proofs/TwistedOperators/ThetaEqB.lean}{GlobalTheory.SquareAgree.algD}
\leanref{Development/AbelianVarietyTheory/Proofs/TwistedOperators/RankAlgebraicDataSchrodingerModelLeAleph0Statement.lean}{GlobalTheory.T.rank_algebraicData_schrodingerModel_le_aleph0}
\leanref{Development/AbelianVarietyTheory/Proofs/TwistedOperators/ConjUMeasurableEmbedding.lean}{GlobalTheory.SquareAgree.thetaInt_eq_modelI}
\leanref{Development/AutomorphicArithmetic/ThetaCorrespondence/ThetaPairLifts/LocalTypeDictionary/IntertwineOfCompatibleMaps.lean}{GlobalTheory.BotUniqT.bottomLevel_ker_left}
\leanref{Development/AutomorphicArithmetic/ThetaCorrespondence/ThetaPairLifts/LocalTypeDictionary/IntertwineOfCompatibleMaps.lean}{GlobalTheory.BotUniqT.bottomLevel_ker_right}
\leanref{Development/AutomorphicArithmetic/ThetaCorrespondence/ThetaPairLifts/LocalTypeDictionary/ThetaKernelMatchOnOppositeParabolic.lean}{GlobalTheory.T.twisted_theta_kernels_agree_up_to_scalar}
\leanref{Development/AbelianVarietyTheory/Proofs/TwistedOperators/ConjUInvMemGfC.lean}{GlobalTheory.SquareAgree.CoinvPass.coinv_passage}
\leanref{Development/AbelianVarietyTheory/Proofs/TwistedOperators/ThetaIntegralModerateGrowthOfBoundedKernelStatement.lean}{GlobalTheory.T.theta_integral_moderate_growth_of_bounded_kernel}
\leanref{Development/AutomorphicArithmetic/ThetaCorrespondence/ThetaPairLifts/LocalTypeDictionary/ExistsMatchingTripleForThetaLift.lean}{GlobalTheory.T.exists_matching_triple_for_theta_lift}
\begin{proof}
For (1), $\beta$ and $\beta/\alpha$ are trivial on $E^\times$, $\beta$
restricts to $\eta^{\operatorname{rk}A}$ on $\mathbf A_F^\times$, and
$\beta/\alpha$ restricts to
$\eta^{\operatorname{rk}A-\operatorname{rk}A'}=\eta^{\operatorname{rk}W}$, by
conditions (a) and (b) of a square triple and $\eta^2=1$. The same holds for
$\alpha$ and $\alpha/\beta$ with $A'$ in place of $A$. The existence of the
Weil representation of a unitary dual pair \cite{Kud94,HKS96}, one of the
published local inputs, gives $\omega$ and $\omega'$.

For (2), let $\nu,\mathcal F_A$ and $\nu',\mathcal F_{A'}$ be the given Haar
measures and fundamental domains, and let $p$ and $q$ be the numbers $s$ and
$s'$ of condition (c) of a square triple. Take the base data $P$, $v_1$,
$m_\alpha$ and $m_\beta$ of Lemma~\ref{th-darboux-reading}(1), the spaces $V$,
$V'$ and $W_\xi$ of Lemma~\ref{th-darboux-reading}(2) and the bijections $c_A$
and $c_{A'}$ of Lemma~\ref{th-darboux-reading}(3). By
Lemma~\ref{th-darboux-reading}(4) the images under $\iota$ of the given
measures and domains are Haar measures and fundamental domains, and
Lemma~\ref{th-darboux-reading}(5) gives good Schrödinger models of
$(V,W_\xi)$ and $(V',W_\xi)$ with them. Lemma~\ref{th-darboux-reading}(6)
gives Weil representations $\varpi$ and $\varpi'$ of these models with
characters $(\beta,\beta/\alpha)$ and $(\alpha,\alpha/\beta)$ that $\Theta$
carries to $\omega$ and $\omega'$. These data form a triple over $E/F$, with
the characters $\alpha$ and $\beta$ and the finite set of places $P$: the
base data and characters have the properties of
Lemma~\ref{th-darboux-reading}(1), $\alpha$ and $\beta$ are Hecke characters
restricting to $\eta^{a'}$ and $\eta^a$ on $\mathbf A_F^\times$ by condition
(a), and $\alpha$ and $\beta$ are unramified above the places outside $P$ by
the definition of $P$; $a+a'=\operatorname{rk}W$ is even by condition (b);
and condition (c) gives $a'<a$, $2(p+q)+1<a'$ and
$2(p+q)+1<(a-a')/2$, the last since $a'+2(2p+2q+1)<a$. The algebraic
functions are carried onto $\mathcal A(-A,W)$ and $\mathcal A(A',W)$ by
Lemma~\ref{th-darboux-reading}(7), and they have countable dimension by
Lemma~\ref{th-darboux-reading}(9). Finally, by
Lemma~\ref{th-darboux-reading}(6) and (8), for $\Theta(\Phi_D)=\Phi$,
\[
I_\omega(\Phi)(h)=\int_{\iota(\mathcal F_A)}\theta_\varpi(g,\kappa(h);\Phi_D)\,d(\iota_*\nu)(g)=I(\Phi_D)(\kappa(h)),
\]
and the same holds for $A'$.

For (3), let $N_p$ and $N'_p$ be as in Lemma~\ref{th-kernel-matching}, and
let $N_0\subset\mathcal A_D$ and $N'_0\subset\mathcal A'_D$ be the sets of
functions with image $0$ under the matched local maps
(Definition~\ref{th-kernel-theta-subspace}(4)). For (a), $N_0$ is a theta
subspace different from $\mathcal A_D$, by
Lemma~\ref{th-kernel-dictionary}(3). By Lemma~\ref{th-kernel-matching}(1),
$N_p$ is a theta subspace different from $\mathcal A_D$, so every element of
$N_p$ has image $0$ by Lemma~\ref{th-kernel-howe}: $N_p\subset N_0$. By the
maximality in Lemma~\ref{th-kernel-matching}(1), $N_0=N_p$ or
$N_0=\mathcal A_D$, and the second is excluded. So $N_0=N_p$: for
$\Phi=\Theta(\Phi_D)$, $p(\Phi)=0$ if and only if the image of $\Phi_D$ is
$0$. In the same way $N'_p$ is a theta subspace of the second kind different
from $\mathcal A'_D$, so $N'_p\subset N'_0$ by Lemma~\ref{th-kernel-howe};
$N'_0$ is a theta subspace of the second kind different from $\mathcal A'_D$
by Lemma~\ref{th-kernel-dictionary}(3), so $N'_0=N'_p$ by
Lemma~\ref{th-kernel-matching}(2).

For (b), let $S\subset\Pi\times\Pi$ be the set of pairs
$(p(\Theta(\Phi_D)),p'(\Theta(\Phi'_D)))$ with $\Phi_D\in\mathcal A_D$ and
$\Phi'_D\in\mathcal A'_D$ having agreeing images. It is a subspace, by the
linearity in Lemma~\ref{th-kernel-dictionary}(3). It is stable under
$(\varrho(h),\varrho(h))$ for $h\in U(W)(\mathbf A_f)$: if $\Phi_D$ and
$\Phi'_D$ have agreeing images, so do $\varpi(1,\kappa(h))\Phi_D$ and
$\varpi'(1,\kappa(h))\Phi'_D$ (Lemma~\ref{th-kernel-dictionary}(4)); these
lie in $\mathcal A_D$ and $\mathcal A'_D$, since by
Lemma~\ref{th-darboux-reading}(6) $\Theta$ carries them to
$\omega(1,h)\Theta(\Phi_D)$ and $\omega'(1,h)\Theta(\Phi'_D)$, which lie in
$\mathcal A(-A,W)$ and $\mathcal A(A',W)$ by condition (a) of a matching; and
$p(\omega(1,h)\Theta(\Phi_D))=\varrho(h)p(\Theta(\Phi_D))$, and the same for
$p'$, by condition (d). In the same way $S$ is stable under
$(\varrho_\tau(Y),\varrho_\tau(Y))$, by Lemma~\ref{th-kernel-dictionary}(4)
and condition (d) for the derivatives along $\exp_\tau(sY)$. If
$(x,0)\in S$, with $x=p(\Theta(\Phi_D))$ and $0=p'(\Theta(\Phi'_D))$, then
the image of $\Phi'_D$ is $0$ by (a), so by linearity $\Phi_D$ and $0$ have
agreeing images, the image of $\Phi_D$ is $0$, and $x=0$ by (a). In the same
way $(0,y)\in S$ implies $y=0$. For every $x=p(\Theta(\Phi_D))$ there is
$\Phi'_D$ with agreeing images (Lemma~\ref{th-kernel-dictionary}(4)), so
$(x,p'(\Theta(\Phi'_D)))\in S$. The space $\Pi$ has countable dimension:
$\Pi=p(\mathcal A(-A,W))$, and $\mathcal A(-A,W)=\Theta(\mathcal A_D)$ with
$\mathcal A_D$ of countable dimension by (2). By the irreducibility of a
matching and Lemma~\ref{th-kernel-graph}, there is $s\neq0$ with $y=sx$ for
every $(x,y)\in S$.

Now let $p(\Phi)=p'(\Phi')$ with $\Phi'_D\in\mathcal S'$. By
Lemma~\ref{th-kernel-dictionary}(4) there is $\Phi_{1,D}\in\mathcal A'_D$
such that $\Phi_D$ and $\Phi_{1,D}$ have agreeing images; put
$\Phi_1=\Theta(\Phi_{1,D})$. Then $(p(\Phi),p'(\Phi_1))\in S$, so
$p'(\Phi_1)=s\,p(\Phi)=s\,p'(\Phi')$, and $p'(\Phi_1-s\Phi')=0$. By (a), the
image of $\Phi_{1,D}-s\Phi'_D$ is $0$, and by linearity $\Phi_D$ and
$\Phi_{1,D}-(\Phi_{1,D}-s\Phi'_D)=s\Phi'_D$ have agreeing images. Since
$s\Phi'_D\in\mathcal S'$, $(\Phi_D,s\Phi'_D)$ is a pair with matching images
by Lemma~\ref{th-kernel-dictionary}(5).

For (c), the place $v_1$ splits completely in $E$ and lies outside the finite
set of places of the triple. That the matched local maps describe the
coinvariants there means in particular that a Schwartz-Bruhat function
$\phi$ at $v_1$ satisfies $r'_{v_1}(\phi)=0$ if and only if
$\phi\in\mathcal R'_{v_1}$, and that every element of $\Pi_{v_1}$ is
$r'_{v_1}(\phi)$ for a function $\phi$ that, read in the linear model at
$v_1$, vanishes on the matrices of rank less than $\operatorname{rk}A'$. Let
$\mathcal R'_D$ be the span of the functions $\varpi'(g,1)\Psi_0-\Psi_0$ with
$g\in U(V')(F_{v_1})$ and $\Psi_0\in\mathcal A'_D$. Every
$\Psi_D\in\mathcal A'_D$ is $\Phi_{0,D}+\rho$ with
$\Phi_{0,D}\in\mathcal A'_D\cap\mathcal S'$ and $\rho\in\mathcal R'_D$. Since
these $\Psi_D$ form a subspace, it suffices to treat a pure tensor
$\Psi_D=f\otimes\bigotimes_v\phi_v$. Choose $\phi_o$ as above with
$r'_{v_1}(\phi_o)=r'_{v_1}(\phi_{v_1})$; then
$\phi_{v_1}-\phi_o\in\mathcal R'_{v_1}$. Let $\Phi_{0,D}$ be the pure tensor
with $\phi_o$ in place of $\phi_{v_1}$; it lies in $\mathcal S'$. Then
$\Psi_D-\Phi_{0,D}$ is the pure tensor with $\phi_{v_1}-\phi_o$ at $v_1$, and
it lies in $\mathcal R'_D$, since the pure tensor with
$\varpi'_{v_1}(g,1)\phi-\phi$ at $v_1$ is $\varpi'(g,1)\Psi_0-\Psi_0$ for
the pure tensor $\Psi_0$ with $\phi$ at $v_1$.

Let $\Psi\in\mathcal A(A',W)$. By Lemma~\ref{th-darboux-reading}(7),
$\Psi=\Theta(\Psi_D)$ with $\Psi_D\in\mathcal A'_D$; write
$\Psi_D=\Phi_{0,D}+\rho$ as above, and put $\Phi_0=\Theta(\Phi_{0,D})$, which
lies in $\mathcal A(A',W)$. For $g\in U(V')(F_{v_1})$ and
$\Psi_0\in\mathcal A'_D$, $g=\iota(g')$ with $g'\in U(A')(\mathbf A_f)$
(Lemma~\ref{th-darboux-reading}(3)), and
$\Theta(\varpi'(g,1)\Psi_0-\Psi_0)=\omega'(g',1)\Theta(\Psi_0)-\Theta(\Psi_0)$
by Lemma~\ref{th-darboux-reading}(6), which lies in $\mathcal R(\omega')$. So
$\Theta(\rho)\in\mathcal R(\omega')$, and $\Psi-\Phi_0=\Theta(\rho)$ by the
linearity of $\Theta$.

For (d), let $\nu'$ and $\mathcal F_{A'}$ be the measure and domain of $I'$
and $I_{\omega'}$, with $\mu'=\iota_*\nu'$ and $\mathcal F'=\iota(\mathcal
F_{A'})$, and let $C$ and $k$ be such that
$|\theta_{\varpi'}(g,h;\Phi'_D)|\leq C\operatorname{ht}(h)^k$ for all $g$ and
$h$. By Lemma~\ref{th-darboux-reading}(6),
$\theta_{\omega'}(g,h;\Phi')=\theta_{\varpi'}(\iota(g),\kappa(h);\Phi'_D)$ for
$\Phi'=\Theta(\Phi'_D)$. Since $\iota$ is a measurable embedding, the
integrability of $g\mapsto\theta_{\varpi'}(g,\kappa(h);\Phi'_D)$ on
$\mathcal F'$ for $\mu'$, which is the convergence of $I'(\Phi'_D)$ at
$\kappa(h)$, gives that of $g\mapsto\theta_{\omega'}(g,h;\Phi')$ on
$\mathcal F_{A'}$ for $\nu'$. There is $C'$ with
$\operatorname{ht}(\kappa(h))\leq C'\operatorname{ht}(h)$ for the adelic
height of $h$, a computation with the entries of $h$; and $\mathcal F_{A'}$
has finite measure, since $U(A')(F)\backslash U(A')(\mathbf A)$ has finite
volume \cite{BHC62}. Hence
\[
\int_{\mathcal F_{A'}}|\theta_{\omega'}(g,h;\Phi')|\,d\nu'(g)\leq\max(C,0)\,C'^k\,\nu'(\mathcal F_{A'})\operatorname{ht}(h)^k .
\]

For (4), the triple over $E/F$ of Lemma~\ref{th-darboux-reading}(10) has
the stated properties: its Gram matrices, characters, character $\psi$ and
Cartan involutions are those of (2), its Schrödinger models are good, and
$\Theta$ carries its Weil representations to $\omega$ and $\omega'$ on all
Schwartz-Bruhat functions. For every triple with these properties,
Lemma~\ref{th-assemble-matching} shows that every family of local maps matched
to it assembles into a matching $(\omega,\omega',\Pi,\varrho,p,p')$ with
$\Pi$, $p$ and $p'$ as there.
% (2): T.exists_placewise_seesaw_triple_over_hermitian (SquareAgree.exists_isOver), with the
%   theta integrals by SquareAgree.thetaInt_eq_modelI and IsOver.theta_eq_A, theta_eq_B, and
%   countable dimension by T.rank_algebraicData_schrodingerModel_le_aleph0.
% (3)(a): BotUniqT.bottomLevel_ker_left (theta_subspace_kernel_maximal,
%   thetaSubspace_contains_polKernel, Lit.HoweDuality.twisted_theta_kernel_condition) and
%   bottomLevel_ker_right (thetaSubspace_eq_ker_projDerivation_of_square,
%   thetaSubgroup_of_ker_derivation, twisted_theta_jet_condition, thetaKernel_projection_isMax).
% (3)(b): SquareAgree.theta_kernels_agree_up_to_scalar, with eq_smul_of_stable_graph,
%   theta_kernel_match_on_parabolic and theta_kernel_match_on_opposite_parabolic (the two halves of
%   (3)(a)), polMatch_linear, polMatch_onto, match_of_polMatch.
% (3)(c): SquareAgree.CoinvPass.coinv_passage with exists_So, tens_mem_relD, conjUInv_mem_GfC.
% (3)(d): ModGrowth.modInt_of_bound with ht_kapX_le, Mem.VolFin.finVol, measurableEmbedding_conjU.
% (4): T.exists_darboux_coordinates_triple and T.exists_matching_triple_for_theta_lift.
\end{proof}

\subsection{Matchings of bottom triples}\label{subsection-th-bottom-matchings}

\begin{lemma}\label{th-bottom-square-facts}
For all sufficiently large data (Definition~\ref{th-large-data}) the following
holds. Let a grid Weil datum, $0\leq i,j<m$, a bottom triple $(A,A',W)$ for
$(i,j)$ (Definition~\ref{th-bottom-triple}) and a tuple $t$ be given, with the
characters $\alpha_t$ and $\beta$ of the square $(i,j)$, and put $a=a(i,j)$
and $a'=a'(i,j)$.
\begin{enumerate}
\item The integer $a+a'$ is even, $a'+2<a$ and $1<a'$.
\item There are $a_\alpha,a_\beta\colon\mathcal T\to\mathbf Q$, which do not
  depend on $t$, such that $\alpha_t$ and $\beta$ have archimedean exponents
  $a_\alpha$ and $a_\beta$ (Definition~\ref{th-first-lift-matching}(2)) and
  $2(a_\alpha(\tau)-a_\beta(\tau))<-(a+a')$ for every $\tau\in\mathcal T$.
\item For every finite place $v$ of $F$, either $\alpha_t$ and $\beta$ are
  unramified at the places over $v$, or $v$ splits in $E$, $\alpha_t/\beta$ is
  ramified at a place over $v$, and $v$ does not divide the different of $F$.
\item For each place $v$ of $F$ below a place of $S_t$ there is an embedding
  $\tau_v\colon E\to F_v$ over $F$ such that every $w\in S_t$ over $v$ is the
  place of $E$ defined by $\tau_v$. The character $\alpha_t/\beta$ is ramified
  at every $w\in S_t$, and $\alpha_t$ and $\beta$ are unramified at every
  place of $E$ that lies over no place of $F$ below $S_t$. Finally
  $\varepsilon_t=\sum_{w\in S_t}1_w$, where $1_w\in\mathbf A_E$ is the
  idempotent of $w$, and $1_w+\bar{1}_w$ is the idempotent of the place of $F$
  below $w$.
\item For the archimedean data $J_\tau=B_\tau=1$ on $A$, $A'$ and $W$, the
  functions $f_t$ and $f'_t$ are archimedean vacua for $(-A,W)$ and $(A',W)$.
\end{enumerate}
\end{lemma}
\leanref{Development/AbelianVarietyTheory/Proofs/DivisorProofs/BottomRanksArithmeticLargeGrids.lean}{GlobalTheory.BotRankSize.bot_arith}
\leanref{Development/AbelianVarietyTheory/Proofs/DivisorProofs/BottomRanksArithmeticLargeGridsStatement.lean}{GlobalTheory.T.bottomRanks_arithmetic_largeGrids}
\leanref{Development/AbelianVarietyTheory/Proofs/DivisorProofs/ExistsBottomInfinityExponentsGapLargeGridsStatement.lean}{GlobalTheory.T.exists_bottomInfinityExponents_gap_largeGrids}
\leanref{Development/AbelianVarietyTheory/Proofs/DivisorProofs/ThetaGridCharactersArchExponentsGapGtRank.lean}{GlobalTheory.SquareExps.alpha_exp}
\leanref{Development/AbelianVarietyTheory/Proofs/DivisorProofs/ThetaGridCharactersArchExponentsGapGtRank.lean}{GlobalTheory.SquareExps.beta_exp}
\leanref{Development/AbelianVarietyTheory/Proofs/DivisorProofs/ChiGridVertexRankAddLtTwoExponentAddCentre.lean}{GlobalTheory.ChiGrid.rankV_add_lt_two_e_add_two_cen}
\leanref{Development/AbelianVarietyTheory/Proofs/DivisorProofs/DefiniteBottomTripleFinitePlacesUnramOrSplitStatement.lean}{GlobalTheory.T.definiteBottomTriple_finitePlaces_unram_or_split}
\leanref{Development/AbelianVarietyTheory/Proofs/DivisorProofs/ExistsEmbeddingSplitPrimeGridLocalTensorPlacesStatement.lean}{GlobalTheory.T.exists_embedding_splitPrime_gridLocalTensorPlaces}
\leanref{Development/AbelianVarietyTheory/Proofs/DivisorProofs/GridLocalTensorPlacesSplitIdempotentSumStatement.lean}{GlobalTheory.T.gridLocalTensorPlaces_splitIdempotent_sum}
\leanref{Development/AbelianVarietyTheory/Proofs/DivisorProofs/SeesawTripleFirstFiniteDatumStandardIsVacuumDatuStatement.lean}{GlobalTheory.T.seesawTriple_firstFiniteDatum_standard_isVacuumDatum}
\leanref{Development/AbelianVarietyTheory/Proofs/DivisorProofs/SeesawTripleSecondFiniteDatumStandardIsVacuumDatStatement.lean}{GlobalTheory.T.seesawTriple_secondFiniteDatum_standard_isVacuumDatum}
\begin{proof}
For (1), put $r=k+m-1$, which is odd, and recall
$\operatorname{rk}V_{i,j}=r+jt+i(r+6)$ (Definition~\ref{th-grid}). Since $m$
is even, $2\lfloor(m-j)/2\rfloor-2\lfloor(m-j-1)/2\rfloor$ is $2$ if $j$ is
even and $0$ if $j$ is odd. So
$a-a'\geq\operatorname{rk}V_{i+1,j+1}-\operatorname{rk}V_{i,j}=t+r+6>2$.
Next $a'\geq\operatorname{rk}V_{i,j}-(m-j)\geq r-m=k-1>1$, since
$k\geq m+3$. Finally $a+a'$ is congruent modulo $2$ to
$\operatorname{rk}V_{i+1,j+1}+\operatorname{rk}V_{i,j}=2r+(2j+1)t+(2i+1)(r+6)$,
which is congruent to $t+r\equiv0$, since $t$ and $r$ are odd.

Parts (2) to (4) are computations from the construction of the level system
and of the grid, and their results are stated. In (2),
$a_\alpha(\tau)=-e_{j+1}(\tau)+\epsilon/2$ and
$a_\beta(\tau)=c_{m-i}(\tau)+\epsilon/2$, where $\epsilon$ is the parity of
$\operatorname{rk}V_{i,j}$, $e_{j+1}(\tau)$ is the entry of the
Harish-Chandra parameter $\lambda(\tau)$, and $c_{m-i}(\tau)$ is the integer
of Definition~\ref{th-grid}(b), the archimedean exponent of $\chi_{m-i}$ at
$\tau$; the
characters $\mu^0_{j+1}$, $\xi_t$ and $\eta_L$ have exponents
$-e_{j+1}(\tau)$, $0$ and $\tfrac12$, by (L2), (L4) and (L3) of
Definition~\ref{th-grid}: for $\xi_t$, which has finite order by (L4), the
archimedean components are characters of finite order of $\mathbf C^\times$,
hence trivial. The integers $c_{m-i}(\tau)$ satisfy
$\operatorname{rk}V_{i+1,j+1}+\operatorname{rk}V_{i,j}<2e_{j+1}(\tau)+2c_{m-i}(\tau)$,
and $a+a'\leq\operatorname{rk}V_{i+1,j+1}+\operatorname{rk}V_{i,j}$, which gives
the inequality. In (3), the second case occurs at the distinguished places
$S_0$ of the grid and at the places below the places of $Q_t$, which are the
places of $F$ below the ramified places of $\xi_t$ by (L5), and the first case
elsewhere.

For (5), $f_t$ is a theta function for $(-A,W)$, obtained from the Fock vacuum
vector for $(-A,W)$ by replacing its component at the finite idempotent
$\varepsilon_t$. At the archimedean places it is therefore annihilated by the
annihilation operators for the standard data, as the Fock vacuum vector is.
The same argument applies to $f'_t$, with $\bar\varepsilon_t$ in place of
$\varepsilon_t$.
\end{proof}

\begin{lemma}\label{th-bottom-split-reading}
Let $D$, $\psi$ and the archimedean embeddings be as in
Definition~\ref{th-first-lift-matching}, let $(A,A',W;\alpha,\beta)$ be a
square triple of Hermitian spaces with archimedean data, and let $\omega$ and
$\omega'$ be Weil data for $(-A,W)$ and $(A',W)$ with characters
$(\beta,\beta/\alpha)$ and $(\alpha,\alpha/\beta)$
(Definition~\ref{th-weil-datum}). Let $\Lambda_A$, $\Lambda_{A'}$ and
$\Lambda_W$ be $\mathcal O_E$-lattices in $E^{\operatorname{rk}A}$,
$E^{\operatorname{rk}A'}$ and $E^{\operatorname{rk}W}$ that are self-dual for
the Hermitian forms of $A$, $A'$ and $W$. Let $S$ be a finite set of finite
places of $E$ such that:
\begin{enumerate}
\item[(i)] for each place $v$ of $F$ below a place of $S$ there is an
  embedding $\tau\colon E\to F_v$ over $F$ such that every $w\in S$ over $v$ is
  the place of $E$ defined by $\tau$;
\item[(ii)] $\alpha/\beta$ is ramified at every $w\in S$;
\item[(iii)] $\alpha$ and $\beta$ are unramified at every place of $E$ that
  lies over no place of $F$ below $S$.
\end{enumerate}
Let $\Sigma$ be the set of places of $F$ below $S$. Then there is a triple over
$E/F$ (Definition~\ref{th-theta-hypotheses}) with the spaces, the characters
$\alpha$ and $\beta$, the character $\psi$ and the Cartan involutions of
Lemma~\ref{th-square-model}(2), whose Weil representations the theta-sum map
$\Theta$ of Lemma~\ref{th-square-model}(2) carries to $\omega$ and $\omega'$,
such that:
\begin{enumerate}
\item the lattices of its two Hermitian spaces are
  $\bar\Lambda_A=\{\bar y:y\in\Lambda_A\}$ and $\bar\Lambda_{A'}$, and the
  lattice of its skew-Hermitian space contains every $y$ with
  $\mathfrak A^{-1}y\subset\Lambda_W$, where
  $\mathfrak A=\mathfrak c\mathcal O_E\cdot\mathfrak b^{-1}$ is the ideal of the
  host (Definition~\ref{th-grid-fock}), for the ideals $\mathfrak b$ of
  $\mathcal O_E$ and $\mathfrak c$ of $\mathcal O_K$ of the host
  (Definition~\ref{hf-host});
\item its finite set of places is $\Sigma$, and its split bases at each
  $v\in\Sigma$ are taken for an embedding $\tau_v\colon E\to F_v$ over $F$,
  where $S$ is the set of the places of $E$ defined by the $\tau_v$ with
  $v\in\Sigma$.
\end{enumerate}
\end{lemma}
\leanref{Development/AutomorphicArithmetic/LocalEmbeddings/BotReading.lean}{GlobalTheory.BotMatch.exists_reading_lat}
\leanref{Development/AutomorphicArithmetic/LocalEmbeddings/BotReading.lean}{GlobalTheory.BotMatch.exists_isHerm_lat}
\leanref{Development/AutomorphicArithmetic/LocalEmbeddings/BotReading.lean}{GlobalTheory.BotMatch.exists_isSkew_lat}
\begin{proof}
A lattice is self-dual for the form of $A$ if and only if it is self-dual for
the form of $-A$. For each real place of $F$ choose an embedding
$\tau\in\mathcal T$ over it; the base of the triple below uses its complex
conjugate $\bar\tau$. By condition (d) of a square triple, $\alpha$ and $\beta$
have archimedean exponents $x(\tau)$ and $x'(\tau)$ with
$2(x(\tau)-x'(\tau))<-\operatorname{rk}W$. At $\bar\tau$ they therefore have
the integer exponents $m_\alpha=-2x$ and $m_\beta=-2x'$, with
$m_\alpha-m_\beta>\operatorname{rk}W$.

Let $T_0$ be the set of places of $F$ below the finite places of $E$ at which
$\alpha$ or $\beta$ is ramified; it is finite, since Hecke characters are
unramified at almost all places. By condition (e) of a square triple, every
$v\in T_0$ splits in $E$, does not divide the different of $F$, and has a
place of $E$ over it at which $\alpha/\beta$ is ramified. By conditions (a) and
(b), $\alpha/\beta$ is trivial on $\mathbf A_F^\times$, so it is ramified at
both places over $v$ if it is ramified at one. By the definition of $T_0$,
$\alpha$ and $\beta$ are unramified at the places of $E$ over places outside
$T_0$. Then $T_0=\Sigma$. Indeed, let $v\in T_0$ and let $\tau_0\colon E\to F_v$
be an embedding over $F$; if $v$ were not below $S$, then by (iii) $\alpha$ and
$\beta$, and so $\alpha/\beta$, would be unramified at the place defined by
$\tau_0$, which lies over $v$. And if $w\in S$ lay over a place outside $T_0$,
then $\alpha$ and $\beta$, and so $\alpha/\beta$, would be unramified at $w$,
against (ii). For $v\in\Sigma$ let $\tau_v$ be an embedding given by (i). Every
$w\in S$ is the place defined by $\tau_v$ for the place $v$ below $w$, and for
$v\in\Sigma$ the place defined by $\tau_v$ lies in $S$, since it equals each of
the places of $S$ over $v$.

By Dirichlet's theorem there is a prime $p\equiv1\pmod{Q_Sqd}$ larger than the
norms of the places of $\Sigma$, and such a prime splits completely in the host
field (Chapter~\ref{chapter-host-fields}). A place $v_1$ of $F$ over $p$ lies
outside $\Sigma$, splits completely in $E$ and has an embedding
$\tau_1\colon E\to F_{v_1}$ over $F$. The character $\psi$, the embeddings
chosen at the real places, $m_\alpha$, $m_\beta$, the set $\Sigma$ with the
$\tau_v$, and $v_1$ with $\tau_1$ form the base of the triple. Its conditions
hold: $E/F$ is a CM extension, $E$ is Galois over $\mathbf Q$ and unramified
over $F$ at the finite places, $\psi$ is standard, every place of $F$ inert in
$E$ does not lie over $2$, has norm larger than $3$ and has $\psi$ unramified
at it, by the properties of the host, and $\psi$ is unramified at every
$v\in\Sigma$, which does not divide the different of $F$.

The three spaces of the triple are those of Lemma~\ref{th-square-model}(2),
with the following lattices and bases. The space $A$ receives the lattice
$\bar\Lambda_A$, which is self-dual for the form of the triple, the Cartan
involutions of its archimedean data, and at each $v\in\Sigma$ a basis of
$F_v^{\operatorname{rk}A}$ that is split for $\tau_v$ and $\bar\Lambda_A$;
such a basis exists since $\tau_v$ identifies $E\otimes_FF_v$ with
$F_v\times F_v$. Its signature is that of condition (c). The space $A'$ with
the negated form receives $\bar\Lambda_{A'}$ and the negated Cartan involutions
in the same way. The skew-Hermitian space receives a lattice built from
$\Lambda_W$ that is self-dual for its trace form and contains every $y$ with
$\mathfrak A^{-1}y\subset\Lambda_W$, and split bases at the places of $\Sigma$.
% exists_isSkew_lat: the lattice is SkewLat.pairLat for Lambda_W and the zero space.

For these spaces and this base, Lemma~\ref{th-darboux-reading}(5) gives good
Schrödinger models of the two pairs, in rational Darboux coordinates with
local vacua, a Haar measure and a fundamental domain, and
Lemma~\ref{th-darboux-reading}(6) gives Weil representations on them, with
characters $(\beta,\beta/\alpha)$ and $(\alpha,\alpha/\beta)$, that $\Theta$
carries to $\omega$ and $\omega'$ on all Schwartz-Bruhat functions.
% TripleDarboux.exists_model (RatDarboux.exists_rational_darboux_of_isGood,
%   VacIndic.eventually_isVac_indicator, T.exists_finitePlace_vacuum_any_darboux_coordinates,
%   Ker.HypModel.exists_haar_fundDom) and TripleDarboux.exists_weil (exists_swap,
%   exists_transport_isWeil). The same two steps build the triple of
%   T.exists_darboux_coordinates_triple behind Lemma th-square-model(4).

The characters of the triple are $\alpha$ and $\beta$, with finite set of
places $\Sigma$. They are Hecke characters with the restrictions to
$\mathbf A_F^\times$ of condition (a) and the integer exponents $m_\alpha$ and
$m_\beta$; $\alpha/\beta$ is ramified at every place over $\Sigma$, $\alpha$ and
$\beta$ are unramified at the places over places outside $\Sigma$, and $v_1$
lies outside $\Sigma$ and splits completely. The conditions on the ranks,
$\operatorname{rk}A'<\operatorname{rk}A$, $\operatorname{rk}W$ even,
$2(s+s')+1<(\operatorname{rk}A-\operatorname{rk}A')/2$ and
$2(s+s')+1<\operatorname{rk}A'$, follow from conditions (b) and (c) of a square
triple, and $\operatorname{rk}W<m_\alpha-m_\beta$ at every real place.
\end{proof}

\begin{lemma}\label{th-bottom-split-assembly}
Assume the published local inputs of the theta correspondence
(Theorem~\ref{lit-bundle-theta}). Let $D$, $\psi$ and the archimedean
embeddings be as in Definition~\ref{th-first-lift-matching}, let
$(A,A',W;\alpha,\beta)$ be a square triple with archimedean data, and let
$\omega$ and $\omega'$ be Weil data for $(-A,W)$ and $(A',W)$ with characters
$(\beta,\beta/\alpha)$ and $(\alpha,\alpha/\beta)$. Let a triple over $E/F$ be
given with the spaces, the characters $\alpha$ and $\beta$, the character
$\psi$ and the Cartan involutions of Lemma~\ref{th-square-model}(2), whose
Weil representations the theta-sum map $\Theta$ carries to $\omega$ and
$\omega'$, and let $\phi^0_v$ and $\phi'^0_v$ be the local vacua of its two
pairs at the finite places $v$. Let local maps be matched to it: linear maps
$p_\infty$ and $p'_\infty$ from the archimedean Schwartz functions of the two
pairs to a space $\Pi_\infty$, and for every finite place $v$ linear maps $p_v$
and $p'_v$ from their Schwartz-Bruhat functions at $v$ to a space $\Pi_v$. Let
$\Pi=\Pi_\infty\otimes\bigotimes'_v\Pi_v$, the restricted tensor product being
taken with respect to the vectors $p_v(\phi^0_v)$. Then there are linear maps
$p$ from the functions on $\mathbb W(-A,W)$ to $\Pi$ and $p'$ from the
functions on $\mathbb W(A',W)$ to $\Pi$ such that
\[
p\Bigl(\Theta\Bigl(\sum_if_i\otimes\bigotimes_v\phi_{i,v}\Bigr)\Bigr)
=\sum_ip_\infty(f_i)\otimes\bigotimes_vp_v(\phi_{i,v})
\]
for every finite family of functions $f_i$ in the archimedean Fock space of
the first pair and Schwartz-Bruhat functions $\phi_{i,v}$ at the finite places
with $\phi_{i,v}=\phi^0_v$ for almost all $v$; the same formula holds for $p'$,
with $p'_\infty$, $p'_v$ and the second pair; and for some operators $\varrho$
on $\Pi$, $(\omega,\omega',\Pi,\varrho,p,p')$ with $p$ and $p'$ restricted to
$\mathcal A(-A,W)$ and $\mathcal A(A',W)$ is a matching of
$(A,A',W;\alpha,\beta)$ (Definition~\ref{th-first-lift-matching}(4)).
\end{lemma}
\leanref{Development/AutomorphicArithmetic/ThetaCorrespondence/ThetaPairLifts/LocalTypeDictionary/FirstFactorMapThetaSumPureTensor.lean}{GlobalTheory.T.tripleProduct_image_thetaSum_pure}
\leanref{Development/AutomorphicArithmetic/ThetaCorrespondence/ThetaPairLifts/LocalTypeDictionary/TripleProductSecondMapThetaPure.lean}{GlobalTheory.T.tripleProduct_secondMap_theta_pure}
\leanref{Development/AbelianVarietyTheory/Proofs/TwistedOperators/ExistsRtensOp.lean}{GlobalTheory.TripleAssembly.PiT}
\leanref{Development/AbelianVarietyTheory/Proofs/TwistedOperators/ExistsRtensOp.lean}{GlobalTheory.TripleAssembly.mapA}
\leanref{Development/AbelianVarietyTheory/Proofs/TwistedOperators/ExistsRtensOp.lean}{GlobalTheory.TripleAssembly.mapB}
\begin{proof}
The map $\Theta$ is injective on the Schwartz-Bruhat functions on the
Schrödinger model of each pair \cite{Wei64}, and the algebraic functions of the
triple (Lemma~\ref{th-square-model}(2)) are Schwartz-Bruhat functions. Let
$l_\infty$ be a linear form on $\Pi_\infty$ and $l_v$ linear forms on the
$\Pi_v$ with $l_v(p_v(\phi^0_v))=1$ for almost all $v$. The restricted tensor
product of the archimedean Schwartz space and the local Schwartz-Bruhat spaces
with respect to the $\phi^0_v$ embeds in the functions on the adelic
Schrödinger model, so there is a linear form on these functions that takes
$f\otimes\bigotimes_v\phi_v$, with $\phi_v=\phi^0_v$ for almost all $v$, to
$l_\infty(p_\infty(f))\prod_vl_v(p_v(\phi_v))$. The forms
$l_\infty\otimes\bigotimes_vl_v$ separate the points of $\Pi$, so these linear
forms together give a linear map $q$ from the functions on the adelic
Schrödinger model to $\Pi$ with
$q(f\otimes\bigotimes_v\phi_v)=p_\infty(f)\otimes\bigotimes_vp_v(\phi_v)$. As
$\Theta$ is injective on a space containing the
algebraic functions, there is a linear map $p$ on the functions on
$\mathbb W(-A,W)$ with $p(\Theta(\Phi_D))=q(\Phi_D)$ for every algebraic
function $\Phi_D$, and the formula follows by linearity. The map $p'$ is built
in the same way.

The operators $\varrho$ and the conditions of a matching for $p$ and $p'$ are
those of the assembly of matched local maps in
Lemma~\ref{th-square-model}(4). That assembly uses of its triple only the
spaces, characters, character $\psi$ and Cartan involutions of
Lemma~\ref{th-square-model}(2) and the fact that $\Theta$ carries its Weil
representations to $\omega$ and $\omega'$, and the triple here has these.
% T.exists_matching_triple_for_theta_lift (TripleAssembly.triple_exists_matching) takes any
%   reading with IsTransp and any IsMatched R Rf; bottom_matching_of_triangle applies its lemmas
%   (T.fockAlgebraicThetaVectors_stable_finiteAdeles_lie, T.triple_algebra_map_surjective,
%   T.triple_assembly_inverse_finite, T.theta_triple_lift_lie_invariant,
%   T.theta_triple_maps_equivariant, T.theta_triple_lift_lie_equivariant,
%   T.theta_triple_lift_irreducible) to the reading of exists_reading_lat. The formula is
%   PureImage.firstFactorMap_thetaSum_pure_tensor: T.thetaTransform_darbouxModel_linear_injective,
%   T.darbouxModelAlgebraicVectors_le_adelicSchwartz, T.weil_model_projection_factorisation.
\end{proof}

\begin{lemma}\label{th-bottom-split-zeta}
Let $k$ be a nonarchimedean local field with ring of integers $\mathcal O$,
uniformizer $\varpi$ and normalized absolute value $|\cdot|$, let $a,a'\geq0$,
put $d=a+a'$, and let $\alpha$ and $\beta$ be unitary characters of $k^\times$
with quotient $\nu=\alpha/\beta$. Let $c\geq1$ and $m\geq0$ be integers such
that $\nu$ is trivial on the units in $1+\varpi^c\mathcal O$ and $\alpha$ and
$\beta$ are trivial on the units in $1+\varpi^m\mathcal O$. For
$X=(X_1\ X_2)\in M_{a\times d}(k)$, with $X_1$ of size $a\times a$, let
$s_A(X)=\nu(\det X_1)^{-1}$ if $X_1\in GL_a(\mathcal O)$ and $X_2$ has entries
in $\varpi^c\mathcal O$, and $s_A(X)=0$ otherwise. For
$Y=(Y_1\ Y_2)\in M_{a'\times d}(k)$, with $Y_2$ of size $a'\times a'$, let
$s_B(Y)=\nu(\det Y_2)^{-1}$ if $Y_2\in GL_{a'}(\mathcal O)$ and $Y_1$ has
entries in $\varpi^c\mathcal O$, and $s_B(Y)=0$ otherwise. For Schwartz-Bruhat
functions $\Xi$ on $M_{a\times d}(k)$ and $\Psi$ on $M_{a'\times d}(k)$ define
functions on $GL_d(k)$ by
\begin{align*}
Z(\Xi)(h)&=\beta(\det h)|\det h|^{-a/2}\sum_n\int_{K\varpi^nK}
\Xi\bigl(gX_0(h^{-1})^{\mathsf T}\bigr)\,\nu(\det g)|\det g|^{d/2}\,dg,\\
Z'(\Psi)(h)&=\alpha(\det h)|\det h|^{a'/2}\int_{GL_{a'}(k)}
\Psi(gY_0h)\,\nu(\det g)|\det g|^{d/2}\,dg,
\end{align*}
where $X_0=(1_a\ 0)$, $Y_0=(0\ 1_{a'})$ and $K=GL_a(\mathcal O)$, the sum runs
over the nonincreasing $n\in\mathbf Z^a$ and is taken to be $0$ if infinitely
many of its terms are nonzero, $\varpi^n$ is the diagonal matrix with entries
$\varpi^{n_1},\ldots,\varpi^{n_a}$, and the Haar measures give $GL_a(\mathcal O)$
and $GL_{a'}(\mathcal O)$ volume $1$. Then for all $g_A\in GL_a(\mathcal O)$,
$g_B\in GL_{a'}(k)$ and $h_0\in GL_d(k)$ there is $\lambda\neq0$ such that
$\Xi(X)=s_A\bigl(g_AX(h_0^{-1})^{\mathsf T}\bigr)$ and
$\Psi(Y)=\lambda\,s_B(g_BYh_0)$ are nonzero Schwartz-Bruhat functions with
$Z(\Xi)=Z'(\Psi)$.
\end{lemma}
\leanref{Development/AutomorphicArithmetic/DivisorFrames/ExistsScaledStandardPairSectionsOfMaxCompactTran.lean}{UpperTheory.exists_scaled_standardPairSections_of_maxCompact_translate}
\leanref{Development/AbelianVarietyTheory/Proofs/DivisorProofs/IsLocPart.lean}{GlobalTheory.StdPairCone.stdA}
\leanref{Development/AbelianVarietyTheory/Proofs/DivisorProofs/IsLocPart.lean}{GlobalTheory.StdPairCone.stdB}
\leanref{Development/AbelianVarietyTheory/Definitions/Divisors/RankEqCardIffExistsFrameMatrixFirstColumnBlock.lean}{GlobalTheory.SplitFlip.bigP}
\leanref{Development/AbelianVarietyTheory/Definitions/Divisors/RankEqCardIffExistsFrameMatrixFirstColumnBlock.lean}{GlobalTheory.SplitFlip.smallP}
\begin{proof}
This is a computation with the two zeta functions, and its result is stated.
For $g_A=1$, $g_B=1$ and $h_0=1$ the zeta functions of $s_A$ and $s_B$ agree
and are not $0$. A change of rows by $g_A\in GL_a(\mathcal O)$ and of columns
by $h_0$ in $\Xi$, together with a change of rows by $g_B$ and of columns by
$h_0$ in $\Psi$, carries an identity $Z(\Xi)=Z'(\Psi)$ to the same identity for
the changed functions once the second is multiplied by a nonzero factor
$\lambda$ depending only on $g_A$, $g_B$ and $h_0$; and these translates of
$s_A$ and $s_B$ are nonzero Schwartz-Bruhat functions.
% T.standardPair_localZetaMatch_ne_zero, T.zetaIdentity_persists_row_column_changes,
%   T.standardPair_members_mem_schwartzBruhat.
\end{proof}

\begin{situation}\label{th-bottom-split-situation}
Let a grid Weil datum with character $\psi$, $0\leq i,j<m$, a bottom triple
$(A,A',W)$ for $(i,j)$ (Definition~\ref{th-bottom-triple}) and a tuple $t$ be
given, with the characters $\alpha_t$ and $\beta$ of the square $(i,j)$ and
$\nu=\nu_{t,i,j}=\alpha_t/\beta$, such that $(A,A',W;\alpha_t,\beta)$ is a
square triple (Definition~\ref{th-first-lift-matching}(3)) and $S_t$ and
$\varepsilon_t$ have the properties of Lemma~\ref{th-bottom-square-facts}(4).
Let archimedean data on $A$, $A'$ and $W$ be given for which $f_t$ and $f'_t$
are archimedean vacua for $(-A,W)$ and $(A',W)$. Put $a=\operatorname{rk}A$,
$a'=\operatorname{rk}A'$ and $d=a+a'$. Let a triple over $E/F$ have the
properties of Lemma~\ref{th-bottom-split-reading} for $(A,A',W;\alpha_t,\beta)$,
the Weil data $\omega_t$ and $\omega'_t$ of the bottom triple, its lattices
$\mathcal O_E^a$, $\mathcal O_E^{a'}$ and $\mathcal O_E^d$, and $S=S_t$, with
the set $\Sigma$ and the embeddings $\tau_v$. Let $\phi^0_v$ and $\phi'^0_v$ be
its local vacua, and let local maps $(\Pi_\infty,p_\infty,p'_\infty)$ and
$(\Pi_v,p_v,p'_v)$ be matched to it. The coordinates of
Definition~\ref{th-bottom-triple}(3) are taken in frames $\beta_A$,
$\beta_{A'}$ and $\beta_W$ of the bottom triple, bases of the completions of
its lattices at the finite places.

For $v\in\Sigma$ let $w_v\in S_t$ be the place defined by $\tau_v$, let
$\alpha_v$, $\beta_v$ and $\nu_v$ be the characters of $F_v^\times$ given by
$\alpha_t$, $\beta$ and $\nu$ at $w_v$ through $\tau_v$, let $c_v=c_t(w_v)$,
and let $s_{A,v}$ and $s_{B,v}$ be the functions $s_A$ and $s_B$ of
Lemma~\ref{th-bottom-split-zeta} for $k=F_v$, $\nu_v$ and $c_v$. For
$x\in\mathbb W(-A,W)$ let $x_v\in M_{d\times a}(F_v)$ be the component at
$w_v$, read through $\tau_v$, of the coordinates of $x$, and for
$y\in\mathbb W(A',W)$ let $\tilde y_v\in M_{a'\times d}(F_v)$ be that of the
dual coordinates of $y$. Let $b$, $b'$ and $b''$ be the split bases at $v$ of
the three spaces of the triple, and $G'$ and $G''$ the Gram matrices of the
second and third spaces read through $\tau_v$. The \emph{coordinate} at $v$ of
an element $u$ of the Schrödinger model of the first pair, an $a\times d$
matrix over $\mathbf A_E$, is $\operatorname{co}_v(u)=b^{-1}u_v(b''^{-1})^{\mathsf T}$,
where $u_v$ is the component of $u$ at $w_v$ read through $\tau_v$; that of an
element $u$ of the Schrödinger model of the second pair is
$\operatorname{co}'_v(u)=-b'^{\mathsf T}G'\bar u_v{G''}^{\mathsf T}b''$, where
$\bar u_v$ is the component of $u$ at the place conjugate to $w_v$, read
through $\tau_v$ composed with complex conjugation. The
coordinate of an element of $\mathbb W(-A,W)$ or $\mathbb W(A',W)$ is that of
the corresponding element of the Schrödinger model. A \emph{linear model} at
$v$ of the first pair is a linear bijection $\iota$ from the Schwartz-Bruhat
functions on its Schrödinger model at $v$ onto those on $M_{a\times d}(F_v)$
such that $\iota(\rho(u)\phi)(X)=\iota(\phi)(X+\operatorname{co}_v(u))$ for $u$
supported at $v$ in the half attached to $w_v$, and
$\iota(\rho(u)\phi)(\operatorname{co}_v(x))=\psi(\langle x,u\rangle)\,\iota(\phi)(\operatorname{co}_v(x))$
for $u$ supported at $v$ in the other half and $x$ supported at $v$ in the half
attached to $w_v$, where $\langle\cdot,\cdot\rangle$ is the symplectic form of
the pair. A linear model at $v$ of the second pair is defined in the same way,
on $M_{a'\times d}(F_v)$, with $\operatorname{co}'_v$ and with the two halves
exchanged.

Among the conditions on matched local maps (Definition~\ref{th-theta-hypotheses})
are the following:
\begin{enumerate}
\item[(M1)] $p_\infty$ and $p'_\infty$ map the archimedean Fock spaces of the two
  pairs onto $\Pi_\infty$;
\item[(M2)] $p_v(\phi^0_v)=p'_v(\phi'^0_v)\neq0$ at every finite place
  $v\notin\Sigma$;
\item[(M3)] at every $v\in\Sigma$ there are linear models $\iota_v$ and
  $\iota'_v$ at $v$ of the two pairs, carrying $\phi^0_v$ and $\phi'^0_v$ to
  the indicator functions of the integral matrices, such that for
  Schwartz-Bruhat functions $\phi$ and $\phi'$ at $v$,
  $p_v(\phi)=p'_v(\phi')$ if and only if $Z(\iota_v\phi)=Z'(\iota'_v\phi')$,
  where $Z$ and $Z'$ are the zeta functions of
  Lemma~\ref{th-bottom-split-zeta} for $k=F_v$, $\alpha_v$ and $\beta_v$.
\end{enumerate}
\end{situation}
\leanref{Development/AbelianVarietyTheory/BlockTheory/DivisorBase/WeilTensorModels.lean}{GlobalTheory.IndefiniteRowSideFullHecke.WeilTensorModels.IsMatched}
\leanref{Development/AbelianVarietyTheory/BlockTheory/DivisorBase/WeilTensorModels.lean}{GlobalTheory.IndefiniteRowSideFullHecke.WeilTensorModels.RealQuot}
\leanref{Development/AbelianVarietyTheory/BlockTheory/DivisorBase/WeilTensorModels.lean}{GlobalTheory.IndefiniteRowSideFullHecke.WeilTensorModels.IsFinMaps}
\leanref{Development/AbelianVarietyTheory/BlockTheory/DivisorBase/WeilTensorModels.lean}{GlobalTheory.IndefiniteRowSideFullHecke.WeilTensorModels.IsZetaAt}
\leanref{Development/AbelianVarietyTheory/BlockTheory/DivisorBase/WeilTensorModels.lean}{GlobalTheory.IndefiniteRowSideFullHecke.WeilTensorModels.Model.IsLin}

\begin{lemma}\label{th-bottom-split-local}
In Situation~\ref{th-bottom-split-situation}, let $v\in\Sigma$. There are
linear models $\iota_v$ and $\iota'_v$ at $v$ of the two pairs of the triple,
nonzero Schwartz-Bruhat functions $\phi_v$ and $\phi'_v$ on their Schrödinger
models at $v$, and $\lambda_v\neq0$, such that $p_v(\phi_v)=p'_v(\phi'_v)$ and,
for all $x\in\mathbb W(-A,W)$ and $y\in\mathbb W(A',W)$,
\[
\iota_v(\phi_v)(\operatorname{co}_v(x))=s_{A,v}(x_v^{\mathsf T}),\qquad
\iota'_v(\phi'_v)(\operatorname{co}'_v(y))=\lambda_v\,s_{B,v}(\tilde y_v).
\]
\end{lemma}
\leanref{Development/AutomorphicArithmetic/LocalEmbeddings/BotLocal.lean}{GlobalTheory.BotMatch.local_pair}
\begin{proof}
Let $\iota_v$ and $\iota'_v$ be the linear models of condition (M3), and let
$\mathcal O_v$ be the ring of integers of $F_v$ and $\varpi$ a uniformizer. The
characters satisfy the hypotheses of Lemma~\ref{th-bottom-split-zeta} for
$k=F_v$, $\alpha_v$, $\beta_v$ and $c=c_v$: $\alpha_v$ and $\beta_v$ are
unitary, $\nu_v=\alpha_v/\beta_v$, and as continuous characters $\alpha_v$ and
$\beta_v$ are trivial on the units in $1+\varpi^m\mathcal O_v$ for some $m$;
$\nu_v$ is trivial on the units in $1+\varpi^{c_v}\mathcal O_v$, since $c_v$ is
the exponent of the conductor of $\nu$ at $w_v$; and $c_v\geq1$, since $\nu$ is
ramified at $w_v$ by Lemma~\ref{th-bottom-square-facts}(4).

Let $\beta_{A,v}$, $\beta_{A',v}$ and $\beta_{W,v}$ be the components of the
frames at $w_v$, read through $\tau_v$, and put
\[
g_A=-\beta_{A,v}^{\mathsf T}\,b,\qquad
g_B=\beta_{A',v}^{-1}(b'^{-1})^{\mathsf T},\qquad
h_0=b''^{-1}\beta_{W,v},
\]
using that the Gram matrices of the bottom triple are identity matrices. The
matrix $g_A$ lies in $GL_a(\mathcal O_v)$: the frame and $b$ are bases of the
completions of the self-dual lattice $\mathcal O_E^a$ and of its conjugate,
which is the lattice of the first space of the triple, and the form pairs these
completions perfectly. Lemma~\ref{th-bottom-split-zeta} gives $\lambda_v\neq0$
and nonzero Schwartz-Bruhat functions
$\Xi(X)=s_{A,v}\bigl(g_AX(h_0^{-1})^{\mathsf T}\bigr)$ and
$\Psi(Y)=\lambda_vs_{B,v}(g_BYh_0)$ with $Z(\Xi)=Z'(\Psi)$. Since $\iota_v$ and
$\iota'_v$ are bijective on Schwartz-Bruhat functions, $\Xi=\iota_v(\phi_v)$ and
$\Psi=\iota'_v(\phi'_v)$ for nonzero Schwartz-Bruhat functions $\phi_v$ and
$\phi'_v$, and $p_v(\phi_v)=p'_v(\phi'_v)$ by (M3). Finally,
$g_A\operatorname{co}_v(x)(h_0^{-1})^{\mathsf T}=x_v^{\mathsf T}$ and
$g_B\operatorname{co}'_v(y)h_0=\tilde y_v$, by multiplying out the definitions of
the coordinates, which gives the two formulas.
% IsZetaAt (hM.fin.zeta), T.seesawCharacter_conductor_at_ramified_place,
%   UpperTheory.localChar_conductor_le_of_heckeExponent,
%   UpperTheory.exists_maxCompact_eq_frame_gram_splitBasis, BotMatch.frMapO_coX and
%   BotMatch.dualCoord_coY. In the Lean g_A, g_B carry the Gram matrices of -A and A'.
\end{proof}

\begin{lemma}\label{th-bottom-split-pure}
In Situation~\ref{th-bottom-split-situation}, for each $v\in\Sigma$ let
$\iota_v$ and $\iota'_v$ be linear models at $v$ of the two pairs of the
triple, let $\phi_v$ and $\phi'_v$ be nonzero Schwartz-Bruhat functions on their
Schrödinger models at $v$, and let $\lambda_v\neq0$, such that
$\iota_v(\phi_v)(\operatorname{co}_v(x))=s_{A,v}(x_v^{\mathsf T})$ and
$\iota'_v(\phi'_v)(\operatorname{co}'_v(y))=\lambda_vs_{B,v}(\tilde y_v)$ for all
$x\in\mathbb W(-A,W)$ and $y\in\mathbb W(A',W)$. For finite places
$v\notin\Sigma$ put $\phi_v=\phi^0_v$ and $\phi'_v=\phi'^0_v$. Then there are
linear maps $L$ from $\mathcal P(-A,W)$ onto the archimedean Fock space of the
first pair of the triple and $L'$ from $\mathcal P(A',W)$ onto that of the
second such that
\[
f_t\otimes P=\Theta\Bigl(L(P)\otimes\bigotimes_v\phi_v\Bigr),\qquad
f'_t\otimes P'=\Theta\Bigl(L'(P')\otimes\bigotimes_v\phi'_v\Bigr)
\]
for all $P\in\mathcal P(-A,W)$ and $P'\in\mathcal P(A',W)$.
\end{lemma}
\leanref{Development/AutomorphicArithmetic/LocalEmbeddings/BotPair.lean}{GlobalTheory.BotMatch.exists_pure_pair}
\leanref{Development/AutomorphicArithmetic/LocalEmbeddings/BotPure.lean}{GlobalTheory.BotMatch.exists_pure}
\leanref{Development/AbelianVarietyTheory/Proofs/DivisorProofs/HeisenbergStableRelationOfVacuaStatement.lean}{GlobalTheory.T.heisenbergStableRelation_of_vacua}
\leanref{Development/AutomorphicArithmetic/DivisorFrames/StdSectionFirstEqProdLocal.lean}{UpperTheory.stdSectionFirst_eq_prod_local}
\leanref{Development/AutomorphicArithmetic/DivisorFrames/StdSectionSecondEqProdLocal.lean}{UpperTheory.stdSectionSecond_eq_prod_local}
\leanref{Development/AutomorphicArithmetic/DivisorFrames/ThetaTransformFockModelEqOfStandardChar.lean}{UpperTheory.theta_transform_fock_model_eq_of_standard_char}
\begin{proof}
The argument is given for $f_t$. For $f'_t$ replace $\varepsilon_t$ by
$\bar\varepsilon_t$, the coordinates by the dual coordinates, and the constant
$1$ in Step~1 by $\kappa'=\prod_{v\in\Sigma}\lambda_v$, which Step~3 divides
out.

\emph{Step 1: the standard function.} By Lemma~\ref{th-bottom-square-facts}(4)
and Lemma~\ref{th-bottom-split-reading}(2),
$\varepsilon_t=\sum_{v\in\Sigma}1_{w_v}$ with distinct places $w_v$, and
$1_{w_v}+\bar1_{w_v}$ is the idempotent of $v$. Let $\Phi_t$ be the standard
function on the $\varepsilon_t$-half of $\mathbb W(-A,W)$ of
Definition~\ref{th-bottom-triple}(3). For $x$ in this half, the coordinates of
$x$ are integral with last $a'$ rows divisible by $\mathfrak f_t$ if and only
if $x_v$ has these properties at every $v\in\Sigma$, and the inverse of $\nu$ at
the determinant of the top block is the product of the local inverses. Hence
\[
\Phi_t(x)=\prod_{v\in\Sigma}s_{A,v}(x_v^{\mathsf T})
=\prod_{v\in\Sigma}\iota_v(\phi_v)(\operatorname{co}_v(x)).
\]
For $f'_t$ the same argument gives
$\prod_v\iota'_v(\phi'_v)(\operatorname{co}'_v(y))=\kappa'\Phi'_t(y)$ on the
$\bar\varepsilon_t$-half of $\mathbb W(A',W)$.

\emph{Step 2: the vacuum.} The map $\Theta$ is a linear bijection from the
Schwartz-Bruhat functions on the Schrödinger model of the first pair onto the
theta functions for $(-A,W)$, and it carries the Heisenberg operators of the
model to the operators $\rho$ of Definition~\ref{th-weil-datum}
\cite{Wei64}. Let $f^0$ be the Fock vacuum vector of the bottom triple from
which $f_t$ is obtained, and write $f^0=\Theta(\Phi^0)$. By
Lemma~\ref{th-bottom-split-reading}(1), at every finite place $v$ the
completion of the lattice of the pair $(-A,W)$ for $\mathcal O_E^a$ and
$\mathcal O_E^d$ lies in the lattice of the Schrödinger model of the triple at
$v$, so $\phi^0_v$ is fixed by the Heisenberg operators of the doubles of its
elements and of its elements in the split halves. Since $f^0$ is a lattice
vacuum for the lattice of the pair and $\Theta$ is injective and equivariant,
$\Phi^0$ and each of its slices at $v$ are fixed by the same operators. This
local lattice is self-dual for $\psi_v$, and the Schwartz-Bruhat functions at
$v$ fixed by these operators form a line \cite{MVW87}, spanned by $\phi^0_v$.
So every slice of $\Phi^0$ at $v$ is a multiple of $\phi^0_v$, and
$\Phi^0=g\otimes\bigotimes_v\phi^0_v$ for an archimedean Schwartz function $g$,
which is not $0$ since $f^0\neq0$.

\emph{Step 3: the theta sum.} Let $\Delta$ be an operator as in
Definition~\ref{th-bottom-triple}(3) that replaces the component on the
$\varepsilon_t$-half: it is linear on the Schwartz-Bruhat functions on the
$\varepsilon_t$-half, carries the indicator function of the lattice of the pair
on that half to $f^0$, and carries translation by $x$ in the
$\varepsilon_t$-half and multiplication by $\psi(\langle\cdot,y\rangle)$ for $y$
in the $\bar\varepsilon_t$-half to $\rho(x)$ and $\rho(y)$; then
$f_t=\Delta(\Phi_t)$. Let $\mathcal R$ be the space of pairs $(F,\Phi)$ of a
function $F$ on $\mathbb W(-A,W)$ and a Schwartz-Bruhat function $\Phi$ on the
Schrödinger model such that $(F,\Theta(\Phi))$ lies in the span of the pairs
$(\Psi,\Delta(\Psi))$. By the properties of $\Delta$ and $\Theta$, $\mathcal R$
is stable under $(F,\Phi)\mapsto(F(\cdot+x),\rho(x)\Phi)$ for $x$ supported at
some $v\in\Sigma$ in the half attached to $w_v$, and under
$(F,\Phi)\mapsto(\psi(\langle\cdot,y\rangle)F,\rho(y)\Phi)$ for $y$ supported at
$v$ in the other half. By the defining properties of $\iota_v$ and the
invariance of $\phi^0_v$, $\iota_v(\phi^0_v)(\operatorname{co}_v(x))$ is
$c_v=\iota_v(\phi^0_v)(0)\neq0$ if $1_{w_v}x$ lies in the local lattice of the
pair at $v$ and $0$ otherwise. With $c=\prod_{v\in\Sigma}c_v$ and $1_t$ the indicator
function of the $\varepsilon_t$-half, the pair
$\bigl(1_t\prod_v\iota_v(\phi^0_v)(\operatorname{co}_v(\cdot)),\,cg\otimes\bigotimes_v\phi^0_v\bigr)$
is $c$ times the pair formed by the indicator function of the lattice on the
$\varepsilon_t$-half and $\Phi^0$, so it lies in $\mathcal R$.

For $\Sigma'\subset\Sigma$ consider the statement that
$\bigl(1_t\prod_{v\in\Sigma}\iota_v(\chi_v)(\operatorname{co}_v(\cdot)),\,cg\otimes\bigotimes_v\chi_v\bigr)\in\mathcal R$
for all Schwartz-Bruhat functions $\chi_v$ with $\chi_v=\phi^0_v$ for
$v\notin\Sigma'$. It holds for $\Sigma'=\emptyset$. Suppose that it holds for
$\Sigma'$, let $v\in\Sigma\setminus\Sigma'$, fix $\chi_{v'}$ for $v'\neq v$,
equal to $\phi^0_{v'}$ outside $\Sigma'$, and let $N$ be the space of
Schwartz-Bruhat functions $\chi$ at $v$ for which the pair with $\chi_v=\chi$
lies in $\mathcal R$. Then $\phi^0_v\in N$. The space $N$ is stable under
$\rho(u)$ for $u$ supported at $v$ in the half attached to $w_v$: $\iota_v$
turns $\rho(u)$ into translation of the coordinate, the coordinates at the
other places of $\Sigma$ do not change, and $\mathcal R$ is stable. It is
stable in the same way under $\rho(u)$ for $u$ supported at $v$ in the other
half, which $\iota_v$ turns into multiplication by $\psi(\langle x,u\rangle)$ at
the coordinate of $x$. These operators generate the Heisenberg operators at
$v$, the Schrödinger representation of a nonarchimedean local field is
irreducible \cite{MVW87}, and $\phi^0_v\neq0$, so $N$ contains every
Schwartz-Bruhat function. Hence the statement holds for $\Sigma'\cup\{v\}$, and
by induction for $\Sigma$. Applied to the family $(\phi_v)$, with Step~1, it
shows that $(\Phi_t,\Theta(cg\otimes\bigotimes_v\phi_v))$ lies in the span of
the pairs $(\Psi,\Delta(\Psi))$, so
$\Theta(cg\otimes\bigotimes_v\phi_v)=\Delta(\Phi_t)=f_t$. Put $f_0=cg\neq0$.
For $f'_t$ this gives $\Theta(c'g'\otimes\bigotimes_v\phi'_v)=\kappa'f'_t$, and
$f'_0=\kappa'^{-1}c'g'$.

\emph{Step 4: the Fock vectors.} The map
$f\mapsto\Theta(f\otimes\bigotimes_v\phi_v)$ on archimedean Schwartz functions
is injective, since $\Theta$ is and $\bigotimes_v\phi_v\neq0$, and it carries
the derivatives along the archimedean directions of the Heisenberg group to
those of the theta functions. The annihilation operators of the archimedean
data kill the archimedean vacuum $f_t$, so the corresponding derivatives kill
$f_0$, and $f_0$ is a Gaussian of the archimedean Fock space of the first pair,
whose Cartan involutions are those of the archimedean data. For a monomial $P$
in the creation operators let $L(P)$ be the result of applying the
corresponding derivatives to $f_0$, and extend $L$ linearly. Then
$\Theta(L(P)\otimes\bigotimes_v\phi_v)=f_t\otimes P$, by the definition of
$f_t\otimes P$. The image of $L$ lies in the Fock space, the smallest space
that contains the Gaussians and is stable under the archimedean derivatives,
and it is that whole space: every Gaussian is a multiple of $f_0$, and the
derivative of $f_t\otimes P$ along any archimedean direction is $f_t\otimes Q$
for some $Q$, so the image of $L$ is stable under the derivatives.
% Step 2: BotMatch.exists_vac (th_laws, localLattice_le_latticeAtPlace, vac_acts,
%   T.schroedinger_fixed_functions_proportional_of_selfDual,
%   T.adelicSchwartz_eq_pureTensor_of_finite_slices). Step 3: graphP in exists_pure, BigDat.lagr_vac,
%   InductLagr.rel_of_vac with T.schroedinger_irreducible_finitePlaceCompletion. Step 4: fock_of
%   (exists_der, dRho_thL, IsGauss, gauss_line, dRho_ptF).
\end{proof}

\begin{lemma}\label{th-bottom-split-matching}
Assume the published local inputs of the theta correspondence
(Theorem~\ref{lit-bundle-theta}) and the local matching
(Definition~\ref{th-theta-hypotheses}). Let a grid Weil datum with character
$\psi$, $0\leq i,j<m$, a bottom triple $(A,A',W)$ for $(i,j)$
(Definition~\ref{th-bottom-triple}) and a tuple $t$ be given, with the
characters $\alpha_t$ and $\beta$ of the square $(i,j)$. Assume that
$(A,A',W;\alpha_t,\beta)$ is a square triple
(Definition~\ref{th-first-lift-matching}(3)), and that the set $S_t$ and the
idempotent $\varepsilon_t$ have the properties of
Lemma~\ref{th-bottom-square-facts}(4). Let archimedean data on $A$, $A'$ and
$W$ be given for which $f_t$ and $f'_t$ are archimedean vacua for $(-A,W)$ and
$(A',W)$. Then $(A,A',W;\alpha_t,\beta)$ has a matching
$(\omega_t,\omega'_t,\Pi,\varrho,p,p')$ such that either $p(f_t\otimes P)=0$
for every $P\in\mathcal P(-A,W)$, or there are $P\in\mathcal P(-A,W)$,
$P'\in\mathcal P(A',W)$ and $\kappa\in\mathbf C$ with
$p(f_t\otimes P)\neq0$ and $p(f_t\otimes P)=\kappa\,p'(f'_t\otimes P')$.
\end{lemma}
\leanref{Development/AutomorphicArithmetic/ThetaCorrespondence/ThetaPairLifts/LocalTypeDictionary/DvBotMain.lean}{GlobalTheory.BotMatch.bottom_matching_of_triangle}
\begin{proof}
% BotMatch.bottom_matching_of_triangle: exists_reading_lat, LocalMatching, local_pair,
%   exists_pure_pair, theta_transform_fock_model_eq_of_standard_char, then mapA and mapB with
%   T.tripleProduct_image_thetaSum_pure and T.tripleProduct_secondMap_theta_pure.
Let a triple over $E/F$, with the set $\Sigma$ and the embeddings $\tau_v$, be
given by Lemma~\ref{th-bottom-split-reading} for $(A,A',W;\alpha_t,\beta)$, the
Weil data $\omega_t$ and $\omega'_t$ of the bottom triple, its lattices
$\mathcal O_E^n$, which are self-dual for the identity Gram matrices, and
$S=S_t$; conditions (i) to (iii) there are properties of $S_t$ in
Lemma~\ref{th-bottom-square-facts}(4). By the local matching, local maps
$(\Pi_\infty,p_\infty,p'_\infty)$ and $(\Pi_v,p_v,p'_v)$ are matched to the
triple. These data are as in Situation~\ref{th-bottom-split-situation}. By
Lemma~\ref{th-bottom-split-assembly} the matched local maps assemble into a
matching $(\omega_t,\omega'_t,\Pi,\varrho,p,p')$ of $(A,A',W;\alpha_t,\beta)$,
with $\Pi=\Pi_\infty\otimes\bigotimes'_v\Pi_v$ and with $p$ and $p'$ given on
theta sums of algebraic functions by the formula there.

For $v\in\Sigma$ let $\phi_v$ and $\phi'_v$, with linear models $\iota_v$ and
$\iota'_v$ and a constant $\lambda_v$, be given by
Lemma~\ref{th-bottom-split-local}, and for finite places $v\notin\Sigma$ put
$\phi_v=\phi^0_v$ and $\phi'_v=\phi'^0_v$. Then $p_v(\phi_v)=p'_v(\phi'_v)$ at
every finite place $v$: for $v\in\Sigma$ by Lemma~\ref{th-bottom-split-local},
and for $v\notin\Sigma$ by condition (M2) of the situation. Let $L$ and $L'$ be
given by Lemma~\ref{th-bottom-split-pure}. For $P\in\mathcal P(-A,W)$ and
$P'\in\mathcal P(A',W)$ the functions $L(P)\otimes\bigotimes_v\phi_v$ and
$L'(P')\otimes\bigotimes_v\phi'_v$ are algebraic, since $L(P)$ and $L'(P')$ lie
in the archimedean Fock spaces, so Lemmas~\ref{th-bottom-split-pure}
and~\ref{th-bottom-split-assembly} give
\[
p(f_t\otimes P)=p_\infty(L(P))\otimes\bigotimes_vp_v(\phi_v),\qquad
p'(f'_t\otimes P')=p'_\infty(L'(P'))\otimes\bigotimes_vp'_v(\phi'_v).
\]
If $p(f_t\otimes P)=0$ for every $P$, the first alternative holds. Otherwise
choose $P$ with $p(f_t\otimes P)\neq0$. By condition (M1), and since $L'$ maps
$\mathcal P(A',W)$ onto the archimedean Fock space of the second pair, there is
$P'$ with $p'_\infty(L'(P'))=p_\infty(L(P))$. As $p_v(\phi_v)=p'_v(\phi'_v)$ for
every $v$, the two displayed vectors are equal, and
$p(f_t\otimes P)=\kappa\,p'(f'_t\otimes P')$ with $\kappa=1$.
\end{proof}

\subsection{Nonvanishing of the first theta lifts}\label{subsection-th-first-lifts}

\begin{lemma}\label{th-first-lift-inputs}
Assume the published local inputs of the theta correspondence
(Theorem~\ref{lit-bundle-theta}), which include Howe duality at the real places
(Theorem~\ref{lit-howe-duality-real}), and the agreement of global and local
theta integrals and the local matching (Definition~\ref{th-theta-hypotheses}).
Let $D$, $\psi$ and the archimedean embeddings be as in
Definition~\ref{th-first-lift-matching}, and let $A$, $A'$ and $W$ be Hermitian
spaces with archimedean data.
\begin{enumerate}
\item Let $a_\alpha,a_\beta\colon\mathcal T\to\mathbf Q$. There are a complex
  vector space $\Pi_\infty$ of countable dimension with linear operators
  $\varrho_{\infty,\tau}(Y)$ for $\tau\in\mathcal T$ and
  $Y\in\mathfrak u_\tau(W)$, such that $0$ and $\Pi_\infty$ are the only
  subspaces stable under all of them, and surjective linear maps
  $q\colon\mathcal P(-A,W)\to\Pi_\infty$ and
  $q'\colon\mathcal P(A',W)\to\Pi_\infty$, with the following property. For
  every square triple $(A,A',W;\alpha,\beta)$ in which $\alpha$ and $\beta$
  have archimedean exponents $a_\alpha$ and $a_\beta$, every matching
  $(\omega,\omega',\Pi,\varrho,p,p')$ of it and every archimedean vacuum $f$ for
  $(-A,W)$, there is a linear map $\varphi\colon\Pi_\infty\to\Pi$ with
  $\varphi\circ\varrho_{\infty,\tau}(Y)=\varrho_\tau(Y)\circ\varphi$ for all
  $\tau$ and $Y$ and with $p(f\otimes P)=\varphi(q(P))$ for all $P$; and the same
  holds for $p'$, $q'$ and the archimedean vacua for $(A',W)$.
\item Every square triple $(A,A',W;\alpha,\beta)$ has a matching.
\item Let $(\omega,\omega',\Pi,\varrho,p,p')$ be a matching of a square triple
  $(A,A',W;\alpha,\beta)$, and let $I=I_\omega$ and $I'=I_{\omega'}$ be theta
  integrals for Haar measures and fundamental domains. There is $c\neq0$ such
  that for every $\Psi\in\mathcal A(A',W)$ there is
  $\Phi_0\in\mathcal A(A',W)$ with $\Psi-\Phi_0\in\mathcal R(\omega')$, with
  the integral defining $I'(\Phi_0)(h)$ absolutely convergent for every $h$ and
  bounded by a constant times a power of an adelic height of $h$, and with
  $I(\Phi)=c\,I'(\Phi_0)$ for every $\Phi\in\mathcal A(-A,W)$ with
  $p(\Phi)=p'(\Phi_0)$.
\item Let $(A,A',W;\alpha,\beta)$ be a square triple, let $\omega$ and
  $\omega'$ be Weil data as in a matching of it, and let Haar measures on
  $U(A)(\mathbf A)$ and $U(A')(\mathbf A)$ and fundamental domains for
  $U(A)(F)$ and $U(A')(F)$ be given. Take the triple over $E/F$ of
  Lemma~\ref{th-square-model}(2) for $\omega$, $\omega'$ and these measures and
  domains, with matched local maps. For $\Phi\in\mathcal A(-A,W)$, every
  algebraic function $\Phi_D$ of the triple with $\Theta(\Phi_D)=\Phi$ and
  every matching $(\omega,\omega',\Pi,\varrho,p,p')$ with these Weil data,
  $p(\Phi)=0$ if and only if the image of $\Phi_D$ under the matched local
  maps is $0$; so whether $p(\Phi)=0$ is the same for all such matchings. The
  same holds for $p'(\Phi')$ with $\Phi'\in\mathcal A(A',W)$. The space $\Pi$ of
  such a matching has countable dimension. These statements hold under the
  published local inputs and the local matching alone.
\item For all sufficiently large data (Definition~\ref{th-large-data}) the
  following holds. Let a grid Weil datum, $0\leq i,j<m$ and a bottom triple
  $(A,A',W)$ for $(i,j)$ (Definition~\ref{th-bottom-triple}) be given, with the
  characters $\alpha_t$ and $\beta$ of the square $(i,j)$. There are
  archimedean data on $A$, $A'$ and $W$ and
  $a_\alpha,a_\beta\colon\mathcal T\to\mathbf Q$ such that $A$, $A'$ and $W$
  are totally positive definite and, for every tuple $t$:
  $(A,A',W;\alpha_t,\beta)$ is a square triple; $\alpha_t$ and $\beta$ have
  archimedean exponents $a_\alpha$ and $a_\beta$; $f_t$ and $f'_t$ are
  archimedean vacua for $(-A,W)$ and $(A',W)$; and the vectors $f_t\otimes P$
  and $f'_t\otimes P'$ formed with these data are those of
  Definition~\ref{th-bottom-triple}, so that $I_t(P,\cdot)=I_{\omega_t}(f_t\otimes P)$
  and $I'_t(P',\cdot)=I_{\omega'_t}(f'_t\otimes P')$ for the measures and
  fundamental domains of the bottom triple.
\item For all sufficiently large data the following holds. Let a grid Weil
  datum, $0\leq i,j<m$, a bottom triple $(A,A',W)$ for $(i,j)$ and a tuple $t$
  be given, and archimedean data on $A$, $A'$ and $W$ for which $f_t$ and
  $f'_t$ are archimedean vacua. Then $(A,A',W;\alpha_t,\beta)$ has a matching
  $(\omega_t,\omega'_t,\Pi,\varrho,p,p')$ such that either
  $p(f_t\otimes P)=0$ for every $P$, or there are $P$, $P'$ and
  $\kappa\in\mathbf C$ with $p(f_t\otimes P)\neq0$ and
  $p(f_t\otimes P)=\kappa\,p'(f'_t\otimes P')$.
\end{enumerate}
\end{lemma}
\leanref{Development/AutomorphicArithmetic/ThetaCorrespondence/ThetaPairLifts/LocalTypeDictionary/TwistedQuotientReal.lean}{GlobalTheory.T.twisted_quotient_real}
\leanref{Development/AutomorphicArithmetic/ThetaCorrespondence/ThetaPairLifts/LocalTypeDictionary/TripleExistsMatching.lean}{GlobalTheory.T.twist_triple_matches}
\leanref{Development/AutomorphicArithmetic/ThetaCorrespondence/ThetaPairLifts/LocalTypeDictionary/TwistedSquareThetaIntegralsEq.lean}{GlobalTheory.T.twisted_square_theta_integrals_eq}
\leanref{Development/AbelianVarietyTheory/Statements/HeckeThetaIntegrals/ThetaIntegralsAgree.lean}{GlobalTheory.SquareSt.ThetaIntegralsAgree}
\leanref{Development/AutomorphicArithmetic/ThetaCorrespondence/ThetaPairLifts/LocalTypeDictionary/IntertwineOfCompatibleMaps.lean}{GlobalTheory.BotUniqT.bottomLevel_ker_left}
\leanref{Development/AutomorphicArithmetic/ThetaCorrespondence/ThetaPairLifts/LocalTypeDictionary/IntertwineOfCompatibleMaps.lean}{GlobalTheory.BotUniqT.bottomLevel_ker_right}
\leanref{Development/AutomorphicArithmetic/ThetaCorrespondence/ThetaPairLifts/LocalTypeDictionary/RankPolarizationPiLeAleph0.lean}{GlobalTheory.SquareAgree.rank_polarization_pi_le_aleph0}
\leanref{Development/AbelianVarietyTheory/Proofs/DivisorProofs/DefiniteBottomDictionaryFramesExponentsKernelStatement.lean}{GlobalTheory.T.definiteBottom_dictionary_frames_exponents_kernel}
\leanref{Development/AutomorphicArithmetic/ThetaCorrespondence/ThetaPairLifts/LocalTypeDictionary/FockFImageEqZeroOrExistsEqSmul.lean}{GlobalTheory.T.fockF_image_eq_zero_or_exists_eq_smul}
\begin{proof}
% (1): RealQuot.twist_quotient_isReal; (2): TripleRead.triple_exists_matching;
% (3): SquareAgree.twisted_theta_integrals_agree; (4): BotUniqT.bottomLevel_ker_left and
% bottomLevel_ker_right, SquareAgree.rank_polarization_pi_le_aleph0, which take the measures
% nu, F, nu1, F1 as data; (5): BotDict.bottom_dictionary; (6):
% T.fockF_image_eq_zero_or_exists_eq_smul. Their steps are followed below.
For (1), let $Q=\mathcal P(-A,W)$ with the operators
$\rho_{\tau,Y}=\partial^{a_\beta}_{W,\tau}(Y)$ and the subspace
$C=\mathcal C_{a_\beta-a_\alpha}(-A,W)$, and let $Q'=\mathcal P(A',W)$ with
$\rho'_{\tau,Y}=\partial^{a_\alpha}_{W,\tau}(Y)$ and
$C'=\mathcal C_{a_\alpha-a_\beta}(A',W)$, both indexed by the pairs $(\tau,Y)$
with $Y\in\mathfrak u_\tau(W)$ (Definition~\ref{th-square-fock-operators}).
These depend only on $A$, $A'$ and $W$ with their archimedean data and on
$a_\alpha$ and $a_\beta$.

Let $(A,A',W;\alpha,\beta)$ be a square triple in which $\alpha$ and $\beta$
have archimedean exponents $a_\alpha$ and $a_\beta$, and let
$(\omega,\omega',\Pi,\varrho,p,p')$ be a matching of it. The characters
$(\beta,\beta/\alpha)$ of $\omega$ have archimedean exponents $a_\beta$ and
$a_\beta-a_\alpha$, and the characters of $\omega'$ have $a_\alpha$ and
$a_\alpha-a_\beta$. For an archimedean vacuum $f$ for $(-A,W)$ let
$\beta_f\colon Q\to\Pi$ be the map $P\mapsto p(f\otimes P)$, and for an
archimedean vacuum $f'$ for $(A',W)$ let $\beta'_{f'}\colon Q'\to\Pi$ be
$P'\mapsto p'(f'\otimes P')$. Let $\sigma_{\tau,Y}=\varrho_\tau(Y)$. Then:
\begin{enumerate}
\item[(i)] each $\beta_f$ is an intertwiner from $(Q,\rho,C)$ and each
  $\beta'_{f'}$ an intertwiner from $(Q',\rho',C')$, in the sense of
  Lemma~\ref{th-square-radical}, and $\varrho(h)\beta_f=\beta_{\omega(1,h)f}$
  for $h\in U(W)(\mathbf A_f)$, by Lemma~\ref{th-square-fock-maps};
\item[(ii)] the images of the $\beta_f$ span $p(\mathcal A(-A,W))$, which is
  $\Pi$ by condition (b) of a matching, and the images of the $\beta'_{f'}$
  span $\Pi$ as well;
\item[(iii)] $(\rho,C)$ and $(\rho',C')$ satisfy the Howe condition, by
  Lemma~\ref{th-square-howe-condition} for $\omega$ and $\omega'$;
\item[(iv)] every subspace of $\Pi$ stable under every $\varrho_\tau(Y)$ has a
  complement with the same property, by Lemma~\ref{th-square-radical}(2) with
  the operators $\kappa_h=\varrho(h)$, since $0$ and $\Pi$ are the only
  subspaces stable under every $\varrho(h)$ and $\varrho_\tau(Y)$.
\end{enumerate}

If no square triple $(A,A',W;\alpha,\beta)$ in which $\alpha$ and $\beta$
have archimedean exponents $a_\alpha$ and $a_\beta$ has a matching, put
$\Pi_\infty=0$ and $q=q'=0$; the property in (1) holds since there is nothing
to check. Otherwise fix such a square triple with a matching, and write
$\beta^0_f$ and $\beta'^0_{f'}$ for its maps. Since the space of this matching
is not $0$, (ii) gives $f_A$ and $f_B$ with $\beta^0_{f_A}\neq0$ and
$\beta'^0_{f_B}\neq0$. The kernel of $\beta^0_{f_A}$ is admissible for
$(\rho,C)$ and different from $Q$, so $R=\operatorname{rad}(\rho,C)$ is
admissible and different from $Q$, and every admissible subspace containing $R$
is $R$ or $Q$, by Lemma~\ref{th-square-radical}(1); in the same way, using
$\beta'^0_{f_B}$, $R'=\operatorname{rad}(\rho',C')\neq Q'$. By (iv) and
Lemma~\ref{th-square-radical}(4), applied to the family $(\beta^0_f)$ and to
$\beta_1=\beta'^0_{f_B}$, there is an isomorphism $e\colon Q/R\to Q'/R'$ that
carries the operators induced by $\partial^{a_\beta}_{W,\tau}(Y)$ to those
induced by $\partial^{a_\alpha}_{W,\tau}(Y)$.

Put $\Pi_\infty=Q/R$, let $\varrho_{\infty,\tau}(Y)$ be the operator induced by
$\partial^{a_\beta}_{W,\tau}(Y)$, which exists since $R$ is admissible, let
$q\colon Q\to\Pi_\infty$ be the quotient map, and let $q'$ be the quotient map
$Q'\to Q'/R'$ followed by $e^{-1}$. Both are surjective. The space
$\Pi_\infty$ has countable dimension, as a quotient of a polynomial ring in
finitely many variables. If a subspace $N$ of $\Pi_\infty$ is stable under
every $\varrho_{\infty,\tau}(Y)$, its inverse image in $Q$ is admissible and
contains $R$, so it is $R$ or $Q$, and $N$ is $0$ or $\Pi_\infty$.

Now let $(A,A',W;\alpha,\beta)$ be any square triple in which $\alpha$ and
$\beta$ have archimedean exponents $a_\alpha$ and $a_\beta$, with a matching
$(\omega,\omega',\Pi,\varrho,p,p')$, and let $f$ be an archimedean vacuum for
$(-A,W)$. By (iv) and Lemma~\ref{th-square-radical}(3),
$R\subset\ker\beta_f$. So $\beta_f=\varphi\circ q$ for the linear map
$\varphi\colon\Pi_\infty\to\Pi$ induced by $\beta_f$, that is,
$p(f\otimes P)=\varphi(q(P))$ for every $P$; and
$\varphi\circ\varrho_{\infty,\tau}(Y)=\varrho_\tau(Y)\circ\varphi$ since
$\beta_f\rho_{\tau,Y}=\varrho_\tau(Y)\beta_f$. In the same way, for an
archimedean vacuum $f'$ for $(A',W)$, $R'\subset\ker\beta'_{f'}$ by
Lemma~\ref{th-square-radical}(3), and $\beta'_{f'}$ induces
$\varphi'_0\colon Q'/R'\to\Pi$, which commutes with the induced operators. The
map $\varphi'=\varphi'_0\circ e$ then satisfies
$p'(f'\otimes P')=\varphi'(q'(P'))$ and
$\varphi'\circ\varrho_{\infty,\tau}(Y)=\varrho_\tau(Y)\circ\varphi'$, since $e$
carries $\varrho_{\infty,\tau}(Y)$ to the operator induced by
$\partial^{a_\alpha}_{W,\tau}(Y)$.

For (2), let $(A,A',W;\alpha,\beta)$ be a square triple. By
Lemma~\ref{th-square-model}(1) there are Weil data $\omega$ and $\omega'$ as
in a matching, and by Lemma~\ref{th-square-model}(4) a triple over $E/F$ for
them. By the local matching it has matched local maps, and by
Lemma~\ref{th-square-model}(4) these assemble into a matching
$(\omega,\omega',\Pi,\varrho,p,p')$.

For (3), let $\nu,\mathcal F$ and $\nu',\mathcal F'$ be the Haar measures and
fundamental domains of $I$ and $I'$. Take the triple over $E/F$ of
Lemma~\ref{th-square-model}(2) for $\omega$, $\omega'$ and these measures and
domains, with the map $\Theta$, and matched local maps for it, which exist by
the local matching. By condition (2) of the agreement of global and local
theta integrals (Definition~\ref{th-theta-hypotheses}), there is $c_0\neq0$
with $I(\Phi_D)=c_0\,I'(\Phi'_D)$ for every pair $(\Phi_D,\Phi'_D)$ with
matching images, and the matched local maps describe the coinvariants at the
places of $F$ that split completely in $E$ and lie outside the finite set of
the triple. Let $s\neq0$ be given by Lemma~\ref{th-square-model}(3)(b), and put
$c=c_0s$. Let $\Psi\in\mathcal A(A',W)$. By
Lemma~\ref{th-square-model}(3)(c) there are $\Phi_0\in\mathcal A(A',W)$ with
$\Psi-\Phi_0\in\mathcal R(\omega')$ and an algebraic function
$\Phi_{0,D}\in\mathcal S'$ with $\Theta(\Phi_{0,D})=\Phi_0$. By condition (1)
of the agreement, which for functions in $\mathcal S'$ gives convergence of the
theta integral and a bound on the theta kernel by a constant times a power of
the height of $h$, these hold for $\Phi_{0,D}$; by
Lemma~\ref{th-square-model}(3)(d), the integral defining $I'(\Phi_0)(h)$
converges absolutely and is bounded by a constant times a power of an adelic
height of $h$. Let $\Phi\in\mathcal A(-A,W)$ with $p(\Phi)=p'(\Phi_0)$, and
let $\Phi_D$ be an algebraic function with $\Theta(\Phi_D)=\Phi$
(Lemma~\ref{th-square-model}(2)). By Lemma~\ref{th-square-model}(3)(b), the
pair $(\Phi_D,s\Phi_{0,D})$ has matching images, so
$I(\Phi_D)=c_0\,I'(s\Phi_{0,D})=c\,I'(\Phi_{0,D})$. By
Lemma~\ref{th-square-model}(2), $I_\omega(\Phi)(h)=I(\Phi_D)(\kappa(h))$ and
$I_{\omega'}(\Phi_0)(h)=I'(\Phi_{0,D})(\kappa(h))$. Hence
$I_\omega(\Phi)=c\,I_{\omega'}(\Phi_0)$.

For (4), the triple of Lemma~\ref{th-square-model}(2) for $\omega$, $\omega'$
and the given measures and domains has matched local maps by the local
matching. Let $(\omega,\omega',\Pi,\varrho,p,p')$ be a matching, let
$\Phi\in\mathcal A(-A,W)$, and let $\Phi_D$ be an algebraic function with
$\Theta(\Phi_D)=\Phi$, which exists by Lemma~\ref{th-square-model}(2). By
Lemma~\ref{th-square-model}(3)(a), $p(\Phi)=0$ if and only if the image of
$\Phi_D$ under the matched local maps is $0$, a condition that does not
involve the matching. The same argument applies to $p'$. The
space $\Pi$ is $p(\mathcal A(-A,W))$ by condition (b) of a matching, and
$\mathcal A(-A,W)$ is the image under $\Theta$ of the algebraic functions of
the triple, which form a space of countable dimension
(Lemma~\ref{th-square-model}(2)). So $\Pi$ has countable dimension. The
argument uses the published local inputs and the local matching only.

For (5), let the data be large enough for Lemma~\ref{th-bottom-square-facts}.
The Gram matrices of $A$, $A'$ and $W$ are identity matrices, so the
archimedean data $J_\tau=B_\tau=1$ at every $\tau$ satisfy
$\bar B_\tau^{\mathsf T}J_\tau B_\tau=\tau(G)$ for each of the three Gram
matrices $G$. With these data $A$, $A'$ and $W$ are totally positive definite.
Let $a_\alpha$ and $a_\beta$ be given by Lemma~\ref{th-bottom-square-facts}(2),
and let $t$ be a tuple. The conditions of a square triple for
$(A,A',W;\alpha_t,\beta)$ hold:
\begin{enumerate}
\item[(a)] $\omega'_t$ is a Weil datum for $(A',W)$ with characters
  $(\alpha_t,\alpha_t/\beta)$, so $\alpha_t$ is trivial on $E^\times$ and
  restricts to $\eta^{\operatorname{rk}A'}$ on $\mathbf A_F^\times$; in the same
  way $\omega_t$ gives this for $\beta$ and $\operatorname{rk}A$;
\item[(b)] $\operatorname{rk}W=a+a'=\operatorname{rk}A+\operatorname{rk}A'$,
  which is even by Lemma~\ref{th-bottom-square-facts}(1);
\item[(c)] with $s=s'=0$ the condition on signatures holds, since the three
  spaces are positive definite, and the inequalities become
  $\operatorname{rk}A'+2<\operatorname{rk}A$ and $1<\operatorname{rk}A'$,
  which are Lemma~\ref{th-bottom-square-facts}(1);
\item[(d)] $\alpha_t$ and $\beta$ have archimedean exponents $a_\alpha(\tau)$
  and $a_\beta(\tau)$ with $2(a_\alpha(\tau)-a_\beta(\tau))<-\operatorname{rk}W$,
  by Lemma~\ref{th-bottom-square-facts}(2);
\item[(e)] is Lemma~\ref{th-bottom-square-facts}(3).
\end{enumerate}
By Lemma~\ref{th-bottom-square-facts}(5), $f_t$ and $f'_t$ are archimedean
vacua for $(-A,W)$ and $(A',W)$. The archimedean data of
Definition~\ref{th-bottom-triple} are these standard data, so the vectors
$f_t\otimes P$ and $f'_t\otimes P'$ formed with them are those of
Definition~\ref{th-bottom-triple}. Since $\omega_t$ and $\omega'_t$ are the Weil
data of the bottom triple, and $I_t$ and $I'_t$ are formed with its measures
and fundamental domains, $I_t(P,\cdot)=I_{\omega_t}(f_t\otimes P)$ and
$I'_t(P',\cdot)=I_{\omega'_t}(f'_t\otimes P')$.

For (6), let the data be large enough for (5) and for
Lemma~\ref{th-bottom-square-facts}. By (5), $(A,A',W;\alpha_t,\beta)$ is a
square triple; this does not involve archimedean data. By
Lemma~\ref{th-bottom-square-facts}(4), $S_t$ and $\varepsilon_t$ have the
properties required in Lemma~\ref{th-bottom-split-matching}, which gives the
matching.
\end{proof}

\begin{lemma}[Dixmier \cite{Dix63}]\label{th-first-lift-dixmier}
Let $\Pi_\infty$ be a complex vector space of countable dimension and
$(\sigma_i)_{i\in I}$ a family of linear operators on it such that $0$ and
$\Pi_\infty$ are the only subspaces stable under every $\sigma_i$. Then every
linear operator $T$ on $\Pi_\infty$ with $T\sigma_i=\sigma_iT$ for all $i$ is a
scalar.
\end{lemma}
\leanref{Development/AbelianVarietyTheory/Proofs/TwistedOperators/DixmierCommutantScalar.lean}{GlobalTheory.dixmier_commutant_scalar}

\begin{lemma}\label{th-first-lift-product}
Let $\Pi_\infty$ and $\Pi$ be complex vector spaces with linear operators
$\sigma_i$ on $\Pi_\infty$ and $\sigma'_i$ on $\Pi$ for $i$ in a set $I$.
Assume that $0$ and $\Pi_\infty$ are the only subspaces of $\Pi_\infty$ stable
under every $\sigma_i$, and that every operator on $\Pi_\infty$ commuting with
every $\sigma_i$ is a scalar. An \emph{intertwiner} is a linear map
$\varphi\colon\Pi_\infty\to\Pi$ with $\varphi\sigma_i=\sigma'_i\varphi$ for
all $i$.
\begin{enumerate}
\item If $\varphi_1,\ldots,\varphi_n$ are linearly independent intertwiners and
  $x_1,\ldots,x_n\in\Pi_\infty$ satisfy $\sum_k\varphi_k(x_k)=0$, then every
  $x_k$ is $0$.
\item Let $q\colon Q\to\Pi_\infty$ and $q'\colon Q'\to\Pi_\infty$ be surjective
  linear maps, and let $\beta$ and $\beta'$ be linear maps from vector spaces
  $S$ and $S'$ to the spaces of linear maps $Q\to\Pi$ and $Q'\to\Pi$. Assume
  that for every $f\in S$ there is an intertwiner $\varphi$ with
  $\beta(f)=\varphi\circ q$, and that the vectors $\beta(f)(P)$ with $f\in S$
  and $P\in Q$ span $\Pi$; and the same for $\beta'$, $S'$ and $q'$. Then there
  are a complex vector space $\Pi_f$, surjective linear maps
  $r\colon S\to\Pi_f$ and $r'\colon S'\to\Pi_f$ and an isomorphism
  $e\colon\Pi_\infty\otimes\Pi_f\to\Pi$ with $\beta(f)(P)=e(q(P)\otimes r(f))$
  and $\beta'(f')(P')=e(q'(P')\otimes r'(f'))$ for all $f\in S$, $P\in Q$,
  $f'\in S'$ and $P'\in Q'$.
\end{enumerate}
\end{lemma}
\leanref{Development/AbelianVarietyTheory/Proofs/TwistedOperators/LsumMapDomainSome.lean}{GlobalTheory.ProductForm.eval_injective}
\leanref{Development/AbelianVarietyTheory/Proofs/TwistedOperators/LsumMapDomainSome.lean}{GlobalTheory.SquareLM.productForm}
\begin{proof}
For (1), suppose there is a counterexample, and take one with the fewest
nonzero $x_k$. Let $s$ be the set of $k$ with $x_k\neq0$, and $k_0\in s$. Let
$R$ be the space of families $(y_k)_{k\in s}$ in $\Pi_\infty$ with
$\sum_{k\in s}\varphi_k(y_k)=0$. Applying $\sigma_i$ to every $y_k$ preserves
$R$, since $\varphi_k\sigma_i=\sigma'_i\varphi_k$. The map
$\pi\colon R\to\Pi_\infty$, $(y_k)\mapsto y_{k_0}$, is injective: a family in
$R$ with $y_{k_0}=0$ has fewer nonzero entries than $(x_k)$, so it is $0$ by
the minimality. The image of $\pi$ is a stable subspace containing
$x_{k_0}\neq0$, so it is $\Pi_\infty$. For $k\in s$, the operator that maps $y$
to the $k$-th entry of $\pi^{-1}(y)$ commutes with every $\sigma_i$, so it is a
scalar $a_k$, and $a_{k_0}=1$. Thus $(a_ky)_{k\in s}$ lies in $R$ for every
$y\in\Pi_\infty$, that is, $\sum_{k\in s}a_k\varphi_k=0$. This contradicts the
linear independence.

For (2), let $F$ be the space of intertwiners $\varphi$ with
$\beta(f)=\varphi\circ q$ for some $f\in S$, and $F'$ the space defined in the
same way by $\beta'$, $S'$ and $q'$. The map
$\mathrm{ev}\colon\Pi_\infty\otimes F\to\Pi$, $x\otimes\varphi\mapsto\varphi(x)$,
is injective by (1), since an element of $\Pi_\infty\otimes F$ can be written
$\sum_kx_k\otimes\varphi_k$ with $(\varphi_k)$ part of a basis of $F$. It is
surjective, since $\beta(f)(P)=\mathrm{ev}(q(P)\otimes\varphi)$ when
$\beta(f)=\varphi\circ q$. The same holds for $F'$.

Next, $F\subset F'$. Let $\varphi\in F$, suppose that $\varphi\notin F'$, and
choose $x$ with $\varphi(x)\neq0$. Since the map $\mathrm{ev}$ for $F'$ is
surjective, $\varphi(x)=\sum_k\varphi'_k(y_k)$ for some $y_k$ and part
$(\varphi'_k)$ of a basis of $F'$. The intertwiners $\varphi$ and $\varphi'_k$
are linearly independent and $\varphi(x)-\sum_k\varphi'_k(y_k)=0$, so $x=0$ by
(1), which is false. In the same way $F'\subset F$, so $F=F'$.

Put $\Pi_f=F$ and $e=\mathrm{ev}$. For $f\in S$ the intertwiner $\varphi$ with
$\beta(f)=\varphi\circ q$ is unique, since $q$ is surjective; let $r(f)$ be
this $\varphi$. Then $r$ is linear, it maps $S$ onto $F$ by the definition of
$F$, and $\beta(f)(P)=r(f)(q(P))=e(q(P)\otimes r(f))$. Define $r'$ in the same
way; it maps $S'$ onto $F'=F$.
\end{proof}

\begin{lemma}\label{th-first-lift-factor}
Assume the hypotheses of Lemma~\ref{th-first-lift-inputs}. Let $D$, $\psi$
and the archimedean embeddings be as in
Definition~\ref{th-first-lift-matching}, and let $A$, $A'$ and $W$ be Hermitian
spaces with archimedean data.
\begin{enumerate}
\item If a square triple $(A,A',W;\alpha,\beta)$ has a matching, then for every
  Weil datum $\omega$ for $(-A,W)$ with characters $(\beta,\beta/\alpha)$ and
  every Weil datum $\omega'$ for $(A',W)$ with characters
  $(\alpha,\alpha/\beta)$ it has a matching with Weil data $\omega$ and
  $\omega'$.
\item Let $a_\alpha,a_\beta\colon\mathcal T\to\mathbf Q$. There are a complex
  vector space $\Pi_\infty$ and surjective linear maps
  $q\colon\mathcal P(-A,W)\to\Pi_\infty$ and
  $q'\colon\mathcal P(A',W)\to\Pi_\infty$ with the following property. Every
  square triple $(A,A',W;\alpha,\beta)$ in which $\alpha$ and $\beta$ have
  archimedean exponents $a_\alpha$ and $a_\beta$ has a matching, and for every
  matching $(\omega,\omega',\Pi,\varrho,p,p')$ of it there are a complex vector
  space $\Pi_f$, linear maps $r$ and $r'$ from the spaces of functions on
  $\mathbb W(-A,W)$ and on $\mathbb W(A',W)$ to $\Pi_f$, each mapping the
  archimedean vacua onto $\Pi_f$, and an isomorphism
  $e\colon\Pi_\infty\otimes\Pi_f\to\Pi$, such that
  $p(f\otimes P)=e(q(P)\otimes r(f))$ for every archimedean vacuum $f$ for
  $(-A,W)$ and $P\in\mathcal P(-A,W)$, and
  $p'(f'\otimes P')=e(q'(P')\otimes r'(f'))$ for every archimedean vacuum $f'$
  for $(A',W)$ and $P'\in\mathcal P(A',W)$.
\end{enumerate}
\end{lemma}
\leanref{Development/AutomorphicArithmetic/ThetaCorrespondence/ThetaPairLifts/LocalTypeDictionary/TripleWeilTransferMatching.lean}{GlobalTheory.T.triple_weil_transfer_matching}
\leanref{Development/AbelianVarietyTheory/Proofs/DivisorProofs/WeilDatumUniqueOnThetaFunctions.lean}{GlobalTheory.WeilUniq.weilDatum_unique}
\leanref{Development/AutomorphicArithmetic/ThetaCorrespondence/ThetaPairLifts/LocalTypeDictionary/TwistedSquareLocalMatching.lean}{GlobalTheory.T.twisted_local_factors_match}
\leanref{Development/AbelianVarietyTheory/Statements/HeckeThetaIntegrals/ThetaIntegralsAgree.lean}{GlobalTheory.SquareSt.LocalMatching}
\begin{proof}
For (1), let $(\omega_0,\omega'_0,\Pi,\varrho,p,p')$ be a matching. The
creation operators preserve theta functions, so $\mathcal A(-A,W)$ consists of
theta functions. Two Weil data for the same pair with the same characters
agree on theta functions. So $\omega(g,h)\Phi=\omega_0(g,h)\Phi$ for all
$(g,h)$ and $\Phi\in\mathcal A(-A,W)$, and likewise for $\omega'$, $\omega'_0$
and $\mathcal A(A',W)$. Conditions (a) to (d) of a matching involve the Weil
data only through these values. So $(\omega,\omega',\Pi,\varrho,p,p')$ is a
matching.

For (2), let $\Pi_\infty$, $\varrho_\infty$, $q$ and $q'$ be given by
Lemma~\ref{th-first-lift-inputs}(1). By Lemma~\ref{th-first-lift-dixmier},
every operator on $\Pi_\infty$ that commutes with every
$\varrho_{\infty,\tau}(Y)$ is a scalar. Let $(A,A',W;\alpha,\beta)$ be a square
triple as in (2). It has a matching by Lemma~\ref{th-first-lift-inputs}(2).
Let $(\omega,\omega',\Pi,\varrho,p,p')$ be a matching of it. The archimedean
vacua for $(-A,W)$ form a vector space $S$, and $f\otimes P$ is linear in $f$;
for $f\in S$ let $\beta_0(f)$ be the linear map $P\mapsto p(f\otimes P)$.
Define $S'$ and $\beta'_0$ in the same way from $(A',W)$ and $p'$. By
Lemma~\ref{th-first-lift-inputs}(1), $\beta_0(f)=\varphi\circ q$ for a linear
map $\varphi$ with $\varphi\circ\varrho_{\infty,\tau}(Y)=\varrho_\tau(Y)\circ\varphi$
for all $\tau$ and $Y$, and likewise for $\beta'_0$ and $q'$. The vectors
$\beta_0(f)(P)$ span $p(\mathcal A(-A,W))$, which is $\Pi$ by condition (b) of
a matching, and the vectors $\beta'_0(f')(P')$ span $\Pi$ as well.
Lemma~\ref{th-first-lift-product}(2), for the operators indexed by the pairs
$(\tau,Y)$, gives $\Pi_f$, surjective linear maps $r\colon S\to\Pi_f$ and
$r'\colon S'\to\Pi_f$ and $e$ with the stated identities. Extend $r$ and $r'$
linearly to all functions.
\end{proof}

\begin{lemma}\label{th-first-lift-compact}
Let $D$, $\psi$ and the archimedean embeddings be as in
Definition~\ref{th-first-lift-matching}, let $S$ and $T$ be Hermitian spaces
with archimedean data, and let $\omega$ be a Weil datum for $(S,T)$ whose
characters have archimedean exponents $a_S$ and $a_T$ at $\tau\in\mathcal T$.
Consider the elements of $U(S)(\mathbf A)$ with component $1$ away from the
place of $\tau$ whose component at $\tau$, written in the orthonormal basis at
$\tau$, is a unitary matrix $h$ commuting with the matrix of signs $J_\tau$;
let $h_+$ and $h_-$ be the diagonal blocks of $h$ for the two signs. Define the
same elements of $U(T)(\mathbf A)$, with blocks $k_+$ and $k_-$. Let $f\neq0$ be
an archimedean vacuum for $(S,T)$. There are integers $x_\pm$ and $y_\pm$,
determined by the numbers of signs of each kind and by $a_S$ and $a_T$, such
that for all such $h$ and $k$ and every $P\in\mathcal P(S,T)$,
\[
\omega(h,k)(f\otimes P)=\det(h_+)^{x_+}\det(h_-)^{x_-}\det(k_+)^{y_+}\det(k_-)^{y_-}\,f\otimes P^{h,k},
\]
where $P^{h,k}$ is obtained from $P$ by the linear substitution of its variables
at $\tau$ given by $\bar h$ and $k$. The total degree of $P^{h,k}$ is at most
that of $P$.
\end{lemma}
\leanref{Development/AbelianVarietyTheory/Proofs/TwistedOperators/CompactSubgroupsActOnFockVectorsAllSignatures.lean}{GlobalTheory.CompactFockAction.compactFockAll}
\leanref{Development/AbelianVarietyTheory/Definitions/TwistedInnerProducts/IsVacFOfIsVacVec.lean}{GlobalTheory.ChainFr.CompactFockF}

\begin{lemma}\label{th-first-lift-relations}
Let $D$, $\psi$ and the archimedean embeddings be as in
Definition~\ref{th-first-lift-matching}, let $M$ and $W$ be Hermitian spaces
with archimedean data, and let $\omega$ be a Weil datum for $(M,W)$. Let $\nu$
be a Haar measure on $U(M)(\mathbf A)$, $\mathcal F$ a fundamental domain for
$U(M)(F)$, and $I_\omega$ the theta integral for them.
\begin{enumerate}
\item If $M$ is totally positive definite, then for $\Phi\in\mathcal A(M,W)$ and
  $h\in U(W)(\mathbf A)$ the function $g\mapsto\theta_\omega(g,h;\Phi)$ is
  integrable on $\mathcal F$.
\item If $M$ is totally positive definite, then $I_\omega$ vanishes on
  $\mathcal R(\omega)$.
\item If $W$ is totally positive definite and the characters of $\omega$ have
  archimedean exponents at every $\tau\in\mathcal T$, then for
  $\Phi\in\mathcal A(M,W)$, $\tau\in\mathcal T$ and $Y\in\mathfrak u_\tau(W)$
  the functions $\omega(1,\exp_\tau(sY))\Phi$, $s\in\mathbf R$, lie in a
  finite-dimensional space.
\end{enumerate}
\end{lemma}
\leanref{Development/AbelianVarietyTheory/Proofs/DivisorProofs/MatchingTripleSignature.lean}{GlobalTheory.DefIrred.integrable_algF}
\leanref{Development/AutomorphicArithmetic/ThetaCorrespondence/ThetaPairLifts/LocalTypeDictionary/ThetaIntegralVanishesOnFiniteRelations.lean}{GlobalTheory.T.theta_integral_vanishes_on_finite_relations}
\leanref{Development/AbelianVarietyTheory/Proofs/TwistedOperators/IntegralRightTranslate.lean}{GlobalTheory.SquareComm.integral_rightTranslate}
\leanref{Development/AbelianVarietyTheory/Proofs/DivisorProofs/OneParamSubgroupFockOrbitFiniteDimensionalStatement.lean}{GlobalTheory.T.oneParamSubgroup_fockOrbit_finiteDimensional}
\begin{proof}
For (1), the creation operators preserve theta functions, so $\Phi$ is a
theta function, and so is $\omega(g,h)\Phi$. By the conditions on a Weil
datum (Definition~\ref{th-weil-datum}(3)), the function
$g\mapsto\theta_\omega(g,h;\Phi)$ is continuous, and for $\gamma\in U(M)(F)$,
$\theta_\omega(\gamma g,h;\Phi)=(\omega(\gamma,1)\omega(g,h)\Phi)(0)=\theta_\omega(g,h;\Phi)$.
Since $M$ is totally definite, $U(M)$ is anisotropic over $F$, and
$U(M)(F)\backslash U(M)(\mathbf A)$ is compact \cite{Bor69}. So the function is
bounded and $\mathcal F$ has finite measure, and the function is integrable on
$\mathcal F$.

For (2), the group $U(M)(\mathbf A)$ is unimodular, so $\nu$ is also right
invariant. Let $u\in U(M)(\mathbf A_f)$ and $\Psi\in\mathcal A(M,W)$. Then
$\theta_\omega(g,h;\omega(u,1)\Psi)=\theta_\omega(gu,h;\Psi)$. The set
$\mathcal Fu$ is a fundamental domain for $U(M)(F)$, and by right invariance
$\int_{\mathcal F}\theta_\omega(gu,h;\Psi)\,d\nu(g)=\int_{\mathcal Fu}\theta_\omega(g,h;\Psi)\,d\nu(g)$,
with integrability on $\mathcal F$ on the left equivalent to integrability on
$\mathcal Fu$ on the right. A function invariant under $U(M)(F)$ on the left
is integrable on one fundamental domain if and only if it is integrable on
another, with the same integral. By (1), the integrals are therefore defined
and $I_\omega(\omega(u,1)\Psi)=I_\omega(\Psi)$. Since the integrals of all the
functions involved converge, $I_\omega$ is linear on their span, and it
vanishes on $\mathcal R(\omega)$.

For (3), by linearity it is enough to treat $\Phi=f\otimes P$ with $f$ an
archimedean vacuum, and $f\neq0$ since otherwise $\Phi=0$. Since $W$ is
totally positive definite, all the signs of the Cartan involution at $\tau$
are equal, and $\exp(sY)$ preserves the positive definite form $\tau(G_W)$; so
in the orthonormal basis at $\tau$ it is a unitary matrix commuting with
$J_\tau$. By Lemma~\ref{th-first-lift-compact} with $h=1$,
$\omega(1,\exp_\tau(sY))(f\otimes P)$ is a multiple of $f\otimes Q$ for a
polynomial $Q$ of total degree at most that of $P$. The vectors $f\otimes Q$
with such $Q$ span a finite-dimensional space.
\end{proof}

\begin{lemma}\label{th-first-lift-comparison}
Assume the hypotheses of Lemma~\ref{th-first-lift-inputs}. Let
$(\omega,\omega',\Pi,\varrho,p,p')$ be a matching of a square triple
$(A,A',W;\alpha,\beta)$ (Definition~\ref{th-first-lift-matching}) in which
$A'$ is totally positive definite, and let $I=I_\omega$ and $I'=I_{\omega'}$ be
theta integrals for Haar measures and fundamental domains. There is $c\neq0$
with $I(\Phi)=c\,I'(\Psi)$ for all $\Phi\in\mathcal A(-A,W)$ and
$\Psi\in\mathcal A(A',W)$ with $p(\Phi)=p'(\Psi)$. In particular $I'(\Psi)=0$
for every $\Psi\in\mathcal A(A',W)$ with $p'(\Psi)=0$.
\end{lemma}
\leanref{Development/AutomorphicArithmetic/ThetaCorrespondence/ThetaPairLifts/LocalTypeDictionary/TwistedSquareThetaIntegralsEq.lean}{GlobalTheory.T.twisted_square_theta_integrals_eq}
\leanref{Development/AutomorphicArithmetic/ThetaCorrespondence/ThetaPairLifts/LocalTypeDictionary/ThetaIntegralVanishesOnFiniteRelations.lean}{GlobalTheory.T.theta_integral_vanishes_on_finite_relations}
\leanref{Development/AutomorphicArithmetic/ThetaCorrespondence/ThetaPairLifts/LocalTypeDictionary/BottomThetaNormEqProdLocalFactors.lean}{GlobalTheory.BottomFactor.bottom_theta_norm_factorization}
\begin{proof}
Let $c\neq0$ be as in Lemma~\ref{th-first-lift-inputs}(3). For
$\Psi\in\mathcal A(A',W)$ let $\Phi_0$ be given for $\Psi$ by
Lemma~\ref{th-first-lift-inputs}(3).
\begin{enumerate}
\item[(i)] $I'(\Psi)=I'(\Phi_0)$. The integral defining $I'(\Psi)$ converges by
  Lemma~\ref{th-first-lift-relations}(1), and the one defining $I'(\Phi_0)$
  converges by Lemma~\ref{th-first-lift-inputs}(3). So
  $I'(\Psi)-I'(\Phi_0)=I'(\Psi-\Phi_0)$, which is $0$ by
  Lemma~\ref{th-first-lift-relations}(2), since
  $\Psi-\Phi_0\in\mathcal R(\omega')$.
\item[(ii)] $p'(\Phi_0)=p'(\Psi)$. By conditions (a) and (c) of a matching,
  $\mathcal R(\omega')$ lies in $\mathcal A(A',W)$ and $p'$ vanishes on it.
\item[(iii)] If $p'(\Psi)=0$, then $I'(\Psi)=0$. Indeed $p'(\Phi_0)=0=p(0)$ by
  (ii), so Lemma~\ref{th-first-lift-inputs}(3) with $\Phi=0$ gives
  $c\,I'(\Phi_0)=I(0)=0$. Since $c\neq0$, $I'(\Phi_0)=0$, and $I'(\Psi)=0$
  by (i).
\item[(iv)] If $p(\Phi)=p'(\Psi)$, then $p(\Phi)=p'(\Phi_0)$ by (ii), so
  $I(\Phi)=c\,I'(\Phi_0)$ by Lemma~\ref{th-first-lift-inputs}(3), and this is
  $c\,I'(\Psi)$ by (i).
\end{enumerate}
Part (iv) is the first assertion and part (iii) the second.
\end{proof}

\begin{remark}\label{th-first-lift-comparison-short}
The second assertion of Lemma~\ref{th-first-lift-comparison} is also the case
$\Phi=0$ of the first: if $p'(\Psi)=0=p(0)$, the first assertion gives
$c\,I'(\Psi)=I(0)=0$, and $c\neq0$. This makes step (iii) of the proof
unnecessary.
\end{remark}

\begin{lemma}\label{th-first-lift-irreducible}
Assume the hypotheses of Lemma~\ref{th-first-lift-inputs}. Let
$(\omega,\omega',\Pi,\varrho,p,p')$ be a matching of a square triple
$(A,A',W;\alpha,\beta)$ (Definition~\ref{th-first-lift-matching}) in which
$A$, $A'$ and $W$ are totally positive definite, and let $I=I_\omega$ and
$I'=I_{\omega'}$ be theta integrals for Haar measures and fundamental domains.
If $I'(\Phi_1)=0$ for some $\Phi_1\in\mathcal A(A',W)$ with $p'(\Phi_1)\neq0$,
then $I'$ vanishes on $\mathcal A(A',W)$.
\end{lemma}
\leanref{Development/AutomorphicArithmetic/ThetaCorrespondence/ThetaPairLifts/LocalTypeDictionary/ThetaLiftVanishesOfIrreducibleDefect.lean}{GlobalTheory.T.theta_lift_vanishes_of_irreducible_defect}
\begin{proof}
Let $K_0$ be the space of functions $\Psi$ on $\mathbb W(A',W)$ such that
$g\mapsto\theta_{\omega'}(g,h;\Psi)$ is integrable on the fundamental domain
for every $h$ and $I'(\Psi)=0$. Let $K=K_0\cap\mathcal A(A',W)$ and
$N=p'(K)$. For $h_1\in U(W)(\mathbf A)$,
$\theta_{\omega'}(g,h;\omega'(1,h_1)\Psi)=\theta_{\omega'}(g,hh_1;\Psi)$, so
$\omega'(1,h_1)$ maps $K_0$ into itself.

The subspace $N$ is stable under $\varrho(h)$ for $h\in U(W)(\mathbf A_f)$.
Indeed, for $\Psi\in K$, $\omega'(1,h)\Psi$ lies in $\mathcal A(A',W)$ by
condition (a) of a matching, hence in $K$, and
$\varrho(h)p'(\Psi)=p'(\omega'(1,h)\Psi)$ by condition (d).

The subspace $N$ is stable under $\varrho_\tau(Y)$ for $Y\in\mathfrak u_\tau(W)$.
Let $\Psi\in K$, and let $\Psi'\in\mathcal A(A',W)$ be the derivative at $s=0$
of $F(s)=\omega'(1,\exp_\tau(sY))\Psi$ (condition (a)); by condition (d),
$p'(\Psi')=\varrho_\tau(Y)p'(\Psi)$. By condition (d) of a square triple, the
characters $\alpha$ and $\alpha/\beta$ of $\omega'$ have archimedean exponents
at every $\tau$. By Lemma~\ref{th-first-lift-relations}(3), the $F(s)$ span a
finite-dimensional space $V$, which lies in $K_0$ by the first paragraph. The
function $\Psi'$ is the pointwise limit of the functions $(F(s)-F(0))/s$ in
$V$ as $s\to0$, and a finite-dimensional space of functions is closed under
pointwise limits. So $\Psi'\in V\subset K_0$, hence $\Psi'\in K$ and
$\varrho_\tau(Y)p'(\Psi)\in N$.

By Lemma~\ref{th-first-lift-relations}(1), $\Phi_1\in K$, so $N$ contains
$p'(\Phi_1)\neq0$. Since $0$ and $\Pi$ are the only stable subspaces, $N=\Pi$.
Let $\Phi_2\in\mathcal A(A',W)$, and choose $\Psi\in K$ with
$p'(\Psi)=p'(\Phi_2)$. Then $p'(\Phi_2-\Psi)=0$, so $I'(\Phi_2-\Psi)=0$ by
Lemma~\ref{th-first-lift-comparison}. By Lemma~\ref{th-first-lift-relations}(1)
the integrals converge, so $I'(\Phi_2)=I'(\Psi)=0$.
\end{proof}

\begin{lemma}\label{th-first-lift-unique}
Assume the published local inputs of the theta correspondence
(Theorem~\ref{lit-bundle-theta}), which include Howe duality at the real places
(Theorem~\ref{lit-howe-duality-real}), and the local matching
(Definition~\ref{th-theta-hypotheses}). Let a grid Weil
datum, $0\leq i,j<m$, a bottom triple $(A,A',W)$ for $(i,j)$
(Definition~\ref{th-bottom-triple}) and a tuple $t$ be given such that
$(A,A',W;\alpha_t,\beta)$ is a square triple, and let $A$, $A'$ and $W$ carry
archimedean data (Definition~\ref{th-first-lift-matching}). Let
$(\omega_t,\omega'_t,\Pi_0,\varrho_0,p_0,p'_0)$ and
$(\omega_t,\omega'_t,\Pi,\varrho,p,p')$ be matchings of
$(A,A',W;\alpha_t,\beta)$ with the Weil data of the bottom triple. There are a
linear isomorphism $E\colon\Pi_0\to\Pi$ and $c\neq0$ with $p=c\,E\circ p_0$ on
$\mathcal A(-A,W)$ and $p'=E\circ p'_0$ on $\mathcal A(A',W)$.
\end{lemma}
\leanref{Development/AutomorphicArithmetic/ThetaCorrespondence/ThetaPairLifts/LocalTypeDictionary/IntertwineOfCompatibleMaps.lean}{GlobalTheory.BotUniqT.bottomLevel_unique_of_triple}
\leanref{Development/AutomorphicArithmetic/ThetaCorrespondence/ThetaPairLifts/LocalTypeDictionary/IntertwineOfCompatibleMaps.lean}{GlobalTheory.T.theta_bottomLevel_uniqueness}
\begin{proof}
By Lemma~\ref{th-first-lift-inputs}(4), applied with the Haar measures and
fundamental domains of the bottom triple, $p_0(\Phi)=0$ if and only if
$p(\Phi)=0$, for $\Phi\in\mathcal A(-A,W)$. Since $p_0$ and $p$ are
surjective, there is a linear isomorphism $E_A\colon\Pi_0\to\Pi$ with
$p=E_A\circ p_0$ on $\mathcal A(-A,W)$. In the same way there is $E_B$ with
$p'=E_B\circ p'_0$ on $\mathcal A(A',W)$.

The map $E_A$ intertwines $\varrho_0$ and $\varrho$. Let $x=p_0(\Phi)$ with
$\Phi\in\mathcal A(-A,W)$, and let $h\in U(W)(\mathbf A_f)$. By conditions (a)
and (d) of a matching, $\omega_t(1,h)\Phi\in\mathcal A(-A,W)$ and
$E_A(\varrho_0(h)x)=E_A(p_0(\omega_t(1,h)\Phi))=p(\omega_t(1,h)\Phi)=\varrho(h)E_A(x)$.
The same argument with the derivative of $\Phi$ along $Y\in\mathfrak u_\tau(W)$
gives $E_A\varrho_{0,\tau}(Y)=\varrho_\tau(Y)E_A$. In the same way $E_B$
intertwines $\varrho_0$ and $\varrho$.

So $E_B^{-1}E_A$ commutes with every $\varrho_0(h)$ and $\varrho_{0,\tau}(Y)$. By
Lemma~\ref{th-first-lift-inputs}(4), $\Pi_0$ has countable dimension, and $0$
and $\Pi_0$ are its only stable subspaces. By
Lemma~\ref{th-first-lift-dixmier}, $E_B^{-1}E_A$ is a scalar $c$, and $c\neq0$
since $E_B^{-1}E_A$ is invertible and $\Pi_0\neq0$. Put $E=E_B$. Then
$p=E_Ap_0=c\,Ep_0$ on $\mathcal A(-A,W)$ and $p'=Ep'_0$ on $\mathcal A(A',W)$.
\end{proof}

\begin{lemma}\label{th-first-lift-finite}
Assume the hypotheses of Lemma~\ref{th-first-lift-inputs}. For all
sufficiently large data (Definition~\ref{th-large-data}) the following holds.
Let a grid Weil datum, $0\leq i,j<m$, a bottom triple $(A,A',W)$ for $(i,j)$
(Definition~\ref{th-bottom-triple}) and a tuple $t$ be given, and archimedean
data on $A$, $A'$ and $W$ for which $f_t$ and $f'_t$ are archimedean vacua
(Definition~\ref{th-first-lift-matching}). Let
$(\omega_t,\omega'_t,\Pi,\varrho,p,p')$ be a matching of
$(A,A',W;\alpha_t,\beta)$. Let $\Pi_\infty$ and $\Pi_f$ be complex vector
spaces, $q\colon\mathcal P(-A,W)\to\Pi_\infty$ and
$q'\colon\mathcal P(A',W)\to\Pi_\infty$ surjective linear maps, $r$ and $r'$
linear maps from the spaces of functions on $\mathbb W(-A,W)$ and on
$\mathbb W(A',W)$ to $\Pi_f$, and $e\colon\Pi_\infty\otimes\Pi_f\to\Pi$ an
isomorphism, such that $p(f\otimes P)=e(q(P)\otimes r(f))$ and
$p'(f'\otimes P')=e(q'(P')\otimes r'(f'))$ for all archimedean vacua $f$ for
$(-A,W)$ and $f'$ for $(A',W)$ and all $P$ and $P'$. Then
$r(f_t)=\kappa\,r'(f'_t)$ for some $\kappa\in\mathbf C$.
\end{lemma}
\leanref{Development/AutomorphicArithmetic/ThetaCorrespondence/ThetaPairLifts/LocalTypeDictionary/ThetaBottomFinitePlaceMatching.lean}{GlobalTheory.T.theta_bottom_finite_place_matching}
\leanref{Development/AutomorphicArithmetic/ThetaCorrespondence/ThetaPairLifts/LocalTypeDictionary/ExistsMatchingFiniteDataBottomLevel.lean}{GlobalTheory.T.theta_bottomLevel_matches_finiteData}
\leanref{Development/AbelianVarietyTheory/Proofs/DivisorProofs/BottomTriangleFiniteFactorExists.lean}{GlobalTheory.BotFinRef.exists_smul_of_factor}
\begin{proof}
Take the data large enough for Lemma~\ref{th-first-lift-inputs}(5) and~(6).
By Lemma~\ref{th-first-lift-inputs}(5), $(A,A',W;\alpha_t,\beta)$ is a square
triple. By Lemma~\ref{th-first-lift-inputs}(6) it has a matching
$(\omega_t,\omega'_t,\Pi_0,\varrho_0,p_0,p'_0)$ such that either
$p_0(f_t\otimes P)=0$ for every $P$, or $p_0(f_t\otimes P)\neq0$ and
$p_0(f_t\otimes P)=\kappa_0\,p'_0(f'_t\otimes P')$ for some $P$, $P'$ and
$\kappa_0$. By Lemma~\ref{th-first-lift-unique} there are an isomorphism
$E\colon\Pi_0\to\Pi$ and $c\neq0$ with $p=c\,Ep_0$ and $p'=Ep'_0$. Put
$e_0=E^{-1}e$, $x=c^{-1}r(f_t)$ and $x'=r'(f'_t)$. Then
$p_0(f_t\otimes P)=e_0(q(P)\otimes x)$ and
$p'_0(f'_t\otimes P')=e_0(q'(P')\otimes x')$ for all $P$ and $P'$.

If $y\in\Pi_\infty$ is nonzero and $y\otimes z=y'\otimes z'$ in
$\Pi_\infty\otimes\Pi_f$, then $z=\ell(y')z'$ for a linear form $\ell$ on
$\Pi_\infty$ with $\ell(y)=1$, as one sees by applying $\ell\otimes1$. In the
first case, $\Pi_\infty\neq0$ since $\Pi_0\neq0$; choose $P$ with
$q(P)\neq0$. Then $q(P)\otimes x=0=q(P)\otimes0$, so $x=0=0\cdot x'$. In the
second case $q(P)\neq0$, since $p_0(f_t\otimes P)\neq0$, and
$q(P)\otimes x=q'(P')\otimes\kappa_0x'$, so $x$ is a multiple of $x'$. In both
cases $r(f_t)=cx$ is a multiple of $r'(f'_t)$.
\end{proof}

\begin{theorem}\label{th-first-lift-nonzero}
Assume the published local inputs of the theta correspondence
(Theorem~\ref{lit-bundle-theta}), the agreement of global and local theta
integrals and the local matching (Definition~\ref{th-theta-hypotheses}), and the
following nonvanishing of $L$-values for all sufficiently large data: for all
$0\leq i,j<m$, every tuple $t$ has an $L$-function $\Lambda$ of
$\nu_{t,i,j}$ with $\Lambda(s)\neq0$ for the integers
$1-e_0(i,j)\leq s\leq h_1(i,j)$, and each of $\nu_{\mu_0,i,j}$,
$\nu_{\mu_0,i,j}\delta$ and $\nu_{\mu_0,i,j}\delta'$ has an $L$-function that
does not vanish at the integers $s\geq1$, where $\delta$ and $\delta'$ are the
characters of (L3) of the level system (Definition~\ref{th-grid}). Then for all
sufficiently large data
(Definition~\ref{th-large-data}), for every grid Weil datum, all $0\leq i,j<m$,
every bottom triple $(A,A',W)$ for $(i,j)$
(Definition~\ref{th-bottom-triple}) and every tuple $t$, there are $P$ and $h$
with $I_t(P,h)\neq0$ and there are $P'$ and $h'$ with $I'_t(P',h')\neq0$.
\end{theorem}
\leanref{Development/AutomorphicArithmetic/ThetaCorrespondence/ThetaPairLifts/LocalTypeDictionary/FirstOccurrenceThetaLiftNeZero.lean}{GlobalTheory.T.first_occurrence_theta_lift_ne_zero}
\leanref{Development/AbelianVarietyTheory/Statements/PlaceShiftStatements/NormSmall.lean}{GlobalTheory.SqTop.LNe}
\leanref{Development/AutomorphicArithmetic/ThetaCorrespondence/ThetaPairLifts/LocalTypeDictionary/FirstOccurrenceThetaLiftNeZero.lean}{GlobalTheory.NonVanBottom.bottom_theta_lift_ne_zero}
\leanref{Development/AutomorphicArithmetic/ThetaCorrespondence/ThetaPairLifts/LocalTypeDictionary/BottomThetaNormEqProdLocalFactors.lean}{GlobalTheory.BottomFactor.bottom_theta_norm_factorization}
\begin{proof}
Take the data large enough for Theorem~\ref{th-rallis-inner-product}, for the
hypothesis on $L$-values and for the factorization below. Fix a grid Weil
datum, $(i,j)$, a bottom triple and a tuple $t$.

By Theorem~\ref{th-rallis-inner-product} there are $P$, a Haar measure $\nu$, a
fundamental domain $\mathcal F$ and $c\neq0$ with
$\int_{\mathcal F}|I_t(P,h)|^2\,d\nu(h)=c\prod_{s=1-e_0}^{h_0}\Lambda(s)$ for
every $L$-function $\Lambda$ of $\nu_{t,i,j}$. By hypothesis some $\Lambda$ does
not vanish at the integers $s$ with $1-e_0\leq s\leq h_1$. Since
$a(i,j)\leq\operatorname{rk}V_{i+1,j+1}$ and
$a'(i,j)\leq\operatorname{rk}V_{i,j}$, we have $h_0\leq h_1$, so the product
is nonzero. Hence the integral is nonzero, and $I_t(P,h)\neq0$ for some $h$.

For the second lift, the two theta integrals of the triple factor through one
space: there are a complex vector space $\Pi_\infty\neq0$ and surjective linear
maps $q$ and $q'$ from the polynomials in the creation operators for $(-A,W)$
and $(A',W)$ to $\Pi_\infty$, and for each tuple $t$ a linear map $J_t$ from
$\Pi_\infty$ to the functions on $U(W)(\mathbf A)$, either $0$ or injective,
such that $I_t(P,\cdot)=\lambda_t\,J_t(q(P))$ for all $P$ and a constant
$\lambda_t$, and $I'_t(P',\cdot)=J_t(q'(P'))$ for all $P'$. We follow Alpöge's
construction. By Lemma~\ref{th-first-lift-inputs}(5) there are archimedean
data on $A$, $A'$ and $W$ and $a_\alpha,a_\beta\colon\mathcal T\to\mathbf Q$
such that $A$, $A'$ and $W$ are totally positive definite and, for every tuple
$t$, $(A,A',W;\alpha_t,\beta)$ is a square triple
(Definition~\ref{th-first-lift-matching}) in which $\alpha_t$ and $\beta$ have
archimedean exponents $a_\alpha$ and $a_\beta$, $f_t$ and $f'_t$ are
archimedean vacua, and $I_t(P,\cdot)=I(f_t\otimes P)$ and
$I'_t(P',\cdot)=I'(f'_t\otimes P')$ for the theta integrals $I=I_{\omega_t}$
and $I'=I_{\omega'_t}$ of the bottom triple. Let $\Pi_\infty$, $q$ and $q'$ be
given for these data by Lemma~\ref{th-first-lift-factor}(2); they do not
depend on $t$.
\begin{enumerate}
\item[(F1)] \emph{Matched local maps.} By Lemma~\ref{th-first-lift-factor} the
  square triple $(A,A',W;\alpha_t,\beta)$ has a matching, and then one
  $(\omega_t,\omega'_t,\Pi,\varrho,p,p')$ with the Weil data of the bottom
  triple. There are a complex vector space $\Pi_f$, linear maps $r$ and $r'$
  and an isomorphism $e\colon\Pi_\infty\otimes\Pi_f\to\Pi$ with
  $p(f\otimes P)=e(q(P)\otimes r(f))$ and
  $p'(f'\otimes P')=e(q'(P')\otimes r'(f'))$ for all archimedean vacua $f$ and
  $f'$. Since $\Pi\neq0$, $\Pi_\infty\neq0$.
\item[(F2)] \emph{Comparison of integrals.} By
  Lemma~\ref{th-first-lift-comparison} there is $c_0\neq0$ with
  $I(\Phi)=c_0\,I'(\Psi)$ whenever $p(\Phi)=p'(\Psi)$, and $I'(\Psi)=0$
  whenever $p'(\Psi)=0$. In particular, if $q'(P')=0$, then
  $p'(f'_t\otimes P')=0$ and $I'_t(P',\cdot)=0$.
\item[(F3)] \emph{Irreducibility.} Either $I'_t(P',\cdot)=0$ for all $P'$, or
  $I'_t(P',\cdot)=0$ implies $q'(P')=0$. Indeed, let $I'_t(P',\cdot)=0$ with
  $q'(P')\neq0$, and let $P''$ be a polynomial. If $r'(f'_t)=0$, then
  $p'(f'_t\otimes P'')=0$, and $I'_t(P'',\cdot)=0$ by (F2). Otherwise
  $p'(f'_t\otimes P')=e(q'(P')\otimes r'(f'_t))\neq0$, and
  Lemma~\ref{th-first-lift-irreducible} gives $I'_t(P'',\cdot)=0$.
\item[(F4)] \emph{The finite parts.} By Lemma~\ref{th-first-lift-finite},
  $r(f_t)=\kappa\,r'(f'_t)$ for some $\kappa\in\mathbf C$. If $q(P)=q'(P')$,
  then $p(f_t\otimes P)=e(q'(P')\otimes\kappa\,r'(f'_t))=\kappa\,p'(f'_t\otimes P')$,
  so by (F2) $I_t(P,\cdot)=c_0\kappa\,I'_t(P',\cdot)$.
\end{enumerate}
By (F2), $J_t(q'(P'))=I'_t(P',\cdot)$ defines a linear map on $\Pi_\infty$,
since $q'$ is surjective. By (F3) it is $0$ or injective. By (F4), with
$\lambda_t=c_0\kappa$ and $P'$ chosen with $q'(P')=q(P)$,
$I_t(P,\cdot)=\lambda_tJ_t(q(P))$. Since $\Pi_\infty\neq0$ and $q'$ is
surjective, $q'(P')\neq0$ for some $P'$.
% The assembly above follows GlobalTheory.BottomFactor.bottom_theta_norm_factorization.

Since $I_t(P,\cdot)\neq0$, the map $J_t$ is not $0$, so it is injective. Choose
$P'$ with $q'(P')\neq0$. Then $I'_t(P',\cdot)=J_t(q'(P'))\neq0$, so
$I'_t(P',h')\neq0$ for some $h'$.
\end{proof}

\section{Seesaws and Bessel periods}\label{section-th-seesaw}

\begin{lemma}\label{th-tensor-rule}
Let $D$ be a definite-side datum and $\psi$ a standard character
(Definition~\ref{th-weil-datum}), and let $S$ and $T$ be Hermitian spaces.
\begin{enumerate}
\item Let $T=A\oplus B$ be an orthogonal decomposition, so that
  $\mathbb W(S,T)=\mathbb W(S,A)\oplus\mathbb W(S,B)$. For functions $f_A$ on
  $\mathbb W(S,A)$ and $f_B$ on $\mathbb W(S,B)$ let $f_A\boxtimes f_B$ be the
  function $(w_1,w_2)\mapsto f_A(w_1)f_B(w_2)$. If $f_A$ and $f_B$ are theta
  functions, then so is $f_A\boxtimes f_B$, and
  $(f_A\boxtimes f_B)(0)=f_A(0)f_B(0)$. If $\omega_T$, $\omega_A$ and
  $\omega_B$ are Weil data for $(S,T)$, $(S,A)$ and $(S,B)$ with characters
  $(\chi_S,\chi_A\chi_B)$, $(\chi_S,\chi_A)$ and $(\chi_S,\chi_B)$, then for
  theta functions $f_A,f_B$, $g\in U(S)(\mathbf A)$, $a\in U(A)(\mathbf A)$ and
  $b\in U(B)(\mathbf A)$,
  \[
  \omega_T\bigl(g,(a,b)\bigr)(f_A\boxtimes f_B)=\omega_A(g,a)f_A\boxtimes\omega_B(g,b)f_B .
  \]
\item Let $S=A\oplus B$ be an orthogonal decomposition, so that
  $\mathbb W(S,T)=\mathbb W(A,T)\oplus\mathbb W(B,T)$, and define
  $f_A\boxtimes f_B$ in the same way. If $f_A$ and $f_B$ are theta functions,
  then so is $f_A\boxtimes f_B$, with value $f_A(0)f_B(0)$ at $0$. If $f$ is a
  theta function for $(S,T)$, then $u\mapsto f(u,v)$ is a theta function for
  $(A,T)$ for every $v$, and $v\mapsto f(u,v)$ is one for $(B,T)$ for every $u$.
  If $\omega_S$, $\omega_A$ and $\omega_B$ are Weil data for $(S,T)$, $(A,T)$
  and $(B,T)$ with characters $(\chi_A\chi_B,\chi_T)$, $(\chi_A,\chi_T)$ and
  $(\chi_B,\chi_T)$, then for every theta function $f$ for $(S,T)$,
  $a\in U(A)(\mathbf A)$, $b\in U(B)(\mathbf A)$, $h\in U(T)(\mathbf A)$ and
  $(u,v)\in\mathbb W(S,T)$,
  \[
  \bigl(\omega_S((a,b),h)f\bigr)(u,v)
  =\Bigl(\omega_A(a,h)\bigl[u'\mapsto\bigl(\omega_B(b,h)f(u',\cdot)\bigr)(v)\bigr]\Bigr)(u).
  \]
\end{enumerate}
\end{lemma}
\leanref{Development/AbelianVarietyTheory/Proofs/DecoratedLatticeProofs/ThetaTensorRuleTargetOrthogonalSum.lean}{GlobalTheory.TensorTargetRule.tensorTarget}
\leanref{Development/AbelianVarietyTheory/BlockTheory/TwistedDiscreteSeries/ConjugateAdjointDagger.lean}{GlobalTheory.IndefiniteRowSideFullHecke.ConjugateAdjointDagger.TensorTarget}
\leanref{Development/AbelianVarietyTheory/Proofs/DecoratedLatticeProofs/ThetaTensorRuleOrthogonalSumSource.lean}{GlobalTheory.TensorSrc.tensorSource}
\leanref{Development/AbelianVarietyTheory/Proofs/DecoratedLatticeProofs/WeilTensorRuleScalarEqOne.lean}{GlobalTheory.TensorScalarOne.tensor_scalar_one}
\begin{proof}
Consider first the decomposition $S=A\oplus B$ of part (2). The pairing
$\langle\ ,\ \rangle$ on $\mathbb W(S,T)$ is the orthogonal sum of the pairings
on $\mathbb W(A,T)$ and $\mathbb W(B,T)$, and the translations $\rho(w)$ act
coordinatewise. So the quasi-periodicity, invariance and smoothness conditions
of Definition~\ref{th-weil-datum}(2) pass from $f_A$ and $f_B$ to
$f_A\boxtimes f_B$, and from $f$ to its slices.

Fix $a$, $b$ and $h$, and let $N$ be the operator on theta functions for $(S,T)$
given by the right side of the displayed formula. Since $\omega_A(a,h)$ and
$\omega_B(b,h)$ preserve theta functions, so does $N$. The relation
$\omega(g,h)\rho(w_0)=\rho(hw_0g^{-1})\omega(g,h)$ of
Definition~\ref{th-weil-datum}(3) for $\omega_A$ and $\omega_B$ gives the same
relation for $N$ with the element $((a,b),h)$. Both $\omega_A(a,h)$ and
$\omega_B(b,h)$ preserve the $L^2$-norm on a fundamental domain, so $N$ does.
Hence $\omega_S((a,b),h)^{-1}N$ commutes with every $\rho(w)$ and preserves the
$L^2$-norm. By Schur's lemma for the translations $\rho(w)$ on theta functions,
it is a scalar of absolute value $1$. So there is $C(a,b,h)$ with $|C|=1$ and
$\omega_S((a,b),h)f=C(a,b,h)\,Nf$ for every theta function $f$.

Since both sides are homomorphisms in $(a,b,h)$, $C$ is multiplicative, and it
is continuous. For rational points $a$, $b$, $h$ the three Weil data satisfy
$(\omega(\gamma,\gamma')f)(0)=f(0)$, and evaluating at $0$ on a product with
$f(0)\neq0$ gives $C(a,b,h)=1$. For $a$, $b$, $h$ supported at a place split in
$E$, Definition~\ref{th-weil-datum}(3) makes the three Weil data act on
invariant linear functionals through the scalars
$r\,\chi_T(x)^{-1}\chi_S(y)$ with $r>0$; the characters multiply, so
$C(a,b,h)$ is a positive real number of absolute value $1$, and $C(a,b,h)=1$.
Thus $a\mapsto C(a,1,1)$ is a continuous unitary character of
$U(A)(\mathbf A)$, trivial on $U(A)(F)$ and on the elements supported at the
split places. Such a character is trivial. It is trivial on the elements with
finite component $1$, by weak approximation at the archimedean places; every
element of $U(A)(\mathbf A)$ is the product of a rational point, an element with
finite component $1$, an element supported at the split places and an element
of an arbitrarily small open subgroup; and the character is continuous. The
same holds for $b\mapsto C(1,b,1)$ and
$h\mapsto C(1,1,h)$, and
$C(a,b,h)=C(a,1,1)C(1,b,1)C(1,1,h)=1$. This proves part (2).

For part (1), let $x\mapsto x^\dagger$ be the adjoint, which identifies
$\mathbb W(S,T)$ with $\mathbb W(T,S)$. The map $f\mapsto\bar f\circ\dagger$
takes theta functions for $(S,T)$ to theta functions for $(T,S)$, and a Weil
datum for $(S,T)$ with characters $(\chi_S,\chi_T)$ to a Weil datum for
$(T,S)$ with characters $(\chi_T,\chi_S)$. It takes $f_A\boxtimes f_B$ for the
decomposition $T=A\oplus B$ of the second space to the product for the same
decomposition of the first space. Part (1) is therefore part (2) for the pair
$(T,S)$.
\end{proof}

\begin{lemma}\label{th-line-comparison}
Let $D$ be a definite-side datum with host element $\xi_0$, let $\psi$ be a
standard character (Definition~\ref{th-weil-datum}), and let $S$ be a Hermitian
space. Let $L$ and $L'$ be Hermitian spaces of rank $1$, and let
$\tau\colon L'\to L$ be a linear map with
$h_L(\tau x,\tau y)=\epsilon\,\xi_0\bar\xi_0\,h_{L'}(x,y)$ for a sign
$\epsilon=\pm1$. For $w\in\mathbb W(S,L')$ let $\kappa(w)=(\tau w)^*$ be the
adjoint of $\tau w$, an element of $\mathbb W(L,S)$.
\begin{enumerate}
\item Let $\epsilon=-1$, and let $\omega$ and $\omega'$ be Weil data for
  $(L,S)$ and $(S,L')$ with characters $(\chi_z,\chi)$ and
  $(\chi',\chi_z^{-1})$. For every theta function $f$ for $(L,S)$, $f\circ\kappa$
  is a theta function for $(S,L')$, and
  $\omega'(h,1)(f\circ\kappa)=(\omega(1,h)f)\circ\kappa$ for
  $h\in U(S)(\mathbf A)$.
\item Let $\epsilon=1$, and let $\omega$ and $\omega'$ be Weil data for
  $(L,S)$ and $(S,L')$ with characters $(\chi_z,\chi)$ and $(\chi',\chi_z)$.
  For every theta function $f$ for $(L,S)$, $\overline{f\circ\kappa}$ is a theta
  function for $(S,L')$, and
  $\omega'(h,1)\overline{f\circ\kappa}=\overline{(\omega(1,h)f)\circ\kappa}$.
\item Let $\epsilon=1$, and let $\zeta\in L$ and $\zeta'\in L'$ be vectors with
  $h(\zeta,\zeta)\neq0$, $h(\zeta',\zeta')\neq0$, $\tau\zeta'=-\zeta$ and
  $h_L(\zeta,\zeta)=\xi_0\bar\xi_0h_{L'}(\zeta',\zeta')$. Then every theta
  function $\Phi_0$ for $(S,L')$ is $\overline{\Xi\circ\kappa}$ for a theta
  function $\Xi$ for $(L,S)$.
\end{enumerate}
\end{lemma}
\leanref{Development/AbelianVarietyTheory/Proofs/DecoratedLatticeProofs/WeilRepSourceLineEqOppositeTargetLineStatement.lean}{GlobalTheory.T.weilRep_sourceLine_eq_opposite_targetLine}
\leanref{Development/AbelianVarietyTheory/Proofs/DecoratedLatticeProofs/WeilRepConjSourceLineEqTargetLine.lean}{GlobalTheory.ConjLine.conj_line}
\leanref{Development/AbelianVarietyTheory/Proofs/TwistedOperators/LinkedLineVectorsInvertibleOnThetaFunctionsStatement.lean}{GlobalTheory.T.linkedLineVectors_invertible_on_theta_functions}
\leanref{Development/AbelianVarietyTheory/Definitions/TwistedInnerProducts/TowerEx.lean}{GlobalTheory.Asc.Lnk}
\begin{proof}
The map $\kappa$ is additive and bijective, since $\tau$ is an isomorphism of
lines. It takes rational points to rational points and preserves the finite
and archimedean components. By the condition on $\tau$ it multiplies the
symplectic pairing by $-\epsilon$:
$\langle\kappa w,\kappa w'\rangle=-\epsilon\langle w,w'\rangle$, up to the
exchange of the two spaces in the pairing.

For (3), put $\Xi=\overline{\Phi_0\circ\kappa^{-1}}$. Complex conjugation
replaces $\psi$ by $\psi^{-1}$ in the quasi-periodicity of
Definition~\ref{th-weil-datum}(2), and $\kappa^{-1}$ changes the sign of the
pairing once more; so $\Xi$ satisfies the quasi-periodicity of a theta function
for $(L,S)$. The invariance under an open subgroup and the smoothness pass
through $\kappa$ because $\kappa$ is continuous and respects the finite and
archimedean components. Then $\Phi_0=\overline{\Xi\circ\kappa}$.

For (1) and (2), the same argument shows that $f\circ\kappa$, respectively
$\overline{f\circ\kappa}$, is a theta function for $(S,L')$. The maps
$h\mapsto[F\mapsto(\omega(1,h)(F\circ\kappa^{-1}))\circ\kappa]$, respectively
with complex conjugations, together with $\omega'$ on the first factor, form a
Weil datum for $(S,L')$ with the characters of $\omega'$: the conditions of
Definition~\ref{th-weil-datum}(3) transport along $\kappa$ and complex
conjugation, and conjugation inverts a unitary character. The transported
Weil datum and $\omega'$ have the same character on $L'$, and their characters
on $S$ restrict to the same character of $\mathbf A_F^\times$. Two Weil data for
the same pair with the same characters agree on theta functions; if the
character of the first space is multiplied by a character $\xi$ trivial on
$\mathbf A_F^\times$, the operator at $(g,h)$ is multiplied by the value of
$\xi$ at a determinant root of $h$, which is $1$ for $h=1$. So the two Weil
data agree at every $(h,1)$, which is the identity in (1) and (2).
\end{proof}

\begin{theorem}[Seesaw identity]\label{th-seesaw-identity}
Let $D$ be a definite-side datum (Definition~\ref{th-definite-datum}) with
host element $\xi_0$, and let $\psi$ be a character of $\mathbf A_F$. Let
$V_0,A,L,S,X,T,L''$ be Hermitian spaces over $E$ with orthogonal
decompositions $V_0=A\oplus L$ and $X=T\oplus L''$, where $L$ and $L''$ have
rank $1$, and let $\tau\colon L''\to L$ be a linear map with
$h_L(\tau x,\tau y)=-\xi_0\bar\xi_0\,h_{L''}(x,y)$. Let
$\chi_A,\chi_z,\chi_S,\chi'_S,\chi_T$ be unitary characters of
$\mathbf A_E^\times$, and let $\omega_1,\omega_A,\omega_z,\omega_X,\omega_T,\omega''_z$
be Weil data (Definition~\ref{th-weil-datum}) for the pairs $(V_0,S)$,
$(A,S)$, $(L,S)$, $(S,X)$, $(S,T)$, $(S,L'')$ with characters
$(\chi_A\chi_z,\chi_S)$, $(\chi_A,\chi_S)$, $(\chi_z,\chi_S)$,
$(\chi'_S,\chi_T\chi_z^{-1})$, $(\chi'_S,\chi_T)$, $(\chi'_S,\chi_z^{-1})$.
Let $\Xi$ be a theta function for $(V_0,S)$ and $\Phi$ one for $(S,T)$. For
$a\in\mathbb W(A,S)$ let $\Xi_a$ be the function
$w\mapsto\Xi(a,(\tau w)^*)$ on $\mathbb W(S,L'')$, where
$\mathbb W(V_0,S)=\mathbb W(A,S)\oplus\mathbb W(L,S)$ and $(\tau w)^*$ is the
adjoint of $\tau w$. Then:
\begin{enumerate}
\item each $\Xi_a$ is a theta function for $(S,L'')$;
\item for all $h_0\in U(A)(\mathbf A)$, $h\in U(S)(\mathbf A)$ and
  $h'\in U(T)(\mathbf A)$,
  \[
  \theta_{\omega_T}(h,h';\Phi)\,\theta_{\omega_1}\bigl((h_0,1),h;\Xi\bigr)
  =\Bigl(\omega_A(h_0,h)\bigl[a\mapsto\theta_{\omega_X}\bigl(h,(h',1);\Phi\boxtimes\Xi_a\bigr)\bigr]\Bigr)(0),
  \]
  where $(h_0,1)\in U(A)\times U(L)\subset U(V_0)$,
  $(h',1)\in U(T)\times U(L'')\subset U(X)$, and $\Phi\boxtimes\Xi_a$ is the
  function $(w_1,w_2)\mapsto\Phi(w_1)\Xi_a(w_2)$ on
  $\mathbb W(S,X)=\mathbb W(S,T)\oplus\mathbb W(S,L'')$.
\end{enumerate}
\end{theorem}
\leanref{Development/AbelianVarietyTheory/Proofs/DecoratedLatticeProofs/LastSeeSawThetaIdentityStatement.lean}{GlobalTheory.T.lastSeeSaw_thetaIdentity}
\leanref{Development/AbelianVarietyTheory/Definitions/LatticeFrames/IsQuasiFam.lean}{GlobalTheory.Unf.SeeLast}
\leanref{Development/AbelianVarietyTheory/Definitions/TwistedInnerProducts/ContinuousResTSymm.lean}{GlobalTheory.MixModel.OrthDec}
\leanref{Development/AbelianVarietyTheory/Proofs/DecoratedLatticeProofs/LastSeeSawThetaIdentity.lean}{GlobalTheory.SeeLastPf.seeLast}
\begin{proof}
The Weil data are standard, so $\psi$ is a standard character. By
Lemma~\ref{th-tensor-rule}(2) for $V_0=A\oplus L$, the function
$v\mapsto\Xi(a,v)$ is a theta function for $(L,S)$ for every $a$. By
Lemma~\ref{th-line-comparison}(1) for $\tau$, with $\omega_z$ and
$\omega''_z$, the function $\Xi_a$ is a theta function for $(S,L'')$, which is
(1), and
\[
\omega''_z(h,1)\Xi_a=\bigl(\omega_z(1,h)\Xi(a,\cdot)\bigr)\circ\kappa ,
\]
where $\kappa(w)=(\tau w)^*$. By Lemma~\ref{th-tensor-rule}(1) for
$X=T\oplus L''$, with $\omega_X$, $\omega_T$ and $\omega''_z$,
\[
\theta_{\omega_X}\bigl(h,(h',1);\Phi\boxtimes\Xi_a\bigr)
=\theta_{\omega_T}(h,h';\Phi)\,\bigl(\omega''_z(h,1)\Xi_a\bigr)(0)
=\theta_{\omega_T}(h,h';\Phi)\,\bigl(\omega_z(1,h)\Xi(a,\cdot)\bigr)(0),
\]
since $\kappa(0)=0$. By Lemma~\ref{th-tensor-rule}(2) for $V_0=A\oplus L$,
with $\omega_1$, $\omega_A$ and $\omega_z$ at $((h_0,1),h)$ and the point
$(0,0)$,
\[
\theta_{\omega_1}\bigl((h_0,1),h;\Xi\bigr)
=\Bigl(\omega_A(h_0,h)\bigl[a\mapsto\bigl(\omega_z(1,h)\Xi(a,\cdot)\bigr)(0)\bigr]\Bigr)(0).
\]
Multiply this by the scalar $\theta_{\omega_T}(h,h';\Phi)$ and move it inside
the linear operator $\omega_A(h_0,h)$. The function of $a$ inside is then the
one in (2), by the previous display.
\end{proof}

The identity expresses in terms of theta kernels the seesaw of
\cite{Kud84} between the pairs $(U(V_0),U(S))$ and $(U(S),U(X))$, with the
line $L$ of $V_0$ matched to the line $L''$ of $X$ through $\tau$.

\begin{lemma}\label{th-matrix-coefficients}
Let $D$ be a definite-side datum, $\psi$ a standard character and $S,T$
Hermitian spaces (Definition~\ref{th-weil-datum}). Write
$\mathbb W=\mathbb W(S,T)$ and
$\mathbb W(E)=M_{\operatorname{rk}T\times\operatorname{rk}S}(E)$ for its
rational points, fix a Haar measure on $\mathbb W$ and a fundamental domain
$\mathcal F$ of volume $1$ for $\mathbb W(E)$, and put
$\langle f,f'\rangle=\int_{\mathcal F}f(w)\overline{f'(w)}\,dw$. Let $f$ and
$f_0$ be theta functions, and put $c(w)=\langle f,\rho(w)f_0\rangle$ for
$w\in\mathbb W$.
\begin{enumerate}
\item Theta functions are bounded, $\rho(w)$ maps theta functions to theta
  functions, and $c$ is continuous and bounded.
\item There is an open subgroup $U\subset\mathbb W$ such that
  $\psi(\langle w,u\rangle)=1$ whenever $c(w)\neq0$ and $u\in U$ has
  archimedean component $0$.
\item For $\tau\in\mathcal T$ and
  $z\in M_{\operatorname{rk}T\times\operatorname{rk}S}(\mathbf C)$ let
  $z_\tau(t)\in\mathbb W$ be the element with component $tz$ at the
  archimedean place of $\tau$ and $0$ at the other places, and put
  $\ell_{\tau,z}(w)=\frac{d}{dt}\psi(\langle z_\tau(t),w\rangle)|_{t=0}$. For
  every $k\geq0$ there is $C$ with $|\ell_{\tau,z}(w)|^k|c(w)|\leq C$ for all
  $w$.
\item There is $C$ with $\sum_{\gamma\in\mathbb W(E)}|c(\gamma+w)|\leq C$ for
  every $w\in\mathbb W$.
\end{enumerate}
\end{lemma}
\leanref{Development/AbelianVarietyTheory/Proofs/TwistedOperators/ContinuousThetaMatrixCoeffHeisenbergTranslateStatement.lean}{GlobalTheory.T.continuous_thetaMatrixCoeff_heisenbergTranslate}
\leanref{Development/AbelianVarietyTheory/Proofs/TwistedOperators/ThetaMatrixCoeffEqZeroUnlessOrthogonalLevelStatement.lean}{GlobalTheory.T.thetaMatrixCoeff_eq_zero_unless_orthogonal_level}
\leanref{Development/AbelianVarietyTheory/Proofs/TwistedOperators/ThetaMatrixCoeffPolynomialDecayRealDirectionsStatement.lean}{GlobalTheory.T.thetaMatrixCoeff_polynomialDecay_realDirections}
\leanref{Development/AbelianVarietyTheory/Proofs/TwistedOperators/RationalTranslateSumBoundedRapidDecayStatement.lean}{GlobalTheory.T.rationalTranslateSum_bounded_rapidDecay}
\begin{proof}
Throughout we use the commutation relation
$\rho(u)\rho(w)=\psi(\langle u,w\rangle)\rho(w)\rho(u)$, which follows from
the definition of $\rho$ and $\langle w,u\rangle=-\langle u,w\rangle$.

For (1): $|\psi|=1$, so the quasi-periodicity of
Definition~\ref{th-weil-datum}(2) makes $|f|$ invariant under translation by
$\mathbb W(E)$. Since $f$ is continuous and $\mathbb W/\mathbb W(E)$ is
compact, $f$ is bounded. The function $\rho(w)f$ is quasi-periodic by the same
computation, it is invariant under $\rho(u)$ for the $u$ with archimedean
component $0$ in the intersection of the open subgroup of $f$ with the kernel
of $u\mapsto\psi(\langle u,w\rangle)$, by the commutation relation, and it is
smooth along the archimedean directions; so it is a theta function. The function
$(w,x)\mapsto f(x)\overline{(\rho(w)f_0)(x)}$ is continuous and bounded by
$\sup|f|\sup|f_0|$, and $\mathcal F$ has volume $1$. By dominated convergence
$c$ is continuous, and $|c|\leq\sup|f|\sup|f_0|$.

For (2): let $U$ be the intersection of the open subgroups of
Definition~\ref{th-weil-datum}(2) for $f$ and $f_0$. Let $u\in U$ have
archimedean component $0$, so that $\rho(u)f=f$ and $\rho(u)f_0=f_0$. The
operator $\rho(u)$ preserves $\langle\ ,\ \rangle$, so by the commutation
relation
\[
c(w)=\langle\rho(u)f,\rho(u)\rho(w)f_0\rangle
=\langle f,\psi(\langle u,w\rangle)\rho(w)f_0\rangle
=\overline{\psi(\langle u,w\rangle)}\,c(w).
\]
If $c(w)\neq0$, then $\psi(\langle u,w\rangle)=1$, hence
$\psi(\langle w,u\rangle)=1$.

For (3): for a theta function $F$ let $\partial F$ be the derivative at $t=0$
of $t\mapsto\rho(z_\tau(t))F$. It exists by the smoothness of theta functions,
$\partial F$ is again a theta function, and $\partial$ is skew-adjoint:
$\langle\partial F,F'\rangle=-\langle F,\partial F'\rangle$. Differentiating
$\rho(z_\tau(t))\rho(w)f_0=\psi(\langle z_\tau(t),w\rangle)\rho(w)\rho(z_\tau(t))f_0$
at $t=0$ gives
$\partial(\rho(w)f_0)=\ell_{\tau,z}(w)\rho(w)f_0+\rho(w)\partial f_0$. Pairing
with $f$ and using skew-adjointness,
\[
\overline{\ell_{\tau,z}(w)}\,\langle f,\rho(w)f_0\rangle
=-\langle\partial f,\rho(w)f_0\rangle-\langle f,\rho(w)\partial f_0\rangle .
\]
So $|\ell_{\tau,z}(w)|\,|c(w)|$ is at most the sum of the absolute values of the
coefficients of the pairs $(\partial f,f_0)$ and $(f,\partial f_0)$. Induction
on $k$, for all pairs of theta functions at once, with the bound of (1) for
$k=0$, gives (3).

For (4): let $|w_\infty|$ be the sum of the absolute values of the entries of
the archimedean components of $w$. For each $\tau$ there are finitely many $z$
and a constant $\kappa$ such that the size of the component of $w$ at $\tau$ is
at most $\kappa\sum_z|\ell_{\tau,z}(w)|$, since $|\ell_{\tau,z}(w)|$ is $2\pi$
times the absolute value of a real linear form in that component. So (3), with
the inequality between the power of a sum and the sum of powers, gives for
every $N$ a constant $C_N$ with $(1+|w_\infty|)^N|c(w)|\leq C_N$. By (2), the
finite components of the $w$ with $c(w)\neq0$ lie in the annihilator of an open
subgroup of the finite part of $\mathbb W$, which is compact. So there is a
compact $Q\subset\mathbb W$ such that $c(w)\neq0$ implies that $w-q$ has finite
component $0$ for some $q\in Q$, and
$|c(w)|\leq C_N\,m(w)$ with $m(w)=(1+|w_\infty|)^{-N}$ for such $w$ and
$m(w)=0$ otherwise. Choose a compact $K$ with $\mathbb W=\mathbb W(E)+K$. For
$N$ large, $b(\xi)=\sup_{k\in K}m(k+\xi)$ is summable over
$\xi\in\mathbb W(E)$: the $\xi$ with $b(\xi)\neq0$ have finite components in a
compact set, hence lie in a lattice of $\mathbb W(E)$, and on that lattice
$m(k+\xi)$ decays like a negative power of the archimedean size of $\xi$,
uniformly in $k$. Writing $w=\gamma_0+k$ with $\gamma_0\in\mathbb W(E)$ and
$k\in K$, we get
$\sum_\gamma|c(\gamma+w)|=\sum_\xi|c(k+\xi)|\leq C_N\sum_\xi b(\xi)$.
\end{proof}

\begin{lemma}\label{th-lattice-vacua}
Let $D$ be a definite-side datum, $\psi$ a standard character and $S,T$
Hermitian spaces (Definition~\ref{th-weil-datum}). Fix a Haar measure on
$\mathbb W=\mathbb W(S,T)$ and a fundamental domain $\mathcal F$ of volume $1$
for its rational points, and put
$\langle f,f'\rangle=\int_{\mathcal F}f(w)\overline{f'(w)}\,dw$.
\begin{enumerate}
\item For theta functions $f$ and $f_0$, the function
  $w\mapsto\langle f,\rho(w)f_0\rangle$ is integrable on $\mathbb W$, and
  \[
  \int_{\mathbb W}\langle f,\rho(w)f_0\rangle\,(\rho(w)f_0)(x)\,dw=\langle f_0,f_0\rangle\,f(x)
  \qquad(x\in\mathbb W).
  \]
\item If $f$ is a nonzero theta function, then $\langle f,f\rangle\neq0$.
\item Let $\Gamma\subset\mathbb W$ be a subgroup, with $\Gamma^\perp$ as in
  Lemma~\ref{th-vacuum-irreducibility}. If $f$ is a lattice vacuum for $\Gamma$
  and $w\in\Gamma^\perp$, then $\rho(w)f$ is a lattice vacuum for $\Gamma$. If
  $f$ and $f_0$ are lattice vacua for $\Gamma$ and $w\notin\Gamma^\perp$, then
  $\langle f,\rho(w)f_0\rangle=0$.
\end{enumerate}
\end{lemma}
\leanref{Development/AbelianVarietyTheory/Proofs/TwistedOperators/MoyalOrthogonalityHeisenbergThetaFunctionsStatement.lean}{GlobalTheory.T.moyal_orthogonality_heisenberg_thetaFunctions}
\leanref{Development/AbelianVarietyTheory/Proofs/TwistedOperators/TameFunctionalVacuumIrreducibility.lean}{GlobalTheory.Irr.thPair_self_ne_zero}
\leanref{Development/AbelianVarietyTheory/Proofs/TwistedOperators/HeisenbergVacuumCommutantPreservesLatticeVacuaStatement.lean}{GlobalTheory.T.heisenbergVacuumCommutant_preserves_latticeVacua}
\begin{proof}
For (1), put $c(w)=\langle f,\rho(w)f_0\rangle$ and write $\mathbb W(E)$ for
the rational points. By Lemma~\ref{th-matrix-coefficients}(1) and~(4), $c$ is
continuous and
$\int_{\mathbb W}|c|=\int_{\mathcal F}\sum_{\gamma}|c(\gamma+w)|\,dw$ is finite,
so $c$ is integrable. Fix $x\in\mathbb W$. The function
$w\mapsto(\rho(w)f_0)(x)$ is continuous and bounded by $\sup|f_0|$, so the
integral over $\mathbb W$ unfolds along $\mathbb W(E)$:
\[
\int_{\mathbb W}c(w)\,(\rho(w)f_0)(x)\,dw
=\int_{\mathcal F}\sum_{\gamma\in\mathbb W(E)}c(\gamma+w)\,(\rho(\gamma+w)f_0)(x)\,dw .\tag{$*$}
\]
Fix $w$ and put $g=\rho(w)f_0$, a theta function. For $\gamma\in\mathbb W(E)$
let $\chi_\gamma(y)=\psi(\langle y,\gamma\rangle)$, a continuous character of
$\mathbb W$ trivial on $\mathbb W(E)$. The quasi-periodicity of $f_0$ gives
$(\rho(\gamma+w)f_0)(y)=\psi(-\tfrac12\langle\gamma,w\rangle)\chi_\gamma(y)g(y)$.
The function $F=f\bar g$ is continuous and invariant under $\mathbb W(E)$, since
the factors of quasi-periodicity of $f$ and $\bar g$ cancel. Hence
\[
c(\gamma+w)=\psi(\tfrac12\langle\gamma,w\rangle)\int_{\mathcal F}\overline{\chi_\gamma(y)}F(y)\,dy,
\qquad
(\rho(\gamma+w)f_0)(x)=\psi(-\tfrac12\langle\gamma,w\rangle)\chi_\gamma(x)g(x).
\]
Put $\Phi(y)=F(y+x)$ and $\hat\Phi(a)=\int_{\mathcal F}\chi_a(y)\Phi(y)\,dy$. The
integral over $\mathcal F$ of an invariant function does not change under
translation, so the product of the two displayed factors is
$g(x)\hat\Phi(-\gamma)$, and $|\hat\Phi(-\gamma)|=|c(\gamma+w)|$ is summable by
Lemma~\ref{th-matrix-coefficients}(4). The characters $\chi_a$, $a\in\mathbb W(E)$,
are the characters of the compact group $\mathbb W/\mathbb W(E)$: they are
multiplicative in $a$, and $y\in\mathbb W(E)$ if and only if $\chi_a(y)=1$ for
all $a\in\mathbb W(E)$ \cite{Tat67}. Fourier inversion for a continuous
function with absolutely summable Fourier coefficients on a compact abelian
group of volume $1$ gives $\sum_a\hat\Phi(a)=\Phi(0)=f(x)\overline{g(x)}$. So
the sum in $(*)$ is $f(x)|g(x)|^2=f(x)|f_0(x+w)|^2$, and
\[
\int_{\mathbb W}c(w)\,(\rho(w)f_0)(x)\,dw=f(x)\int_{\mathcal F}|f_0(x+w)|^2\,dw
=\langle f_0,f_0\rangle\,f(x),
\]
since $|f_0|^2$ is invariant under $\mathbb W(E)$. This is the orthogonality
relation of Moyal for the Heisenberg group \cite[Chapter~1]{Fol89}, in the
model of quasi-periodic functions, where the integral over the group becomes
an integral over $\mathbb W$ and the inner product an integral over
$\mathcal F$.

For (2): $|\psi|=1$, so the quasi-periodicity of $f$ makes $|f|$ invariant under
rational translations. If $\langle f,f\rangle=0$, the continuous function $f$
vanishes almost everywhere on $\mathcal F$, hence on $\mathbb W$, hence
everywhere.

For (3), recall that $\rho(u)\rho(w)=\psi(\langle w,u\rangle)\rho(w)\rho(u)$ up
to the order of the arguments of the pairing, and that the operators $\rho(u)$
preserve $\langle\ ,\ \rangle$. If $w\in\Gamma^\perp$ and $u\in\Gamma$, the
scalar for $\rho(2u)$, and for $\rho(u)$ when $u$ lies in the half attached to a
split place, is $1$; so $\rho(w)f$ has the invariance of a lattice vacuum. If
$w\notin\Gamma^\perp$, there is such an operator $\rho(u')$, with
$u'=2u$ or $u'=u$, fixing $f$ and $f_0$ and with
$c=\psi(\langle u',w\rangle)\neq1$. Then
$\langle f,\rho(w)f_0\rangle=\langle\rho(u')f,\rho(u')\rho(w)f_0\rangle
=\bar c\,\langle f,\rho(w)f_0\rangle$, so $\langle f,\rho(w)f_0\rangle=0$.
\end{proof}

\begin{lemma}\label{th-vacuum-irreducibility}
Let $D$ be a definite-side datum, $\psi$ a standard character and $S,T$
Hermitian spaces, and let $\Gamma\subset\mathbb W(S,T)$ be a subgroup
(Definition~\ref{th-weil-datum}). Let $\Gamma^\perp$ be the set of $w$ with
$\psi(\langle w,2u\rangle)=1$ for $u\in\Gamma$ and $\psi(\langle w,u\rangle)=1$
for those $u\in\Gamma$ in the half of $\mathbb W$ attached to a split place.
Call a map $\ell$ from functions on $\mathbb W$ to $\mathbf C$ \emph{tame} if
it is additive and $\mathbf C$-linear on theta functions, and for every theta
function $f$ the function $w\mapsto\ell(\rho(w)f)$ is continuous and bounded
and $\ell(\int\alpha(w)\rho(w)f\,dw)=\int\alpha(w)\ell(\rho(w)f)\,dw$ for
integrable $\alpha$ for which the left side is a theta function. Let $\ell$ be
tame and let $f_0\neq0$ be a lattice vacuum for $\Gamma$ with
$\ell(\rho(w)f_0)=0$ for all $w\in\Gamma^\perp$. Then $\ell(f)=0$ for every
lattice vacuum $f$ for $\Gamma$.
\end{lemma}
\leanref{Development/AbelianVarietyTheory/Proofs/TwistedOperators/TameFunctionalVacuumIrreducibilityStatement.lean}{GlobalTheory.T.tame_functional_vacuum_irreducibility}
\leanref{Development/AbelianVarietyTheory/Definitions/TwistedInnerProducts/TowerEx.lean}{GlobalTheory.Asc.Irr}
\leanref{Development/AbelianVarietyTheory/Definitions/TwistedInnerProducts/TowerEx.lean}{GlobalTheory.Asc.Tame}
\leanref{Development/AbelianVarietyTheory/Proofs/TwistedOperators/TameFunctionalVacuumIrreducibility.lean}{GlobalTheory.Irr.irr}
\begin{proof}
Let $f$ be a lattice vacuum for $\Gamma$, and use the pairing of
Lemma~\ref{th-lattice-vacua} for a Haar measure and a fundamental domain of
volume $1$. Put $\alpha(w)=\langle f,\rho(w)f_0\rangle$, which is integrable by
Lemma~\ref{th-lattice-vacua}(1). By the same lemma,
$\int\alpha(w)\rho(w)f_0\,dw=\langle f_0,f_0\rangle f$, a theta function. Since
$\ell$ is tame,
\[
\langle f_0,f_0\rangle\,\ell(f)=\int_{\mathbb W}\alpha(w)\,\ell(\rho(w)f_0)\,dw .
\]
For $w\in\Gamma^\perp$ the factor $\ell(\rho(w)f_0)$ vanishes by hypothesis.
For $w\notin\Gamma^\perp$ the factor $\alpha(w)$ vanishes by
Lemma~\ref{th-lattice-vacua}(3). So the integral is $0$. Since
$\langle f_0,f_0\rangle\neq0$ by Lemma~\ref{th-lattice-vacua}(2), $\ell(f)=0$.
\end{proof}

\begin{definition}\label{th-heights}
Let $D$ be a definite-side datum with host field $E/F$ and set $\mathcal T$ of
complex embeddings (Definition~\ref{th-definite-datum}), and let $V$ be a
Hermitian space of rank $n$ (Definition~\ref{th-spaces}).
\begin{enumerate}
\item For $g\in GL_n(\mathbf A_E)$ and a place $v$ of $F$ let $\|g\|_v$ be the
  operator norm of $g_v$ at an archimedean $v$, and the largest absolute value
  of the entries of $g_v$ at a finite $v$, the absolute value on
  $E\otimes_FF_v$ being the spectral norm. The \emph{height} of
  $g\in U(V)(\mathbf A)$ is
  $\operatorname{ht}_0(g)=\prod_v\max(1,\|g\|_v,\|g^{-1}\|_v)$, and its
  \emph{height on the quotient} is
  $\operatorname{ht}(g)=\inf_{\gamma\in U(V)(F)}\operatorname{ht}_0(\gamma g)$.
\item For $a\in U(V)(\mathbf A)$ with finite component $1$ put
  $\|a\|_\infty=\sum_{\tau\in\mathcal T}\sum_{i,j}|\tau(a_{ij})|$.
\item The \emph{size} of $x\in\mathbf A_E^n$ is
  $|x|_{\mathbf A}=\prod_v\max_i|x_i|_v$, with the norms of (1). For a nonzero
  $\xi\in E^n$ and $g\in U(V)(\mathbf A)$ the vector $g^{-1}\xi$ is the vector
  $\xi$ \emph{seen from} $g$.
\end{enumerate}
\end{definition}
\leanref{Development/AbelianVarietyTheory/Definitions/TwistedInnerProducts/ExistsHeightLeTwoMulQuotHeight.lean}{GlobalTheory.Grow.height}
\leanref{Development/AbelianVarietyTheory/Definitions/TwistedInnerProducts/ExistsHeightLeTwoMulQuotHeight.lean}{GlobalTheory.Grow.quotHeight}
\leanref{Development/AbelianVarietyTheory/Definitions/Foundations/EventuallyNormFinLeOne.lean}{GlobalTheory.AdelicU.heightGL}
\leanref{Development/AbelianVarietyTheory/Proofs/TwistedOperators/VsizeNonneg.lean}{GlobalTheory.HtOrth.vsize}
\leanref{Development/AbelianVarietyTheory/Proofs/TwistedOperators/VsizeNonneg.lean}{GlobalTheory.HtOrth.seen}

The height is at least $1$, submultiplicative, and finite because $g_v$ is
integral and unitary at almost every $v$. Since $\operatorname{ht}_0$ is
continuous and $U(V)(F)$ is discrete, for every $g$ there is a rational
$\gamma$ with $\operatorname{ht}_0(\gamma g)\leq2\operatorname{ht}(g)$.

\begin{lemma}\label{th-archimedean-growth}
Let $D$ be a definite-side datum, $\psi$ a standard character, $S,T$
Hermitian spaces and $\omega$ a Weil datum for $(S,T)$
(Definition~\ref{th-weil-datum}). There is $N$ such that for every theta
function $f$ for $(S,T)$ and all $x\in U(S)(\mathbf A)$ and
$y\in U(T)(\mathbf A)$ there is $C$ with
\[
|(\omega(xa,yb)f)(w)|\leq C\bigl((1+\|a\|_\infty)(1+\|b\|_\infty)\bigr)^N
\]
for all $a\in U(S)(\mathbf A)$ and $b\in U(T)(\mathbf A)$ with finite
component $1$ and all $w\in\mathbb W(S,T)$, with the notation of
Definition~\ref{th-heights}.
\end{lemma}
\leanref{Development/AbelianVarietyTheory/Proofs/IwasawaDecomposition/WeilRepThetaFunctionPolynomialGrowthRealPoints.lean}{GlobalTheory.Mem.LiftBound.om_growth}
\leanref{Development/AbelianVarietyTheory/Proofs/IwasawaDecomposition/ExistsOscillatorActionUnitaryPolynomialBoundStatement.lean}{GlobalTheory.T.hasDerivAt_oscillatorAction_unitaryPath_quadratic}
\leanref{Development/AbelianVarietyTheory/Proofs/IwasawaDecomposition/WeakWeilRepEqUnitSmulOscillatorRealPointsStatement.lean}{GlobalTheory.T.weakWeilRep_eq_unit_smul_oscillator_real_points}
\begin{proof}
Let $\mathbb W_\infty$ be the real vector space of elements of
$\mathbb W(S,T)$ with finite component $0$, with the symplectic form
$\operatorname{Im}$ of the Hermitian form attached to $\langle\ ,\ \rangle$,
and let $G_\infty$ be the group of pairs $(a,b)$ with finite components $1$,
acting on $\mathbb W_\infty$ by the symplectic maps $m(a,b)\colon w\mapsto bwa^{-1}$.

\emph{The oscillator representation.} Let $B$ be the Banach space of bounded
continuous functions on $\mathbb W(S,T)$ with the supremum norm. Choose a complex
structure $J$ on $\mathbb W_\infty$ compatible with the symplectic form, with
the Euclidean norm $|u|$ it defines, and let $P$ be the Gaussian average
$PF=\int_{\mathbb W_\infty}e^{-\pi|u|^2/2}\rho(u)F\,du$. Alpöge defines, for
$g\in G_\infty$ and $F\in B$,
\[
\Omega(g)F=c(g)\int_{\mathbb W_\infty}\rho(m(g)w)\int_{\mathbb W_\infty}e^{-\pi|u|^2/2}\rho(m(g)u)\,P\rho(-w)F\,du\,dw,
\]
where, in coordinates in which the Hermitian form whose imaginary part is the
symplectic form is a diagonal sign matrix $s$,
$c(g)=\det\bigl(\tfrac12(1+s)m(g)+\tfrac12(1-s)\bigr)$, the determinant of the
block of $m(g)$ on the positive part. This is the oscillator representation in
the Fock model \cite[Chapter~4]{Fol89}: on
smooth vectors $F$, those for which $w\mapsto\rho(w)F$ is smooth, $\Omega$
is a homomorphism and $\Omega(g)\rho(w)=\rho(m(g)w)\Omega(g)$. The integral
is bounded by $\|\Omega(g)F\|\leq|c(g)|K(F)$, with
$K(F)=\int e^{-\pi|u|^2/2}du\int\|P\rho(-w)F\|\,dw$, which is finite for smooth
$F$; and $|c(g)|\leq L(1+\sum_{i,j}|m(g)_{ij}|)^n$, with $n$ the dimension of
$\mathbb W_\infty$ over $\mathbf C$ for $J$ and a constant $L$, since $c(g)$ is a
determinant of size $n$ whose entries are linear in those of $m(g)$.

\emph{Comparison with $\omega$.} For $(a,b)\in G_\infty$ there is $c$ with
$|c|=1$ and $\omega(a,b)F=c\,\Omega(a,b)F$ for every theta function $F$. Both
sides satisfy the relation $\omega(g,h)\rho(w)=\rho(hwg^{-1})\omega(g,h)$ of
Definition~\ref{th-weil-datum}(3) and preserve the $L^2$-norm on a fundamental
domain, so this is the argument with Schur's lemma for the translations of
Lemma~\ref{th-tensor-rule}.

\emph{The bound.} Write $x=k\,a_0$ and $y=k'b_0$ with $k,k'$ of archimedean
component $1$ and $a_0,b_0$ of finite component $1$; elements with
archimedean component $1$ commute with those with finite component $1$. Then
\[
\omega(xa,yb)f=\omega(a_0a,b_0b)\,\omega(k,k')f .
\]
The function $f'=\omega(k,k')f$ is a theta function, hence smooth for the
archimedean translations and in $B$. By the comparison,
$|(\omega(a_0a,b_0b)f')(w)|\leq\|\Omega(a_0a,b_0b)f'\|\leq|c(a_0a,b_0b)|K(f')$.
The sum of the absolute values of the entries of $m(a',b')$ is at most
$K(1+\|a'\|_\infty)(1+\|b'\|_\infty)$ for a constant $K$, and
$1+\|a_0a\|_\infty\leq(1+\|a_0\|_\infty)(1+\|a\|_\infty)$. So
$|c(a_0a,b_0b)|\leq C'((1+\|a\|_\infty)(1+\|b\|_\infty))^n$ with $C'$
depending on $a_0$ and $b_0$, and the lemma holds with $N=n$.
\end{proof}

\begin{definition}\label{th-sized-lattice}
Let $E$ be a CM field. For $a\in E$ let $\operatorname{house}(a)$ be the largest
absolute value of the images of $a$ under the embeddings $E\to\mathbf C$. Let
$V$ be a finite-dimensional $E$-vector space with a nondegenerate Hermitian form
$h$, linear in the first variable, with $h(y,x)=\overline{h(x,y)}$ (for the
form $h$ of Definition~\ref{th-spaces}, take $(x,y)\mapsto h(y,x)$). Let
$\Lambda\subset V$ be a finitely generated subgroup, $N$ an integer, $\nu$ a
real function on $V$ and $B$ a real number. Then $\Lambda$ is
\emph{$(N,B)$-sized} for $\nu$ if:
\begin{enumerate}
\item $\Lambda$ is stable under multiplication by the algebraic integers of $E$,
  and $h(x,y)$ is an algebraic integer for all $x,y\in\Lambda$;
\item $Ny\in\Lambda$ for every $y\in V$ such that $h(y,x)$ is an algebraic
  integer for every $x\in\Lambda$;
\item $\nu\geq0$, and $\nu(x+y)\leq\nu(x)+\nu(y)$,
  $\nu(ax)\leq\operatorname{house}(a)\nu(x)$ and
  $\operatorname{house}(h(x,y))\leq\nu(x)\nu(y)$ for all $a\in E$ and
  $x,y\in V$;
\item $B\geq2$, and $V$ is spanned over $E$ by the $x\in\Lambda$ with
  $\nu(x)\leq B$.
\end{enumerate}
\end{definition}
\leanref{Development/AbelianVarietyTheory/Proofs/HeightProofs/IsLieDataCircle.lean}{GlobalTheory.SmoothForms.Sized}

Since $B\geq2$, a constant $K$ is at most $B^{k_0}$ for every $k_0$ with
$2^{k_0}\geq K$. The bounds below absorb constants into powers of $B$ in this
way. An algebraic integer $a\neq0$ of $E$ with $\operatorname{house}(a)\leq U$,
where $U\geq1$, satisfies $|\sigma(a)|\,U^{[E:\mathbf Q]}\geq1$ for every embedding
$\sigma$, since $|N_{E/\mathbf Q}(a)|\geq1$ is the product of $|\sigma(a)|$ and
$[E:\mathbf Q]-1$ factors at most $U$. Hence
$\operatorname{house}(a^{-1})\leq U^{[E:\mathbf Q]}$.

\begin{lemma}\label{th-anisotropic-frame}
Let $E$ be a CM field. For every $n$ there is $k$ such that for every
$E$-vector space $V$ of dimension $n$ with a Hermitian form, and every
$(N,B)$-sized $\Lambda\subset V$ for a function $\nu$
(Definition~\ref{th-sized-lattice}), there are $v_1,\ldots,v_n\in\Lambda$ with
$h(v_i,v_j)=0$ for $i\neq j$, $h(v_i,v_i)\neq0$ and $\nu(v_i)\leq B^k$.
\end{lemma}
\leanref{Development/AbelianVarietyTheory/Proofs/HeightProofs/IsLieDataCircle.lean}{GlobalTheory.SmoothForms.FrameSmall}
\leanref{Development/AbelianVarietyTheory/Proofs/HeightProofs/OrthogonalAnisotropicFrameBoundedCMStatement.lean}{GlobalTheory.T.orthogonalAnisotropicFrame_bounded_CM}
\leanref{Development/AbelianVarietyTheory/Proofs/HeightProofs/OrthogonalAnisotropicFrameBoundedCM.lean}{GlobalTheory.SmoothForms.frameSmall}
\leanref{Development/AbelianVarietyTheory/Proofs/HeightProofs/OrthogonalAnisotropicFrameBoundedCM.lean}{GlobalTheory.SmoothForms.exists_small_anisotropic_mem}
\leanref{Development/AbelianVarietyTheory/Proofs/HeightProofs/OrthogonalAnisotropicFrameBoundedCM.lean}{GlobalTheory.SmoothForms.exists_small_frame_mem}
\begin{proof}
We prove more. Put $k_0=0$ and $k_{n+1}=10k_n+4$. Let $W\subset V$ be a subspace
of dimension $n$ on which $h$ is nondegenerate, and $C\geq2$ such that $W$ is
spanned by the $x\in\Lambda\cap W$ with $\nu(x)\leq C$, the \emph{spanning
vectors}. We show that there are $v_1,\ldots,v_n\in\Lambda\cap W$, pairwise
orthogonal, with $h(v_i,v_i)\neq0$ and $\nu(v_i)\leq C^{k_n}$. The lemma is the
case $W=V$ and $C=B$.

\emph{An anisotropic vector.} If $W\neq0$, there is $v\in\Lambda\cap W$ with
$h(v,v)\neq0$ and $\nu(v)\leq C^4$. Indeed, some spanning vector $x$ is nonzero,
and since $h$ is nondegenerate on $W$, some spanning vector $y$ has
$h(x,y)\neq0$. If $h(x,x)\neq0$ or $h(y,y)\neq0$, take $v=x$ or $v=y$.
Otherwise put $c=h(x,y)$ and $v=x+cy$, which lies in $\Lambda$ because $c$ is an
algebraic integer. Then $h(v,v)=\bar c\,h(x,y)+c\,h(y,x)=2c\bar c\neq0$, and
$\nu(v)\leq\nu(x)+\operatorname{house}(c)\nu(y)\leq C+C^3\leq C^4$.

\emph{Induction on $n$.} The case $n=0$ is empty. Let $W$ have dimension $n+1$,
and take $v$ as above. Let $W'=\{x\in W:h(x,v)=0\}$, of dimension $n$, and
$P(x)=h(v,v)x-h(x,v)v$. Then $P$ maps $W$ into $W'$ and $\Lambda$ into
$\Lambda$, and $P(x)=h(v,v)x$ for $x\in W'$. Since $h(v,v)\neq0$, $W'$ is spanned
by the vectors $P(x)$ for spanning vectors $x$ of $W$, and these satisfy
\[
\nu(P(x))\leq\operatorname{house}(h(v,v))\nu(x)+\operatorname{house}(h(x,v))\nu(v)
\leq C^8C+C\,C^4C^4\leq C^{10}.
\]
The form is nondegenerate on $W'$: if $y\in W'$ is orthogonal to $W'$, then
$\overline{h(v,v)}\,h(y,x)=h(y,P(x))=0$ for every $x\in W$, so $y=0$. By
induction, with $C^{10}$ in place of $C$, there are pairwise orthogonal
$w_1,\ldots,w_n\in\Lambda\cap W'$ with $h(w_i,w_i)\neq0$ and
$\nu(w_i)\leq C^{10k_n}$. With $v$ they are the required vectors, since $v$ is
orthogonal to $W'$ and $C^4$ and $C^{10k_n}$ are at most $C^{10k_n+4}$.
\end{proof}

\begin{lemma}\label{th-short-vector}
Let $E$ be a CM field. For every $n\geq1$ and every nonzero integer $N$ there are
$C$ and $k$ such that for every $E$-vector space $V$ of dimension $n$ with a
Hermitian form, and every $(N,B)$-sized $\Lambda\subset V$ for a function $\nu$
(Definition~\ref{th-sized-lattice}), there is $v\in\Lambda$ with $v\neq0$,
$\operatorname{house}(h(v,v))\leq C$ and $\nu(v)\leq B^k$.
\end{lemma}
\leanref{Development/AbelianVarietyTheory/Proofs/HeightProofs/IsLieDataCircle.lean}{GlobalTheory.SmoothForms.ShortSmall}
\leanref{Development/AbelianVarietyTheory/Proofs/HeightProofs/ShortLatticeVectorBoundedSizeOfOrthogonalFrameStatement.lean}{GlobalTheory.T.shortLatticeVectorBoundedSize_of_orthogonalFrame}
\leanref{Development/AbelianVarietyTheory/Proofs/HeightProofs/ShortLatticeVectorBoundedSizeOfOrthogonalFrame.lean}{GlobalTheory.SmoothForms.Short.shortSmall}
\leanref{Development/AbelianVarietyTheory/Proofs/HeightProofs/ShortLatticeVectorBoundedSizeOfOrthogonalFrame.lean}{GlobalTheory.SmoothForms.Short.exists_short_coord}
\leanref{Development/AbelianVarietyTheory/PublishedTheorems/AdelicLattices₂/ExistsShortOfTraceDet.lean}{GlobalTheory.HermClass.det_ne_zero_and_abs_le}
\leanref{Development/AbelianVarietyTheory/PublishedTheorems/AdelicLattices₂/ExistsShortOfTraceDet.lean}{GlobalTheory.HermClass.exists_int_vecMul_lt}
\begin{proof}
Let $F$ be the maximal totally real subfield of $E$ and $d=[E:\mathbf Q]$. Fix a
nonzero algebraic integer $\delta\in E$ with $\bar\delta=-\delta$, and a nonzero
integer $D_E$ such that $D_Ea$ is an algebraic integer whenever $a\in E$ and
$\operatorname{Tr}_{E/\mathbf Q}(at)\in\mathbf Z$ for every algebraic integer
$t$ of $E$. Put $M=2ND_E$ and $t_0=|M|+1$, let $k_0$ be the exponent of
Lemma~\ref{th-anisotropic-frame} for $n$, and fix $m$ with
$2^m\geq2nt_0+1+\operatorname{house}(\delta)$. We take $C=2nt_0^2$.

\emph{An orthogonal basis over $F$.} By Lemma~\ref{th-anisotropic-frame} there
are pairwise orthogonal $e_1,\ldots,e_n\in\Lambda$ with $h(e_i,e_i)\neq0$ and
$\nu(e_i)\leq B^{k_0}$. Regard $V$ as an $F$-vector space of dimension $2n$ with
the symmetric $F$-bilinear form $q(x,y)=h(x,y)+h(y,x)$. The $2n$ vectors
$e_i$ and $\delta e_i$, which we call $g$, lie in $\Lambda$ and are pairwise
orthogonal for $q$: for instance
$q(e_i,\delta e_i)=(\bar\delta+\delta)h(e_i,e_i)=0$. Since
$q(g,g)=2h(g,g)\neq0$, they form an $F$-basis of $V$; write $c_g(x)\in F$ for
the coordinates of $x$ in it. Also
$\nu(g)\leq(1+\operatorname{house}(\delta))B^{k_0}\leq B^{m+k_0}$.

\emph{The discriminant.} The group $\Lambda$ is free and spans $V$ over
$\mathbf Q$, since it spans $V$ over $E$ and is stable under the algebraic
integers. Let $b_1,\ldots,b_r$ be a basis of $\Lambda$ over $\mathbf Z$, so
$r=dn$, and let $T$ be the rational $r\times r$ matrix with
$T_{ij}=\operatorname{Tr}_{F/\mathbf Q}q(b_i,b_j)=\operatorname{Tr}_{E/\mathbf Q}h(b_i,b_j)$.
If $x\in\mathbf Q^r$ and $xT\in\mathbf Z^r$, then $y=\sum_ix_ib_i$ has
$\operatorname{Tr}_{E/\mathbf Q}h(y,w)\in\mathbf Z$ for $w\in\Lambda$. Applied to
$\bar sw$ for algebraic integers $s$, this gives
$\operatorname{Tr}_{E/\mathbf Q}(s\,h(y,w))\in\mathbf Z$, so $D_Eh(y,w)$ is an
algebraic integer for every $w\in\Lambda$. By
Definition~\ref{th-sized-lattice}(2), $ND_Ey\in\Lambda$, and $Mx$ is integral.
Hence $T$ is invertible: if $xT=0$, then $Mcx$ is integral for every rational
$c$, so $x=0$. The rows of $T^{-1}$ are the $x$ with $xT$ a unit vector, so
$MT^{-1}$ is an integer matrix, and $|\det(MT^{-1})|\geq1$ gives
$0<|\det T|\leq|M|^r$.

\emph{Minkowski's theorem.} Let $\varphi$ run through the $[F:\mathbf Q]$ real
embeddings of $F$. Since
$\operatorname{Tr}_{F/\mathbf Q}q(x,y)=\sum_\varphi\sum_g\varphi(c_g(x))\varphi(c_g(y))\varphi(q(g,g))$,
the matrix $T$ is $A'\Delta A'^{\mathsf T}$, where $A'$ has entries
$A'_{i,(\varphi,g)}=\varphi(c_g(b_i))$ and $\Delta$ is diagonal with entries
$\varphi(q(g,g))$. Put $A=A'|\Delta|^{1/2}$. Then $(\det A)^2=|\det T|$, which
is nonzero and less than $t_0^{2r}$. By Minkowski's linear forms theorem
\cite{Cas59}, a real $r\times r$ matrix $A$ with $0<|\det A|<t_0^r$ has a nonzero
$z\in\mathbf Z^r$ with $|\sum_iz_iA_{i,(\varphi,g)}|<t_0$ for every column
$(\varphi,g)$. The vector $v=\sum_iz_ib_i\in\Lambda$ is nonzero, and the
inequalities say $\varphi(c_g(v))^2|\varphi(q(g,g))|<t_0^2$ for all $\varphi$ and
$g$. Summing over the $2n$ vectors $g$ gives
$|\varphi(q(v,v))|\leq2nt_0^2$. As $h(v,v)=q(v,v)/2$ lies in $F$,
$\operatorname{house}(h(v,v))\leq C$.

\emph{The size of $v$.} Each $q(g,g)$ is a nonzero algebraic integer with
$\operatorname{house}(q(g,g))\leq2\nu(g)^2\leq B^{1+2(m+k_0)}$. Put
$U=B^{(1+2(m+k_0))d}$. Then $|\varphi(q(g,g))|U\geq1$, so
$\varphi(c_g(v))^2\leq t_0^2U$ and $\operatorname{house}(c_g(v))\leq t_0U$.
Hence
\[
\nu(v)\leq\sum_g\operatorname{house}(c_g(v))\,\nu(g)\leq2nt_0U\,B^{m+k_0}\leq
B^m\,U\,B^{m+k_0},
\]
and the lemma holds with $k=2m+k_0+(1+2(m+k_0))d$.
\end{proof}

\begin{lemma}\label{th-anisotropic-vector}
Let $E$ be a CM field. For every $n\geq1$ and every nonzero integer $N$ there are
a finite set $S\subset E^\times$ and $k$ such that for every $E$-vector space $V$
of dimension $n$ with a Hermitian form, and every $(N,B)$-sized $\Lambda\subset V$
for a function $\nu$ (Definition~\ref{th-sized-lattice}), there is $v\in\Lambda$
with $h(v,v)\in S$ and $\nu(v)\leq B^k$.
\end{lemma}
\leanref{Development/AbelianVarietyTheory/Proofs/HeightProofs/IsLieDataCircle.lean}{GlobalTheory.SmoothForms.AnisoSmall}
\leanref{Development/AbelianVarietyTheory/Proofs/HeightProofs/AnisotropicFiniteLengthsOfShortLatticeVectorStatement.lean}{GlobalTheory.T.anisotropicFiniteLengths_of_shortLatticeVector}
\leanref{Development/AbelianVarietyTheory/Proofs/HeightProofs/AnisotropicFiniteLengthsOfShortLatticeVector.lean}{GlobalTheory.AnisoSize.anisoSmall_of_shortSmall}
\leanref{Development/AbelianVarietyTheory/Proofs/HeightProofs/AnisotropicFiniteLengthsOfShortLatticeVector.lean}{GlobalTheory.AnisoSize.exists_small_generator}
\leanref{Development/AbelianVarietyTheory/Proofs/HeightProofs/AnisotropicFiniteLengthsOfShortLatticeVector.lean}{GlobalTheory.AnisoSize.exists_small_lift}
\leanref{Development/AbelianVarietyTheory/Proofs/HeightProofs/AnisotropicFiniteLengthsOfShortLatticeVector.lean}{GlobalTheory.AnisoSize.exists_near_integer}
\leanref{Development/AbelianVarietyTheory/PublishedTheorems/AdelicLattices₂/ExistsIntMulIsIntegralOfTrace.lean}{GlobalTheory.HermClass.exists_finite_residues}
\begin{proof}
Let $d=[E:\mathbf Q]$, $d_E$ the discriminant of $E$ and $\mathcal O_E$ its ring of
integers. We fix the following, which depend only on $E$, $n$ and $N$.
\begin{enumerate}
\item[(a)] $C$ and $k_1$ as in Lemma~\ref{th-short-vector}.
\item[(b)] The integer $M=\lceil|d_E|^{1/2}\rceil!$. Every nonzero ideal $J$ of
  $\mathcal O_E$ contains $j\neq0$ with
  $\operatorname{house}(j)\leq|d_E|^{1/2}\,\mathrm NJ$ and $j\mid Ma$ for every
  $a\in J$. Indeed, Minkowski's theorem \cite[Chapter~I, \S6]{Neu99} gives $j\in J$,
  $j\neq0$, with $\operatorname{house}(j)^d\leq|d_E|^{1/2}\mathrm NJ$, and the
  right side is at least $1$. Then $(j)=JI$ with
  $\mathrm NI\leq|N_{E/\mathbf Q}(j)|/\mathrm NJ\leq|d_E|^{1/2}$, so $\mathrm NI$
  divides $M$ and $M\in I$, and $MJ\subset IJ=(j)$.
\item[(c)] A constant $C_2$ such that every $c\in E$ has an algebraic integer $s$
  with $\operatorname{house}(c-s)\leq C_2$: round the coordinates of $c$ in an
  integral basis.
\item[(d)] A finite set $R$ of representatives of $L/2NML$, where $L$ is the
  group of algebraic integers $a\in E$ with $\bar a=a$, which is finitely
  generated.
\item[(e)] The set $S$ of the nonzero elements of $R$, of $\{2NM\}$, and of the
  algebraic integers of $E$ all of whose conjugates have absolute value at most
  $C$. It is finite, since the characteristic polynomials of the last have
  bounded integer coefficients.
\end{enumerate}
Let $\Lambda$ be given. Every constant above is at most a fixed power of $B$, and
each bound below is a power $B^{k'}$ with $k'$ depending only on $E$, $n$ and
$N$. Let $v$ be as in Lemma~\ref{th-short-vector}. If $h(v,v)\neq0$, then
$h(v,v)$ is an algebraic integer whose conjugates are at most $C$, so it lies in
$S$, and we take $v$. Suppose now $h(v,v)=0$.

\emph{The ideal of $v$.} The values $h(x,v)$, $x\in\Lambda$, form an ideal $J$ of
$\mathcal O_E$, since $\Lambda$ is stable under $\mathcal O_E$. As $v\neq0$ and
$V$ is spanned by the $x\in\Lambda$ with $\nu(x)\leq B$, one of these, $x_0$, has
$a=h(x_0,v)\neq0$. Then $\operatorname{house}(a)\leq B^{1+k_1}$ and
$\mathrm NJ\leq|N_{E/\mathbf Q}(a)|\leq B^{(1+k_1)d}$. Take $j\in J$ as in (b),
and $w_0\in\Lambda$ with $h(w_0,v)=j$.

\emph{A small lift.} There is $w\in\Lambda$ with $h(w,v)=j$ and $\nu(w)$ at most
a power of $B$. Let $x_1,\ldots,x_n\in\Lambda$ with $\nu(x_i)\leq B$ be a basis of
$V$, and put $p_i=ax_i-h(x_i,v)x_0\in\Lambda$, so that $h(p_i,v)=0$ and
$\nu(p_i)\leq2B^{2+k_1}$. Writing $w_0=\sum_ic_ix_i$ with $c_i\in E$, one has
$j=\sum_ic_ih(x_i,v)$ and so
$w_0=(j/a)x_0+\sum_i(c_i/a)p_i$. Choose algebraic integers $s_i$ with
$\operatorname{house}(c_i/a-s_i)\leq C_2$ and put $w=w_0-\sum_is_ip_i\in\Lambda$.
Then $h(w,v)=j$ and $w=(j/a)x_0+\sum_i(c_i/a-s_i)p_i$, whose size is bounded by
$\operatorname{house}(j)\operatorname{house}(a^{-1})B+nC_2\max_i\nu(p_i)$, a power
of $B$ by the bound for $\operatorname{house}(a^{-1})$ after
Definition~\ref{th-sized-lattice}.

\emph{A hyperbolic pair.} Put $v'=N\,\overline{(M/j)}\,v$. For $x\in\Lambda$,
$h(\overline{(M/j)}\,v,x)=\overline{(M/j)h(x,v)}$ is an algebraic integer,
because $j$ divides $Mh(x,v)$; by Definition~\ref{th-sized-lattice}(2),
$v'\in\Lambda$. Moreover $h(w,v')=NM$, $h(v',v')=0$, and
$\nu(v')\leq\operatorname{house}(NM)\operatorname{house}(j^{-1})\nu(v)$ is at most
a power of $B$. For $s\in E$ with $\bar s=s$,
\[
h(w-sv',w-sv')=h(w,w)-2NMs .
\]

\emph{The norm.} Since $h(w,w)\in L$, write $h(w,w)=\rho+2NMy$ with $\rho\in R$
and $y\in L$. If $\rho\neq0$, put $s=y$; then $h(w-sv',w-sv')=\rho\in S$. If
$\rho=0$, put $s=y-1$; then $h(w-sv',w-sv')=2NM\in S$. In both cases
$\operatorname{house}(s)\leq\operatorname{house}((2NM)^{-1})\bigl(\nu(w)^2+\operatorname{house}(\rho)\bigr)+1$,
so $\nu(w-sv')\leq\nu(w)+\operatorname{house}(s)\nu(v')$ is at most a power of
$B$, and $w-sv'\in\Lambda$ is the required vector.
\end{proof}

\begin{lemma}\label{th-orthogonal-frame}
Let $E$ be a CM field. For every $n$ and every nonzero integer $N$ there are a
finite set $S_0\subset E^\times$ and $k$ such that for every $E$-vector space $V$
of dimension $n$ with a Hermitian form, and every $(N,B)$-sized $\Lambda\subset V$
for a function $\nu$ (Definition~\ref{th-sized-lattice}), there are
$v_1,\ldots,v_n\in\Lambda$ with $h(v_i,v_j)=0$ for $i\neq j$,
$h(v_i,v_i)\in S_0$ and $\nu(v_i)\leq B^k$.
\end{lemma}
\leanref{Development/AbelianVarietyTheory/Proofs/HeightProofs/IsLieDataCircle.lean}{GlobalTheory.SmoothForms.OrthoSmall}
\leanref{Development/AbelianVarietyTheory/Proofs/HeightProofs/OrthogonalFrameOfAnisotropicFiniteLengthsStatement.lean}{GlobalTheory.T.orthogonalFrame_of_anisotropicFiniteLengths}
\leanref{Development/AbelianVarietyTheory/Proofs/HeightProofs/OrthogonalFrameOfAnisotropicFiniteLengths.lean}{GlobalTheory.SmoothForms.orthoSmall_of_anisoSmall}
\begin{proof}
By induction on $n$; for $n=0$ there is nothing to prove. Let $S_1$ and $k_1$ be
as in Lemma~\ref{th-anisotropic-vector} for $n+1$ and $N$, let $M\neq0$ be an
integer such that $M/a$ is an algebraic integer for every $a\in S_1$, and let
$S_2$ and $k_2$ be as in the lemma for $n$ and $MN$. We show that $S_0=S_1\cup S_2$
and $k=k_1+4(k_1+1)k_2$ work for $n+1$.

Let $\Lambda\subset V$ be given, with $\dim V=n+1$. By
Lemma~\ref{th-anisotropic-vector} there is $v\in\Lambda$ with $a=h(v,v)\in S_1$
and $\nu(v)\leq B^{k_1}$; then $a\neq0$ and $\bar a=a$. Let
$W=\{x\in V:h(x,v)=0\}$, of dimension $n$, and $p(x)=ax-h(x,v)v$, a map from $V$
onto $W$ with $p(\Lambda)\subset\Lambda$ and $p(x)=ax$ for $x\in W$; for $y\in W$,
$h(y,p(x))=a\,h(y,x)$. Put $\Lambda'=\Lambda\cap W$ and $T=B^{k_1+1}$. Then
$\Lambda'$ is $(MN,T^4)$-sized for $\nu$ in $W$:
\begin{enumerate}
\item the form is nondegenerate on $W$, since $y\in W$ orthogonal to $W$ has
  $a\,h(y,x)=h(y,p(x))=0$ for every $x\in V$; conditions (1) and (3) are
  inherited from $\Lambda$;
\item if $y\in W$ and $h(y,x')$ is an algebraic integer for all $x'\in\Lambda'$,
  then $h(y,p(x))=a\,h(y,x)$ is one for all $x\in\Lambda$, so $Nay\in\Lambda$, and
  $MNy=(M/a)(Nay)\in\Lambda\cap W$;
\item for $x\in\Lambda$ with $\nu(x)\leq B$,
  $\nu(p(x))\leq\operatorname{house}(a)\nu(x)+\operatorname{house}(h(x,v))\nu(v)\leq T^2T+T\,T\,T\leq T^4$,
  and these $p(x)$ span $W$, since $p$ maps $V$ onto $W$.
\end{enumerate}
By induction there are pairwise orthogonal $u_1,\ldots,u_n\in\Lambda'$ with
$h(u_i,u_i)\in S_2$ and $\nu(u_i)\leq T^{4k_2}$. The vectors $v,u_1,\ldots,u_n$
are pairwise orthogonal, with norms in $S_0$ and sizes at most
$B^{k_1+4(k_1+1)k_2}$.
\end{proof}

\begin{lemma}[Reduction]\label{th-reduction}
Let $D$ be a definite-side datum and $V$ a Hermitian space, with the notation
of Definition~\ref{th-heights}. There is $M$ such that for every open subgroup
$U\subset U(V)(\mathbf A)$ there are a finite set $\Sigma\subset U(V)(\mathbf A)$
and $C\geq0$ with the following property: every $g\in U(V)(\mathbf A)$ satisfies
$\gamma g=x\,a\,k$ for some $\gamma\in U(V)(F)$, $x\in\Sigma$, $a$ with finite
component $1$ and $1+\|a\|_\infty\leq C\operatorname{ht}_0(g)^M$, and $k\in U$
with archimedean component $1$.
\end{lemma}
\leanref{Development/AbelianVarietyTheory/Proofs/IwasawaDecomposition/AdelicHeightGrowthOfCosetGrowthUnitaryPair.lean}{GlobalTheory.Mem.LiftGrow.reduction_unif}
\leanref{Development/AbelianVarietyTheory/Proofs/HeightProofs/UnitaryReductionPolynomialHeightBoundStatement.lean}{GlobalTheory.T.unitary_reduction_polynomial_height_bound}
\leanref{Development/AbelianVarietyTheory/Proofs/HeightProofs/ExistsFiniteCosetsAdelicUnitaryBoundedDenominatoStatement.lean}{GlobalTheory.T.exists_finite_cosets_adelicUnitary_bounded_denominator}
\begin{proof}
Let $L_0$ be the lattice of $V$, and for $g\in U(V)(\mathbf A)$ let
$L_g=E^n\cap g\widehat L_0$, a lattice of $V$, self-dual up to the fixed
scalar of the Gram matrix.

\emph{Step 1: denominators.} There are $C_1$ and $M_1$ such that for every $g$
there is an integer $1\leq d\leq C_1\operatorname{ht}_0(g)^{M_1}$ with
$de_i\in L_g$ for the standard basis vectors $e_i$. Indeed $d$ can be taken to
be a fixed integer times a common denominator of the coordinates of the
entries of the finite component of $g^{-1}$, and the product over the finite
places of the local denominators is at most a power of the finite part of
$\operatorname{ht}_0(g)$.

\emph{Step 2: orthogonal frames.} Fix nonzero integers $N_0$ and $N$ such that,
for every $g$, the form $(x,y)\mapsto N_0\,x^{\mathsf T}G_V\bar y$ is integral on
$L_g$ and $N$ times the dual of $L_g$ for it lies in $L_g$, and a constant $K$
for which $\nu(x)=K\sum_i\operatorname{house}(x_i)$ satisfies
Definition~\ref{th-sized-lattice}(3). Since $de_i\in L_g$ by Step 1, $L_g$ is
$(N,B)$-sized for this form and $\nu$ (Definition~\ref{th-sized-lattice}) with
$B=2Kd$. By
Lemma~\ref{th-orthogonal-frame} there are a finite set $S_0\subset E^\times$ and
$k$, independent of $g$, such that $L_g$ contains a basis $v_1,\ldots,v_n$ of
$V$, orthogonal for this form, with norms in $S_0$ and $\nu(v_i)\leq B^k$.

\emph{Step 3: a rational element.} Apply Step 2 to $L_g$ and $d$ from Step 1.
Since $S_0$ is finite, the diagonal Gram matrices of the frames run through a
finite set, and for each of them fix one frame of $V$ with that Gram matrix.
The $E$-linear map $\gamma'$ taking the frame of $L_g$ to the fixed frame of
the same Gram matrix is an isometry, so it lies in $U(V)(F)$, and its
archimedean size is at most a constant times that of the frame of $L_g$,
hence at most $C_2\operatorname{ht}_0(g)^{M_1k}$. Put $\gamma=\gamma'$. The
lattice $L_{\gamma g}=\gamma L_g$ is spanned by a fixed frame up to
the bounded index of a frame in its lattice, so the finite components of
$\gamma g$ and $(\gamma g)^{-1}$ have entries with denominators dividing a fixed
integer $e$.

\emph{Step 4: finitely many cosets.} The finite elements of $U(V)(\mathbf A)$
whose entries and those of their inverses have denominators dividing $e$ form a
compact set. It is covered by finitely many cosets $xU$ with $x$ of archimedean
component $1$; let $\Sigma$ be the set of these $x$. Write $\gamma g=k_0a$ with
$k_0$ of archimedean component $1$ and $a$ of finite component $1$, and
$k_0=xk$ with $x\in\Sigma$ and $k\in U$. Since $a$ commutes with $k$,
$\gamma g=xak$. Finally $\|a\|_\infty=\|\gamma g\|_\infty\leq
c\|\gamma\|_\infty\|g\|_\infty$, and $\|g\|_\infty$ is at most a constant times
$\operatorname{ht}_0(g)$, so $1+\|a\|_\infty\leq C\operatorname{ht}_0(g)^{M_1k+1}$.
\end{proof}

\begin{lemma}\label{th-height-comparison}
Let $D$ be a definite-side datum and $V$ a Hermitian space of rank $n$, with the
notation of Definition~\ref{th-heights}.
\begin{enumerate}
\item There are $C$ and $b$ such that if $0<t\leq1$, $g\in U(V)(\mathbf A)$ and
  $|g^{-1}\xi|_{\mathbf A}\geq t$ for every nonzero $\xi\in E^n$, then
  $\operatorname{ht}(g)\,t^b\leq C$.
\item There are $c>0$ and $N$ with
  $|g^{-1}\xi|_{\mathbf A}\operatorname{ht}(g)^N\geq c$ for all $g$ and all
  nonzero $\xi\in E^n$.
\item Let $S=A\oplus B$ be an orthogonal decomposition. There are $C$ and $a$
  with $\operatorname{ht}(h)\leq C\operatorname{ht}((h,1))^a$ and
  $\operatorname{ht}((h,1))\leq C\operatorname{ht}(h)^a$ for
  $h\in U(A)(\mathbf A)$.
\item Let $S=S_0\oplus H^\rho$ be a hyperbolic decomposition with flag group
  $N_S$ (Definition~\ref{th-step}) and embedding $\iota\colon U(S_0)\to U(S)$.
  There are $C$ and $a$ with $\operatorname{ht}(h)\leq C\operatorname{ht}(n\iota(h))^a$
  and $\operatorname{ht}(n\iota(h))\leq C\operatorname{ht}(h)^a$ for
  $n\in N_S(\mathbf A)$ and $h\in U(S_0)(\mathbf A)$.
\end{enumerate}
\end{lemma}
\leanref{Development/AbelianVarietyTheory/Proofs/TwistedOperators/UnitaryQuotientHeightLeOfNoShortVectorStatement.lean}{GlobalTheory.T.unitaryQuotientHeight_le_of_no_short_vector}
\leanref{Development/AbelianVarietyTheory/Proofs/TwistedOperators/LeAdelicVectorSizeProductMulUnitaryQuotientHeighStatement.lean}{GlobalTheory.T.le_adelicVectorSizeProduct_mul_unitaryQuotientHeight_pow}
\leanref{Development/AbelianVarietyTheory/Proofs/TwistedOperators/UnitaryQuotientHeightOrthogonalSumComparableStatement.lean}{GlobalTheory.T.unitaryQuotientHeight_orthogonalSum_comparable}
\leanref{Development/AbelianVarietyTheory/Proofs/TwistedOperators/UnitaryQuotientHeightComparableWittDecompositionStatement.lean}{GlobalTheory.T.unitaryQuotientHeight_comparable_wittDecomposition}
\begin{proof}
Part (1) is Mahler's compactness criterion \cite{Bor69} in the polynomial form
that reduction theory gives. A rational translate of $g$ lies in a Siegel set,
on which the height is bounded by a power of the torus coordinates, and these
are bounded by a power of $1/t$ when no nonzero rational vector seen from $g$
has size below $t$. Alpöge proves this form with his own reduction theory, as
the note explains.

For (2), choose a rational $\gamma$ with
$\operatorname{ht}_0(\gamma g)\leq2\operatorname{ht}(g)$. Then
$g^{-1}\xi=(\gamma g)^{-1}\gamma\xi$, and $\gamma\xi$ is a nonzero vector of
$E^n$. At each place
$\max_i|(\gamma\xi)_i|_v\leq c_v\|\gamma g\|_v\max_i|(g^{-1}\xi)_i|_v$, with
$c_v=1$ at the finite places, and by the product formula
$\prod_v\max_i|(\gamma\xi)_i|_v\geq1$. Hence
$|g^{-1}\xi|_{\mathbf A}\geq c\operatorname{ht}_0(\gamma g)^{-1}\geq
\tfrac c2\operatorname{ht}(g)^{-1}$.

For the upper bound in (3), choose a rational $\gamma$ with
$\operatorname{ht}_0(\gamma h)\leq2\operatorname{ht}(h)$. The element
$(\gamma,1)$ is rational, and the height of $(g,1)$ is at most
$C_2\operatorname{ht}_0(g)^{a_2}$, since the entries of $(g,1)$ are linear in
those of $g$ with rational coefficients. So
$\operatorname{ht}((h,1))\leq\operatorname{ht}_0((\gamma h,1))\leq
C_2(2\operatorname{ht}(h))^{a_2}$. For the lower bound, put
$H=\operatorname{ht}((h,1))$ and let $\eta\in E^{\operatorname{rk}A}$ be nonzero.
The vector $\eta$ seen from $h$ and its image $\iota\eta$ in $S$ seen from
$(h,1)$ satisfy $|(h,1)^{-1}\iota\eta|_{\mathbf A}\leq C_3|h^{-1}\eta|_{\mathbf A}$.
By (2) for $S$, $c\leq|(h,1)^{-1}\iota\eta|_{\mathbf A}H^N$, so
$|h^{-1}\eta|_{\mathbf A}\geq t=c/(KH^N)$ with $K=\max(C_3,c,1)$, and $t\leq1$.
By (1) for $A$, $\operatorname{ht}(h)\leq C_1t^{-b}=C_1(K/c)^bH^{Nb}$.

For (4): the embedding satisfies $\operatorname{ht}_0(\iota(h))\leq
C\operatorname{ht}_0(h)$; every $n\in N_S(\mathbf A)$ is $\gamma^{-1}n_1$ with
$\gamma\in N_S(F)$ and $\operatorname{ht}_0(n_1)\leq C_1$, since
$N_S(F)\backslash N_S(\mathbf A)$ is compact; and $\iota(U(S_0)(F))$ normalizes
$N_S$. Given $n$ and $h$, choose a rational $\delta$ with
$\operatorname{ht}_0(\delta h)\leq2\operatorname{ht}(h)$ and then $\gamma\in N_S(F)$
with $n'=\gamma\iota(\delta)n\iota(\delta)^{-1}$ of height at most $C_1$. With
$h'=\delta h$ we have $n'\iota(h')=\gamma\iota(\delta)\,n\iota(h)$, so the two have
the same height on the quotient, and
$\operatorname{ht}(n\iota(h))\leq\operatorname{ht}_0(n')\operatorname{ht}_0(\iota(h'))
\leq2CC_1\operatorname{ht}(h)$. For the lower bound, let $e_1,\ldots,e_\rho$ be the
isotropic vectors of the decomposition and $\eta\in E^{\operatorname{rk}S_0}$
nonzero, and choose a rational $\gamma$ with
$\operatorname{ht}_0(\gamma n\iota(h))\leq2\operatorname{ht}(n\iota(h))$. Some
$(\rho+1)\times(\rho+1)$ minor $\Delta$ of the matrix $\gamma(e_1,\ldots,e_\rho,\iota\eta)$
is a nonzero element of $E$. Since $n$ fixes the $e_a$ modulo the span of the
earlier ones and fixes $\iota(S_0)$ modulo the span of the $e_a$, $\Delta$ is
also that minor of $\gamma n\iota(h)\,(e_1,\ldots,e_\rho,\iota(h^{-1}\eta))$. By
the product formula $|\Delta|_{\mathbf A}\geq1$, and expanding the minor,
$|\Delta|_{\mathbf A}\leq C\operatorname{ht}_0(\gamma n\iota(h))^{\rho+1}|h^{-1}\eta|_{\mathbf A}$.
Hence $|h^{-1}\eta|_{\mathbf A}\geq c'\operatorname{ht}(n\iota(h))^{-\rho-1}$
for every nonzero $\eta$, and (1) for $S_0$ gives
$\operatorname{ht}(h)\leq C\operatorname{ht}(n\iota(h))^{(\rho+1)b}$.
\end{proof}

\begin{lemma}\label{th-weil-estimates}
Let $D$ be a definite-side datum, $\psi$ a standard character, $S,T$
Hermitian spaces, $\omega$ a Weil datum for $(S,T)$
(Definition~\ref{th-weil-datum}) and $f$ a theta function for $(S,T)$. For a
Hermitian space $V$ write $\operatorname{ht}$ for the height on
$U(V)(F)\backslash U(V)(\mathbf A)$ (Definition~\ref{th-heights}).
\begin{enumerate}
\item For $g\in U(S)(\mathbf A)$, $h\in U(T)(\mathbf A)$ and
  $u,x\in\mathbb W(S,T)$,
  $(\omega(g,h)\rho(u)f)(x)=\psi(\tfrac12\langle x,hug^{-1}\rangle)\,(\omega(g,h)f)(x+hug^{-1})$.
\item The function $(g,h,x)\mapsto(\omega(g,h)f)(x)$ is continuous.
\item There are $C$ and $A$ with
  $|(\omega(g,h)f)(x)|\leq C\,(\operatorname{ht}(g)\operatorname{ht}(h))^A$ for
  all $g$, $h$ and $x$.
\item Let $Y$ be a second countable space with an $s$-finite Borel measure $m$,
  let $F_y$ ($y\in Y$) be theta functions for $(S,T)$, and let $(g,h)$ be
  fixed. Suppose that $(y,x)\mapsto F_y(x)$ and
  $(y,x)\mapsto(\omega(g,h)F_y)(x)$ are continuous, that both are bounded in
  absolute value by $\beta(y)$ for an integrable function $\beta$, and that
  $F=\int_YF_y\,dm(y)$ is a theta function. Then
  $\omega(g,h)F=\int_Y\omega(g,h)F_y\,dm(y)$.
\item Let $X=T\oplus L$ be an orthogonal decomposition. There are $C$ and $a$
  with $\operatorname{ht}((h,1))\leq C\operatorname{ht}(h)^a$ and
  $\operatorname{ht}(h)\leq C\operatorname{ht}((h,1))^a$ for
  $h\in U(T)(\mathbf A)$. Let $V=V_0\oplus H^\rho$ be a hyperbolic decomposition
  with flag group $N_V$ (Definition~\ref{th-step}) and
  $\iota\colon U(V_0)\to U(V)$ the embedding. There are $C$ and $a$ with
  $\operatorname{ht}(n\iota(g))\leq C\operatorname{ht}(g)^a$ and
  $\operatorname{ht}(g)\leq C\operatorname{ht}(n\iota(g))^a$ for
  $n\in N_V(\mathbf A)$ and $g\in U(V_0)(\mathbf A)$.
\end{enumerate}
\end{lemma}
\leanref{Development/AbelianVarietyTheory/Proofs/DecoratedLatticeProofs/WeilDatumThetaJointlyContinuous.lean}{GlobalTheory.Unf.jointCont_all}
\leanref{Development/AbelianVarietyTheory/Proofs/TwistedOperators/WeilRepresentationThetaFunctionSupLeQuotientHeigStatement.lean}{GlobalTheory.T.weilRepresentation_thetaFunction_sup_le_quotientHeight}
\leanref{Development/AbelianVarietyTheory/Proofs/DecoratedLatticeProofs/WeilOperatorIntegralCommOfBoundedFiniteMeasure.lean}{GlobalTheory.CommIntBdd.om_integral_comm_dom}
\leanref{Development/AbelianVarietyTheory/Proofs/TwistedOperators/WeilKernelFunctionalSuperpositionIsTame.lean}{GlobalTheory.KerTame.om_heis_apply}
\begin{proof}
Part (1) is the relation $\omega(g,h)\rho(u)=\rho(hug^{-1})\omega(g,h)$ of
Definition~\ref{th-weil-datum}(3), evaluated at $x$.

For (2), fix $(g_0,h_0,x_0)$. For $(g,h)$ near $(g_0,h_0)$ the theta functions
$\omega(g,h)f$ are invariant under one open subgroup of the elements with
archimedean component $0$, and their derivatives along the archimedean
directions are bounded uniformly. Write $x=x_0+x_\infty+u$ with $x_\infty$
archimedean and $u$ finite and small. The invariance removes $u$ up to a
factor of $\psi$ that is continuous in $x$, the derivative bound controls the
translation by $x_\infty$, and $(\omega(g,h)f)(x_0)$ is continuous in $(g,h)$ by
Definition~\ref{th-weil-datum}(3).

For (3) we first bound by the heights $\operatorname{ht}_0$ of
Definition~\ref{th-heights}. For rational $\gamma,\gamma'$ and a theta function
$F$,
\[
(\omega(\gamma,\gamma')F)(x)=(\rho(x)\omega(\gamma,\gamma')F)(0)
=(\omega(\gamma,\gamma')\rho(\gamma'^{-1}x\gamma)F)(0)=F(\gamma'^{-1}x\gamma),
\]
by Definition~\ref{th-weil-datum}(3). Put
$K(g,h)=\sup_x|(\omega(g,h)f)(x)|$, which is finite by
Lemma~\ref{th-matrix-coefficients}(1). Then:
\begin{enumerate}
\item[(a)] $K(\gamma g,\gamma'h)=K(g,h)$ for rational $\gamma,\gamma'$, by the
  last display applied to $F=\omega(g,h)f$.
\item[(b)] There are open subgroups $U$ and $U'$ of $U(S)(\mathbf A)$ and
  $U(T)(\mathbf A)$ with $\omega(k,k')f=f$ for $k\in U$ and $k'\in U'$ of
  archimedean component $1$, so $K(gk,hk')=K(g,h)$. Indeed, let $c_F(w)=\langle F,\rho(w)F\rangle$
  be the diagonal matrix coefficient of Lemma~\ref{th-matrix-coefficients}. By
  part (2) of that lemma, $c_f(w)\neq0$ only if the finite component of $w$
  lies in a compact set, and $f$ is invariant under the translations by an open
  subgroup of the elements with archimedean component $0$; so
  $c_f(kwk'^{-1})=c_f(w)$ for all $w$ and all $(k,k')$ of archimedean
  component $1$ in a neighbourhood of $1$, for which $kwk'^{-1}-w$ lies in that
  subgroup when $w$ is in the compact set. By the relation of
  Definition~\ref{th-weil-datum}(3) and since $\omega(k,k')$ preserves the
  pairing, $\omega(k,k')f$ has the same diagonal matrix coefficient as $f$, so
  $\omega(k,k')f=c\,f$ with $|c|=1$: by the orthogonality relation of
  Lemma~\ref{th-lattice-vacua}(1), a theta function is determined by its
  diagonal matrix coefficient up to such a scalar. Taking the neighbourhood so
  small that $|(\omega(k,k')f)(w)/f(w)-1|<\tfrac12$ at a point $w$ with
  $f(w)\neq0$, and $U$, $U'$ open subgroups inside it, the powers $c^n$ all
  satisfy $|c^n-1|<\tfrac12$, so $c=1$.
\item[(c)] By Lemma~\ref{th-archimedean-growth}, $K(xa,yb)\leq
  C_{x,y}((1+\|a\|_\infty)(1+\|b\|_\infty))^N$ for $a,b$ of finite component $1$.
\end{enumerate}
By Lemma~\ref{th-reduction} for $U$ and $U'$, write $\gamma g=xak$ and
$\gamma'h=ybk'$ with $x\in\Sigma$, $y\in\Sigma'$ and
$1+\|a\|_\infty\leq C\operatorname{ht}_0(g)^M$,
$1+\|b\|_\infty\leq C'\operatorname{ht}_0(h)^{M'}$. By (a), (b) and (c),
\[
|(\omega(g,h)f)(x')|\leq K(xa,yb)\leq\max_{x,y}C_{x,y}\,(CC')^N\bigl(\operatorname{ht}_0(g)\operatorname{ht}_0(h)\bigr)^{(M+M')N}
\]
for all $x'$. Now choose rational $\gamma,\gamma'$ with
$\operatorname{ht}_0(\gamma g)\leq2\operatorname{ht}(g)$ and
$\operatorname{ht}_0(\gamma'h)\leq2\operatorname{ht}(h)$. Then
$|(\omega(g,h)f)(x)|=|(\omega(\gamma g,\gamma'h)f)(\gamma'x\gamma^{-1})|\leq
4^AC(\operatorname{ht}(g)\operatorname{ht}(h))^A$, where $C$ and $A$ are the
constants of the last display.

For (4), let $q=\omega(g,h)F-\int_Y\omega(g,h)F_y\,dm(y)$. It is continuous and
quasi-periodic, since each $\omega(g,h)F_y$ is a theta function. Let $f'$ be a
theta function and put $f'_0=\omega(g,h)^{-1}f'$. Since $\omega(g,h)$ preserves
the $L^2$-norm on a fundamental domain, it preserves the pairing
$\langle\ ,\ \rangle$, and by Fubini, justified by the bound $\beta$,
\[
\langle\omega(g,h)F,f'\rangle=\langle F,f'_0\rangle=\int_Y\langle F_y,f'_0\rangle\,dm(y)
=\int_Y\langle\omega(g,h)F_y,f'\rangle\,dm(y)=\Bigl\langle\int_Y\omega(g,h)F_y\,dm(y),f'\Bigr\rangle .
\]
So $q$ is orthogonal to every theta function. A continuous quasi-periodic
function orthogonal to every theta function is $0$, so $q=0$.

Part (5) is Lemma~\ref{th-height-comparison}(3) and~(4).
\end{proof}

\begin{lemma}\label{th-kernel-tame}
Let $D$ be a definite-side datum, $\psi$ a standard character, $S,T$
Hermitian spaces and $\omega$ a Weil datum for $(S,T)$
(Definition~\ref{th-weil-datum}). Let $Y$ be a second countable space with an
$s$-finite Borel measure $m$, let $\lambda$ be a measurable function on $Y$, and let
$g\colon Y\to U(S)(\mathbf A)$, $h\colon Y\to U(T)(\mathbf A)$ and
$w\colon Y\to\mathbb W(S,T)$ be continuous. Suppose that
$y\mapsto|\lambda(y)|\,(\operatorname{ht}(g(y))\operatorname{ht}(h(y)))^k$ is
integrable for every $k\geq0$. Then for every theta function $f$ the function
$y\mapsto\lambda(y)\,(\omega(g(y),h(y))f)(w(y))$ is integrable, and
\[
\ell(f)=\int_Y\lambda(y)\,(\omega(g(y),h(y))f)(w(y))\,dm(y)
\]
is tame (Lemma~\ref{th-vacuum-irreducibility}).
\end{lemma}
\leanref{Development/AbelianVarietyTheory/Proofs/TwistedOperators/WeilKernelFunctionalSuperpositionIsTameStatement.lean}{GlobalTheory.T.weilKernelFunctional_superposition_isTame}
\begin{proof}
Write $\omega_y=\omega(g(y),h(y))$. Let $f$ be a theta function, and take $C$
and $A$ from Lemma~\ref{th-weil-estimates}(3). By part (1) of that lemma and
$|\psi|=1$,
\[
|\lambda(y)\,(\omega_y\rho(u)f)(w(y))|\leq C\,|\lambda(y)|\,(\operatorname{ht}(g(y))\operatorname{ht}(h(y)))^A
\]
for all $y$ and all $u\in\mathbb W(S,T)$, and the right side is integrable by
hypothesis. With $u=0$ this gives the integrability. The map $\ell$ is additive
and linear on theta functions because $\omega_y$ and the integral are linear.
For each $y$ the integrand at $\rho(u)f$ is continuous in $u$, by parts (1)
and (2) of Lemma~\ref{th-weil-estimates}, so $u\mapsto\ell(\rho(u)f)$ is
continuous by dominated convergence, and it is bounded by the integral of the
right side.

For the last condition of tameness, let $\alpha$ be integrable for a measure
$m'$ on $\mathbb W(S,T)$ such that $F=\int\alpha(u)\rho(u)f\,dm'(u)$ is a theta
function. The functions $F_u=\alpha(u)\rho(u)f$ and $\omega_yF_u$ are bounded
by $|\alpha(u)|$ times $\sup|f|$ and $\sup|\omega_yf|$, by part (1), so part (4)
of Lemma~\ref{th-weil-estimates} gives
$(\omega_yF)(w(y))=\int\alpha(u)(\omega_y\rho(u)f)(w(y))\,dm'(u)$. The double
integrand $\lambda(y)\alpha(u)(\omega_y\rho(u)f)(w(y))$ is bounded by
$C|\lambda(y)|(\operatorname{ht}(g(y))\operatorname{ht}(h(y)))^A|\alpha(u)|$, which
is integrable on $Y\times\mathbb W(S,T)$. Exchanging the integrals gives
$\ell(F)=\int\alpha(u)\ell(\rho(u)f)\,dm'(u)$.
\end{proof}

\begin{theorem}\label{th-seesaw-period-nonvanishing}
Let $D$ be a definite-side datum (Definition~\ref{th-definite-datum}) and
$\psi$ a character of $\mathbf A_F$. Let $W$ and $W_0$ be Hermitian spaces
with an isometry $\iota_0\colon W_0\to W$, with a Haar measure and a
fundamental domain $\mathcal F_{W_0}$ for $U(W_0)(F)$, and let $X=T\oplus L''$
be an orthogonal decomposition with $L''$ of rank $1$, with a Haar measure and
a fundamental domain $\mathcal F_T$ for $U(T)(F)$; write
$\iota''\colon U(T)\to U(X)$ for $h\mapsto(h,1)$. Let
$\Gamma\subset\mathbb W(W,T)$, $\Gamma_L\subset\mathbb W(W_0,-L'')$ and
$\Gamma_X\subset\mathbb W(W_0,X)$ be subgroups such that for
$w=(w_1,w_2)\in\Gamma_X^\perp$ the components satisfy
$\iota_0(w_1)\in\Gamma^\perp$ and $w_2\in\Gamma_L^\perp$, and such that there
are nonzero lattice vacua $f_1$ for $\Gamma$ and $f_2$ for $\Gamma_L$ with
$(f_1\circ\iota_0)\boxtimes\bar f_2$ a lattice vacuum for $\Gamma_X$. Let
$\omega$, $\omega_X$, $\omega_z$ be Weil data for $(W,T)$, $(W_0,X)$ and
$(W_0,-L'')$ with characters $(\chi_S,\chi_T)$, $(\chi_S,\chi_T\chi_z^{-1})$
and $(\chi_S^{-1},\chi_z)$ (Definition~\ref{th-weil-datum}). Let $\varphi$ on
$U(W_0)(\mathbf A)$ and $u$ on $U(T)(\mathbf A)$ be rapidly decreasing: left
invariant under the rational points, continuous, and with
$|\varphi(g)|\operatorname{ht}(g)^N$ bounded for every $N$, where
$\operatorname{ht}$ is the height on the quotient. Let $\Xi$ be a lattice
vacuum for $\Gamma_X$ with
\[
\int_{\mathcal F_T}\theta_\Xi(\varphi)\bigl(\iota''(h)\bigr)\,u(h)\,dh\neq0 .
\]
Then there are lattice vacua $\Phi$ for $\Gamma$ and $\Phi_0$ for $\Gamma_L$
with
\[
\int_{\mathcal F_{W_0}}\varphi(g)\,\theta_\Phi(u)\bigl(\iota_0(g)\bigr)\,
\overline{\theta_{\omega_z}(g,1;\Phi_0)}\,dg\neq0 .
\]
\end{theorem}
\leanref{Development/AbelianVarietyTheory/Proofs/TwistedOperators/LastSeesawPeriodNonvanishingStatement.lean}{GlobalTheory.T.lastSeesaw_periodNonvanishing}
\leanref{Development/AbelianVarietyTheory/Definitions/TwistedInnerProducts/TowerEx.lean}{GlobalTheory.Asc.UpLast}
\leanref{Development/AbelianVarietyTheory/Proofs/TwistedOperators/LastSeesawPeriodNonvanishing.lean}{GlobalTheory.LastSeesaw.per_zero}
\leanref{Development/AbelianVarietyTheory/Definitions/TwistedInnerProducts/ExistsHeightLeTwoMulQuotHeight.lean}{GlobalTheory.Grow.RapidDec}
\begin{proof}
For a theta function $F$ for $(W_0,X)$ put
\[
\ell(F)=\int_{\mathcal F_{W_0}\times\mathcal F_T}\varphi(g)\,u(h)\,
\theta_{\omega_X}\bigl(g,\iota''(h);F\bigr)\,dg\,dh .
\]
This is the map of Lemma~\ref{th-kernel-tame} on
$Y=\mathcal F_{W_0}\times\mathcal F_T$ with $\lambda(g,h)=\varphi(g)u(h)$, the points
$(g,\iota''(h))$ and $w=0$. Its integrability hypothesis holds: by
Lemma~\ref{th-weil-estimates}(5),
$\operatorname{ht}(\iota''(h))\leq C\operatorname{ht}(h)^a$, so
$|\lambda(g,h)|(\operatorname{ht}(g)\operatorname{ht}(\iota''(h)))^k$ is bounded by
$C^k$ times $|\varphi(g)|\operatorname{ht}(g)^k\cdot|u(h)|\operatorname{ht}(h)^{ak}$,
which is bounded because $\varphi$ and $u$ are rapidly decreasing, and the
fundamental domains have finite volume. So $\ell$ is tame. By Fubini,
$\ell(\Xi)$ is the integral in the hypothesis, so $\ell(\Xi)\neq0$.

Let $\Phi$ be a theta function for $(W,T)$ and $\Phi_0$ one for $(W_0,-L'')$,
and write $\Phi\circ\iota_0$ for the pullback of $\Phi$ to $\mathbb W(W_0,T)$.
The pullback of $\omega$ along $\iota_0$ is a Weil datum for $(W_0,T)$ with
characters $(\chi_S,\chi_T)$, and complex conjugation turns $\omega_z$ into a
Weil datum for $(W_0,L'')$ with characters $(\chi_S,\chi_z^{-1})$ that sends
$\bar\Phi_0$ to $\overline{\omega_z(g,1)\Phi_0}$. By
Lemma~\ref{th-tensor-rule} applied to $X=T\oplus L''$, the function
$(\Phi\circ\iota_0)\boxtimes\bar\Phi_0$ is a theta function for $(W_0,X)$ and
\[
\theta_{\omega_X}\bigl(g,\iota''(h);(\Phi\circ\iota_0)\boxtimes\bar\Phi_0\bigr)
=\theta_\omega\bigl(\iota_0(g),h;\Phi\bigr)\,\overline{\theta_{\omega_z}(g,1;\Phi_0)} .
\]
Integrating over $h$ first,
$\ell\bigl((\Phi\circ\iota_0)\boxtimes\bar\Phi_0\bigr)$ is the period in the
conclusion.

Suppose that this period vanishes for all lattice vacua $\Phi$ for $\Gamma$
and $\Phi_0$ for $\Gamma_L$. Put $f_0=(f_1\circ\iota_0)\boxtimes\bar f_2$,
which is a nonzero lattice vacuum for $\Gamma_X$. For
$w=(w_1,w_2)\in\Gamma_X^\perp$ the translate $\rho(w)f_0$ is
$(\rho(\iota_0w_1)f_1\circ\iota_0)\boxtimes\overline{\rho(w_2)f_2}$. Since
$\iota_0(w_1)\in\Gamma^\perp$ and $w_2\in\Gamma_L^\perp$, the translates
$\rho(\iota_0w_1)f_1$ and $\rho(w_2)f_2$ are again lattice vacua for $\Gamma$
and $\Gamma_L$. So $\ell(\rho(w)f_0)$ is one of the vanishing periods, and
$\ell(\rho(w)f_0)=0$ for all $w\in\Gamma_X^\perp$.
Lemma~\ref{th-vacuum-irreducibility} gives $\ell(\Xi)=0$, a contradiction.
\end{proof}

The period in the conclusion is the Fourier-Jacobi period of $\varphi$ against
the theta lift of $u$ at depth zero, the case in which the unipotent part of
the period is trivial.

\begin{definition}\label{th-periods}
Let $D$ be a definite-side datum and $\psi$ a character of $\mathbf A_F$. Let
$V$ and $V_0$ be Hermitian spaces with an orthogonal decomposition
$V=V_0\oplus H$, where $H$ is a sum of $\rho$ hyperbolic planes with isotropic
bases $(e_a)$, $(e'_a)$, and let $N\subset U(V)$ be the unipotent radical of
the stabilizer of the flag spanned by the $e_a$, with a Haar measure and a
fundamental domain $\mathcal N$ for $N(F)$ of volume $1$.
\begin{enumerate}
\item Let $V_0=A\oplus L$ with $L$ a line spanned by $\zeta$. The
  \emph{Bessel period} of functions $\varphi$ on $U(A)(\mathbf A)$ and $f$ on
  $U(V)(\mathbf A)$ is
  \[
  \mathcal B(\varphi,f)=\int_{\mathcal F_A}\int_{\mathcal N}\varphi(h)\,
  f(n\,\iota(h))\,\overline{\psi_{\zeta}(n)}\,dn\,dh,
  \]
  where $\iota\colon U(A)\to U(V)$ is the inclusion and $\psi_\zeta$ is the
  generic character of $N$ attached to the flag and $\zeta$.
\item Let $L'$ be a line spanned by $\zeta'$ and $\omega_z$ a Weil datum for
  $(V_0,L')$. The \emph{Fourier-Jacobi period} of $\varphi$ on
  $U(V_0)(\mathbf A)$ and $f$ on $U(V)(\mathbf A)$ against a theta function
  $\Phi_0$ for $(V_0,L')$ is
  \[
  \mathcal J(\varphi,f;\Phi_0)=\int_{\mathcal F_{V_0}}\int_{\mathcal N}\varphi(h)\,
  f(n\,\iota(h))\,\overline{\theta_J(n,h;\Phi_0)}\,\overline{\psi_N(n)}\,dn\,dh,
  \]
  where $\theta_J(n,h;\Phi_0)$ is the Jacobi theta kernel of $\omega_z$ at
  $\zeta'$ and $\psi_N$ the character of the flag.
\end{enumerate}
For a theta grid with grid Weil datum (Definition~\ref{th-grid}), the
\emph{top Bessel period} $\mathcal B^+(\varphi,F)$ of $\varphi$ on
$U(Q_0)(\mathbf A)$ and $F$ on $U(V_{m,0})(\mathbf A)$ is the Bessel period
for $Q_0=P_m\subset V_{m,0}$, whose complement is a line and $(m-2)/2$
hyperbolic planes.
\end{definition}
\leanref{Development/AbelianVarietyTheory/Definitions/TwistedInnerProducts/WittSmallGroupEmbeddingMemPointwiseStabilizerIso.lean}{GlobalTheory.MixPer.besPer}
\leanref{Development/AbelianVarietyTheory/Definitions/TwistedInnerProducts/WittSmallGroupEmbeddingMemPointwiseStabilizerIso.lean}{GlobalTheory.MixPer.fjSrc}
\leanref{Development/AbelianVarietyTheory/Definitions/Divisors/BplusOfEquivariant.lean}{GlobalTheory.FunObj.Bplus}

\begin{definition}\label{th-core-chain}
Let $D$ be a definite-side datum with a level system, a theta grid and a grid
Weil datum (Definition~\ref{th-grid}), and fix Haar measures and fundamental
domains $\mathcal F_{V_{i,0}}$ and $\mathcal F_{P_i}$ for $U(V_{i,0})(F)$ and
$U(P_i)(F)$, $0\leq i\leq m$.
\begin{enumerate}
\item \emph{Row data} $(\varrho,\Phi)$ consist of, for each $0\leq i<m$, a
  coset vector $\Phi_i$ of the pair $(V_{i+1,0},V_{i,0})$ and a polynomial
  $\varrho_i$ in the creation operators at the archimedean places; the
  corresponding theta function $f_i$ of the pair is $\varrho_i$ applied to
  $\Phi_i$. For a function $u$ on $U(Y_0)(\mathbf A)=U(V_{0,0})(\mathbf A)$
  the \emph{partial row lifts} are $\theta^{(0)}(u)=u$ and, for $n<m$,
  \[
  \theta^{(n+1)}(u)(g)=\int_{\mathcal F_{V_{n,0}}}\theta_{\omega}(g,h;f_n)\,
  \theta^{(n)}(u)(h)\,dh\qquad\bigl(g\in U(V_{n+1,0})(\mathbf A)\bigr),
  \]
  where $\omega$ is the Weil datum of the pair $(-V_{n+1,0},V_{n,0})$ and
  $U(V_{n+1,0})=U(-V_{n+1,0})$ (Definition~\ref{th-weil-datum}).
\item A \emph{column chain of core forms} is a tuple $(\Psi_0,\ldots,\Psi_m)$
  of rapidly decreasing functions $\Psi_i$ on $U(P_i)(\mathbf A)$ such that for
  each $i<m$ there is a theta function $f'_i$ for the pair $(-P_{i+1},P_i)$,
  a lattice vacuum for the depth lattice of the pair (the sublattice of its
  lattice cut out by the modulus of the grid at the distinguished places),
  with
  \[
  \Psi_i(g)=\int_{\mathcal F_{P_{i+1}}}\theta_\omega(h,g;f'_i)\,\Psi_{i+1}(h)\,dh,
  \]
  where $\omega$ is the Weil datum of the pair.
\end{enumerate}
\end{definition}
\leanref{Development/AbelianVarietyTheory/Definitions/TwistedInnerProducts/OffQOfSizeRows.lean}{GlobalTheory.SepRows.IsColChain}
\leanref{Development/AbelianVarietyTheory/Definitions/TwistedInnerProducts/OffQOfSizeRows.lean}{GlobalTheory.SepRows.part}
\leanref{Development/AbelianVarietyTheory/Definitions/TwistedInnerProducts/OffQOfSizeRows.lean}{GlobalTheory.SepRows.RapidRow}
\leanref{Development/AbelianVarietyTheory/Definitions/TwistedInnerProducts/OffQOfSizeRows.lean}{GlobalTheory.SepRows.depLat}
\leanref{Development/AbelianVarietyTheory/Definitions/Divisors/BasePlaceNotMemChiGridBadPlaces.lean}{GlobalTheory.ChiGrid.GridData.Weil.datum}
\leanref{Development/AbelianVarietyTheory/Definitions/Divisors/BplusOfEquivariant.lean}{GlobalTheory.FunObj.thetaChi}

\section{One step of the tower}\label{section-th-steps}

\begin{definition}\label{th-step}
Let $D$ be a definite-side datum and $\psi$ a standard character
(Definition~\ref{th-weil-datum}). For Hermitian spaces $S$ and $S_0$ and an
integer $\rho\geq0$, a \emph{hyperbolic decomposition} $S=S_0\oplus H^\rho$ is
an isometric embedding of $S_0$ in $S$ together with vectors
$e_1,\ldots,e_\rho,e'_1,\ldots,e'_\rho$ of $S$, orthogonal to $S_0$, with
$h(e_a,e_b)=h(e'_a,e'_b)=0$ and $h(e'_a,e_b)=\delta_{ab}$. Write $N_S$ for the
group $N$ of Definition~\ref{th-periods} for this decomposition, $U_\rho$ for
the group of upper unitriangular matrices in $GL_\rho$, and
$\psi_U(u)=\psi\bigl(\operatorname{Tr}_{\mathbf A_E/\mathbf A_F}\sum_{a<\rho}u_{a,a+1}\bigr)$
for $u\in U_\rho(\mathbf A_E)$.
\begin{enumerate}
\item \emph{Bessel data} of depth $\rho$ consist of a hyperbolic decomposition
  $V=V_0\oplus H^\rho$, an orthogonal decomposition $V_0=A\oplus L$ with $L$ of
  rank $1$, a vector $\zeta\in L$ with $h(\zeta,\zeta)\neq0$, a Haar measure and
  a fundamental domain for $U(A)(F)$, and a Haar measure on $N_V(\mathbf A)$
  with a fundamental domain of volume $1$ for $N_V(F)$. For functions $\varphi$
  on $U(A)(\mathbf A)$ and $f$ on $U(V)(\mathbf A)$, $\mathcal B_V(\varphi,f)$
  is the Bessel period of Definition~\ref{th-periods}(1).
\item \emph{Fourier-Jacobi data} of depth $\rho$ consist of a hyperbolic
  decomposition $W=W_0\oplus H^\rho$, a Haar measure and a fundamental domain for
  $U(W_0)(F)$, and a Haar measure on $N_W(\mathbf A)$ with a fundamental domain
  of volume $1$ for $N_W(F)$. For a Hermitian line $L'$ with a vector $\zeta'$,
  a Weil datum $\omega_z$ for $(W_0,L')$, functions $\varphi$ on
  $U(W_0)(\mathbf A)$ and $f$ on $U(W)(\mathbf A)$ and a theta function
  $\Phi_0$ for $(W_0,L')$, $\mathcal J_W(\varphi,f;\Phi_0)$ is the
  Fourier-Jacobi period of Definition~\ref{th-periods}(2).
\item A \emph{step} from $S=S_0\oplus H^\rho$ to $T=T_0\oplus H^\sigma$ consists
  of unitary characters $\chi_S,\chi_T$ of $\mathbf A_E^\times$, Weil data
  $\omega$ for $(S,T)$ and $\omega_0$ for $(S_0,T_0)$, both with characters
  $(\chi_S,\chi_T)$, a Haar measure and a fundamental domain for $U(T)(F)$, and
  Haar measures with fundamental domains of volume $1$ for the rational points
  of $M_{\operatorname{rk}T\times\rho}(\mathbf A_E)$,
  $M_{\sigma\times\operatorname{rk}S_0}(\mathbf A_E)$ and $U_\rho(\mathbf A_E)$.
  For a function $G$ on $U(T)(\mathbf A)$ and a theta function $\Phi$ for
  $(S,T)$, $\theta_\Phi(G)$ is the theta lift of $G$ to $U(S)(\mathbf A)$ by
  $\omega$ (Definition~\ref{th-weil-datum}(4)).
\item The \emph{flat map} of a step at points
  $x\in M_{\operatorname{rk}T\times\rho}(\mathbf A_E)$ and
  $y\in M_{\operatorname{rk}S_0\times\sigma}(\mathbf A_E)$ sends a theta
  function $\Phi$ for $(S,T)$ to the function
  \[
  \Phi^\flat(w)=\int_{U_\rho(\mathbf A_E)}\overline{\psi_U(u)}\,
  \Phi_{xu,y}(w)\,du\qquad\bigl(w\in\mathbb W(S_0,T_0)\bigr),
  \]
  the integral over the whole group $U_\rho(\mathbf A_E)$ against the Haar
  measure of (3), where $\Phi_{x,y}$ is the mixed model of $\Phi$ at $(x,y)$. It is obtained in
  two stages. First, in the coordinates of $\mathbb W(S,T)$ attached to the two
  isotropic halves of $H^\rho\subset S$, translate $\Phi$ by $x$ in the first
  half and by $w'\in\mathbb W(S_0,T)$, integrate the translates over the
  fundamental domain of $M_{\operatorname{rk}T\times\rho}(\mathbf A_E)$ in the
  second half, and take the value at $0$; this is a function of $w'$. Second,
  apply the same construction to this function on the side of $T$: in the
  coordinates attached to the isotropic halves of $H^\sigma\subset T$,
  integrate over the fundamental domain of
  $M_{\sigma\times\operatorname{rk}S_0}(\mathbf A_E)$ in the first half,
  translate by the adjoint of $y$ in the second half and by
  $w\in\mathbb W(S_0,T_0)$, and take the value at $0$.
\end{enumerate}
A step from Bessel data $V=V_0\oplus H^\rho$ to Fourier-Jacobi data
$W=W_0\oplus H^\rho$ of the same depth is taken with the points $x_0$, the
matrix of the vectors $e_1,\ldots,e_\rho$ of $W$, and $y_0$, the matrix with
last column $-\zeta$ and other columns $0$. A step from Fourier-Jacobi data
$W=W_0\oplus H^{\sigma+1}$ to Bessel data $V=V_0\oplus H^\sigma$ is taken with
the points $x_1$, the matrix with columns the vectors $e_1,\ldots,e_\sigma$ of
$V$ followed by $\zeta$, and $y=0$.
\end{definition}
\leanref{Development/AbelianVarietyTheory/Definitions/TwistedInnerProducts/WittDecompProjAnisotropicAnisotropicEmbedding.lean}{GlobalTheory.MixModel.HypDec}
\leanref{Development/AbelianVarietyTheory/Definitions/TwistedInnerProducts/TowerEx.lean}{GlobalTheory.Asc.BesL}
\leanref{Development/AbelianVarietyTheory/Definitions/TwistedInnerProducts/TowerEx.lean}{GlobalTheory.Asc.JacL}
\leanref{Development/AbelianVarietyTheory/Definitions/LatticeFrames/IsThetaFam.lean}{GlobalTheory.Unf.Step}
\leanref{Development/AbelianVarietyTheory/Definitions/TwistedInnerProducts/WittSmallGroupEmbeddingMemPointwiseStabilizerIso.lean}{GlobalTheory.MixPer.flatAt}
\leanref{Development/AbelianVarietyTheory/Definitions/TwistedInnerProducts/WittSmallGroupEmbeddingMemPointwiseStabilizerIso.lean}{GlobalTheory.MixPer.evalAt}
\leanref{Development/AbelianVarietyTheory/Definitions/TwistedInnerProducts/TowerEx.lean}{GlobalTheory.Asc.flatS}

\begin{definition}\label{th-step-local}
Let a step from $S=S_0\oplus H^\rho$ to $T=T_0\oplus H^\sigma$ with Weil datum
$\omega$ and flat map $\Phi\mapsto\Phi^\flat$ at $(x,y)$ be given
(Definition~\ref{th-step}). Let $\iota\colon\mathbb W(S_0,T_0)\to\mathbb W(S,T)$
be the inclusion through $S_0\subset S$ and $T_0\subset T$, and let $R_x$ be the
group of $r\in U(T)(\mathbf A)$ with $rx=xu$ for some $u\in U_\rho(\mathbf A_E)$.
\emph{Local data} for the step consist of a set $\mathcal S$ of functions on
$\mathbb W(S,T)$, a subgroup $\Gamma_0\subset\mathbb W(S_0,T_0)$, and a
compact subgroup $K$ of the group of elements of $U(T)(\mathbf A)$ with
archimedean component $1$, open in that group, such that:
\begin{enumerate}
\item every $\Phi\in\mathcal S$ is a theta function;
\item for $\Phi\in\mathcal S$, $\rho(w)\Phi\in\mathcal S$ for every
  $w\in\mathbb W(S,T)$ with finite component $0$, and
  $\rho(\iota w_0)\Phi\in\mathcal S$ for every $w_0\in\Gamma_0^\perp$
  (Lemma~\ref{th-vacuum-irreducibility});
\item $\Phi^\flat$ is a lattice vacuum for $\Gamma_0$ for every $\Phi\in\mathcal S$;
\item $\omega(1,k)\Phi=\Phi$ for $\Phi\in\mathcal S$ and $k\in K$;
\item if $\Phi\in\mathcal S$, $h\in U(T)(\mathbf A)$ and
  $(\omega(1,h)\Phi)^\flat\neq0$, then $h=rak$ with $r\in R_x$, $a$ of finite
  component $1$ and $k\in K$;
\item $\Phi^\flat\neq0$ for some $\Phi\in\mathcal S$.
\end{enumerate}
\end{definition}
\leanref{Development/AbelianVarietyTheory/Definitions/TwistedInnerProducts/TowerEx.lean}{GlobalTheory.Asc.Loc}
\leanref{Development/AbelianVarietyTheory/Definitions/LatticeFrames/KerBF.lean}{GlobalTheory.Unf.Rgrp}

\begin{definition}\label{th-step-periods}
Let $D$ be a definite-side datum and $\psi$ a standard character.
\begin{enumerate}
\item Let a step from Bessel data $V=V_0\oplus H^\rho$, $V_0=A\oplus L$, to
  Fourier-Jacobi data $W=W_0\oplus H^\rho$ be given, with Weil datum $\omega_0$
  for $(V_0,W_0)$ (Definition~\ref{th-step}). Let $Y\in\mathbb W(V_0,W)$ be the
  element attached to $y_0$ in the isotropic half of $W$ spanned by
  $e'_1,\ldots,e'_\rho$. For $n\in N_W(\mathbf A)$ let $p(n)$ be the component in
  $\mathbb W(V_0,W_0)$ of $n^{-1}Y$, put
  $c(n)=\tfrac12\langle Y,n^{-1}Y\rangle$, and let
  $j(n)\Xi=\psi(c(n))\,\rho(p(n))\Xi$ for functions $\Xi$ on
  $\mathbb W(V_0,W_0)$. For functions $\varphi$ on $U(A)(\mathbf A)$ and $G$ on
  $U(W)(\mathbf A)$ and a theta function $\Xi$ for $(V_0,W_0)$, the
  \emph{Fourier-Jacobi functional} is
  \[
  \ell_{\varphi,G}(\Xi)=\int_{\mathcal F_{W_0}}\int_{\mathcal N_W}G(nh)\,
  \overline{\psi_N(n)}\int_{\mathcal F_A}\varphi(h_0)\,
  \bigl(\omega_0((h_0,1),1)\,j(n)\,\omega_0(1,h)\Xi\bigr)(0)\,dh_0\,dn\,dh ,
  \]
  with the fundamental domains and measures of the Bessel and Fourier-Jacobi
  data and the character $\psi_N$ of Definition~\ref{th-periods}.
\item Let a step from Fourier-Jacobi data $W=W_0\oplus H^{\sigma+1}$ to Bessel
  data $V=V_0\oplus H^\sigma$, $V_0=A\oplus L$, be given, with Weil datum
  $\omega_0$ for $(W_0,V_0)$. Fix a Haar measure on $\mathbb W(W_0,L)$ with a
  fundamental domain $\mathcal F_L$ of volume $1$, a Weil datum $\omega_z$ for
  $(W_0,L)$ and a theta function $\Phi_0$ for $(W_0,L)$. For a function
  $\varphi$ on $U(W_0)(\mathbf A)$ and a theta function $\Xi$ for $(W_0,V_0)$,
  the \emph{handed-down function} is the function on $U(A)(\mathbf A)$
  \[
  \varphi_\Xi(h_1)=\int_{\mathcal F_{W_0}}\varphi(h_0)\int_{\mathcal F_L}
  \bigl(\omega_0(h_0,(h_1,1))\Xi\bigr)(0,v)\,
  \overline{\bigl(\omega_z(h_0,1)\Phi_0\bigr)(v)}\,dv\,dh_0 ,
  \]
  where $\mathbb W(W_0,V_0)=\mathbb W(W_0,A)\oplus\mathbb W(W_0,L)$.
\end{enumerate}
\end{definition}
\leanref{Development/AbelianVarietyTheory/Definitions/TwistedInnerProducts/TowerEx.lean}{GlobalTheory.Asc.JacL.perAll}
\leanref{Development/AbelianVarietyTheory/Definitions/TwistedInnerProducts/WittSmallGroupEmbeddingMemPointwiseStabilizerIso.lean}{GlobalTheory.MixPer.fjAll}
\leanref{Development/AbelianVarietyTheory/Definitions/TwistedInnerProducts/WittSmallGroupEmbeddingMemPointwiseStabilizerIso.lean}{GlobalTheory.MixPer.jacT}
\leanref{Development/AbelianVarietyTheory/Definitions/TwistedInnerProducts/TowerEx.lean}{GlobalTheory.Asc.handP}

\begin{lemma}\label{th-step-functionals}
Let $D$ be a definite-side datum and $\psi$ a standard character. Let $T$ be a
Hermitian space with a Haar measure and a fundamental domain for $U(T)(F)$, and
let $G$ be a rapidly decreasing function on $U(T)(\mathbf A)$
(Theorem~\ref{th-seesaw-period-nonvanishing}).
\begin{enumerate}
\item Let $V=V_0\oplus H^\rho$, $V_0=A\oplus L$, be Bessel data, $\omega$ a Weil
  datum for $(V,T)$ and $\varphi$ a rapidly decreasing function on
  $U(A)(\mathbf A)$ (Definition~\ref{th-step}). Then
  $\Phi\mapsto\mathcal B_V(\varphi,\theta_\Phi(G))$ is tame
  (Lemma~\ref{th-vacuum-irreducibility}).
\item Let $W=W_0\oplus H^\rho$ be Fourier-Jacobi data, $\omega$ a Weil datum for
  $(W,T)$, $L'$ a line with a vector $\zeta'$, $\omega_z$ a Weil datum for
  $(W_0,L')$, $\varphi$ a rapidly decreasing function on $U(W_0)(\mathbf A)$ and
  $\Phi_0$ a theta function for $(W_0,L')$. Then
  $\Phi\mapsto\mathcal J_W(\varphi,\theta_\Phi(G);\Phi_0)$ is tame.
\item Let a step from Bessel data $V=V_0\oplus H^\rho$, $V_0=A\oplus L$ with
  vector $\zeta$, to Fourier-Jacobi data $W=W_0\oplus H^\rho$ be given, with
  characters $(\chi_V,\chi_W)$, let $T=W$, and let $\varphi$ be rapidly
  decreasing on $U(A)(\mathbf A)$. Then $\Xi\mapsto\ell_{\varphi,G}(\Xi)$ is tame
  (Definition~\ref{th-step-periods}(1)). Let moreover $L'$ be a line with a
  vector $\zeta'$, $h(\zeta',\zeta')\neq0$, and $\tau\colon L'\to L$ a linear
  map with $\tau\zeta'=-\zeta$ and
  $h(\zeta,\zeta)=\xi_0\bar\xi_0\,h(\zeta',\zeta')$; let
  $\chi_V=\chi_A\chi_z$, and let $\omega_A$ and $\omega_z$ be Weil data for
  $(A,W_0)$ and $(W_0,L')$ with characters $(\chi_A,\chi_W)$ and
  $(\chi',\chi_z)$. Then for theta functions $\Xi_1$ for $(A,W_0)$ and $\Xi_2$
  for $(L,W_0)$,
  \[
  \ell_{\varphi,G}(\Xi_1\boxtimes\Xi_2)=\mathcal J_W\bigl(\theta_{\Xi_1}(\varphi),G;\overline{\Xi_2\circ\kappa}\bigr),
  \]
  where $\theta_{\Xi_1}(\varphi)$ is the theta lift of $\varphi$ to
  $U(W_0)(\mathbf A)$ by $\omega_A$, the Fourier-Jacobi period is taken with
  $\zeta'$ and $\omega_z$, and $\kappa(w)=(\tau w)^*$ as in
  Lemma~\ref{th-line-comparison}.
\item Let a step from Fourier-Jacobi data $W=W_0\oplus H^{\sigma+1}$ to Bessel
  data $V=V_0\oplus H^\sigma$, $V_0=A\oplus L$, be given, with characters
  $(\chi_W,\chi_V)$, let $T=V$, let $\omega_z$ be a Weil datum for $(W_0,L)$, let
  $\varphi$ be rapidly decreasing on $U(W_0)(\mathbf A)$, and let $\Phi_0$ be a
  theta function for $(W_0,L)$. Then $\Xi\mapsto\mathcal B_V(\varphi_\Xi,G)$ is
  tame (Definition~\ref{th-step-periods}(2)). If moreover
  $\chi_V=\chi_A\chi_z$, and $\omega_A$ and $\omega_z$ are Weil data for
  $(W_0,A)$ and $(W_0,L)$ with characters $(\chi_W,\chi_A)$ and
  $(\chi_W,\chi_z)$, then for theta functions $\Xi_1$ for $(W_0,A)$ and
  $\Phi'$ for $(W_0,L)$,
  \[
  \varphi_{\Xi_1\boxtimes\Phi'}=\langle\Phi',\Phi_0\rangle\,\theta_{\Xi_1}(\varphi),
  \]
  where $\langle\Phi',\Phi_0\rangle=\int_{\mathcal F_L}\Phi'\,\overline{\Phi_0}$
  and $\theta_{\Xi_1}(\varphi)$ is the theta lift of $\varphi$ to
  $U(A)(\mathbf A)$ by $\omega_A$.
\item For a step from $S=S_0\oplus H^\rho$ to $T=T_0\oplus H^\sigma$, every
  theta function $\Phi$ for $(S,T)$ and every $w\in\mathbb W(S_0,T_0)$,
  $(\rho(\iota w)\Phi)^\flat=\rho(w)\Phi^\flat$, with $\iota$ as in
  Definition~\ref{th-step-local}.
\end{enumerate}
\end{lemma}
\leanref{Development/AbelianVarietyTheory/Proofs/TwistedOperators/ThetaLiftPeriodsIsTameThetaFunctionalStatement.lean}{GlobalTheory.T.thetaLift_periods_isTameThetaFunctional}
\leanref{Development/AbelianVarietyTheory/Proofs/TwistedOperators/FourierJacobiPeriodAllDataIsTameThetaFunctionalStatement.lean}{GlobalTheory.T.fourierJacobiPeriod_all_data_isTameThetaFunctional}
\leanref{Development/AbelianVarietyTheory/Proofs/TwistedOperators/PeriodAllDataProductsEqFourierJacobiPeriodStatement.lean}{GlobalTheory.T.periodAllData_products_eq_fourierJacobiPeriod}
\leanref{Development/AbelianVarietyTheory/Proofs/TwistedOperators/BesselPeriodHandedDownIsTameThetaFunctionalStatement.lean}{GlobalTheory.T.besselPeriod_handedDown_isTameThetaFunctional}
\leanref{Development/AbelianVarietyTheory/Proofs/TwistedOperators/BesselHandedFunctionProductEqCoreLiftMulPairingStatement.lean}{GlobalTheory.T.besselHandedFunction_product_eq_coreLift_mul_pairing}
\leanref{Development/AbelianVarietyTheory/Proofs/DecoratedLatticeProofs/FlatMapIntertwinesCoreHeisenbergOperators.lean}{GlobalTheory.FlatHeis.flatAt_heisO}
\begin{proof}
\emph{Tameness, in (1) to (4).} Each map is the map $\ell$ of
Lemma~\ref{th-kernel-tame} for a measure space $Y$ of finite measure, a weight
$\lambda$, group elements and points, which we list; by Fubini's theorem, which the
bounds below justify, $\ell$ agrees with the period on theta functions. In each
case Lemma~\ref{th-weil-estimates}(5) and~(3) and the rapid decrease of
$\varphi$ and $G$ bound $|\lambda(y)|(\operatorname{ht}\operatorname{ht})^k$ on $Y$
for every $k$, which is the hypothesis of Lemma~\ref{th-kernel-tame}. Write
$\mathcal N$ for the fundamental domain of volume $1$ for the rational points
of the flag group of the decomposition at hand.

In (1), $\theta_\Phi(G)(g)=\int_{\mathcal F_T}(\omega(g,t)\Phi)(0)G(t)\,dt$, so
the period is $\ell$ on $Y=\mathcal F_A\times\mathcal N\times\mathcal F_T$ with
$\lambda(h,n,t)=\varphi(h)\overline{\psi_\zeta(n)}G(t)$, the elements
$(n\iota(h),t)$ and the point $0$. Here
$\operatorname{ht}(n\iota(h))\leq C\operatorname{ht}(\iota(h))^a\leq
C'\operatorname{ht}(h)^{a'}$ by the two comparisons of
Lemma~\ref{th-weil-estimates}(5), and $|\psi_\zeta|=1$. In (2) the same holds
with $Y=\mathcal F_{W_0}\times\mathcal N\times\mathcal F_T$ and
$\lambda(h,n,t)=\varphi(h)\overline{\theta_J(n,h;\Phi_0)}\,\overline{\psi_N(n)}G(t)$,
where $|\theta_J(n,h;\Phi_0)|=|(\omega_z(h,1)\Phi_0)(p)|$ for a point $p$
depending on $n$, which is at most $C\operatorname{ht}(h)^a$ by
Lemma~\ref{th-weil-estimates}(3).

In (3), with $\omega_0$ the Weil datum of the step for $(V_0,W_0)$, the map
$\ell_{\varphi,G}$ is $\ell$ on
$Y=\mathcal F_{W_0}\times\mathcal N_W\times\mathcal F_A$ with
$\lambda(h,n,h_0)=G(nh)\overline{\psi_N(n)}\varphi(h_0)\psi(c(n))$, the elements
$((h_0,1),h)$ and the point $p(n)$ of Definition~\ref{th-step-periods}(1):
since $(h_0,1)$ fixes $p(n)$, which lies in $\mathbb W(L,W_0)$, the
equivariance of $\omega_0$ gives
$(\omega_0((h_0,1),1)j(n)\omega_0(1,h)\Xi)(0)=\psi(c(n))(\omega_0((h_0,1),h)\Xi)(p(n))$.
The bound uses $\operatorname{ht}((h_0,1))\leq C\operatorname{ht}(h_0)^a$ and
$\operatorname{ht}(h)\leq C\operatorname{ht}(nh)^a$ from
Lemma~\ref{th-weil-estimates}(5). In (4), the composite
$\Xi\mapsto\mathcal B_V(\varphi_\Xi,G)$ is $\ell$ on
$Y=\mathcal F_A\times\mathcal N_V\times\mathcal F_{W_0}\times\mathcal F_L$ with
$\lambda(h_1,n,h_0,v)=G(n\iota(h_1))\overline{\psi_\zeta(n)}\varphi(h_0)\overline{(\omega_z(h_0,1)\Phi_0)(v)}$,
the elements $(h_0,(h_1,1))$ and the point $(0,v)$; the factor
$|(\omega_z(h_0,1)\Phi_0)(v)|$ is at most $C\operatorname{ht}(h_0)^a$ by
Lemma~\ref{th-weil-estimates}(3).

\emph{The product formula in (3).} A Weil datum $\omega_L$ for $(L,W_0)$ with
characters $(\chi_z,\chi_W)$ exists, since $\chi_z$ and $\chi_W$ are characters
of Weil data and have the required restrictions to $\mathbf A_F^\times$. By
Lemma~\ref{th-tensor-rule}(2) for $V_0=A\oplus L$ and $\chi_V=\chi_A\chi_z$,
\[
\bigl(\omega_0((h_0,1),h)(\Xi_1\boxtimes\Xi_2)\bigr)(u,v)=\bigl(\omega_A(h_0,h)\Xi_1\bigr)(u)\,\bigl(\omega_L(1,h)\Xi_2\bigr)(v).
\]
The point $p(n)$ has component $0$ in $\mathbb W(A,W_0)$, since $y_0$ is
supported on $\zeta\in L$; call its component in $\mathbb W(L,W_0)$ $p_L(n)$.
The point of the Jacobi theta kernel of $\omega_z$ at $\zeta'$
(Definition~\ref{th-periods}(2)) is mapped by $\kappa$ to $p_L(n)$, because
$\tau\zeta'=-\zeta$, and its central term is $-c(n)$, because
$h(\zeta,\zeta)=\xi_0\bar\xi_0\,h(\zeta',\zeta')$. By
Lemma~\ref{th-line-comparison}(2) for $\omega_L$ and $\omega_z$,
$\omega_z(h,1)\overline{\Xi_2\circ\kappa}=\overline{(\omega_L(1,h)\Xi_2)\circ\kappa}$.
Substituting, the inner integral over $h_0$ in $\ell_{\varphi,G}(\Xi_1\boxtimes\Xi_2)$
becomes $\theta_{\Xi_1}(\varphi)(h)$ times the complex conjugate of the Jacobi
theta kernel of $\overline{\Xi_2\circ\kappa}$ at $(n,h)$. The resulting integral is
the Fourier-Jacobi period.

\emph{The product formula in (4).} By Lemma~\ref{th-tensor-rule}(1) for the
decomposition $V_0=A\oplus L$ of the second space and $\chi_V=\chi_A\chi_z$,
\[
\omega_0\bigl(h_0,(h_1,1)\bigr)(\Xi_1\boxtimes\Phi')=\omega_A(h_0,h_1)\Xi_1\boxtimes\omega_z(h_0,1)\Phi' .
\]
Evaluating at $(0,v)$ and integrating against
$\overline{(\omega_z(h_0,1)\Phi_0)(v)}$ over $\mathcal F_L$ gives
$(\omega_A(h_0,h_1)\Xi_1)(0)\,\langle\omega_z(h_0,1)\Phi',\omega_z(h_0,1)\Phi_0\rangle$.
The Weil operators preserve the $L^2$-norm on a fundamental domain, hence the
pairing, so this is $(\omega_A(h_0,h_1)\Xi_1)(0)\langle\Phi',\Phi_0\rangle$.
Integrating against $\varphi(h_0)$ gives
$\langle\Phi',\Phi_0\rangle\,\theta_{\Xi_1}(\varphi)(h_1)$.

\emph{Part (5).} The mixed model is built from translations along the
hyperbolic parts and evaluation at $0$. The element $\iota w$ is orthogonal to
the hyperbolic parts in the pairing, so $\rho(\iota w)$ commutes with these
translations, and the evaluation at $0$ of the translate by $\iota w$ is the
value at $w$ of the translate in $\mathbb W(S_0,T_0)$, with the same central
factor. Hence $(\rho(\iota w)\Phi)_{x,y}=\rho(w)\Phi_{x,y}$ for all $x,y$, and
integrating over $u$ gives (5).
\end{proof}

\begin{lemma}[Mixed model]\label{th-mixed-model}
Let $D$ be a definite-side datum and $\psi$ a standard character. Let
$V=V_0\oplus H^\rho$ be a hyperbolic decomposition with flag group $N_V$
(Definition~\ref{th-step}), with the elements $m(a)$, $n_B(b)$ and $n_C(c)$ of
$U(V)(\mathbf A)$ and the embedding $\iota\colon U(V_0)\to U(V)$, and let $W$ be
a Hermitian space. Fix a Haar measure on $M_{\operatorname{rk}W\times\rho}(\mathbf A_E)$
with a fundamental domain $\mathcal F$ of volume $1$ for the rational points.
For $x\in M_{\operatorname{rk}W\times\rho}(\mathbf A_E)$ write $X(x)$ and $Y(x)$
for the elements of $\mathbb W(V,W)$ supported on the $e_a$ and on the $e'_a$
with coefficient matrix $x$, and for a theta function $f$ for $(V,W)$ let $f_x$
be the function on $\mathbb W(V_0,W)$ given by
\[
f_x(w)=\int_{\mathcal F}\bigl(\rho(Y(a))\rho(X(x))\rho(w)f\bigr)(0)\,da .
\]
\begin{enumerate}
\item For $a,x\in M_{\operatorname{rk}W\times\rho}(\mathbf A_E)$ and $w,w'\in\mathbb W(V_0,W)$:
  $(\rho(X(a))f)_x=f_{x+a}$;
  $(\rho(Y(a))f)_x=\psi(\langle X(x),Y(a)\rangle)f_x$;
  $(\rho(w')f)_x=\rho(w')f_x$; $(x,w)\mapsto f_x(w)$ is continuous; each $f_x$
  is a theta function for $(V_0,W)$; and
  $\sum_{x\in M_{\operatorname{rk}W\times\rho}(E)}f_x(w)=f(w)$, the sum
  converging.
\item Let $\omega$ and $\omega_0$ be Weil data for $(V,W)$ and $(V_0,W)$ with the
  same characters $(\chi_V,\chi_W)$ (Definition~\ref{th-weil-datum}). Then for
  $h\in U(W)(\mathbf A)$, $g_0\in U(V_0)(\mathbf A)$, $a\in GL_\rho(\mathbf A_E)$,
  $b\in M_{\rho\times\operatorname{rk}V_0}(\mathbf A_E)$ and skew-Hermitian
  $c$:
  \begin{align*}
  (\omega(1,h)f)_x&=\omega_0(1,h)\,f_{h^{-1}x},&
  (\omega(\iota(g_0),1)f)_x&=\omega_0(g_0,1)\,f_x,\\
  (\omega(m(a),1)f)_x&=\chi_W(\det a)^{-1}|\det a|_{\mathbf A_E}^{\operatorname{rk}W/2}f_{xa},&
  (\omega(n_B(b),1)f)_x&=\rho(xb)\,f_x,\\
  (\omega(n_C(c),1)f)_x&=\psi\bigl(\tfrac12\langle X(x),Y(xc)\rangle\bigr)f_x .
  \end{align*}
\end{enumerate}
\end{lemma}
\leanref{Development/AbelianVarietyTheory/Definitions/TwistedInnerProducts/WittDecompSymplInclIsotropicInclAnisotropic.lean}{GlobalTheory.MixModel.HypDec.mixS}
\leanref{Development/AbelianVarietyTheory/Proofs/DecoratedLatticeProofs/HeisenbergOperatorThetaFunctionalMixedSourceModeStatement.lean}{GlobalTheory.T.heisenbergOperator_thetaFunctional_mixedSourceModel}
\leanref{Development/AbelianVarietyTheory/Proofs/DecoratedLatticeProofs/MixedModelFormulasWittSourceProofStatement.lean}{GlobalTheory.T.mixedModelFormulasWittSource_proof}
\leanref{Development/AbelianVarietyTheory/BlockTheory/TwistedDiscreteSeries/SkewGramHalfModule.lean}{GlobalTheory.IndefiniteRowSideFullHecke.SkewGramHalfModule.FormulasS}
\leanref{Development/AbelianVarietyTheory/Proofs/DecoratedLatticeProofs/WeilOperatorsWittParabolicOriginToPointStatement.lean}{GlobalTheory.T.weil_operators_witt_parabolic_origin_to_point}
\leanref{Development/AbelianVarietyTheory/Proofs/DecoratedLatticeProofs/WeilRepresentationValueAtOriginLeviFactorStatement.lean}{GlobalTheory.T.weilRepresentation_value_at_origin_leviFactor}
\leanref{Development/AbelianVarietyTheory/Proofs/DecoratedLatticeProofs/WittMixedModelUnipotentFormulasStatement.lean}{GlobalTheory.T.wittMixedModel_unipotent_formulas}
\leanref{Development/AbelianVarietyTheory/Proofs/DecoratedLatticeProofs/AdelicGLCircleCharacterTrivialOnDetOneStatement.lean}{GlobalTheory.T.adelic_GL_circle_character_trivial_on_det_one}
\leanref{Development/AbelianVarietyTheory/Proofs/DecoratedLatticeProofs/HeckeCharacterEqOneOfTrivialSplitPlacesStatement.lean}{GlobalTheory.T.heckeCharacter_eq_one_of_trivial_splitPlaces}
\begin{proof}
For (1): the first three identities follow from the definition of $\rho$, since
the span of the $e_a$, the span of the $e'_a$ and $\mathbb W(V_0,W)$ are
orthogonal for $\langle\ ,\ \rangle$ except for the pairing between the first
two; the change of variables $a\mapsto a$ in $\mathcal F$ is allowed because the
integrand is invariant under rational translations, by quasi-periodicity. The
function $a\mapsto(\rho(Y(a))\rho(w)f)(0)$ is invariant under
$M_{\operatorname{rk}W\times\rho}(E)$, since the span of the $e'_a$ is isotropic,
and its Fourier coefficients on the compact quotient are the $f_x(w)$ for
rational $x$; summing them at $a=0$ is Fourier inversion, the coefficients being
absolutely summable by the argument of Lemma~\ref{th-matrix-coefficients}(4).

For (2): by (1), $f_x=(\rho(X(x))f)_0$, and by
Definition~\ref{th-weil-datum}(3),
$\rho(X(x))\,\omega(g,h)=\omega(g,h)\,\rho(h^{-1}X(x)g)$ on theta functions.
Hence, for every $(g,h)$,
\[
(\omega(g,h)f)_x=\bigl(\omega(g,h)\,\rho(h^{-1}X(x)g)f\bigr)_0 .\tag{M}
\]
Here $h^{-1}X(x)=X(h^{-1}x)$ and $X(x)\,m(a)\iota(g_0)=X(xa)$, and a
computation in the coordinates of the hyperbolic decomposition gives
\[
X(x)\,n_B(b)=X(x)+xb-Y\bigl(\tfrac12xbb^+\bigr),\qquad X(x)\,n_C(c)=X(x)+Y(xc),
\]
where $xb\in\mathbb W(V_0,W)$ and $b^+=G_{V_0}^{-1}\bar b^{\mathsf T}$. We use
Lemma~\ref{th-mixed-origin} below, whose proof uses only part (1).

\emph{The group of $W$.} By (M) and Lemma~\ref{th-mixed-origin}(1) for the
theta function $\rho(X(h^{-1}x))f$,
$(\omega(1,h)f)_x=\omega_0(1,h)(\rho(X(h^{-1}x))f)_0=\omega_0(1,h)f_{h^{-1}x}$.

\emph{The Levi factor and the core.} By (M) and
Lemma~\ref{th-mixed-origin}(2) for $\rho(X(xa))f$,
\[
(\omega(m(a)\iota(g_0),1)f)_x=\alpha(a)\,|\det a|_{\mathbf A_E}^{\operatorname{rk}W/2}\,\xi(g_0)\,\omega_0(g_0,1)f_{xa}
\]
with the characters $\alpha$ and $\xi$ of that lemma. It remains to show
$\xi=1$ and $\alpha(a)=\chi_W(\det a)^{-1}$; the formulas for $\iota(g_0)$ and
for $m(a)$ are the cases $a=1$ and $g_0=1$.
\begin{enumerate}
\item[(a)] $\xi=1$. The character $\xi$ is continuous, trivial on $U(V_0)(F)$
  and on the elements supported at the places of $F$ split in $E$. A
  continuous unitary character of the adelic points of a unitary group with
  these properties is trivial, as in the proof of
  Lemma~\ref{th-tensor-rule}.
\item[(b)] $\alpha$ is trivial on the matrices of determinant $1$. Let
  $i\neq j$, $t\in\mathbf A_E$, $u=1+tE_{ij}$, and let $d$ be the diagonal
  matrix with entry $2$ at $i$ and $1$ elsewhere. Then $du\,d^{-1}=u^2$, so
  $\alpha(u)=\alpha(u)^2$ and $\alpha(u)=1$. The ring $\mathbf A_E$ has stable
  rank one, so the matrices of determinant $1$ over it are products of such
  $u$.
\item[(c)] Let $\delta(t)$ be the diagonal matrix with entries
  $t,1,\ldots,1$. The character $t\mapsto\alpha(\delta(t))\chi_W(t)$ of
  $\mathbf A_E^\times$ is continuous and trivial on $E^\times$, since $\alpha$
  is trivial on $GL_\rho(E)$ and $\chi_W$ on $E^\times$, and it is trivial on
  the elements supported at the split places, by
  Lemma~\ref{th-mixed-origin}(2). A Hecke character of $E$ with this last
  property is trivial: by strong approximation every idele is the product of
  an idele with finite component $1$ and an element of the group $H$
  generated by $E^\times$, the ideles supported at the split places and the
  ideles congruent to $1$ modulo a given integer; the ideles with finite
  component $1$ lie in the closure of $H$; and the character is trivial on an
  open subgroup of the last kind. So $\alpha(\delta(t))=\chi_W(t)^{-1}$, and by
  (b), $\alpha(a)=\alpha\bigl(a\,\delta(\det a)^{-1}\bigr)\,\alpha\bigl(\delta(\det a)\bigr)=\chi_W(\det a)^{-1}$.
\end{enumerate}

\emph{The unipotent radical.} The elements $Y(\tfrac12xbb^+)$, $X(x)$ and $xb$
are pairwise orthogonal for $\langle\ ,\ \rangle$: the pairing of the first two
is a multiple of $\langle xb,xb\rangle$, which is $0$ since the form is
alternating, and the others vanish by the orthogonality of the decomposition. So
$\rho(X(x)n_B(b))=\rho(-Y(\tfrac12xbb^+))\,\rho(X(x))\,\rho(xb)$. By (M),
Lemma~\ref{th-mixed-origin}(3), the identity $(\rho(Y(a'))F)_0=F_0$ of (1) and
then (1) twice,
\[
(\omega(n_B(b),1)f)_x=\bigl(\rho(X(x)n_B(b))f\bigr)_0
=\bigl(\rho(X(x))\rho(xb)f\bigr)_0=(\rho(xb)f)_x=\rho(xb)f_x .
\]
For $n_C(c)$ put $q=\tfrac12\langle X(x),Y(xc)\rangle$. The translations
satisfy $\rho(X(x))\rho(Y(xc))=\psi(q)\,\rho(X(x)+Y(xc))$, and by (1),
$(\rho(X(x))\rho(Y(xc))f)_0=(\rho(Y(xc))f)_x=\psi(2q)f_x$. So by (M) and
Lemma~\ref{th-mixed-origin}(3),
$(\omega(n_C(c),1)f)_x=(\rho(X(x)+Y(xc))f)_0=\psi(q)f_x$.
\end{proof}

\begin{remark}\label{th-mixed-model-literature}
Formulas of the shape of Lemma~\ref{th-mixed-model}(2) are classical for the
mixed models of the Weil representation over a local field; the metaplectic
representation of a $p$-adic symplectic group and its models are treated in
\cite{MVW87}. The proof above uses only the conditions of
Definition~\ref{th-weil-datum}: the operators at the origin of the mixed model
are determined up to characters by their norms and matrix coefficients
(Lemma~\ref{th-mixed-origin}), and the characters are fixed by their values at
rational points and at the places of $F$ split in $E$.
\end{remark}

\begin{lemma}[Mixed model at the origin]\label{th-mixed-origin}
Let $D$ be a definite-side datum and $\psi$ a standard character. Let
$V=V_0\oplus H^\rho$ be a hyperbolic decomposition with the elements $m(a)$,
$n_B(b)$ and $n_C(c)$ of $U(V)(\mathbf A)$ and the embedding
$\iota\colon U(V_0)\to U(V)$ (Definition~\ref{th-step}), let $W$ be a Hermitian
space, and let $\omega$ and $\omega_0$ be Weil data for $(V,W)$ and $(V_0,W)$
with the same characters $(\chi_V,\chi_W)$ (Definition~\ref{th-weil-datum}).
Fix a Haar measure on $M_{\operatorname{rk}W\times\rho}(\mathbf A_E)$ with a
fundamental domain $\mathcal F$ of volume $1$ for the rational points, and for
a theta function $f$ for $(V,W)$ let $f_0$ be its mixed model at $0$
(Lemma~\ref{th-mixed-model}). Let $P_V$ be the subgroup of
$U(V)(\mathbf A)$ generated by the elements $g$ with $ge_a=e_a$ for every $a$
and $gv-v$ in the span of the $e_a$ for every $v$ orthogonal to all the $e_a$.
\begin{enumerate}
\item For every theta function $f$ for $(V,W)$ and $h\in U(W)(\mathbf A)$,
  $(\omega(1,h)f)_0=\omega_0(1,h)f_0$.
\item There are continuous homomorphisms $\alpha$ from $GL_\rho(\mathbf A_E)$
  and $\xi$ from $U(V_0)(\mathbf A)$ to the complex numbers of absolute value
  $1$, trivial on $GL_\rho(E)$ and on $U(V_0)(F)$, such that
  \[
  (\omega(m(a)\iota(g_0),1)f)_0=\alpha(a)\,|\det a|_{\mathbf A_E}^{\operatorname{rk}W/2}\,\xi(g_0)\,\omega_0(g_0,1)f_0
  \]
  for every theta function $f$, $a\in GL_\rho(\mathbf A_E)$ and
  $g_0\in U(V_0)(\mathbf A)$. Moreover $\xi(g_0)=1$ and
  $\alpha(a)=\chi_W(\det a)^{-1}$ whenever $g_0$ and $a$ are supported at a
  place of $F$ split in $E$.
\item $(\omega(n,1)f)_0=f_0$ for every theta function $f$ and every $n\in P_V$;
  this applies to $n=n_B(b)$ and $n=n_C(c)$.
\end{enumerate}
\end{lemma}
\leanref{Development/AbelianVarietyTheory/Proofs/DecoratedLatticeProofs/MixedCoordinateOriginWeilRepTargetEquivariantStatement.lean}{GlobalTheory.T.mixedCoordinate_origin_weilRep_target_equivariant}
\leanref{Development/AbelianVarietyTheory/Proofs/DecoratedLatticeProofs/WeilActionLeviOriginEqCharScalarStatement.lean}{GlobalTheory.T.weilAction_levi_origin_eq_charScalar}
\leanref{Development/AbelianVarietyTheory/Proofs/DecoratedLatticeProofs/LeviFactorCharactersSplitPlaceStatement.lean}{GlobalTheory.T.leviFactor_characters_splitPlace}
\leanref{Development/AbelianVarietyTheory/Proofs/DecoratedLatticeProofs/MixedCoordinateOriginUnipotentRadicalInvariantStatement.lean}{GlobalTheory.T.mixedCoordinate_origin_unipotentRadical_invariant}
\leanref{Development/AbelianVarietyTheory/Proofs/TwistedOperators/RestrictedThetaCoefficientEqIsotropicIntegralStatement.lean}{GlobalTheory.T.restrictedTheta_coefficient_eq_isotropic_integral}
\leanref{Development/AbelianVarietyTheory/Proofs/IwasawaDecomposition/ThetaFunctionEqUnitSmulOfDiagonalCoefficientEq.lean}{GlobalTheory.ThetaCoefficient.exists_unit_smul_of_coefficient_eq}
\leanref{Development/AbelianVarietyTheory/Proofs/TwistedOperators/WittTowerRestrictionIntertwinesWeilUpToCharacterStatement.lean}{GlobalTheory.T.wittTowerRestriction_intertwines_weil_upToCharacter}
\leanref{Development/AbelianVarietyTheory/Proofs/TwistedOperators/RestrictionWeilDataTwistCharacterTrivialSplitPlaStatement.lean}{GlobalTheory.T.restriction_weilData_twistCharacter_trivial_splitPlaces}
\leanref{Development/AbelianVarietyTheory/Proofs/TwistedOperators/RestrictionWittTowerIntertwinesLeviWeilRepStatement.lean}{GlobalTheory.T.restriction_wittTower_intertwines_levi_weilRep}
\leanref{Development/AbelianVarietyTheory/Proofs/TwistedOperators/LeviElementScalesRestrictionNormSqAbsDetStatement.lean}{GlobalTheory.T.leviElement_scales_restrictionNormSq_absDet}
\leanref{Development/AbelianVarietyTheory/Proofs/TwistedOperators/IsotropicMeanChangesByScalarOfTrivialActionStatement.lean}{GlobalTheory.T.isotropicMean_changes_by_scalar_of_trivial_action}
\leanref{Development/AbelianVarietyTheory/Proofs/TwistedOperators/UnipotentLeviMatrixActionIdentityStatement.lean}{GlobalTheory.T.unipotent_levi_matrix_action_identity}
\begin{proof}
If $\rho=0$, then $V=V_0$, $f_0=f$, the elements $m(a)$, $n_B(b)$ and
$n_C(c)$ are $1$, and $\omega$ and $\omega_0$ agree on theta functions, since
two Weil data for the same pair with the same characters agree on theta
functions; the lemma holds with $\alpha=\xi=1$. Let $\rho\geq1$. By
Lemma~\ref{th-mixed-model}(1), $f_0$ is a theta function for $(V_0,W)$,
$(\rho(w')f)_0=\rho(w')f_0$ for $w'\in\mathbb W(V_0,W)$, and
$(\rho(Y(a'))f)_0=f_0$ for every $a'$.

\emph{Norms and coefficients.} Fix Haar measures on $\mathbb W(V,W)$ and
$\mathbb W(V_0,W)$ with fundamental domains $\mathcal F_V$ and $\mathcal F_0$
of volume $1$ for their rational points, write $\|u\|$ for the $L^2$-norm on
these domains, and put $c_u(w)=\int_{\mathcal F_0}(\rho(w)u)\bar u$ for a theta
function $u$ for $(V_0,W)$ and $w\in\mathbb W(V_0,W)$. Four facts are used.
\begin{enumerate}
\item[(C1)] There is a Haar measure $dy$ on
  $M_{\operatorname{rk}W\times\rho}(\mathbf A_E)$ such that, for theta
  functions $f,g$ for $(V,W)$ and $w\in\mathbb W(V_0,W)$,
  $\int_{\mathcal F_0}(\rho(w)f_0)\overline{g_0}
  =\int\int_{\mathcal F_V}(\rho(Y(y)+w)f)\,\bar g\,dy$, the outer integral over
  all of $M_{\operatorname{rk}W\times\rho}(\mathbf A_E)$. This is a computation:
  unfold the integrals over $\mathcal F$ defining $f_0$ and $g_0$ and use the
  quasi-periodicity of $f$ and $g$.
\item[(C2)] If $u$ and $v$ are theta functions for $(V_0,W)$ with $c_u=c_v$,
  then $v=cu$ with $|c|=1$.
\item[(C3)] Every theta function for $(V_0,W)$ is $f_0$ for some theta function
  $f$ for $(V,W)$.
\item[(C4)] If a linear operator $M$ on the theta functions for $(V_0,W)$
  commutes with every $\rho(w)$ and preserves the norm, then it is a scalar of
  absolute value $1$; this is Schur's lemma for the translations, as in the
  proof of Lemma~\ref{th-tensor-rule}.
\end{enumerate}
The rational points act on theta functions by transport:
$(\omega(\gamma,\gamma')F)(\gamma'w\gamma^{-1})=F(w)$ for rational
$(\gamma,\gamma')$, by the value at $0$ and the relation of
Definition~\ref{th-weil-datum}(3).

\emph{Part (1).} Fix $h$. Since $\omega(1,h)$ and $\omega_0(1,h)$ preserve the
norms, the relation of Definition~\ref{th-weil-datum}(3) and (C1) express
$c_{(\omega(1,h)f)_0}(w)$ as the integral over $y$ of the coefficient of $f$ at
$Y(h^{-1}y)+h^{-1}w$, and $c_{\omega_0(1,h)f_0}(w)$ as the same integral at
$Y(y)+h^{-1}w$. The change of variables $y\mapsto h^{-1}y$ has modulus $1$,
since $|\det h|_{\mathbf A_E}=1$, so the two coefficients agree, and by (C2)
$(\omega(1,h)f)_0=c\,\omega_0(1,h)f_0$ with $|c|=1$. Comparing $f$ with a fixed
$f'$ with $f'_0\neq0$, through the coefficients of $f+f'$, shows that
$c=c(h)$ does not depend on $f$. Since $\omega$ and $\omega_0$ are
homomorphisms, $c$ is a homomorphism, and it is continuous by the continuity
of the Weil data. It is trivial on $U(W)(F)$: for rational $h$, the action by
transport and the invariance of the integral over $\mathcal F$ under
$y\mapsto h^{-1}y$ give $(\omega(1,h)f)_0(hw)=f_0(w)=(\omega_0(1,h)f_0)(hw)$.
It is trivial on the elements $h$ supported at a place $v$ of $F$ split in $E$.
For this, take a linear functional $\ell_0$ on the theta functions for
$(V_0,W)$ invariant under the half of $\mathbb W(V_0,W)$ attached to $v$, and
let $\ell(f)$ be the average of $\ell_0((\rho(u)f)_0)$ over the elements $u$ of
the corresponding half of the span of the $Y(y)$; then $\ell$ is invariant
under the half of $\mathbb W(V,W)$ attached to $v$ and nonzero on some $f$. By
the condition at split places of Definition~\ref{th-weil-datum}(3), $\omega(1,h)$
multiplies $\ell$ by $r\chi_W(y_h)$ and $\omega_0(1,h)$ multiplies $\ell_0$ by
$r_0\chi_W(y_h)$, with $r,r_0>0$ and $y_h$ the determinant of $h$ at $v$, and
the change of variables in the average contributes a factor $\kappa>0$. So
$r=c(h)\,r_0\,\kappa$, and $c(h)$, a positive number of absolute value $1$, is
$1$. A continuous unitary character of $U(W)(\mathbf A)$ trivial on $U(W)(F)$
and on the elements supported at the split places is trivial, as in the proof
of Lemma~\ref{th-tensor-rule}. So $c=1$.

\emph{Part (2).} Fix $a$ and $g_0$, and put $m=m(a)\iota(g_0)$. A computation
with (C1) and the change of variables $y\mapsto ya$ gives
$\|(\omega(m,1)f)_0\|^2=|\det a|_{\mathbf A_E}^{\operatorname{rk}W}\|f_0\|^2$.
So by (C3) there is a well defined operator $M$ on the theta functions for
$(V_0,W)$ with
$M(f_0)=|\det a|_{\mathbf A_E}^{-\operatorname{rk}W/2}\,\omega_0(g_0,1)^{-1}(\omega(m,1)f)_0$.
It preserves the norm, by the computation and the unitarity of
$\omega_0(g_0,1)$, and it commutes with the translations, since
$(\rho(w)f)_0=\rho(w)f_0$ and both $\omega(m,1)$ and $\omega_0(g_0,1)$ move
$\rho(w)$ to $\rho(wg_0^{-1})$. By (C4), $M$ is a scalar $c(a,g_0)$ of absolute
value $1$, which gives the formula. Put $\alpha(a)=c(a,1)$ and
$\xi(g_0)=c(1,g_0)$. Since some $f_0$ is nonzero,
$c(a,g_0)=\alpha(a)\xi(g_0)$ and $\alpha$ and $\xi$ are homomorphisms; they are
continuous by the continuity of the Weil data, and trivial at rational points
by the action by transport, as in part (1). At a split place, the same
comparison of the condition at split places for $\omega$ and $\omega_0$, applied
to $m$, whose determinant there is read from $\det a$ and $\det g_0$, shows that
$\xi(g_0)$ and $\alpha(a)\chi_W(\det a)$ are positive numbers of absolute value
$1$, hence equal to $1$.

\emph{Part (3).} Let $n$ be a generator of $P_V$. Then $Y(y)n=Y(y)$ for every
$y$, and $wn-w$ lies in the span of the $Y(y)$ for every $w\in\mathbb W(V_0,W)$.
The comparison of coefficients of part (1), with the translation of $y$ by an
element depending on $w$ in place of $y\mapsto h^{-1}y$, gives
$(\omega(n,1)f)_0=c(n)f_0$ with $|c(n)|=1$, and then the same holds on all of
$P_V$ with a multiplicative $c$. Let $s\in U(V)(F)$ be the rational element
with $se_a=2e_a$ and $se'_a=\tfrac12e'_a$ for every $a$ and $s=1$ on $V_0$. By
the action by transport, $(\omega(s,1)f)_0$ is the mean of $f$ over the
$Y(\tfrac12y)$, which equals the mean over the $Y(y)$ because the integrand is
invariant under rational translations; so $(\omega(s,1)f)_0=f_0$. Every
generator $n$ of $P_V$ is a commutator $sms^{-1}m^{-1}$ with $m\in P_V$, so
$(\omega(n,1)f)_0=c(m)c(m)^{-1}f_0=f_0$, and $c=1$ on $P_V$. The elements
$n_B(b)$ and $n_C(c)$ fix the $e_a$ and move the vectors orthogonal to the
$e_a$ by elements of their span, so they lie in $P_V$.
\end{proof}

\begin{lemma}[Unfolding]\label{th-unfolding}
Let $D$ be a definite-side datum and $\psi$ a standard character. Let
$V=V_0\oplus H^\rho$ be a hyperbolic decomposition with flag group $N_V$
(Definition~\ref{th-step}), $W$ a Hermitian space and $\omega$ a Weil datum for
$(V,W)$. Fix Haar measures with fundamental domains for $U(W)(F)$ and, of
volume $1$, for $N_V(F)$, for $U_\rho(E)$ in $U_\rho(\mathbf A_E)$, for the
rational points of $M_{\rho\times\operatorname{rk}V_0}(\mathbf A_E)$ and for
those of $M_{\operatorname{rk}W\times\rho}(\mathbf A_E)$. Use the following
coordinates on $N_V(\mathbf A)$: for $u\in U_\rho(\mathbf A_E)$, $m(u)$ sends
the matrix $(e_1,\ldots,e_\rho)$ to $(e_1,\ldots,e_\rho)u$, fixes $V_0$ and
sends $(e'_1,\ldots,e'_\rho)$ to $(e'_1,\ldots,e'_\rho)(u^*)^{-1}$; the elements
$n_B(b)$, $b\in M_{\rho\times\operatorname{rk}V_0}(\mathbf A_E)$, and $n_C(c)$,
$c$ a skew-Hermitian $\rho\times\rho$ matrix, make up the unipotent radical of
the stabilizer of the span of the $e_a$, the $n_C(c)$ forming its centre; and
every element of $N_V(\mathbf A)$ is $n_C(c)n_B(b)m(u)$. For
$x\in M_{\operatorname{rk}W\times\rho}(\mathbf A_E)$ write $X(x)$ and $Y(x)$
for the elements of $\mathbb W(V,W)$ supported on the $e_a$ and on the $e'_a$
with coefficient matrix $x$, and for a theta function $\Psi$ for $(V,W)$ let
$\Psi_x$ be the function on $\mathbb W(V_0,W)$ given by the first stage of the
flat map in Definition~\ref{th-step}(4):
$\Psi_x(w')=\int(\rho(Y(a))\rho(X(x))\rho(w')\Psi)(0)\,da$, the integral over
the fundamental domain for $M_{\operatorname{rk}W\times\rho}(E)$.

Let $x_0\in M_{\operatorname{rk}W\times\rho}(E)$ have linearly independent
columns, and let $G_0=x_0^*G_Wx_0$ satisfy $u^*G_0u=G_0$ for
$u\in U_\rho(\mathbf A_E)$. Let $\Theta$ be a function on $N_V(\mathbf A)$ such
that:
\begin{enumerate}
\item $\Theta$ is continuous and left invariant under $N_V(F)$;
\item $\Theta(n_C(c)n)=\overline{\psi(q(c))}\,\Theta(n)$, where
  $q(c)=\tfrac12\langle X(x_0),Y(x_0c)\rangle$;
\item for every $j$ and every $a\in E^\rho$ with $a_j=1$, $a_i=0$ for $i>j$ and
  $G_0a=0$, there are an element $t\in N_V(\mathbf A)$ that fixes
  $e_1,\ldots,e_j$ and is the transvection along $\sum_ia_ie_i$ and a vector of
  $V\otimes\mathbf A$ orthogonal to it, and $c\neq1$, with
  $\Theta(ntn')=c\,\Theta(nn')$ for all $n,n'\in N_V(\mathbf A)$.
\end{enumerate}
For a theta function $\Psi$ for $(V,W)$ put
\[
\ell_\Theta(\Psi)=\int_{U_\rho(\mathbf A_E)}\int\Theta(n_B(b)m(u))\,\Psi_{x_0u}(x_0b)\,db\,du ,
\]
the outer integral over the whole group $U_\rho(\mathbf A_E)$ against the fixed
Haar measure, and the inner integral over the fundamental domain for
$M_{\rho\times\operatorname{rk}V_0}(E)$. Let $G$ be rapidly decreasing on
$U(W)(\mathbf A)$, $\Phi$ a theta function for $(V,W)$, and $\mathcal B$ a
fundamental domain for $U(W)(F)\cap R_{x_0}$, with $R_{x_0}$ as in
Definition~\ref{th-step-local}. Then $h\mapsto G(h)\ell_\Theta(\omega(1,h)\Phi)$
is integrable on $\mathcal B$, and
\[
\int_{N_V(F)\backslash N_V(\mathbf A)}\Theta(n)\,\theta_\Phi(G)(n)\,dn
=\int_{\mathcal B}G(h)\,\ell_\Theta(\omega(1,h)\Phi)\,dh .
\]
\end{lemma}
\leanref{Development/AbelianVarietyTheory/Proofs/DecoratedLatticeProofs/UnfoldOverFlagUnipotentThetaLift.lean}{GlobalTheory.Unf.unfold_over_N_lift}
\leanref{Development/AbelianVarietyTheory/Proofs/DecoratedLatticeProofs/UnipotentMeanThetaLiftEqOrbitIntegralStatement.lean}{GlobalTheory.T.unipotentMean_thetaLift_eq_orbitIntegral}
\leanref{Development/AbelianVarietyTheory/Proofs/DecoratedLatticeProofs/MixedThetaSeriesTermsLawsUnipotentRationalStatement.lean}{GlobalTheory.T.mixedThetaSeries_terms_laws_unipotent_rational}
\leanref{Development/AbelianVarietyTheory/Proofs/DecoratedLatticeProofs/FlagUnitaryOrbitActionTransitiveWithStabilizerStatement.lean}{GlobalTheory.T.flagUnitaryOrbitAction_transitive_with_stabilizer}
\leanref{Development/AbelianVarietyTheory/Proofs/DecoratedLatticeProofs/FlagUnipotentFundamentalDomainLayersIntegralStatement.lean}{GlobalTheory.T.flagUnipotent_fundamentalDomain_layers_integral}
\begin{proof}
For a theta function $\Psi$ for $(V,W)$ and
$x\in M_{\operatorname{rk}W\times\rho}(E)$ put $t_\Psi(x)=\Psi_x(0)$.

\emph{Step 1: the terms.} In the notation of Lemma~\ref{th-mixed-model},
$\Psi_x$ is the mixed model of $\Psi$ at $x$, so $t_\Psi(x)=\Psi_x(0)$. Let
$\omega_0$ be a Weil datum for $(V_0,W)$ with the characters of $\omega$; one
exists, by the existence of Weil data with given characters. The
following hold for every theta function $\Psi$.
\begin{enumerate}
\item[(a)] $\Psi(0)=\sum_xt_\Psi(x)$, the sum over
  $x\in M_{\operatorname{rk}W\times\rho}(E)$ converging. This is
  Lemma~\ref{th-mixed-model}(1) at $w=0$.
\item[(b)] $t_{\omega(n_C(c)n_B(b)m(u),1)\Psi}(x)=\psi(q_x(c))\,\Psi_{xu}(xb)$
  for $u\in U_\rho(\mathbf A_E)$, where $q_x(c)=\tfrac12\langle X(x),Y(xc)\rangle$.
  Indeed, by the formulas of Lemma~\ref{th-mixed-model}(2) for $n_C(c)$, $n_B(b)$
  and $m(u)$ in turn, the left side is
  $\psi(q_x(c))\,(\rho(xb)(\omega(m(u),1)\Psi)_x)(0)=\psi(q_x(c))\,(\omega(m(u),1)\Psi)_x(xb)$,
  and $(\omega(m(u),1)\Psi)_x=\Psi_{xu}$ because $\det u=1$.
\item[(c)] Every $\gamma\in N_V(F)$ is $m(u)n_B(b)n_C(c)$ with $u$, $b$, $c$
  rational, and $t_{\omega(\gamma,1)\Psi}(x)=t_\Psi(xu)$. Indeed, by the same
  formulas the left side is $\psi(q_{xu}(c))\Psi_{xu}(xub)$; the scalar is $1$
  because $\psi$ is trivial on $F$ and $q_{xu}(c)\in F$ for rational $xu$ and $c$;
  and $\Psi_{xu}(xub)=\Psi_{xu}(0)$ by the quasi-periodicity of the theta
  function $\Psi_{xu}$ at the rational point $xub$. For $\delta\in U(W)(F)$,
  $t_{\omega(1,\delta)\Psi}(x)=(\omega_0(1,\delta)\Psi_{\delta^{-1}x})(0)=t_\Psi(\delta^{-1}x)$
  by Lemma~\ref{th-mixed-model}(2) and Definition~\ref{th-weil-datum}(3) at
  rational points.
\item[(d)] $(n,h)\mapsto t_{\omega(n,h)\Psi}(x)$ is continuous, and there are
  $C$ and $A$ with
  $\sum_x|t_{\omega(n,h)\Psi}(x)|\leq C\operatorname{ht}(h)^A$ for
  $n\in N_V(\mathbf A)$ and $h\in U(W)(\mathbf A)$.
\end{enumerate}
The continuity in (d) follows from the joint continuity of
Lemma~\ref{th-weil-estimates}(2), the integral defining $\Psi_x(0)$ being over
the compact set $\mathcal F$ of a function invariant under rational
translations. For the bound we use three facts.
\begin{enumerate}
\item[(d1)] There is $k$ such that for every neighbourhood $K_0$ of $0$ in
  $\mathbb W(V,W)$ there is $C_0$ with the following property: if $n\geq1$ is an
  integer and $F$ is a theta function with $\rho(nu)F=F$ for every $u\in K_0$ of
  archimedean component $0$, then $\sum_x|F_x(w)|\leq C_0n^k\|F\|_k$ for every
  $w$, where $\|F\|_k$ is the sum of the supremum norms of the derivatives of $F$
  of order at most $k$ along the archimedean directions. The $F_x(w)$ are the
  Fourier coefficients of a function on the compact quotient of
  $M_{\operatorname{rk}W\times\rho}(\mathbf A_E)$; the invariance puts the
  finite components of the $x$ with $F_x(w)\neq0$ in $\frac1n$ times a fixed
  lattice, and integration by parts along the archimedean directions gives
  $|F_x(w)|\leq C\|F\|_k(1+|x_\infty|)^{-k}$, which is summable over such $x$
  with sum at most a constant times $n^k\|F\|_k$ for $k$ large.
\item[(d2)] There are a neighbourhood $K_0$ of $0$ and $m\geq1$ such that for all
  $g,h$ there is an integer $1\leq n\leq m\operatorname{ht}_0(g)\operatorname{ht}_0(h)$
  with $\rho(nu)\,\omega(g,h)\Psi=\omega(g,h)\Psi$ for $u\in K_0$ of archimedean
  component $0$. Indeed $\Psi$ is invariant under an open subgroup $U_\Psi$, so
  $\omega(g,h)\Psi$ is invariant under $hU_\Psi g^{-1}$ by
  Definition~\ref{th-weil-datum}(3), and $h^{-1}(nK_0)g\subset U_\Psi$ when $n$
  clears the denominators of $g^{-1}$ and $h$, which are bounded by their
  heights.
\item[(d3)] For every $k$ there are $c$ and $A$ with
  $\|\omega(g,h)\Psi\|_k\leq c(\operatorname{ht}_0(g)\operatorname{ht}_0(h))^A$.
  A derivative of $\omega(g,h)\Psi$ along an archimedean direction is, by the
  relation of Definition~\ref{th-weil-datum}(3), $\omega(g,h)$ applied to the
  derivative of $\Psi$ along the direction transported by $(g,h)$, whose size is
  polynomial in the archimedean norms of $g$ and $h$; and
  Lemma~\ref{th-weil-estimates}(3), applied to the derivatives of $\Psi$, which
  are theta functions, bounds the supremum norms.
\end{enumerate}
By (d1) to (d3), $\sum_x|t_{\omega(g,h)\Psi}(x)|\leq C(\operatorname{ht}_0(g)\operatorname{ht}_0(h))^A$.
Now let $n\in N_V(\mathbf A)$ and $h\in U(W)(\mathbf A)$. Write $n=\gamma n_1$
with $\gamma\in N_V(F)$ and $n_1$ in a fixed compact set, which is possible
because $N_V(F)\backslash N_V(\mathbf A)$ is compact, and $h=\delta h_1$ with
$\delta\in U(W)(F)$ and $\operatorname{ht}_0(h_1)\leq2\operatorname{ht}(h)$. By (c),
the terms of $\omega(\gamma n_1,\delta h_1)\Psi=\omega(\gamma,\delta)\omega(n_1,h_1)\Psi$
are those of $\omega(n_1,h_1)\Psi$ at the points $\delta^{-1}xu$, and
$x\mapsto\delta^{-1}xu$ is a bijection of the rational matrices. So the sum of
their absolute values is at most $C(\operatorname{ht}_0(n_1)\cdot2\operatorname{ht}(h))^A$,
which is (d).

\emph{Step 2: the orbit.} Let $O$ be the set of
$x\in M_{\operatorname{rk}W\times\rho}(E)$ with linearly independent columns
and $x^*G_Wx=G_0$. By (d), $\Theta$ is bounded, being continuous and invariant
under $N_V(F)$ with compact quotient, and $G$ is rapidly decreasing, so the
integrals below converge absolutely and may be exchanged. For each $h$,
\[
\int\Theta(n)\,(\omega(n,1)\Psi)(0)\,dn=\int\Theta(n)\sum_{x\in O}t_{\omega(n,1)\Psi}(x)\,dn
\]
over $N_V(F)\backslash N_V(\mathbf A)$, with $\Psi=\omega(1,h)\Phi$. By (a) it
suffices to show that the sum over the rational $x\notin O$ has integral $0$.
For a set $P$ of rational matrices put
$F_P(n)=\sum_{x\in P}\Theta(n)t_{\omega(n,1)\Psi}(x)$; by (d) with $h=1$ and the
boundedness of $\Theta$, $\sum_x|\Theta(n)t_{\omega(n,1)\Psi}(x)|$ is bounded
uniformly in $n$, so sums and integrals may be exchanged below. If $P$ is
stable under $x\mapsto xu$ for the unitriangular $u\in U_\rho(E)$, then $F_P$ is
left invariant under $N_V(F)$, by condition (1) and (c).

\emph{The terms with $x^*G_Wx\neq G_0$.} Let $P$ be the set of these $x$; it is
stable under $x\mapsto xu$, since $u^*G_0u=G_0$. Let $Z$ be the centre of
$N_V$, the group of the $n_C(c)$; it is closed and normal, and $Z(F)\backslash
Z(\mathbf A)$ is compact. By (b) and condition (2), for $x\in P$,
\[
\Theta(n_C(c)n)\,t_{\omega(n_C(c)n,1)\Psi}(x)=a_x(c)\,\Theta(n)\,t_{\omega(n,1)\Psi}(x),\qquad
a_x(c)=\overline{\psi(q(c))}\,\psi(q_x(c)),
\]
where we use that $\omega(n_C(c)n,1)=\omega(n_C(c),1)\omega(n,1)$ and the formula
of Lemma~\ref{th-mixed-model}(2) for $n_C(c)$. The map $c\mapsto a_x(c)$ is a
continuous character, trivial on the rational $c$ because $\psi$ is trivial on
$F$, and nontrivial: $q_x(c)-q(c)$ depends on $c$ through
$x^*G_Wx-G_0$, by an $\mathbf A_F$-linear form in the skew-Hermitian matrix $c$
built from the trace and $\xi_0$, which is nonzero when $x^*G_Wx\neq G_0$, and
$\psi$ is nontrivial. Integrating over
$N_V(F)\backslash N_V(\mathbf A)$ first over $Z(F)\backslash Z(\mathbf A)$,
\[
\int F_P(n)\,dn=\int\sum_{x\in P}\Theta(n)t_{\omega(n,1)\Psi}(x)\int_{Z(F)\backslash Z(\mathbf A)}a_x(c)\,dc\,dn=0 .
\]

\emph{The terms with $x^*G_Wx=G_0$ and dependent columns.} For $1\leq j\leq\rho$ let
$P_j$ be the set of these $x$ whose first relation is at $j$: there is $a\in E^\rho$
with $a_j=1$, $a_i=0$ for $i>j$ and $xa=0$, and no nonzero relation supported on
the indices $i<j$. The vector $a=a(x)$ is then unique. Unitriangular $u$
preserve $P_j$ (with $a(xu)=u^{-1}a(x)$), so $F_{P_j}$ is invariant under
$N_V(F)$. Let $N_j\subset N_V$ be the subgroup of elements fixing
$e_1,\ldots,e_j$, which is closed and normal with $N_j(F)\backslash N_j(\mathbf A)$
compact. For $\gamma\in N_j(F)$, the block $u$ of $\gamma$ fixes the vectors
supported on the first $j$ indices, so $a(xu)=a(x)$, and the partial sums
$F_{P_j,a}$ over the $x\in P_j$ with $a(x)=a$ are invariant under $N_j(F)$.
For each $a$ occurring, $G_0a=x^*G_Wxa=0$, so condition (3) gives $t\in N_j(\mathbf A)$
and $c\neq1$ with $\Theta(mtn)=c\,\Theta(mn)$; and for $m\in N_j(\mathbf A)$
and $x$ with $xa=0$, $t_{\omega(mtn,1)\Psi}(x)=t_{\omega(mn,1)\Psi}(x)$, because
the transvection $t$ moves only along $\sum_ia_ie_i$ and the vector orthogonal
to it, which $x$ does not see. Hence $F_{P_j,a}(mtn)=c\,F_{P_j,a}(mn)$ for
$m\in N_j(\mathbf A)$. Integrating over $N_V(F)\backslash N_V(\mathbf A)$ first
over $N_j(F)\backslash N_j(\mathbf A)$, the inner integral of $F_{P_j,a}(mn)$ in
$m$ is unchanged by the right translation $m\mapsto mt$ and multiplied by
$c\neq1$, so it vanishes; summing over $a$ gives $\int F_{P_j}=0$.
Inserting $\theta_\Phi(G)(n)=\int G(h)(\omega(n,h)\Phi)(0)\,dh$ and exchanging
the integrals, the left side of the lemma is
\[
\int\!\!\int\Theta(n)\,G(h)\sum_{x\in O}t_{\omega(n,h)\Phi}(x)\,dn\,dh,
\]
over $N_V(F)\backslash N_V(\mathbf A)$ and $U(W)(F)\backslash U(W)(\mathbf A)$.

\emph{Step 3: unfolding the orbit.} The group $N_V(F)\times U(W)(F)$ acts on
$O$ through (c), and the action is transitive by Witt's theorem. The stabilizer
of $x_0$ is the product of the rational points of the subgroup $N_0$ of $N_V$
of the elements $n_B(b)n_C(c)$, which fix $x_0$ because $u=1$ and are the only
ones to do so because $x_0$ has linearly independent columns, and of
$U(W)(F)\cap R_{x_0}$. The integrand
$\Theta(n)G(h)t_{\omega(n,h)\Phi}(x)$ is invariant under the simultaneous
action of $N_V(F)\times U(W)(F)$ on $(n,h)$ and on $x$, by condition (1), the
invariance of $G$ and (c). Unfolding the sum over the orbit, the integral of
Step 2 equals
\[
\int_{\mathcal A}\int_{\mathcal B}\Theta(n)\,G(h)\,t_{\omega(n,h)\Phi}(x_0)\,dh\,dn,
\]
where $\mathcal A$ is a fundamental domain for $N_0(F)$ in $N_V(\mathbf A)$.

\emph{Step 4: the layers.} By (b) and condition (2), the function
$n\mapsto\Theta(n)t_{\omega(n,h)\Phi}(x_0)$ is invariant under left
multiplication by $n_C(c)$. For such functions the integral over $\mathcal A$ is
the iterated integral of the value at $n_B(b)m(u)$ over $u\in U_\rho(\mathbf A_E)$
and over $b$ in the fundamental domain for $M_{\rho\times\operatorname{rk}V_0}(E)$;
the variable $u$ runs over the whole group, since $N_0$ has $u=1$. By (b),
$t_{\omega(n_B(b)m(u),h)\Phi}(x_0)=(\omega(1,h)\Phi)_{x_0u}(x_0b)$. So the
integral of Step 3 is $\int_{\mathcal B}G(h)\ell_\Theta(\omega(1,h)\Phi)\,dh$,
and the absolute convergence of Step 2 gives the integrability.
\end{proof}

\begin{lemma}\label{th-bessel-unfolding}
Let $D$ be a definite-side datum and $\psi$ a standard character, and let a
step from Bessel data $V=V_0\oplus H^\rho$, $V_0=A\oplus L$ with vector
$\zeta$, to Fourier-Jacobi data $W=W_0\oplus H^\rho$ be given, with its flat map
at $(x_0,y_0)$ (Definition~\ref{th-step}). For $n\in N_V(\mathbf A)$ write
$n e_q=e_q+\sum_{p<q}u_{pq}(n)e_p$ and $n\zeta=\zeta+\sum_pb_p(n)e_p+(\text{a vector of }V_0)$,
so that the Bessel character is
$\psi_\zeta(n)=\psi\bigl(\operatorname{Tr}_{\mathbf A_E/\mathbf A_F}(\sum_{p<\rho}u_{p,p+1}(n)+b_\rho(n))\bigr)$.
\begin{enumerate}
\item The matrix $x_0$ of $e_1,\ldots,e_\rho\in W$ has linearly independent columns
  and $x_0^*G_Wx_0=0$; the group $R_{x_0}$ of Definition~\ref{th-step-local} is
  the Jacobi group $N_W\rtimes U(W_0)$; and for $n\in N_W(\mathbf A)$ and
  $h\in U(W_0)(\mathbf A)$, $n\iota(h)$ is rational if and only if $n$ and $h$
  are.
\item $\Theta=\overline{\psi_\zeta}$ satisfies conditions (1) to (3) of
  Lemma~\ref{th-unfolding} for $x_0$ and $G_0=0$.
\item $\ell_\Theta(\Psi)=\Psi^\flat(0)$ for every theta function $\Psi$ for
  $(V,W)$, with $\ell_\Theta$ as in Lemma~\ref{th-unfolding}.
\end{enumerate}
\end{lemma}
\leanref{Development/AbelianVarietyTheory/Proofs/TwistedOperators/BesselLevelUnfoldingCharacterAndFlatMapStatement.lean}{GlobalTheory.T.besselLevel_unfolding_character_and_flatMap}
\leanref{Development/AbelianVarietyTheory/Proofs/DecoratedLatticeProofs/UnfoldingBasePointFourierJacobiGroupStatement.lean}{GlobalTheory.T.unfolding_base_point_fourierJacobi_group}
\leanref{Development/AbelianVarietyTheory/Definitions/TwistedInnerProducts/WittSmallGroupEmbeddingMemPointwiseStabilizerIso.lean}{GlobalTheory.MixPer.psiBes}
\begin{proof}
For (1): the $e_a$ are isotropic and pairwise orthogonal, so the Gram matrix
$x_0^*G_Wx_0$ is $0$, and they are linearly independent. An element $r$ with
$rx_0=x_0u$ for a unitriangular $u$ stabilizes the flag of the $e_a$ and acts on
it through $u$; it therefore also stabilizes the orthogonal of the flag, and
its action on the quotient, which is $W_0$, is an element $h\in U(W_0)$, so that
$r\iota(h)^{-1}\in N_W$. Conversely $N_W$ and $\iota(U(W_0))$ satisfy
$rx_0=x_0u$. The decomposition $r=n\iota(h)$ is unique and its factors are given
by polynomials in the entries of $r$ with coefficients in $F$, so $r$ is
rational if and only if $n$ and $h$ are.

For (2): $\psi_\zeta$ is a continuous character of $N_V(\mathbf A)$, trivial
on $N_V(F)$ because $\psi$ is trivial on $F$; this gives condition (1). The
elements $n_C(c)$ of the centre have $u_{p,p+1}=0$ and $b_\rho=0$, so
$\psi_\zeta(n_C(c))=1$; and $q(c)=\frac12\langle X(x_0),Y(x_0c)\rangle$ depends on
$c$ through $x_0^*G_Wx_0=0$, so $\psi(q(c))=1$. This gives condition (2). For
condition (3), let $j$ and $a$ be given, with $a_j=1$ and $a_i=0$ for $i>j$, and
choose $s\in\mathbf A_E$ with $\psi(\operatorname{Tr}s)\neq1$, which exists since
$\psi$ is nontrivial. There are a vector $v$ of $V\otimes\mathbf A$ orthogonal to
$v_a=\sum_ia_ie_i$ and the transvection $t$ along $v_a$ and $v$, an element of
$N_V(\mathbf A)$ fixing $e_1,\ldots,e_j$, with
$\sum_pu_{p,p+1}(t)+b_\rho(t)=s$. If $j<\rho$, the vector $v$ is built from
$e'_{j+1}$, so that $t$ moves $e_{j+1}$ by a multiple of $v_a$ and the
superdiagonal entry $u_{j,j+1}(t)$ is $s$, the others vanishing because $v$ is
orthogonal to $e_1,\ldots,e_j$ and $a_i=0$ for $i>j$. If $j=\rho$, it is built
from $\zeta$, so that $b_\rho(t)=s$; this uses $h(\zeta,\zeta)\neq0$ to make $v$
orthogonal to $v_a$. Since $\psi_\zeta$ is a
character, $\Theta(ntn')=\overline{\psi(\operatorname{Tr}s)}\,\Theta(nn')$, and
$c=\overline{\psi(\operatorname{Tr}s)}\neq1$.

For (3): for $u\in U_\rho(\mathbf A_E)$ and $b\in M_{\rho\times\operatorname{rk}V_0}(\mathbf A_E)$,
$\Theta(n_B(b)m(u))=\overline{\psi_U(u)}\,\overline{\psi_\zeta(n_B(b))}$, because
$\psi_\zeta(m(u))=\psi_U(u)$ and $\psi_\zeta$ is a character. The value
$\psi_\zeta(n_B(b))$ depends only on the component of $b$ along $\zeta$ in its
last row. The second stage of the flat map at $y_0$ integrates over the points
$x_0b$ of the half of $\mathbb W(V_0,W)$ attached to the $e_a$ of $W$, after
translating by the adjoint of $y_0$, whose only nonzero column is $-\zeta$; by
Lemma~\ref{th-mixed-model}(1) this translation multiplies the value at $x_0b$
by $\overline{\psi_\zeta(n_B(b))}$. So the integrand of the second stage at $b$
is $\overline{\psi_\zeta(n_B(b))}\,\Psi_{x_0u}(x_0b)$, which is the integrand of
$\ell_\Theta$ at $(b,u)$ without the factor $\overline{\psi_U(u)}$.
So the inner integral of $\ell_\Theta$ over $b$ is the second stage of the flat
map applied to $\Psi_{x_0u}$ and evaluated at $0$, and the outer integral over
$u$ against $\overline{\psi_U(u)}$ gives $\Psi^\flat(0)$.
\end{proof}

\begin{lemma}\label{th-fourier-strips}
Let $D$ be a definite-side datum, $\psi$ a standard character,
$V=V_0\oplus H^\rho$ a hyperbolic decomposition and $W$ a Hermitian space, with
$X(x)$ and $Y(x)$ as in Lemma~\ref{th-mixed-model}, and write
$\mathcal M=M_{\operatorname{rk}W\times\rho}(\mathbf A_E)$. Say that a set
$S\subset\mathcal M$ is \emph{read on the archimedean components} if $z\in S$
and $z'\in\mathcal M$ with the same archimedean components imply $z'\in S$.
\begin{enumerate}
\item Let $(Y,m)$ be a measure space, $F$ an integrable function on $Y$ and
  $x\colon Y\to\mathcal M$ a measurable map with
  $\int_Y\psi(\langle X(x(y)),Y(c)\rangle)F(y)\,dm(y)=0$ for every
  $c\in\mathcal M$ with finite component $0$. Then $\int_{x^{-1}(S)}F\,dm=0$ for
  every measurable $S$ read on the archimedean components.
\item Let $x_0\in M_{\operatorname{rk}W\times\rho}(E)$ have linearly independent
  columns, let $R=R_{x_0}$ (Definition~\ref{th-step-local}), and let $C$ be a
  compact set of elements of $U(W)(\mathbf A)$ with finite component $1$. The
  \emph{strip} $R(\mathbf A)\,C\,U(W)(\mathbf A)_f$, with $U(W)(\mathbf A)_f$ the
  elements with archimedean component $1$, is measurable, and there is a
  measurable $S\subset\mathcal M$ read on the archimedean components such that
  for all $h\in U(W)(\mathbf A)$ and $u\in U_\rho(\mathbf A_E)$,
  $h^{-1}x_0u\in S$ if and only if $h$ lies in the strip.
\end{enumerate}
\end{lemma}
\leanref{Development/AbelianVarietyTheory/Proofs/TwistedOperators/FourierUniquenessRealComponentsAdelicHomStatement.lean}{GlobalTheory.T.fourier_uniqueness_realComponents_adelicHom}
\leanref{Development/AbelianVarietyTheory/Proofs/TwistedOperators/AdelicStripSetMeasurableReadOnRealPartStatement.lean}{GlobalTheory.T.adelicStripSet_measurable_read_on_realPart}
\leanref{Development/AbelianVarietyTheory/Definitions/TwistedInnerProducts/PartFixed.lean}{GlobalTheory.Asc.strip}
\begin{proof}
For (1), let $P\colon\mathcal M\to\mathbf C^{I}$, $I$ the set of triples
$(i,j,\tau)$, send $z$ to its archimedean coordinates $\tau(z_{ij})$, and let
$s\colon\mathbf C^I\to\mathcal M$ be a continuous section with finite
components $0$, so that $P(s(w))=w$. Since $S$ is read on the archimedean
components, $z\in S$ if and only if $s(P(z))\in S$, so
$x^{-1}(S)=(P\circ x)^{-1}(s^{-1}(S))$ with $s^{-1}(S)$ measurable. Let $\nu$ be
the image of the complex measure $F\,dm$ under $P\circ x$, a finite measure on
the real vector space $\mathbf C^I$. For every real linear form $l$ on
$\mathbf C^I$ there is $c$ with finite component $0$ and
$\exp(i\,l(P(z)))=\psi(\langle X(z),Y(c)\rangle)$ for all $z$: at the archimedean
places $\psi$ is $\exp(-2\pi i\sum_\tau\tau(\cdot))$, and the pairing
$\langle X(z),Y(c)\rangle$ is a nondegenerate real bilinear form in the
archimedean coordinates of $z$ and $c$. So the Fourier transform of $\nu$
vanishes, hence $\nu=0$ by the uniqueness theorem for Fourier transforms of
finite measures on $\mathbf R^d$ \cite[1.3.6]{Rud62}, and
$\int_{x^{-1}(S)}F\,dm=\nu(s^{-1}(S))=0$.

For (2), choose a rational $y$ with $yx_0=1$, and let $e_\infty\in\mathbf A_E$ be
the idempotent of the archimedean places. Let $Q$ be the set of $w\in\mathcal M$
with $x_0yw=w$ and $yw$ upper triangular with diagonal entries $e_\infty$, and
$S$ the set of $z$ with $e_\infty cz\in Q$ for some $c\in C$. Since $C$ is
compact and $Q$ closed, $S$ is closed, hence measurable; it is read on the
archimedean components because $e_\infty cz$ depends only on the archimedean
component of $z$ when $c$ has finite component $1$. If $h=rcb$ is in the strip,
then $c\,h^{-1}x_0u=b^{-1}r^{-1}x_0u$, since $c$ commutes with $b$; its
archimedean component is that of $r^{-1}x_0u=x_0vu$ with $v$ unitriangular, which
lies in $Q$. Conversely let $h^{-1}x_0u\in S$ with $c\in C$, and write $h=ab$ with
$a$ archimedean and $b$ finite. Then the archimedean component of
$(ca^{-1})x_0u$ lies in $Q$, so $(ca^{-1})x_0=x_0v$ for a unitriangular $v$, that
is $ca^{-1}\in R(\mathbf A)$, and $h=(ca^{-1})^{-1}\,c\,b$ is in the strip.
Taking $u=1$, the strip is the preimage of $S$ under the continuous map
$h\mapsto h^{-1}x_0$, hence measurable.
\end{proof}

\begin{proposition}\label{th-vanishing-descends}
Let $D$ be a definite-side datum and $\psi$ a standard character.
\begin{enumerate}
\item Let a step from Bessel data $V=V_0\oplus H^\rho$, $V_0=A\oplus L$, to
  Fourier-Jacobi data $W=W_0\oplus H^\rho$ with its flat map at $(x_0,y_0)$, and
  local data $(\mathcal S,\Gamma_0,K)$ for it, be given
  (Definitions~\ref{th-step} and~\ref{th-step-local}). Let $\varphi$ on
  $U(A)(\mathbf A)$ and $G$ on $U(W)(\mathbf A)$ be rapidly decreasing, with
  $G(gk)=G(g)$ for $k\in K$. If $\mathcal B_V(\varphi,\theta_\Phi(G))=0$ for
  every $\Phi\in\mathcal S$, then $\ell_{\varphi,G}(\Phi^\flat)=0$ for every
  $\Phi\in\mathcal S$.
\item Let a step from Fourier-Jacobi data $W=W_0\oplus H^{\sigma+1}$ to Bessel
  data $V=V_0\oplus H^\sigma$, $V_0=A\oplus L$ with vector $\zeta$, with its flat
  map at $(x_1,0)$, and local data $(\mathcal S,\Gamma_0,K)$ for it, be given.
  Fix a Haar measure on $\mathbb W(W_0,L)$ with a fundamental domain of volume
  $1$, a Weil datum $\omega_z$ for $(W_0,L)$ and a theta function $\Phi_0$ for
  $(W_0,L)$. Let $\varphi$ on $U(W_0)(\mathbf A)$ and $G$ on $U(V)(\mathbf A)$
  be rapidly decreasing, with $G(gk)=G(g)$ for $k\in K$. If
  $\mathcal J_W(\varphi,\theta_\Phi(G);\Phi_0)=0$ for every $\Phi\in\mathcal S$,
  the period taken with $\zeta$ and $\omega_z$, then
  $\mathcal B_V(\varphi_{\Phi^\flat},G)=0$ for every $\Phi\in\mathcal S$
  (Definition~\ref{th-step-periods}(2)).
\end{enumerate}
\end{proposition}
\leanref{Development/AbelianVarietyTheory/Proofs/TwistedOperators/VanishingDescendsBesselToFourierJacobiStatement.lean}{GlobalTheory.T.vanishing_descends_bessel_to_fourierJacobi}
\leanref{Development/AbelianVarietyTheory/Proofs/TwistedOperators/VanishingDescendsFourierJacobiToBesselStatement.lean}{GlobalTheory.T.vanishing_descends_fourierJacobi_to_bessel}
\leanref{Development/AbelianVarietyTheory/Definitions/TwistedInnerProducts/TowerEx.lean}{GlobalTheory.Asc.VanBF}
\leanref{Development/AbelianVarietyTheory/Definitions/TwistedInnerProducts/TowerEx.lean}{GlobalTheory.Asc.VanFB}
\leanref{Development/AbelianVarietyTheory/Proofs/TwistedOperators/VanishingDescendsBesselToFourierJacobi.lean}{GlobalTheory.Asc.vanBF}
\leanref{Development/AbelianVarietyTheory/Proofs/TwistedOperators/VanishingDescendsFourierJacobiToBessel.lean}{GlobalTheory.Asc.VanJac.vanFB}
\begin{proof}
We prove (1). The proof of (2) has the same steps: Lemma~\ref{th-unfolding}
applies to the character of the Fourier-Jacobi period and the point $x_1$ of
Definition~\ref{th-step}, with the Bessel group $N_V\rtimes U(A)$ of $V$ in
place of the Jacobi group of $W$, the function
$h\mapsto G(h)\,\varphi_{(\omega(1,h)\Phi)^\flat}$ in place of $f$ below, and the
two periods exchanged.

Fix $\Phi\in\mathcal S$ and put, for $h\in U(W)(\mathbf A)$,
\[
f(h)=G(h)\int_{\mathcal F_A}\varphi(h_0)\,\bigl(\omega(\iota(h_0),h)\Phi\bigr)^\flat(0)\,dh_0,
\]
where $\iota(h_0)$ is the image of $(h_0,1)$ in $U(V)$. Let $R$ be the Jacobi
group $N_W\rtimes U(W_0)$ of $W$, which is the group $R_{x_0}$ of
Definition~\ref{th-step-local}.

\emph{Unfolding.} Put $\Theta=\overline{\psi_\zeta}$ on $N_V(\mathbf A)$, and
let $x_0$ be the matrix of the vectors $e_1,\ldots,e_\rho$ of $W$. By
Lemma~\ref{th-bessel-unfolding}, its Gram matrix is $0$, $\Theta$ satisfies
conditions (1) to (3) of Lemma~\ref{th-unfolding} for $x_0$, for every theta
function $\Psi$, $\ell_\Theta(\Psi)=\Psi^\flat(0)$ for the flat map at
$(x_0,y_0)$, and $R_{x_0}=R$.
For fixed $h_0$, $\theta_\Phi(G)(n\iota(h_0))=\theta_{\omega(\iota(h_0),1)\Phi}(G)(n)$,
so Lemma~\ref{th-unfolding} for $\omega(\iota(h_0),1)\Phi$ gives
\[
\int\overline{\psi_\zeta(n)}\,\theta_\Phi(G)(n\iota(h_0))\,dn
=\int_{\mathcal B}G(h)\,\bigl(\omega(\iota(h_0),h)\Phi\bigr)^\flat(0)\,dh ,
\]
with $\mathcal B$ a fundamental domain for $R(F)=U(W)(F)\cap R$. The integral
over $\mathcal B$ of $|G(h)|$ times the integral of the absolute value of the
integrand of the flat map is at most $C\operatorname{ht}(\iota(h_0))^A$,
hence at most $C'\operatorname{ht}(h_0)^{A'}$ by
Lemma~\ref{th-weil-estimates}(5). Since $\varphi$ is rapidly decreasing, we may
integrate against $\varphi(h_0)$ over $\mathcal F_A$ and exchange the integrals:
$\mathcal B_V(\varphi,\theta_\Phi(G))=\int_{\mathcal B}f(h)\,dh$, absolutely
convergent. The same holds with $\mathcal B\cap S$ in place of $\mathcal B$ for
every measurable set $S$ of $U(W)(\mathbf A)$ invariant under left
multiplication by $R(\mathbf A)$, the unfolded integral then being taken over
the $(h,h_0,u,b)$ with $h\in S$.

\emph{Strips.} For $c\in M_{\operatorname{rk}W\times\rho}(\mathbf A_E)$ with
finite component $0$, the theta function $\rho(Y(c))\Phi$, translated in the
span of the $e'_a$ of $V$, lies in $\mathcal S$ by
Definition~\ref{th-step-local}(2), so its Bessel period vanishes. By
Lemma~\ref{th-weil-estimates}(1) and the definition of the first stage of the
flat map, replacing $\Phi$ by $\rho(Y(c))\Phi$ multiplies the unfolded
integrand at $(h,h_0,u,b)$ by
$\psi(\langle X(h^{-1}x_0u),Y(c)\rangle)$. So the integrals of the unfolded
integrand against all the characters $\psi(\langle X(\cdot),Y(c)\rangle)$ of
the archimedean component of $h^{-1}x_0u$ vanish. By
Lemma~\ref{th-fourier-strips}(1), applied to the unfolded integrand on the
space of the $(h,h_0,u,b)$ with $x(h,h_0,u,b)=h^{-1}x_0u$, its integral over the
points with $h^{-1}x_0u$ in a measurable set read on the archimedean components
vanishes. For a compact set $C$ of elements of $U(W)(\mathbf A)$ with finite
component $1$, Lemma~\ref{th-fourier-strips}(2) shows that the strip
$S=R(\mathbf A)\,C\,U(W)(\mathbf A)_f$ is measurable and is the set of $h$ with
$h^{-1}x_0u$ in such a set. Hence $\int_{\mathcal B\cap S}f=0$ for every strip
$S$.

\emph{Invariance.} $f$ is left invariant under $R(F)$: $G$ is invariant under
the rational points, the flat map of $\omega(1,nh')\Phi$ is the translate of the
flat map of $\Phi$ under the action of the Jacobi group on theta functions for
$(V_0,W_0)$, with the character $\psi_N$ on $N_W$, and at rational points this
action fixes the value at $0$ by Definition~\ref{th-weil-datum}(3) for
$\omega_0$.

\emph{The mean.} For $h\in U(W)(\mathbf A)$, write the integral of
$r\mapsto f(rh)$ over $R(F)\backslash R(\mathbf A)$ with $r=nh'$, $n\in N_W$,
$h'\in U(W_0)$, by the Haar measure on the semidirect product. The same
equivariance of the flat map turns $f(nh'h)$ into the integrand of
$\ell_{\varphi,G(\cdot\,h)}$ evaluated at $(\omega(1,h)\Phi)^\flat$. So the mean
of $f$ at $h$ is $\ell_{\varphi,G(\cdot\,h)}\bigl((\omega(1,h)\Phi)^\flat\bigr)$;
it converges absolutely and is continuous in $h$.

\emph{Conclusion.} Write $T(h)$ for the mean of $f$ at $h$, which is
continuous. It is left invariant under $R(\mathbf A)$, as the integral over
$R(F)\backslash R(\mathbf A)$ of a function invariant under $R(F)$, and right
invariant under $K$ by Definition~\ref{th-step-local}(4) and the invariance of
$G$. By Definition~\ref{th-step-local}(5), $T(h)\neq0$ implies
$h\in R(\mathbf A)\,U(W)(\mathbf A)_\infty\,K$, where $U(W)(\mathbf A)_\infty$ is
the group of elements with finite component $1$. Let $C$ be a set of elements
with finite component $1$ and $x=rcb$ an element of its strip, with $r\in R$,
$c\in C$ and $b$ of archimedean component $1$. If $T(x)\neq0$, write $x=qak$
with $q\in R$, $a$ archimedean and $k\in K$. Then
$(ac^{-1})(kb^{-1})=q^{-1}r\in R$, the first factor archimedean and the second
finite, and an element of $R$ with this decomposition has both factors in $R$;
so $ac^{-1}\in R$ and $T(x)=T(a)=T(c)$.

Suppose that $T(1)\neq0$. The set $\mathcal U$ of $h$ with
$\operatorname{Re}(T(h)/T(1))>\tfrac12$ is an open neighbourhood of $1$. Choose
a compact neighbourhood $C$ of $1$ in $U(W)(\mathbf A)_\infty$ inside
$\mathcal U$, and
let $S$ be its strip. By the last paragraph, on $S$ the function
$\operatorname{Re}(T/T(1))$ is $\geq0$, and $>\tfrac12$ where $T\neq0$. By
Strips, $\int_{\mathcal B\cap S}f=0$, so the integral of $\mathbf 1_ST$ over
$R(\mathbf A)\backslash U(W)(\mathbf A)$ vanishes. Choose a measurable
$\beta\geq0$ on $U(W)(\mathbf A)$ with $\int_{R(\mathbf A)}\beta(rx)\,dr=1$
for every $x$. Then $\int\beta\,\mathbf 1_ST/T(1)=0$ over $U(W)(\mathbf A)$,
and the real part of the integrand is $\geq0$, so it vanishes almost
everywhere. Hence $\beta=0$ almost everywhere on $A=S\cap\{T\neq0\}$, which is
left invariant under $R(\mathbf A)$, and integrating over $R(\mathbf A)$ shows
that $A$ has measure $0$. But $A$ contains the open set of products $ck$ with
$c$ in the interior of $C$ and $k\in K$, since $T(ck)=T(c)$ and
$\operatorname{Re}(T(c)/T(1))>\tfrac12$. This open set contains $1$, so it has
positive measure, a contradiction. Hence
$\ell_{\varphi,G}(\Phi^\flat)=T(1)=0$.
\end{proof}

\begin{proposition}\label{th-step-nonvanishing}
Let $D$ be a definite-side datum with host element $\xi_0$, and $\psi$ a
standard character.
\begin{enumerate}
\item \emph{From Fourier-Jacobi to Bessel periods.} Let a step from Bessel data
  $V=V_0\oplus H^\rho$, $V_0=A\oplus L$ with vector $\zeta$, to Fourier-Jacobi
  data $W=W_0\oplus H^\rho$, with characters $(\chi_V,\chi_W)$, and local data
  $(\mathcal S,\Gamma_0,K)$ for it be given (Definitions~\ref{th-step}
  and~\ref{th-step-local}). Let $L'$ be a line with a vector $\zeta'$,
  $h(\zeta',\zeta')\neq0$, and $\tau\colon L'\to L$ a linear map with
  $\tau\zeta'=-\zeta$ and $h(\zeta,\zeta)=\xi_0\bar\xi_0\,h(\zeta',\zeta')$.
  Let $\chi_A,\chi_z,\chi'$ be characters with $\chi_V=\chi_A\chi_z$, and let
  $\omega_A$ and $\omega_z$ be Weil data for $(A,W_0)$ and $(W_0,L')$ with
  characters $(\chi_A,\chi_W)$ and $(\chi',\chi_z)$. Let $\varphi$ on
  $U(A)(\mathbf A)$ and $G$ on $U(W)(\mathbf A)$ be rapidly decreasing, with
  $G(gk)=G(g)$ for $k\in K$, and let $\Xi_1$ and $\Xi_2$ be theta functions for
  $(A,W_0)$ and $(L,W_0)$ such that $\Xi_1\boxtimes\Xi_2$ is a lattice vacuum
  for $\Gamma_0$. If
  $\mathcal J_W\bigl(\theta_{\Xi_1}(\varphi),G;\overline{\Xi_2\circ\kappa}\bigr)\neq0$,
  with the notation of Lemma~\ref{th-step-functionals}(3), then
  $\mathcal B_V(\varphi,\theta_\Phi(G))\neq0$ for some $\Phi\in\mathcal S$.
\item \emph{From Bessel to Fourier-Jacobi periods.} Let a step from
  Fourier-Jacobi data $W=W_0\oplus H^{\sigma+1}$ to Bessel data
  $V=V_0\oplus H^\sigma$, $V_0=A\oplus L$ with vector $\zeta$, with characters
  $(\chi_W,\chi_V)$, and local data $(\mathcal S,\Gamma_0,K)$ for it be given.
  Let $\chi_A,\chi_z$ be characters with $\chi_V=\chi_A\chi_z$, and let
  $\omega_A$ and $\omega_z$ be Weil data for $(W_0,A)$ and $(W_0,L)$ with
  characters $(\chi_W,\chi_A)$ and $(\chi_W,\chi_z)$. Let $\varphi$ on
  $U(W_0)(\mathbf A)$ and $G$ on $U(V)(\mathbf A)$ be rapidly decreasing, with
  $G(gk)=G(g)$ for $k\in K$, and let $\Xi_1$ be a theta function for $(W_0,A)$
  and $\Phi_0\neq0$ one for $(W_0,L)$ such that $\Xi_1\boxtimes\Phi_0$ is a
  lattice vacuum for $\Gamma_0$. If
  $\mathcal B_V(\theta_{\Xi_1}(\varphi),G)\neq0$, with $\theta_{\Xi_1}(\varphi)$
  the lift by $\omega_A$, then
  $\mathcal J_W(\varphi,\theta_\Phi(G);\Phi_0)\neq0$ for some $\Phi\in\mathcal S$,
  the period taken with $\zeta$ and $\omega_z$.
\end{enumerate}
\end{proposition}
\leanref{Development/AbelianVarietyTheory/Proofs/TwistedOperators/FourierJacobiPeriodGivesBesselPeriodOfLiftStatement.lean}{GlobalTheory.T.fourierJacobi_period_gives_bessel_period_of_lift}
\leanref{Development/AbelianVarietyTheory/Proofs/TwistedOperators/BesselPeriodGivesFourierJacobiPeriodOfLiftStatement.lean}{GlobalTheory.T.bessel_period_gives_fourierJacobi_period_of_lift}
\leanref{Development/AbelianVarietyTheory/Definitions/TwistedInnerProducts/TowerEx.lean}{GlobalTheory.Asc.UpBF}
\leanref{Development/AbelianVarietyTheory/Definitions/TwistedInnerProducts/TowerEx.lean}{GlobalTheory.Asc.UpFB}
\leanref{Development/AbelianVarietyTheory/Proofs/TwistedOperators/FourierJacobiPeriodGivesBesselPeriodOfLift.lean}{GlobalTheory.UpBessel.upBF}
\leanref{Development/AbelianVarietyTheory/Proofs/TwistedOperators/BesselPeriodGivesFourierJacobiPeriodOfLift.lean}{GlobalTheory.UpFromBessel.upFB}
\begin{proof}
For (1), suppose that $\mathcal B_V(\varphi,\theta_\Phi(G))=0$ for every
$\Phi\in\mathcal S$. Put $\ell=\ell_{\varphi,G}$, a tame map on theta functions
for $(V_0,W_0)$ by Lemma~\ref{th-step-functionals}(3). By
Proposition~\ref{th-vanishing-descends}(1), $\ell(\Phi^\flat)=0$ for every
$\Phi\in\mathcal S$. Choose $\Phi_1\in\mathcal S$ with $\Phi_1^\flat\neq0$, by
Definition~\ref{th-step-local}(6); it is a lattice vacuum for $\Gamma_0$ by
Definition~\ref{th-step-local}(3). For $w\in\Gamma_0^\perp$, the function
$\rho(\iota w)\Phi_1$ lies in $\mathcal S$ by Definition~\ref{th-step-local}(2),
and $(\rho(\iota w)\Phi_1)^\flat=\rho(w)\Phi_1^\flat$ by
Lemma~\ref{th-step-functionals}(5). Hence $\ell(\rho(w)\Phi_1^\flat)=0$ for all
$w\in\Gamma_0^\perp$. By Lemma~\ref{th-vacuum-irreducibility}, $\ell$ vanishes
on every lattice vacuum for $\Gamma_0$, in particular on $\Xi_1\boxtimes\Xi_2$.
By Lemma~\ref{th-step-functionals}(3),
$\ell(\Xi_1\boxtimes\Xi_2)=\mathcal J_W(\theta_{\Xi_1}(\varphi),G;\overline{\Xi_2\circ\kappa})$,
which is nonzero by hypothesis. This is a contradiction.

For (2), suppose that $\mathcal J_W(\varphi,\theta_\Phi(G);\Phi_0)=0$ for every
$\Phi\in\mathcal S$. Fix a Haar measure on $\mathbb W(W_0,L)$ with a fundamental
domain of volume $1$, and put $\ell(\Xi)=\mathcal B_V(\varphi_\Xi,G)$, with the
handed-down function for $\Phi_0$. It is tame by
Lemma~\ref{th-step-functionals}(4), and $\ell(\Phi^\flat)=0$ for every
$\Phi\in\mathcal S$ by Proposition~\ref{th-vanishing-descends}(2). As in (1),
Lemma~\ref{th-vacuum-irreducibility} shows that $\ell$ vanishes on every lattice
vacuum for $\Gamma_0$, in particular on $\Xi_1\boxtimes\Phi_0$. By
Lemma~\ref{th-step-functionals}(4),
$\varphi_{\Xi_1\boxtimes\Phi_0}=\langle\Phi_0,\Phi_0\rangle\,\theta_{\Xi_1}(\varphi)$,
and the Bessel period is linear in its first argument, so
$\langle\Phi_0,\Phi_0\rangle\,\mathcal B_V(\theta_{\Xi_1}(\varphi),G)=0$. Since
$\Phi_0\neq0$, $\langle\Phi_0,\Phi_0\rangle\neq0$ by
Lemma~\ref{th-lattice-vacua}(2), and $\mathcal B_V(\theta_{\Xi_1}(\varphi),G)=0$,
a contradiction.
\end{proof}

\section{The ascent along the tower}\label{section-th-ascent}

\begin{lemma}\label{th-grid-conjugation}
Let $D$ be a definite-side datum in which $p_{de}$ and $p'_{de}$ are prime to
$Q_Sq$, with a level system with base tuple, a theta grid, a grid Weil datum
$(\omega_p)_p$ and measures as in Definition~\ref{th-core-chain}, and let
$0\leq i<m$. Let $\omega^{\mathrm{row}}_i$ and $\omega^{\mathrm{core}}_i$ be the
families of operators on functions on $\mathbb W(V_{i+1,0},V_{i,0})$ and
$\mathbb W(P_{i+1},P_i)$ given by
$\omega^{\mathrm{row}}_i(g,h)f=\overline{\omega_p(g,h)\bar f}$ for the grid pair
$p=(V_{i+1,0},V_{i,0})$, and in the same way for $(P_{i+1},P_i)$; here the
first space of the pair carries the negated form, which has the same unitary
group. Let $(\chi_{S,i},\chi_{T,i})$ be the characters of the grid pair
$(V_{i+1,0},V_{i,0})$. Then:
\begin{enumerate}
\item $\omega^{\mathrm{row}}_i$ is a Weil datum for $(V_{i+1,0},V_{i,0})$ with
  characters $(\chi_{S,i}^{-1},\chi_{T,i}^{-1})$, and $\omega^{\mathrm{core}}_i$
  is one for $(P_{i+1},P_i)$ with characters
  $(\chi_{m-i}^{-1},\chi_{m-i}^{-1})$;
\item if $i$ is odd, then $\chi_{S,i}^{-1}=\chi_{m-i}^{-1}\eta_L$ and
  $\chi_{T,i}=\chi_{m-i}$; if $i$ is even, then
  $\chi_{T,i}^{-1}=\chi_{m-i}^{-1}\eta_L$ and $\chi_{S,i}=\chi_{m-i}$;
\item a function $f$ is a theta function for $(-P_{i+1},P_i)$ and a lattice
  vacuum for the depth lattice of the core pair if and only if $\bar f$ is a
  theta function for $(P_{i+1},P_i)$ and a lattice vacuum for the same lattice;
\item if $\Psi_i$ is the lift of $\Psi_{i+1}$ by $f'_i$ as in
  Definition~\ref{th-core-chain}(2), then $\bar\Psi_i$ is the theta lift
  $\theta_{\bar f'_i}(\bar\Psi_{i+1})$ from $U(P_{i+1})$ to $U(P_i)$ by
  $\omega^{\mathrm{core}}_i$;
\item for row data $(\varrho,\Phi)$ and a function $u$ on $U(Y_0)(\mathbf A)$,
  $\overline{\theta^{(i+1)}(u)}$ is the theta lift
  $\theta_{\bar f_i}\bigl(\overline{\theta^{(i)}(u)}\bigr)$ from $U(V_{i,0})$ to
  $U(V_{i+1,0})$ by $\omega^{\mathrm{row}}_i$, where $f_i$ is the theta function
  of the row data at $i$;
\item $\theta^{(i)}(u)$ depends only on the row data $(\varrho_j,\Phi_j)$ with
  $j<i$.
\end{enumerate}
\end{lemma}
\leanref{Development/AbelianVarietyTheory/Proofs/TwistedOperators/GridConjugationWeilDataInverseCharacterStatement.lean}{GlobalTheory.T.gridConjugation_weilData_inverseCharacter}
\leanref{Development/AbelianVarietyTheory/Definitions/TwistedInnerProducts/PartFixed.lean}{GlobalTheory.Asc.ConjRead}
\leanref{Development/AbelianVarietyTheory/Definitions/TwistedInnerProducts/TowerEx.lean}{GlobalTheory.Asc.cjOm}
\leanref{Development/AbelianVarietyTheory/Definitions/TwistedInnerProducts/GenInhabited.lean}{GlobalTheory.Asc.FamilyModel.conjRead}
\begin{proof}
For (1): complex conjugation preserves theta functions up to the exchange of
$\psi$ with $\psi^{-1}$, which the negation of the form of the first space
undoes, and it inverts unitary characters; so the conjugate of a Weil datum
for $(-S,T)$ with characters $(\chi_S,\chi_T)$ is a Weil datum for $(S,T)$
with characters $(\chi_S^{-1},\chi_T^{-1})$. The characters of the core pair are
$(\chi_{m-i},\chi_{m-i})$, because the ranks of the core spaces are even
(Definition~\ref{th-grid}). For (2): by Definition~\ref{th-grid} the characters
of the row pair are $\chi_{m-i}$ times $\eta_L^{-1}$ for each space of odd rank;
since $r$ is odd, $\operatorname{rk}V_{i,0}=r+i(r+6)$ is odd if and only if $i$
is even. Part (3) holds because complex conjugation preserves the invariance
conditions of a lattice vacuum. For (4) and (5), conjugate the integral
defining the lift and use
$\overline{\omega_p(g,h)f(0)}=(\omega^{\mathrm{row}}_i(g,h)\bar f)(0)$, and the
same for the core pairs. Part (6) holds because $\theta^{(i)}(u)$ is built from
$u$ by the lifts at $0,\ldots,i-1$.
\end{proof}

\begin{lemma}\label{th-lattice-signs}
Let $D$ be a definite-side datum, $\psi$ a standard character and $S,T$
Hermitian spaces with self-dual lattices $\Lambda_S$ and $\Lambda_T$, and let
$\Lambda\subset\mathbb W(S,T)$ be the closure of the lattice of the pair,
consisting of elements with archimedean component $0$
(Definition~\ref{th-weil-datum}).
\begin{enumerate}
\item An element $u$ with archimedean component $0$ lies in $\Lambda$ if and only
  if $\psi(\langle u,u'\rangle)=1$ for every $u'\in\Lambda$.
\item For every $u\in\Lambda$ there is $\varepsilon=\pm1$ with
  $\rho(u)f=\varepsilon f$ for every lattice vacuum $f$ for $\Lambda$.
\end{enumerate}
\end{lemma}
\leanref{Development/AbelianVarietyTheory/Proofs/AdelicLatticeProofs/SelfDualLatticePairingDualStatement.lean}{GlobalTheory.T.self_dual_lattice_pairing_dual}
\leanref{Development/AbelianVarietyTheory/Proofs/AdelicLatticeProofs/SelfDualLatticePairingDual.lean}{GlobalTheory.LatVac.signGroup}
\begin{proof}
Part (1) is the self-duality of $\Lambda_S$ and $\Lambda_T$, place by place, for
the pairing $\operatorname{Tr}(\xi_0\operatorname{tr}(\cdot))$ and the standard
character $\psi$.

For (2), let $G$ be the set of $u\in\Lambda$ for which some $\varepsilon$ with
$\varepsilon^2=1$ satisfies $\rho(u)f=\varepsilon f$ for every lattice vacuum
$f$ for $\Lambda$. It is a subgroup: for $u,u'\in\Lambda$,
$\rho(u)\rho(u')=\psi(\frac12\langle u,u'\rangle)\rho(u+u')$ up to the order of
the arguments, and $\psi(\frac12\langle u,u'\rangle)^2=\psi(\langle u,u'\rangle)=1$
by (1). It contains $2u$ for every $u\in\Lambda$, and every $u\in\Lambda$ lying
in the half of $\mathbb W$ attached to a split place, with $\varepsilon=1$, by
the definition of a lattice vacuum. The places of $F$ over $2$ split in $E$, a
property of the host field, so $1=\sum_{v\mid2}(e_v+\bar e_v)+2\beta$ in the
finite adeles, with $e_v$ the idempotent of a place of $E$ over $v$, $\bar e_v$
its conjugate, and $\beta$ integral. Hence every $u\in\Lambda$ is
$\sum_v(e_vu+\bar e_vu)+2\beta u$, a sum of elements of $\Lambda$ in split
halves and of an element of $2\Lambda$, so $u\in G$.
\end{proof}

\begin{lemma}\label{th-coset-vectors}
In the situation of Lemma~\ref{th-grid-conjugation}, let $0\leq i<m$, and write
$S=V_{i+1,0}$ and $T=V_{i,0}$, with self-dual lattices $\Lambda_S$ and
$\Lambda_T$ of the grid. Let $\Lambda\subset\mathbb W(S,T)$ be the lattice of
the row pair, $\Gamma\subset\Lambda$ its depth lattice, $\mathfrak m$ the
modulus of the grid, $M=\operatorname{Hom}(\Lambda_S,\Lambda_T)$, and $e$ the
sum of the idempotents of $\mathbf A_E$ of the chosen places of $E$ over the
distinguished places, with conjugate $\bar e$ (Definition~\ref{th-grid}(W6)
and (W7)). For a class $\xi$ in $M/\mathfrak mM$ let $\Phi_\xi$ be the coset
vector of $\xi$ in the grid Weil datum (Definition~\ref{th-grid}(W7)): the
image, under every place model at $e$ of the vacuum vector $f_0$ of the pair
(Definition~\ref{th-grid-fock}(7)), of the indicator function of the set of
$x$ in the $e$-half with $x-ex_0$ in the closure of $\mathfrak mM$ for some
$x_0\in M$ representing $\xi$.
\begin{enumerate}
\item $\Phi_\xi\neq0$; it is a theta function for $(-S,T)$, a Fock vacuum vector
  at every archimedean place and a lattice vacuum for $\Gamma$; and
  $\bar\Phi_\xi$ is a theta function for $(S,T)$ and a lattice vacuum for
  $\Gamma$.
\item For $w\in\Gamma^\perp$ with archimedean component $0$ there are a class
  $\xi'$ and $c\in\mathbf C$ with $\rho(w)\bar\Phi_\xi=c\,\bar\Phi_{\xi'}$.
\end{enumerate}
\end{lemma}
\leanref{Development/AbelianVarietyTheory/Proofs/TwistedOperators/CosetVectorsNonzeroDepthFockVacuaStatement.lean}{GlobalTheory.T.coset_vectors_nonzero_depth_fock_vacua}
\leanref{Development/AbelianVarietyTheory/Proofs/TwistedOperators/HeisenbergCommutantPermutesFirstRowCosetVectorsStatement.lean}{GlobalTheory.T.heisenbergCommutant_permutes_firstRowCosetVectors}
\begin{proof}
For (1): $f_0$ has norm $1$, and the indicator function takes the value $1$ at
$ex_0$, so replacing a component of $f_0$ by it gives a nonzero function. The
replacement preserves quasi-periodicity, the invariance under an open subgroup
and smoothness, so $\Phi_\xi$ is a theta function. It changes only a finite
component, and the annihilation operators act at the archimedean places, so
$\Phi_\xi$ is a Fock vacuum vector.

For the lattice vacuum, write $u=eu+\bar eu+r$ with $r=u-eu-\bar eu$, so that
$er=\bar er=0$. Three kinds of translations fix $\Phi_\xi$.
\begin{enumerate}
\item[(A)] $\rho(v)$ for $v$ in the $e$-half and in the closure of
  $\mathfrak mM$: it acts on the $e$-component by translation, which preserves
  the indicator function of the coset.
\item[(B)] $\rho(v)$ for $v\in\Lambda$ in the $\bar e$-half: it acts on the
  $e$-component by multiplication by $t\mapsto\psi(\langle t,v\rangle)$, which
  is $1$ on the coset because the coset lies in $\Lambda$ and $\Lambda$ is
  self-dual.
\item[(C)] $\rho(r)$ for $r$ with $er=\bar er=0$ and $\rho(r)f_0=f_0$: such a
  translation does not touch the replaced component.
\end{enumerate}
Let $u\in\Gamma$. Then $eu$ lies in the closure of $\mathfrak mM$, by the
definition of the depth lattice, and $eu,\bar eu,r\in\Lambda$. The vectors
$2\bar eu$ and $2r$ are orthogonal, and $\psi(\tfrac12\langle2eu,2\bar eu+2r\rangle)=1$
because $\Lambda$ is self-dual, so
$\rho(2u)=\rho(2eu)\rho(2\bar eu)\rho(2r)$. Since $f_0$ is a lattice vacuum for
$\Lambda$, $\rho(2r)f_0=f_0$, and (A), (B), (C) give $\rho(2u)\Phi_\xi=\Phi_\xi$.
If $u$ lies in the half attached to a split place, so do $eu$, $\bar eu$ and
$r$, these are pairwise orthogonal, $\rho(r)f_0=f_0$, and
$\rho(u)\Phi_\xi=\Phi_\xi$ in the same way. Complex conjugation takes theta
functions for $(-S,T)$ to theta functions for $(S,T)$, with
$\rho(u)\bar f=\overline{\rho_{-S}(u)f}$ for the translations $\rho_{-S}$ of
the pair $(-S,T)$; so $\bar\Phi_\xi$ is a lattice vacuum for $\Gamma$.

For (2), decompose $w=ew+\bar ew+w'$ with $w'=w-ew-\bar ew$; all three have
archimedean component $0$.
\begin{enumerate}
\item[(i)] For $u\in\Lambda$ put $u_0=u-eu-\bar eu$. Since $1-e-\bar e$ is
  fixed by complex conjugation, $\langle w',u\rangle=\langle w,u_0\rangle$. Write
  $u_0=\sum_{v\mid2}(e_vu_0+\bar e_vu_0)+2\beta u_0$ as in the proof of
  Lemma~\ref{th-lattice-signs}. Each $e_vu_0$, $\bar e_vu_0$ and $\beta u_0$ lies
  in $\Lambda$ with $e$-component $0$, hence in $\Gamma$; the first two lie in
  split halves, so $w\in\Gamma^\perp$ gives $\psi(\langle w,e_vu_0\rangle)=
  \psi(\langle w,\bar e_vu_0\rangle)=\psi(\langle w,2\beta u_0\rangle)=1$. Hence
  $\psi(\langle w',u\rangle)=1$ for all $u\in\Lambda$, and $w'\in\Lambda$ by
  Lemma~\ref{th-lattice-signs}(1). By Lemma~\ref{th-lattice-signs}(2),
  $\rho(w')f_0=c_0f_0$ with $c_0=\pm1$, $f_0$ being a lattice vacuum for
  $\Lambda$. Since $w'$ has $e$- and $\bar e$-components $0$, it does not touch
  the replaced component, and $\rho(w')\Phi_\xi=c_0\Phi_\xi$.
\item[(ii)] $\rho(\bar ew)$ multiplies the $e$-component by
  $t\mapsto\psi(\langle t,\bar ew\rangle)$. For $t,t'$ in the coset, $t-t'$ lies
  in the $e$-half and in the closure of $\mathfrak mM$, hence in $\Gamma$, so
  $\psi(\langle t-t',w\rangle)=1$ and the factor is constant on the coset:
  $\rho(\bar ew)\Phi_\xi=c_2\Phi_\xi$.
\item[(iii)] For $u\in\Lambda$, $\langle ew,u\rangle=\langle w,\bar eu\rangle$,
  and $\bar eu$ lies in $\Gamma$ and in the $\bar e$-half, which is attached to
  split places; so $\psi(\langle ew,u\rangle)=1$, and $ew\in\Lambda$ by
  Lemma~\ref{th-lattice-signs}(1): $ew$ lies in the closure of $M$ at the
  chosen places and in the $e$-half. Choose bases of $\Lambda_S$ and $\Lambda_T$
  and write $ew$ in coordinates, with entries in the closure of $\mathcal O_E$ at
  the chosen places. By the Chinese remainder theorem there is $x_1\in M$ whose
  coordinates are congruent to them modulo $\mathfrak m$ at the chosen places, so
  that $ew-ex_1$ lies in the closure of $\mathfrak mM$. Translation by $ew$ maps
  the indicator function of the coset of $x_0$ to that of the coset of
  $x_0-x_1$, so $\rho(ew)\Phi_\xi=\Phi_{\xi'}$.
\end{enumerate}
By the commutation relation of $\rho$, $\rho(w)\Phi_\xi$ is a multiple of
$\Phi_{\xi'}$, and conjugating gives (2).
\end{proof}

\begin{lemma}\label{th-invariant-form}
Let $D$ be a definite-side datum, $\psi$ a standard character and $S,T$
Hermitian spaces (Definition~\ref{th-weil-datum}). Let $e\in\mathbf A_E$ be the
idempotent of a place of $E$ over a place of $F$ split in $E$, with conjugate
$\bar e$, and call the elements $x$ of $\mathbb W(S,T)$ with $x=ex$ the
\emph{$e$-half}. Let $\Lambda\subset\mathbb W(S,T)$ be a subgroup such that for
every $y\notin\Lambda$ in the $e$-half there is $l\in\Lambda$ in the $\bar e$-half
with $\psi(\langle l,y\rangle)\neq1$. Let $f\neq0$ be a function on
$\mathbb W(S,T)$ with $\rho(u)f=f$ for every $u\in\Lambda$ in the $e$-half or in
the $\bar e$-half. Then:
\begin{enumerate}
\item there is a linear form $\kappa$ on the functions on $\mathbb W(S,T)$ with
  $\kappa(f)=1$ such that $\kappa(\rho(y)f)$ is real and $\geq0$ for every $y$
  in the $e$-half;
\item there is a linear form $\ell$ on the functions on $\mathbb W(S,T)$ with
  $\ell(f)\neq0$ and $\ell(\rho(y)\varphi)=\ell(\varphi)$ for every $y$ in the
  $e$-half and every theta function $\varphi$.
\end{enumerate}
\end{lemma}
\leanref{Development/AbelianVarietyTheory/Proofs/SatakeParameters/ExistsInvariantForm.lean}{GlobalTheory.SplitHecke.exists_count_form'}
\leanref{Development/AbelianVarietyTheory/Proofs/SatakeParameters/ExistsInvariantForm.lean}{GlobalTheory.SplitHecke.exists_invariant_form'}
\leanref{Development/AbelianVarietyTheory/Proofs/AdelicLatticeProofs/NotMemSpanOfCommute.lean}{GlobalTheory.SplitScalar.notMem_span_of_commute}
\begin{proof}
Both halves are isotropic for $\langle\ ,\ \rangle$, since $e\bar e=0$. Hence
$\rho(u)\rho(u')=\rho(u+u')$ for $u,u'$ in the same half, and the
translations by one half commute with each other. For all $w,w'$,
$\rho(w)\rho(w')=\psi(\langle w,w'\rangle)\rho(w')\rho(w)$. We use the
following fact from linear algebra.
\begin{enumerate}
\item[$(*)$] Let $\mathcal G$ be a set of pairwise commuting linear operators on
  a complex vector space, $X$ a set of vectors each of which is killed by some
  element of $\mathcal G$, and $f$ a vector with $Bf\neq0$ for every finite
  product $B$ of elements of $\mathcal G$. Then $f$ is not in the span of $X$.
\end{enumerate}
Indeed, every vector of the span is killed by a finite product of elements of
$\mathcal G$: if $Bx=0$ and $B'x'=0$, then $BB'$ kills $x+x'$, since
$BB'=B'B$.

For (1), let $X_1$ be the set of the $\rho(y)f$ with $y\notin\Lambda$ in the
$e$-half, and $\mathcal G_1$ the set of the operators $\rho(l)-z$ with
$l\in\Lambda$ in the $\bar e$-half and $z\in\mathbf C$, $z\neq1$. These operators
commute. Given $\rho(y)f\in X_1$, take $l$ as in the hypothesis and
$z=\psi(\langle l,y\rangle)$. Then $\rho(l)\rho(y)f=z\,\rho(y)\rho(l)f=z\,\rho(y)f$,
so $\rho(l)-z$ kills $\rho(y)f$. Each $\rho(l)-z$ in $\mathcal G_1$ multiplies
$f$ by $1-z\neq0$, so every finite product multiplies $f$ by a nonzero constant.
By $(*)$, $f$ is not in the span of $X_1$, and there is a linear form $\kappa$
that vanishes on $X_1$ with $\kappa(f)=1$. For $y$ in the $e$-half,
$\kappa(\rho(y)f)$ is $\kappa(f)=1$ if $y\in\Lambda$, and $0$ if $y\notin\Lambda$.

For (2), let $X_2$ be the set of the differences $\rho(y)\varphi-\varphi$ with $y$
in the $e$-half and $\varphi$ a theta function, and $\mathcal G_2$ the set of the
averaging operators $A_{y,N}=\sum_{j=0}^N\rho(jy)$ with $y$ in the $e$-half and
$N\geq0$. These operators commute. Given $\rho(y)\varphi-\varphi\in X_2$, there is
an open subgroup $U\subset\mathbb W(S,T)$ such that $\rho(u)\varphi=\varphi$ for
$u\in U$ with archimedean component $0$
(Definition~\ref{th-weil-datum}(2)). The element $y$ has archimedean component
$0$, so $(N+1)y\in U$ for some $N\geq0$, and
\[
A_{y,N}(\rho(y)\varphi-\varphi)=\sum_{j=0}^N\bigl(\rho((j+1)y)\varphi-\rho(jy)\varphi\bigr)
=\rho((N+1)y)\varphi-\varphi=0 .
\]
Let $\kappa$ be as in (1). For every finite product $B$ of elements of
$\mathcal G_2$, $\kappa(B\rho(y)f)\geq0$ for all $y$ in the $e$-half and
$\kappa(Bf)>0$. This holds for $B=1$ by (1), and passes from $B$ to $BA_{y',N}$
because
\[
\kappa(BA_{y',N}\rho(y)f)=\sum_{j=0}^N\kappa\bigl(B\rho(jy'+y)f\bigr)\geq0,\qquad
\kappa(BA_{y',N}f)=\sum_{j=0}^N\kappa\bigl(B\rho(jy')f\bigr)\geq\kappa(Bf)>0,
\]
the term $j=0$ of the second sum being $\kappa(Bf)$. So $Bf\neq0$ for every such
$B$. By $(*)$, $f$ is not in the span of $X_2$, and there is a linear form $\ell$
that vanishes on $X_2$ with $\ell(f)\neq0$; it has the invariance of (2).
\end{proof}

\begin{lemma}\label{th-split-separation}
Let $D$ be a definite-side datum, $\psi$ a standard character and $S,T$
Hermitian spaces. Let $e\in\mathbf A_E$ be the idempotent of a place of $E$ over
a place of $F$ split in $E$, and call the elements $x$ of $\mathbb W(S,T)$ with
$x=ex$ the \emph{$e$-half}. Let $f$ be a theta function such that
$\ell(\rho(x)f)=0$ for every $x$ in the $e$-half and every linear map $\ell$ from
theta functions to $\mathbf C$ with $\ell(\rho(u)\varphi)=\ell(\varphi)$ for all
theta functions $\varphi$ and all $u$ in the $\bar e$-half. Then $f=0$.
\end{lemma}
\leanref{Development/AbelianVarietyTheory/Proofs/InertiaEigenvalues/SplitPlaceInvariantFormsSeparateThetaFunctions.lean}{GlobalTheory.DownLift.SplitSep.invForms_separate}
\begin{proof}
Let $\Lambda$ be the closure of the lattice of a pair of self-dual lattices of
$S$ and $T$, as in Lemma~\ref{th-lattice-signs}. The two halves are isotropic
and $e\Lambda,\bar e\Lambda\subset\Lambda$. Since $f$ is a theta function, there
is $k\geq1$ with $\rho(ku)f=f$ for every $u\in\Lambda$ in either half. Let $Y$ be
a set of representatives of the $y$ in the $\bar e$-half with $ky\in\Lambda$
modulo $k(\Lambda\cap\bar e\text{-half})$, and $X$ one of the $x$ in the
$e$-half with $kx\in\Lambda$ modulo $k(\Lambda\cap e\text{-half})$; both are
finite.

\emph{The averages vanish.} Let $x\in X$ and
$F_x=\sum_{y\in Y}\rho(y)\rho(x)f$. Let $\Lambda'$ be the group of the $y$ in the
$\bar e$-half with $ky\in\Lambda$ together with the $ku$, $u\in\Lambda$ in the
$e$-half. Then $F_x$ is fixed by $\rho(u)$ for $u\in\Lambda'$: in the
$\bar e$-half, $\rho(u)$ permutes the representatives up to translations by
$k(\Lambda\cap\bar e\text{-half})$, which commute with $\rho(x)$ because
$\psi(\langle ku,x\rangle)=1$ for $kx\in\Lambda$, and fix $f$; in the $e$-half,
$\rho(ku)$ commutes with $\rho(y)$ and $\rho(x)$ for the same reason and fixes
$f$. Suppose $F_x\neq0$. By the self-duality of $\Lambda$, every $y$ in the
$\bar e$-half outside $\Lambda'$ pairs nontrivially with some element of
$\Lambda'$ in the $e$-half. Lemma~\ref{th-invariant-form}(2), for $\Lambda'$, the
idempotent $\bar e$ and $F_x$, gives a linear map $\ell$ with $\ell(F_x)\neq0$
and $\ell(\rho(y)\varphi)=\ell(\varphi)$ for every $y$ in the $\bar e$-half and
every theta function $\varphi$. The function $\rho(x)f$ is a theta function, so
$\ell(F_x)=\sum_{y\in Y}\ell(\rho(y)\rho(x)f)=|Y|\,\ell(\rho(x)f)=0$ by
hypothesis, a contradiction. So $F_x=0$.

\emph{The average over $X$.} For $y$ in the $\bar e$-half and $x$ in the
$e$-half, $\rho(-x)\rho(y)\rho(x)=\psi(\langle y,x\rangle)\rho(y)$, the halves
being isotropic. For each $y\in Y$, $\sum_{x\in X}\psi(\langle y,x\rangle)$ is
$|X|$ if $y\in k\Lambda$ and $0$ otherwise, by the orthogonality of characters
of the finite group of the $x$ modulo $k\Lambda$. Exactly one $y\in Y$ lies in
$k\Lambda$, and for it $\rho(y)f=f$. Hence
$0=\sum_{x\in X}\rho(-x)F_x=|X|\,f$, and $f=0$.
\end{proof}

\begin{lemma}\label{th-nonsplit-vacua}
Let $D$ be a definite-side datum, $\psi$ a standard character and $S,T$
Hermitian spaces with self-dual lattices, with closure $\Lambda$ of the lattice
of the pair, and let $\omega$ be a Weil datum for $(S,T)$ with characters
$(\chi_S,\chi_T)$ (Definition~\ref{th-weil-datum}). Let $v$ be a finite place of
$F$ not split in $E$ at which $\chi_S$ and $\chi_T$ are unramified, and let
$g\in U(S)(\mathbf A)$ and $h\in U(T)(\mathbf A)$ have components $1$ away from
$v$ and stabilize the lattices at $v$. If $f$ is a theta function with
$\rho(u)f=f$ for every $u\in\Lambda$ supported at $v$, then $\omega(g,h)f=f$.
\end{lemma}
\leanref{Development/AbelianVarietyTheory/Proofs/DivisorProofs/WeilDatumFixesVacuumNonSplitPlaceUnramifiedStatement.lean}{GlobalTheory.T.weilDatum_fixes_vacuum_nonSplitPlace_unramified}
\leanref{Development/AbelianVarietyTheory/Proofs/DivisorProofs/ExistsWeilDatumPowersUnramifiedEtaTildeStatement.lean}{GlobalTheory.T.exists_weilDatum_powers_unramifiedEtaTilde}
\leanref{Development/AbelianVarietyTheory/Proofs/DivisorProofs/ExistsWeilDatumTwistHeckeCharactersStatement.lean}{GlobalTheory.T.exists_weilDatum_twist_heckeCharacters}
\begin{proof}
Let $\tilde\eta$ be a Hecke character of $E$, unramified at every finite place,
whose restriction to $\mathbf A_F^\times$ is $\eta$, with archimedean components
of exponent $\frac12$; Alpöge's datum provides one, and its existence forces
$E/F$ to be unramified at the finite places, so that $v$ is inert. There is a
Weil datum
$\tilde\omega$ for $(S,T)$ with characters $(\tilde\eta^{\operatorname{rk}S},\tilde\eta^{\operatorname{rk}T})$
that fixes $f$ under $(g,h)$ as in the statement at every nonsplit place: it is
the restricted tensor product of the local Weil representations given by
Kudla's splittings for $\tilde\eta$ \cite{Kud94}, and at an inert place where
$\tilde\eta$ is unramified and the lattices are self-dual for the pairing the
local representation fixes the characteristic function of the lattice under its
stabilizer \cite{Kud94}.

The characters $\xi_S=\chi_S\tilde\eta^{-\operatorname{rk}S}$ and
$\xi_T=\chi_T\tilde\eta^{-\operatorname{rk}T}$ are Hecke characters trivial on
$\mathbf A_F^\times$. Twisting $\tilde\omega$ by them gives a Weil datum
$\omega_1$ with characters $(\chi_S,\chi_T)$ and
$\omega_1(g,h)=\xi_T(x)^{-1}\xi_S(y)\,\tilde\omega(g,h)$ whenever $x/\bar x=\det g$
and $y/\bar y=\det h$. Two Weil data with the same characters agree on theta
functions, so $\omega(g,h)f=\omega_1(g,h)f$. At the nonsplit place $v$, $\det g$
and $\det h$ are units of norm $1$, and by Hilbert's Theorem~90 for the
unramified extension $E_v/F_v$ there are units $x$ and $y$ of $E_v$, extended
by $1$ elsewhere, with $x/\bar x=\det g$ and
$y/\bar y=\det h$. The characters $\chi_S$, $\chi_T$ and $\tilde\eta$ are
unramified at $v$, so $\xi_T(x)=\xi_S(y)=1$, and
$\omega(g,h)f=\tilde\omega(g,h)f=f$.
\end{proof}

\begin{lemma}\label{th-depth-fixed}
In the situation of Lemma~\ref{th-grid-conjugation}, for a space $X$ of the grid
let $K_X$ be the subgroup of $U(X)(\mathbf A)$ of elements with archimedean
component $1$ whose finite component lies in the congruence subgroup of the
lattice of $X$ in the grid for the modulus of the grid. Let $0\leq i<m$.
\begin{enumerate}
\item If $f$ is a theta function for $(-V_{i+1,0},V_{i,0})$ and a lattice vacuum
  for the depth lattice of the row pair, then $\omega_p(k,k')f=f$ for
  $k\in K_{V_{i+1,0}}$ and $k'\in K_{V_{i,0}}$. The same holds for the core pairs
  $(-P_{i+1},P_i)$, with $K_{P_{i+1}}$ and $K_{P_i}$.
\item For all row data $(\varrho,\Phi)$, the theta function $f_i$ of the row data
  at $i$ is a lattice vacuum for the depth lattice of the row pair.
\item For all row data, every function $u$ on $U(Y_0)(\mathbf A)$ and every
  $k\in K_{V_{i+1,0}}$, $\theta^{(i+1)}(u)(gk)=\theta^{(i+1)}(u)(g)$.
\end{enumerate}
\end{lemma}
\leanref{Development/AbelianVarietyTheory/Proofs/TwistedOperators/DependenceFixedStatement.lean}{GlobalTheory.T.dependence_fixed}
\leanref{Development/AbelianVarietyTheory/Definitions/TwistedInnerProducts/OffQOfSizeRows.lean}{GlobalTheory.SepRows.DepFixed}
\leanref{Development/AbelianVarietyTheory/Definitions/TwistedInnerProducts/OffQOfSizeRows.lean}{GlobalTheory.SepRows.DatumDep}
\leanref{Development/AbelianVarietyTheory/Definitions/TwistedInnerProducts/OffQOfSizeRows.lean}{GlobalTheory.SepRows.lev}
\leanref{Development/AbelianVarietyTheory/Definitions/TwistedInnerProducts/PartFixed.lean}{GlobalTheory.Asc.PartFixed}
\leanref{Development/AbelianVarietyTheory/Definitions/TwistedInnerProducts/GenInhabited.lean}{GlobalTheory.Asc.FamilyModel.partFixed}
\begin{proof}
For (1), we treat the row pair; the core pairs are the same. Write $S$ and $T$
for the two spaces of the pair, the first with the negated form, $\Lambda$ for
the lattice of the pair, $\Gamma$ for its depth lattice, $\chi_S$ and $\chi_T$
for the characters of the pair, and $S_0$ for the set of distinguished places,
with the chosen place $w_v$ of $E$ over each $v\in S_0$ and the exponent
$d(v)$ of $w_v$ in the modulus. Let $k\in K_S$ and $k'\in K_T$.

\emph{The characters.} $\chi_S$ and $\chi_T$ are unramified at the places of
$F$ outside $S_0$, and at each $v\in S_0$ they are trivial on the units of
$E_{w_v}$ congruent to $1$ modulo $w_v^{d(v)}$. This holds for $\chi_{m-i}$ by
conditions (c), (d) and (e) of Definition~\ref{th-grid}, since $d(v_0)\geq1$
and $d(v)=2$ for the other $v\in S_0$; it holds for $\eta_L$, which is
unramified at every finite place by (L3) of Definition~\ref{th-grid}; and hence
it holds for $\chi_{m-i}\eta_L^{-1}$.

\emph{Localization.} The function $(g,h)\mapsto\omega(g,h)f$ is continuous, and
$k$ and $k'$ have archimedean component $1$. So it suffices to show
$\omega(k_v,k'_v)f=f$ for every finite place $v$ of $F$, where $k_v$ and $k'_v$
agree with $k$ and $k'$ at $v$ and are $1$ elsewhere: then
$\omega(k,k')f=f$, as the limit of the finite products of these components.

\emph{A place $v\notin S_0$.} The characters are unramified at $v$, $k_v$ and
$k'_v$ stabilize the lattices at $v$, and every element of $\Lambda$ supported
at $v$ lies in $\Gamma$, since the depth lattice differs from $\Lambda$ only at
the chosen places. So $f$ is a lattice vacuum for the lattice of the pair at
$v$. If $v$ is not split in $E$, then $v$ is not over $2$, since the places over
$2$ split; so every $u\in\Lambda$ supported at $v$ is $2u'$ with $u'\in\Gamma$,
$\rho(u)f=f$, and Lemma~\ref{th-nonsplit-vacua} gives $\omega(k_v,k'_v)f=f$. If
$v$ splits, the computation of the next paragraph applies with the characters
unramified at $v$, the lattice $\Lambda$ at $v$ in place of $\Gamma$, and
$x$, $y$ units, so that $\chi_T(x)=\chi_S(y)=1$.

\emph{A place $v\in S_0$.} The place $v$ splits in $E$. Let $e$ be the
idempotent of $\mathbf A_E$ of $w_v$; the idempotent of $v$ is $e+\bar e$, and
the $e$-half and the $\bar e$-half of $\mathbb W(S,T)$ at $v$ are isotropic.
Let $x$ and $y$ be the determinants of $k_v$ and $k'_v$ at $w_v$. We use the
following consequence of the condition at split places in
Definition~\ref{th-weil-datum}(3). \emph{Suppose that (i) $u\in\Lambda$ if and
only if $k'_vuk_v^{-1}\in\Lambda$; (ii) $\rho(k'_vuk_v^{-1}-u)f=f$ for
$u\in\Lambda$ in the $e$-half; and (iii) $\rho(u)f=f$ for $u\in\Lambda$ in the
$\bar e$-half. Then $\omega(k_v,k'_v)f=\chi_T(x)^{-1}\chi_S(y)f$.} Indeed, by
Definition~\ref{th-weil-datum}(3) there is $r>0$ such that $\omega(k_v,k'_v)$
multiplies every linear functional $\ell$ invariant under the translations of
the corresponding half by $r\chi_T(x)^{-1}\chi_S(y)$. Let
$F=\omega(k_v,k'_v)f-r\chi_T(x)^{-1}\chi_S(y)f$, a theta function. For such
$\ell$ and $z=k'_vwk_v^{-1}$ in the $e$-half,
$\ell(\rho(z)F)=r\chi_T(x)^{-1}\chi_S(y)\,(\ell(\rho(w)f)-\ell(\rho(z)f))$ by the
equivariance of $\omega$. If $w\in\Lambda$, then
$\rho(z)f=\rho(w)\rho(z-w)f=\rho(w)f$ by (ii), the half being isotropic. If
$w\notin\Lambda$, then $z\notin\Lambda$ by (i), and both terms vanish: the
self-duality of $\Lambda$ gives $l\in\Lambda$ in the other half with
$\psi(\langle l,w\rangle)\neq1$, and comparing
$\ell(\rho(l)\rho(w)f)$ with $\ell(\rho(w)\rho(l)f)$, using (iii), gives
$\ell(\rho(w)f)=0$. So $\ell(\rho(z)F)=0$ for all such $\ell$ and $z$, and
$F=0$ by Lemma~\ref{th-split-separation}. Hence
$\omega(k_v,k'_v)f=r\chi_T(x)^{-1}\chi_S(y)f$; since $\omega(k_v,k'_v)$
preserves the $L^2$-norm on a fundamental domain and $f\neq0$, $r=1$.

The three conditions hold. Condition (i) holds because $k_v$ and $k'_v$
stabilize $\Lambda_S$ and $\Lambda_T$. For (ii), $k_v$ and $k'_v$ are congruent
to $1$ modulo $w_v^{d(v)}$, so $k'_vuk_v^{-1}-u$ lies in the closure of
$\mathfrak m\operatorname{Hom}(\Lambda_S,\Lambda_T)$ and in the $e$-half,
hence in $\Gamma$; the hypothesis that $p_{de}$ and $p'_{de}$ are prime to
$Q_Sq$ gives the integrality at $w_v$ of the elements that compare $\Lambda$
with $\operatorname{Hom}(\Lambda_S,\Lambda_T)$. Then $f$, a lattice vacuum for
$\Gamma$, is fixed by $\rho(k'_vuk_v^{-1}-u)$, which lies in the half attached
to a split place. For (iii), an element of $\Lambda$ in the $\bar e$-half lies
in $\Gamma$, because the depth lattice constrains only the components at the
chosen places. Finally $x$ and $y$ are congruent to $1$ modulo $w_v^{d(v)}$, so
$\chi_T(x)=\chi_S(y)=1$ by the condition on the characters, and
$\omega(k_v,k'_v)f=f$.

For (2), $f_i$ is $\varrho_i$ applied to the coset vector $\Phi_i$, which is a
lattice vacuum for the depth lattice by Lemma~\ref{th-coset-vectors}(1). The
creation operators act at the archimedean places, and the elements of the
depth lattice have archimedean component $0$, so the translations by these
elements commute with $\varrho_i$. Hence $f_i$ is again a lattice vacuum for
the depth lattice.

For (3), by definition
$\theta^{(i+1)}(u)(g)=\int\theta_{\omega_p}(g,h;f_i)\,\theta^{(i)}(u)(h)\,dh$.
For $k\in K_{V_{i+1,0}}$,
$\omega_p(gk,h)f_i=\omega_p(g,h)\,\omega_p(k,1)f_i=\omega_p(g,h)f_i$ by (1)
and (2). So the integrands at $gk$ and at $g$ agree.
\end{proof}

\begin{lemma}\label{th-fock-translates}
Let $D$ be a definite-side datum, $\psi$ a standard character and $S,T$
Hermitian spaces with Cartan involutions and orthonormal bases at the
archimedean places (Definition~\ref{th-weil-datum}). Let $\ell$ be a tame map on
theta functions for $(S,T)$ (Lemma~\ref{th-vacuum-irreducibility}), and let $f$
be a Fock vacuum vector. If $\ell(f\otimes P)=0$ for every polynomial $P$ in the
creation operators, then $\ell(\rho(w)f)=0$ for every $w\in\mathbb W(S,T)$ with
finite component $0$.
\end{lemma}
\leanref{Development/AbelianVarietyTheory/Proofs/TwistedOperators/TameFunctionalKillsFockVacuumRealTranslatesStatement.lean}{GlobalTheory.T.tameFunctional_kills_fockVacuum_realTranslates}
\leanref{Development/AbelianVarietyTheory/Proofs/TwistedOperators/FockGaussianReadingDifferentiableStatement.lean}{GlobalTheory.T.fockGaussianReading_differentiable}
\leanref{Development/AbelianVarietyTheory/Proofs/TwistedOperators/FockGaussianJetsAtOriginEqFockVectorsStatement.lean}{GlobalTheory.T.fock_gaussian_jets_at_origin_eq_fock_vectors}
\leanref{Development/AbelianVarietyTheory/Proofs/TwistedOperators/FockRealPointTransportContinuousIntegralStatement.lean}{GlobalTheory.T.fockRealPoint_transport_continuous_integral}
\begin{proof}
The Cartan involutions and orthonormal bases give complex coordinates on the
elements with finite component $0$: a continuous real-linear map
$z\mapsto w(z)$ from
$V_{\mathbf C}=\bigoplus_{\tau\in\mathcal T}M_{\operatorname{rk}T\times\operatorname{rk}S}(\mathbf C)$
onto the elements of $\mathbb W(S,T)$ with finite component $0$, built from the
archimedean embeddings and the bases. For an integrable $g$ on $V_{\mathbf C}$,
the image of $g(z)\,dz$ under $z\mapsto w(z)$ is an $s$-finite measure
$\alpha\,dm$ on $\mathbb W(S,T)$ with $\alpha$ integrable, and
$\int g(z)H(w(z))\,dz=\int\alpha H\,dm$ for every $H$; so an integral
$\int g(z)\,\rho(w(z))f\,dz$ is an integral of the translates $\rho(w)f$ against
an integrable function, as in the definition of a tame map. Let
$q(z)=\sum_\tau c_\tau|z_\tau|^2$, with
$c_\tau=2\pi|\tau(\xi_0)|>0$, be the quadratic form of the Fock model and
$b(z_0,z)$ its polarization, linear in $\bar z$. Put $F(z)=\ell(\rho(w(z))f)$,
which is continuous and bounded since $\ell$ is tame. For
$x\in\mathbb W(S,T)$ put $R_x(z)=e^{q(z)/2}(\rho(w(z))f)(x)$.
\begin{enumerate}
\item[(a)] $R_x$ is an entire function of $z$, with
  $|R_x(z)|\leq\sup|f|\,e^{q(z)/2}$. The bound holds because
  $|(\rho(w)f)(x)|=|f(x+w)|$. For holomorphy, let $J$ be the complex structure
  of the coordinates. Since $f$ is annihilated by the annihilation operators,
  the function $G(w)=(\rho(w)f)(x)$ satisfies
  $dG_w(Jy)=i\,dG_w(y)+c\,\operatorname{Im}h(w,y)G(w)$ for the symplectic form
  $\operatorname{Im}h$ of the coordinates and a constant $c$, where the second
  term comes from the commutation relation of the translations. The Gaussian
  $E(z)=e^{q(z)/2}$ satisfies the same equation with the opposite inhomogeneous
  term, and the symplectic form in the coordinates is
  $-\sum_\tau c_\tau\operatorname{Im}(z_\tau\bar y_\tau)$. So the product
  $R_x=(G\circ w)\,E$ satisfies the Cauchy-Riemann equations.
\item[(b)] For every monomial $z^\mu$ there is a constant $C_\mu$, not depending
  on $x$ or on $f$, with $\partial^\mu R_x(0)=C_\mu(f\otimes z^\mu)(x)$, where
  $z^\mu$ is read as a monomial in the creation operators. Since $R_x$ is
  holomorphic, its derivative in a coordinate $z_k$ is half the derivative along
  the real direction minus $i$ times the derivative along the $J$-rotated
  direction; on $\rho(w(z))f$ this combination is the creation operator, up to
  a constant, the Gaussian factor cancelling the central term of the commutation
  relation. Induction on $|\mu|$ gives (b).
\item[(c)] For an entire function $R$ with $|R|\leq Ce^{q/2}$,
  $\int\overline{z^\mu}e^{-q(z)}R(z)\,dz=K_\mu\,\partial^\mu R(0)$ for constants
  $K_\mu$, and $R(z_0)=K'\int e^{b(z_0,z)}e^{-q(z)}R(z)\,dz$ \cite{Bar61}.
\end{enumerate}
By (a), (b) and (c),
$\int\overline{z^\mu}e^{-q(z)/2}(\rho(w(z))f)(x)\,dz=K_\mu C_\mu(f\otimes z^\mu)(x)$
for every $x$. The left side is an integral of translates of $f$ against an
integrable function, so tameness gives
$\int\overline{z^\mu}e^{-q(z)/2}F(z)\,dz=K_\mu C_\mu\,\ell(f\otimes z^\mu)=0$.
Expanding $e^{b(z_0,z)}$ in its power series in $\bar z$, which converges
absolutely against $e^{-q/2}|F|$, gives
$\int e^{b(z_0,z)}e^{-q(z)/2}F(z)\,dz=0$ for every $z_0$. By the second formula
of (c) applied to $R_x$,
$(\rho(w(z_0))f)(x)=K'e^{-q(z_0)/2}\int e^{b(z_0,z)}e^{-q(z)/2}(\rho(w(z))f)(x)\,dz$,
again an integral of translates of $f$. Tameness gives
$\ell(\rho(w(z_0))f)=K'e^{-q(z_0)/2}\int e^{b(z_0,z)}e^{-q(z)/2}F(z)\,dz=0$.
\end{proof}

\begin{lemma}\label{th-fock-span}
In the situation of Lemma~\ref{th-grid-conjugation}, let $0\leq i<m$, and let
$\ell$ be a tame map on functions on $\mathbb W(V_{i+1,0},V_{i,0})$
(Lemma~\ref{th-vacuum-irreducibility}). For a coset vector $\kappa$ of the row
pair and a polynomial $P$ in the creation operators, let $f_{\kappa,P}$ be $P$
applied to $\kappa$. If $\ell(\bar f_{\kappa,P})=0$ for all $\kappa$ and $P$,
then $\ell(f)=0$ for every lattice vacuum $f$ for the depth lattice of the row
pair.
\end{lemma}
\leanref{Development/AbelianVarietyTheory/Proofs/TwistedOperators/FockConjugatesSpanDepthDataStatement.lean}{GlobalTheory.T.fockConjugates_span_depth_data}
\leanref{Development/AbelianVarietyTheory/Definitions/TwistedInnerProducts/TowerEx.lean}{GlobalTheory.Asc.FockTot}
\leanref{Development/AbelianVarietyTheory/Proofs/TwistedOperators/FockConjugatesSpanDepthData.lean}{GlobalTheory.ConjDataTotal.fockTot}
\leanref{Development/AbelianVarietyTheory/Proofs/TwistedOperators/HeisenbergCommutantPermutesFirstRowCosetVectorsStatement.lean}{GlobalTheory.T.heisenbergCommutant_permutes_firstRowCosetVectors}
\leanref{Development/AbelianVarietyTheory/Proofs/TwistedOperators/TameFunctionalKillsFockVacuumRealTranslatesStatement.lean}{GlobalTheory.T.tameFunctional_kills_fockVacuum_realTranslates}
\begin{proof}
Let $\Gamma$ be the depth lattice of the row pair. By
Lemma~\ref{th-coset-vectors}, the coset vectors are nonzero Fock vacuum
vectors at every archimedean place for the Cartan involutions of the grid Weil
datum, their conjugates $\bar\kappa$ are lattice vacua for $\Gamma$, and for a
fixed coset vector $\kappa_0$ and $w\in\Gamma^\perp$ with archimedean component
$0$, $\rho(w)\bar\kappa_0$ is a multiple of some $\bar\kappa$.

First, $\ell(\rho(w)\bar\kappa)=0$ for every coset vector $\kappa$ and every $w$
with finite component $0$. Indeed, $F\mapsto\overline{\ell(\bar F)}$ is tame on
theta functions for $(-V_{i+1,0},V_{i,0})$ and vanishes on $P$ applied to the
Fock vacuum vector $\kappa$ for every $P$. By Lemma~\ref{th-fock-translates}
for the pair $(-V_{i+1,0},V_{i,0})$ it vanishes on the translates of $\kappa$
by the elements with finite component $0$. Complex conjugation takes these
translates to the translates $\rho(w)\bar\kappa$ for the pair
$(V_{i+1,0},V_{i,0})$, so $\ell(\rho(w)\bar\kappa)=0$.

Next let $w\in\Gamma^\perp$, and write $w=w_\infty+w_f$ with archimedean and
finite components. Since $\Gamma$ has archimedean component $0$, $w_f$ lies in
$\Gamma^\perp$, and $\rho(w)=\rho(w_\infty)\rho(w_f)$ because the two components
pair to $0$. So $\rho(w)\bar\kappa_0$ is $\rho(w_\infty)$ applied to a
multiple of some $\bar\kappa$, and $\ell(\rho(w)\bar\kappa_0)=0$ by the first
step and linearity. Since $\bar\kappa_0\neq0$ is a lattice vacuum for $\Gamma$,
Lemma~\ref{th-vacuum-irreducibility} gives $\ell(f)=0$ for every lattice vacuum
$f$ for $\Gamma$.
\end{proof}

\begin{definition}\label{th-tower-decomposition}
Let $D$ be a definite-side datum in which $p_{de}$ and $p'_{de}$ are prime to
$Q_Sq$, with host element $\xi_0$, a level system with base tuple, a theta grid,
a grid Weil datum and measures as in Definition~\ref{th-core-chain}. A
\emph{tower decomposition} of the grid consists of:
\begin{enumerate}
\item a Hermitian line $L_z$ and an orthogonal decomposition $P_0=Y_0\oplus L_z$;
\item vectors $\zeta_\rho\in-L_z$ with $h(\zeta_\rho,\zeta_\rho)\neq0$ for
  $\rho\geq0$, and linear maps $\tau_\rho$ of $-L_z$ with
  $\tau_\rho\zeta_\rho=-\zeta_{\rho+1}$ and
  $h(\zeta_{\rho+1},\zeta_{\rho+1})=\xi_0\bar\xi_0\,h(\zeta_\rho,\zeta_\rho)$;
\item for each $\rho$ with $2\rho+2\leq m$, Bessel data $B_\rho$ of depth $\rho$,
  $V_{2\rho+2,0}=V'_\rho\oplus H^\rho$ with $V'_\rho=P_{2\rho+2}\oplus(-L_z)$ and
  vector $\zeta_{\rho+1}$, and Fourier-Jacobi data $J_\rho$ of depth $\rho$,
  $V_{2\rho+1,0}=P_{2\rho+1}\oplus H^\rho$ (Definition~\ref{th-step}), with
  the measures on $U(P_{2\rho+2})$ and $U(P_{2\rho+1})$ of
  Definition~\ref{th-core-chain}.
\end{enumerate}
It \emph{fits the ends} if the embedding $U(Y_0)\to U(P_0)$ given by (1) is the
embedding $\iota_0$ of the grid, and if, for $2\rho+2=m$ and rapidly decreasing
$\varphi$ on $U(P_m)(\mathbf A)$ and $f$ on $U(V_{m,0})(\mathbf A)$,
$\mathcal B_{B_\rho}(\bar\varphi,\bar f)\neq0$ implies
$\mathcal B^+(\varphi,f)\neq0$ (Definition~\ref{th-periods}). With the notation
of Lemmas~\ref{th-grid-conjugation} and~\ref{th-depth-fixed}, it is \emph{good}
for subgroups $\Gamma_L(\rho)\subset\mathbb W(P_{2\rho+1},-L_z)$ if:
\begin{enumerate}\setcounter{enumi}{3}
\item for $2\rho+2\leq m$ and every step from $B_\rho$ to $J_\rho$ with Weil datum
  $\omega^{\mathrm{row}}_{2\rho+1}$ on $(V_{2\rho+2,0},V_{2\rho+1,0})$, there are
  local data $(\mathcal S,\Gamma_0,K_{V_{2\rho+1,0}})$ for it
  (Definition~\ref{th-step-local}) such that every $\Phi\in\mathcal S$ is a
  lattice vacuum for the depth lattice of the row pair, and
  $\Xi_1\boxtimes\Xi_2$ is a lattice vacuum for $\Gamma_0$ whenever $\Xi_1$ is a
  lattice vacuum for the depth lattice of the core pair
  $(P_{2\rho+2},P_{2\rho+1})$ and $\Xi_2$ is a theta function for
  $(-L_z,P_{2\rho+1})$ with $\overline{\Xi_2\circ\kappa_\rho}$ a lattice vacuum
  for $\Gamma_L(\rho)$, where $\kappa_\rho(w)=(\tau_\rho w)^*$;
\item for $2\sigma+4\leq m$ and every step from $J_{\sigma+1}$ to $B_\sigma$ with
  Weil datum $\omega^{\mathrm{row}}_{2\sigma+2}$ on
  $(V_{2\sigma+3,0},V_{2\sigma+2,0})$, there are local data
  $(\mathcal S,\Gamma_0,K_{V_{2\sigma+2,0}})$ for it such that every
  $\Phi\in\mathcal S$ is a lattice vacuum for the depth lattice of the row pair,
  and $\Xi_1\boxtimes\Phi_0$ is a lattice vacuum for $\Gamma_0$ whenever $\Xi_1$
  is a lattice vacuum for the depth lattice of the core pair
  $(P_{2\sigma+3},P_{2\sigma+2})$ and $\Phi_0$ is a lattice vacuum for
  $\Gamma_L(\sigma+1)$;
\item for $2\rho+2\leq m$, $\Gamma_L(\rho)$ has a nonzero lattice vacuum;
\item the subgroups $\Gamma$, $\Gamma_L(0)$ and $\Gamma_X$, the depth lattice of
  the row pair $(V_{1,0},V_{0,0})$, $\Gamma_L(0)$, and the depth lattice of the
  core pair $(P_1,P_0)$, satisfy the hypotheses on the subgroups in
  Theorem~\ref{th-seesaw-period-nonvanishing} for $W=V_{1,0}$, $W_0=P_1$,
  $T=Y_0$ and $X=P_0=Y_0\oplus L_z$.
\end{enumerate}
\end{definition}
\leanref{Development/AbelianVarietyTheory/Definitions/TwistedInnerProducts/TowerEx.lean}{GlobalTheory.Asc.Frames}
\leanref{Development/AbelianVarietyTheory/Definitions/TwistedInnerProducts/TowerEx.lean}{GlobalTheory.Asc.Frames.Ends}
\leanref{Development/AbelianVarietyTheory/Definitions/TwistedInnerProducts/TowerEx.lean}{GlobalTheory.Asc.Frames.Good}
\leanref{Development/AbelianVarietyTheory/Definitions/TwistedInnerProducts/TowerEx.lean}{GlobalTheory.Asc.LastLat}

\begin{definition}\label{th-standard-pieces}
Let $D$ be a definite-side datum with host element $\xi_0$, $\psi$ a standard
character, and a step from $S=S_0\oplus H^\rho$ to $T=T_0\oplus H^\sigma$ with
Weil datum $\omega$ (Definition~\ref{th-step}). A matrix over $\mathbf A_E$ is
\emph{finite} if its archimedean component is $0$, and \emph{supported at} a set
$P$ of finite places of $F$ if it is finite and its components at the places
outside $P$ are $0$. For $a\in M_{\operatorname{rk}T\times\rho}(\mathbf A_E)$
write $X(a)$ and $Y(a)$ for the elements of $\mathbb W(S,T)$ supported on the
$e_i$ and on the $e'_i$ of $S$ with coefficient matrix $a$; for
$\xi\in M_{\sigma\times\operatorname{rk}S_0}(\mathbf A_E)$ write $X_T(\xi)$ and
$Y_T(\xi)$ for the elements of $\mathbb W(S_0,T)$ supported on the $e_j$ and on
the $e'_j$ of $T$ with coefficient matrix $\xi$; write
$\iota_S\colon\mathbb W(S_0,T)\to\mathbb W(S,T)$ and
$\iota_T\colon\mathbb W(S_0,T_0)\to\mathbb W(S_0,T)$ for the inclusions, so that
$\iota=\iota_S\iota_T$ is the inclusion of Definition~\ref{th-step-local}; and for
$y\in M_{\operatorname{rk}S_0\times\sigma}(\mathbf A_E)$ put
$y^*=Y_T(-\xi_0^{-1}\bar y^{\mathsf T}G_{S_0})$. Let
$x\in M_{\operatorname{rk}T\times\rho}(\mathbf A_E)$ and
$y\in M_{\operatorname{rk}S_0\times\sigma}(\mathbf A_E)$.
\begin{enumerate}
\item \emph{Pieces} for the step are subgroups $A_1,A_2$ of
  $M_{\operatorname{rk}T\times\rho}(\mathbf A_E)$, $B_1,B_2$ of
  $M_{\sigma\times\operatorname{rk}S_0}(\mathbf A_E)$ and $\Gamma_0$ of
  $\mathbb W(S_0,T_0)$. They are \emph{admissible} at $x$ if:
  \begin{enumerate}
  \item[(a)] $A_1$, $A_2$, $B_1$ and $B_2$ consist of finite matrices and are
    compact and open in the group of finite matrices;
  \item[(b)] a finite $a$ lies in $A_1$ if and only if
    $\psi(\langle X(a),Y(a')\rangle)=1$ for every $a'\in A_2$, a finite $a'$ lies
    in $A_2$ if and only if this holds for every $a\in A_1$, and the same holds
    for $B_1$, $B_2$ and $\psi(\langle X_T(\xi),Y_T(\eta)\rangle)$;
  \item[(c)] $\psi_U(u)=1$ for every $u\in U_\rho(\mathbf A_E)$ with $u-1$
    finite and $xu-x\in A_1$;
  \item[(d)] $\Gamma_0$ consists of finite elements and has a nonzero lattice
    vacuum.
  \end{enumerate}
\item The pieces \emph{split} a subgroup $\Gamma\subset\mathbb W(S,T)$ at $(x,y)$
  if every $u\in\Gamma$ is
  $X(a)+Y(a')+\iota_S\bigl(X_T(\xi)+Y_T(\eta)\bigr)+\iota w$ for some $a\in A_1$,
  $a'\in A_2$, $\xi\in B_1$, $\eta\in B_2$ and $w\in\Gamma_0$ with
  $\psi(\langle X(x),Y(a')\rangle)=1$ and $\psi(\langle y^*,X_T(\xi)\rangle)=1$.
\item The \emph{class} of the pieces at $(x,y)$ is the set of theta functions
  $\Phi$ for $(S,T)$ with
  \[
  \rho(X(a))\Phi=\Phi,\quad
  \rho(Y(a'))\Phi=\psi(\langle X(x),Y(a')\rangle)\Phi,\quad
  \rho(\iota_SX_T(\xi))\Phi=\psi(\langle y^*,X_T(\xi)\rangle)\Phi,\quad
  \rho(\iota_SY_T(\eta))\Phi=\Phi
  \]
  for $a\in A_1$, $a'\in A_2$, $\xi\in B_1$ and $\eta\in B_2$, and with
  $\rho(2\iota w)\Phi=\Phi$ for every $w\in\Gamma_0$ and $\rho(\iota w)\Phi=\Phi$
  for every $w\in\Gamma_0$ such that $\iota w$ lies in the half of
  $\mathbb W(S,T)$ attached to a split place.
\item An element $h\in U(T)(\mathbf A)$ satisfies the \emph{unfolding conditions}
  at a set $P$ of finite places if:
  \begin{enumerate}
  \item[(a)] there are $u\in U_\rho(\mathbf A_E)$ and $a\in A_1$ such that
    $h^{-1}xu-x-a$ has component $0$ at every place of $P$;
  \item[(b)] $\psi_U(1+N)=1$ for every strictly upper triangular
    $N\in M_\rho(\mathbf A_E)$ supported at $P$ with $h^{-1}xN\in A_1$;
  \item[(c)] $\psi(\langle y^*,X_T(\xi')\rangle)=\psi(\langle y^*,X_T(\xi)\rangle)$
    whenever $\xi\in B_1$ is supported at $P$ and $hX_T(\xi)=X_T(\xi')$.
  \end{enumerate}
\end{enumerate}
\end{definition}
\leanref{Development/AbelianVarietyTheory/Definitions/TwistedInnerProducts/IsAuxiliaryVectorWittDec.lean}{GlobalTheory.AscFin.Pieces}
\leanref{Development/AbelianVarietyTheory/Definitions/TwistedInnerProducts/IsAuxiliaryVectorWittDec.lean}{GlobalTheory.AscFin.Pieces.Is}
\leanref{Development/AbelianVarietyTheory/Definitions/TwistedInnerProducts/IsAuxiliaryVectorWittDec.lean}{GlobalTheory.AscFin.Splits}
\leanref{Development/AbelianVarietyTheory/Definitions/TwistedInnerProducts/IsAuxiliaryVectorWittDec.lean}{GlobalTheory.AscFin.cls}
\leanref{Development/AbelianVarietyTheory/Definitions/TwistedInnerProducts/IsAuxiliaryVectorWittDec.lean}{GlobalTheory.AscFin.Cond}

\begin{lemma}\label{th-standard-class}
Let $D$ be a definite-side datum, $\psi$ a standard character, a step from
$S=S_0\oplus H^\rho$ to $T=T_0\oplus H^\sigma$ with Weil datum $\omega$, and
$x\in M_{\operatorname{rk}T\times\rho}(\mathbf A_E)$ and
$y\in M_{\operatorname{rk}S_0\times\sigma}(\mathbf A_E)$. Let pieces
$(A_1,A_2,B_1,B_2,\Gamma_0)$ for the step be admissible at $x$
(Definition~\ref{th-standard-pieces}).
\begin{enumerate}
\item If the pieces split a subgroup $\Gamma\subset\mathbb W(S,T)$ at $(x,y)$,
  every member of their class at $(x,y)$ is a lattice vacuum for $\Gamma$.
\end{enumerate}
Let moreover $x$ have entries in $E$ and linearly independent columns, and let
$\Phi\mapsto\Phi^\flat$ be the flat map of the step at $(x,y)$
(Definition~\ref{th-step}(4)).
\begin{enumerate}\setcounter{enumi}{1}
\item If $\Phi$ is in the class, then so are $\rho(w)\Phi$ for every
  $w\in\mathbb W(S,T)$ with finite component $0$ and $\rho(\iota w)\Phi$ for
  every $w\in\Gamma_0^\perp$ (Lemma~\ref{th-vacuum-irreducibility}), and
  $\Phi^\flat$ is a lattice vacuum for $\Gamma_0$.
\item Some member $\Phi$ of the class has $\Phi^\flat\neq0$.
\item If $\Phi$ is in the class, $h\in U(T)(\mathbf A)$ and
  $(\omega(1,h)\Phi)^\flat\neq0$, then $h$ satisfies the unfolding conditions at
  the set of all finite places.
\end{enumerate}
\end{lemma}
\leanref{Development/AbelianVarietyTheory/Proofs/TwistedOperators/DoubleMixedModelSupportedClassSubsetSplitLatticeStatement.lean}{GlobalTheory.T.doubleMixedModelSupportedClass_subset_splitLatticeVacua}
\leanref{Development/AbelianVarietyTheory/Proofs/TwistedOperators/DoubleMixedModelClassStableFlatVacuaStatement.lean}{GlobalTheory.T.doubleMixedModelClass_stable_flat_vacua}
\leanref{Development/AbelianVarietyTheory/Proofs/TwistedOperators/ExistsUnfoldingFlatMapNeZeroFullRankStatement.lean}{GlobalTheory.T.exists_unfolding_flatMap_ne_zero_fullRank}
\leanref{Development/AbelianVarietyTheory/Proofs/TwistedOperators/UnfoldingConditionsOfFlatMapWeilActionNeZeroStatement.lean}{GlobalTheory.T.unfoldingConditions_of_flatMap_weilAction_ne_zero}
\leanref{Development/AbelianVarietyTheory/Proofs/DecoratedLatticeProofs/FlatThetaMapsConvergeThetaFunctions.lean}{GlobalTheory.Unf.Flat.flatLaw}
\leanref{Development/AbelianVarietyTheory/Proofs/TwistedOperators/FlatMapHeisenbergAndHomTransformationStatement.lean}{GlobalTheory.T.flatMap_heisenberg_and_hom_transformation}
\leanref{Development/AbelianVarietyTheory/Proofs/TwistedOperators/ExistsAdelicBumpPeriodicUnderAdjointLatticeStatement.lean}{GlobalTheory.T.exists_adelicBump_periodic_under_adjointLattice}
\leanref{Development/AbelianVarietyTheory/Proofs/TwistedOperators/ContinuousReprSymm.lean}{GlobalTheory.AscFin.exists_bump_le}
\leanref{Development/AbelianVarietyTheory/Proofs/TwistedOperators/MixedModelTargetBumpMulCoreThetaExistsStatement.lean}{GlobalTheory.T.mixedModelTarget_bump_mul_coreTheta_exists}
\leanref{Development/AbelianVarietyTheory/Proofs/TwistedOperators/ExistsThetaFunctionMixedCoordinateBumpMulCoreStatement.lean}{GlobalTheory.T.exists_thetaFunction_mixedCoordinate_bump_mul_core}
\leanref{Development/AbelianVarietyTheory/Proofs/TwistedOperators/UnfoldingConditionsFromSecondMixedCoordinateStatement.lean}{GlobalTheory.T.unfolding_conditions_from_second_mixedCoordinate}
\leanref{Development/AbelianVarietyTheory/Proofs/TwistedOperators/ClassConditionsIsotropicHomFromFirstMixedCoordinStatement.lean}{GlobalTheory.T.class_conditions_isotropicHom_from_firstMixedCoordinate}
\leanref{Development/AbelianVarietyTheory/Proofs/TwistedOperators/ClassConditionsCoreFromMixedCoordinateFunctionStatement.lean}{GlobalTheory.T.class_conditions_core_from_mixedCoordinateFunction}
\leanref{Development/AbelianVarietyTheory/Proofs/TwistedOperators/GenericCharacterEqRealPartCharOnBumpSupportStatement.lean}{GlobalTheory.T.genericCharacter_eq_realPartChar_on_bumpSupport}
\leanref{Development/AbelianVarietyTheory/Proofs/TwistedOperators/FlatMapProductDatumEqUnitriangularIntegralBumpStatement.lean}{GlobalTheory.T.flatMap_productDatum_eq_unitriangularIntegral_bump}
\begin{proof}
We use two rules for the translations of Definition~\ref{th-weil-datum}: for
$w,w'\in\mathbb W(S,T)$,
$\rho(w)\rho(w')=\psi(\tfrac12\langle w,w'\rangle)\rho(w+w')$ and
$\rho(w)\rho(w')=\psi(\langle w,w'\rangle)\rho(w')\rho(w)$. The relations of the
two hyperbolic decompositions give the orthogonality used below: the elements
$X(a)$, $Y(a')$, $\iota_SX_T(\xi)$, $\iota_SY_T(\eta)$ and $\iota w$ are pairwise
orthogonal for $\langle\ ,\ \rangle$, except the pairs $(X(a),Y(a'))$ and
$(\iota_SX_T(\xi),\iota_SY_T(\eta))$, and
$\langle\iota w,\iota w'\rangle=\langle w,w'\rangle$.

\emph{Part (1).} Let $\Phi$ be in the class and $u\in\Gamma$, and write
$u=u_1+\dots+u_5$ with $u_1=X(a)$, $u_2=Y(a')$, $u_3=\iota_SX_T(\xi)$,
$u_4=\iota_SY_T(\eta)$ and $u_5=\iota w$ as in
Definition~\ref{th-standard-pieces}(2). By the conditions of the class,
$\rho(u_1)$ and $\rho(u_4)$ fix $\Phi$, and so do $\rho(u_2)$ and $\rho(u_3)$,
since the splitting makes their characters $1$; and $\rho(2u_5)$ fixes $\Phi$.
Hence every $\rho(2u_i)$ fixes $\Phi$. The scalars
$\psi(\tfrac12\langle 2u_i,2u_j\rangle)=\psi(\langle u_i,u_j\rangle)^2$ are $1$,
by the orthogonality and by the duality of $A_1$ with $A_2$ and of $B_1$ with
$B_2$. So $\rho(2u)\Phi=\Phi$. If $u$ lies in the half of $\mathbb W(S,T)$
attached to a split place, so do the $u_i$, which are obtained from $u$ by
multiplication with matrices on either side; the half is isotropic, so
$\rho(u)=\rho(u_1)\cdots\rho(u_5)$, and $\rho(u_5)$ fixes $\Phi$ by the condition
of the class at split places. So $\rho(u)\Phi=\Phi$, and $\Phi$ is a lattice
vacuum for $\Gamma$.

\emph{Part (2).} The translates of a theta function are theta functions. If
$\rho(c)$ commutes with all the translations in the definition of the class,
then $\rho(c)\Phi$ is in the class; by the second rule, this holds when
$\psi(\langle c,c'\rangle)=1$ for all these translations $\rho(c')$. If $c$ has
finite component $0$, then $\langle c,c'\rangle=0$ for every finite $c'$, the
pairing being a sum of local terms, and all the $c'$ are finite by the
admissibility. If $c=\iota w$ with $w\in\Gamma_0^\perp$, then $c$ is orthogonal
to the $X(a)$, $Y(a')$, $\iota_SX_T(\xi)$ and $\iota_SY_T(\eta)$, and
$\psi(\langle\iota w,2\iota w'\rangle)=\psi(\langle w,2w'\rangle)=1$ for
$w'\in\Gamma_0$, and $\psi(\langle\iota w,\iota w'\rangle)=1$ for $w'\in\Gamma_0$
in a split half. For the flat map, $\Phi^\flat$ is a theta function for
$(S_0,T_0)$, the integral defining it converging; this is a computation in the
mixed models of Lemma~\ref{th-mixed-model}, which uses that $x$ has entries in
$E$ and linearly independent columns. By Lemma~\ref{th-step-functionals}(5),
$\rho(w)\Phi^\flat=(\rho(\iota w)\Phi)^\flat$ for $w\in\mathbb W(S_0,T_0)$. Since
$\rho(2\iota w)\Phi=\Phi$ for $w\in\Gamma_0$, and $\rho(\iota w)\Phi=\Phi$ for
$w\in\Gamma_0$ in a split half, whose image $\iota w$ lies in the same half,
$\Phi^\flat$ is a lattice vacuum for $\Gamma_0$.

\emph{Part (3).} Let $g\neq0$ be a lattice vacuum for $\Gamma_0$, and let $L$ be
a matrix over $E$ with $Lx=1$. The following functions exist; each is a
construction in the mixed models of Lemma~\ref{th-mixed-model}.
\begin{enumerate}
\item[(a)] A continuous compactly supported function $\beta_2$ on
  $M_{\operatorname{rk}S_0\times\sigma}(\mathbf A_E)$, smooth along the
  archimedean directions and invariant under an open subgroup of finite
  translations, with $\beta_2(y)\neq0$, invariant under the translations by the
  $b$ with $b^*\in Y_T(B_2)$, and vanishing outside the points $y+r+b$ with $r$
  of finite component $0$ and $b^*\in Y_T(B_2)$.
\item[(b)] A function $\beta_1$ on $M_{\operatorname{rk}T\times\rho}(\mathbf A_E)$
  with the same regularity and values in $[0,1]$, with $\beta_1(x)=1$, invariant
  under the translations by $A_1$, and vanishing outside the points $x+r+a$ with
  $a\in A_1$ and $r$ of finite component $0$ such that the real part of
  $\overline{\psi(\operatorname{Tr}s(Lr))}$ is positive, where $s(M)$ is the sum
  of the superdiagonal entries $M_{i,i+1}$ of a matrix $M$. The points $x+r$ with
  this property form a neighbourhood of $x$.
\item[(c)] A theta function $g'$ for $(S_0,T)$ whose mixed model along the
  hyperbolic part of $T$ is $\beta_2(y')g$ at every point $y'$.
\item[(d)] A theta function $\Phi$ for $(S,T)$ whose mixed model along the
  hyperbolic part of $S$ is $\Phi_{x'}=\beta_1(x')g'$ for every $x'$.
\end{enumerate}
Then $\Phi$ is in the class. The conditions for $A_1$ and $A_2$ follow from
Lemma~\ref{th-mixed-model}(1) and the invariance and support of $\beta_1$, by
the duality of $A_1$ and $A_2$; those for $B_1$ and $B_2$ follow in the same
way from the second mixed model and the invariance and support of $\beta_2$; and
those for $\Gamma_0$ from the fact that $g$ is a lattice vacuum for $\Gamma_0$.
By the definition of the flat map,
\[
\Phi^\flat(w)=I\cdot\beta_2(y)\,g(w),\qquad
I=\int\overline{\psi_U(u)}\,\beta_1(xu)\,du,
\]
the integral over $u$ being that of Definition~\ref{th-step}(4), against a Haar
measure of $U_\rho(\mathbf A_E)$. If $\beta_1(xu)\neq0$, then $xu=x+r+a$ as in
(b). The archimedean component $u_\infty$ of $u$ then has $u_\infty-1=Lr$, and
its finite component $u_f$ has $xu_f-x\in A_1$, so $\psi_U(u_f)=1$ by the
admissibility, and $\psi_U(u)=\psi_U(u_\infty)=\psi(\operatorname{Tr}s(Lr))$. So
the real part of the integrand is nonnegative. It is continuous and compactly
supported in $u$, since $x$ has linearly independent columns, and positive at
$u=1$, where $\beta_1(x)=1$. Hence $I$ has positive real part, and
$\Phi^\flat\neq0$ because $\beta_2(y)\neq0$ and $g\neq0$.

\emph{Part (4).} Put $F=\omega(1,h)\Phi$, a theta function with $F^\flat\neq0$.
Two computations in the mixed models, by Lemma~\ref{th-mixed-model}(1), give for
every theta function $f$:
(i) $(\rho(X(xN))f)^\flat=\psi_U(1+N)f^\flat$ for every strictly upper triangular
$N\in M_\rho(\mathbf A_E)$, since $\rho(X(xN))$ moves the point $xu$ of the first
mixed model to $x(u+N)$, and $u\mapsto u+N$ preserves the Haar measure of
$U_\rho(\mathbf A_E)$ and multiplies $\psi_U$ by $\psi_U(1+N)$; and
(ii) $(\rho(\iota_SX_T(\xi))f)^\flat=\psi(\langle y^*,X_T(\xi)\rangle)f^\flat$
for every $\xi$, by the same rule in the second mixed model. Recall that
$\omega(1,h)\rho(w)=\rho(hw)\omega(1,h)$ (Definition~\ref{th-weil-datum}(3)).
\begin{enumerate}
\item[(a)] Since $F^\flat\neq0$, the mixed model $F_{xu}$ does not vanish at some
  point $w_0$, for some $u\in U_\rho(\mathbf A_E)$. For $a'\in A_2$,
  $\rho(Y(ha'))F=\omega(1,h)\rho(Y(a'))\Phi=\psi(\langle X(x),Y(a')\rangle)F$,
  while Lemma~\ref{th-mixed-model}(1) gives
  $(\rho(Y(ha'))F)_{xu}=\psi(\langle X(xu),Y(ha')\rangle)F_{xu}$. Comparing at
  $w_0$, and using $\langle X(xu),Y(ha')\rangle=\langle X(h^{-1}xu),Y(a')\rangle$,
  gives $\psi(\langle X(h^{-1}xu-x),Y(a')\rangle)=1$. Let $a$ be the finite
  component of $h^{-1}xu-x$. Since $a'$ is finite,
  $\psi(\langle X(a),Y(a')\rangle)=1$ for all $a'\in A_2$, so $a\in A_1$, and
  $h^{-1}xu-x-a$ has finite component $0$.
\item[(b)] Let $N$ be finite and strictly upper triangular with $h^{-1}xN\in A_1$.
  Then $\rho(X(xN))F=\omega(1,h)\rho(X(h^{-1}xN))\Phi=F$, so (i) gives
  $\psi_U(1+N)F^\flat=F^\flat$, and $\psi_U(1+N)=1$.
\item[(c)] Let $\xi\in B_1$ and $hX_T(\xi)=X_T(\xi')$. Then
  $\rho(\iota_SX_T(\xi'))F=\omega(1,h)\rho(\iota_SX_T(\xi))\Phi
  =\psi(\langle y^*,X_T(\xi)\rangle)F$, and (ii) gives
  $\psi(\langle y^*,X_T(\xi')\rangle)F^\flat=\psi(\langle y^*,X_T(\xi)\rangle)F^\flat$.
\end{enumerate}
These are the unfolding conditions at the set of all finite places.
\end{proof}

\begin{lemma}\label{th-standard-local-data}
In the situation of Lemma~\ref{th-standard-class}, let $x$ have entries in $E$ and
linearly independent columns, let the pieces split a subgroup
$\Gamma\subset\mathbb W(S,T)$ at $(x,y)$, and let $K$ be a compact subgroup of the
group of elements of $U(T)(\mathbf A)$ with archimedean component $1$, open in
that group. Suppose that $\omega(1,k)f=f$ for every lattice vacuum $f$ for
$\Gamma$ and every $k\in K$, and that every $h\in U(T)(\mathbf A)$ such that $hk$
satisfies the unfolding conditions at all finite places for every $k\in K$ is
$h=rak'$ with $r\in R_x$ (Definition~\ref{th-step-local}), $a$ of finite
component $1$ and $k'\in K$. Let $\mathcal S$ be the class of the pieces at
$(x,y)$. Then $(\mathcal S,\Gamma_0,K)$ are local data for the step
(Definition~\ref{th-step-local}), and every member of $\mathcal S$ is a lattice
vacuum for $\Gamma$.
\end{lemma}
\leanref{Development/AbelianVarietyTheory/Proofs/TwistedOperators/StandardAuxiliaryFramesGoodRowGrids.lean}{GlobalTheory.AscFin.GoodStd.loc_cls}
\begin{proof}
Every member of $\mathcal S$ is a lattice vacuum for $\Gamma$ by
Lemma~\ref{th-standard-class}(1). Conditions (1) to (6) of
Definition~\ref{th-step-local} hold as follows. Condition (1) is part of the
definition of the class. Conditions (2) and (3) are
Lemma~\ref{th-standard-class}(2), and condition (6) is
Lemma~\ref{th-standard-class}(3). Condition (4) holds because the members of
$\mathcal S$ are lattice vacua for $\Gamma$. For condition (5), let
$\Phi\in\mathcal S$ and $(\omega(1,h)\Phi)^\flat\neq0$. For $k\in K$,
$\omega(1,hk)\Phi=\omega(1,h)\omega(1,k)\Phi=\omega(1,h)\Phi$ by condition (4),
so $(\omega(1,hk)\Phi)^\flat\neq0$, and $hk$ satisfies the unfolding conditions
at all finite places by Lemma~\ref{th-standard-class}(4). By hypothesis $h=rak'$
as required.
\end{proof}

\begin{lemma}\label{th-standard-pieces-grid}
Let $D$, the level system, the theta grid, the grid Weil datum and the tower
decomposition be as in Lemma~\ref{th-standard-tower}, and let $\mathfrak m$ be the
modulus of the grid. For $0\leq i<m$ let $\Gamma_i$ be the depth lattice of the
row pair $(V_{i+1,0},V_{i,0})$, and $\Gamma'_i$ the lattice of the pair for the
lattice of $V_{i+1,0}$ in the grid and $\mathfrak m\bar{\mathfrak m}^{-1}$ times
the lattice of $V_{i,0}$. For a step from $S=S_0\oplus H^\rho$ to
$T=T_0\oplus H^\sigma$ with $(S,T)=(V_{i+1,0},V_{i,0})$, the \emph{pieces cut
out by the grid} are, in the notation of Definition~\ref{th-standard-pieces},
\[
A_1=\{a:X(a)\in\Gamma_i\},\quad
B_1=\{\xi:\iota_SX_T(\xi)\in\Gamma'_i\},\quad
B_2=\{\eta:\iota_SY_T(\eta)\in\Gamma'_i\},\quad
\Gamma_0=\{w:\iota w\in\Gamma_i\},
\]
and the closed subgroup $A_2$ generated by the finite $a'$ with
$\psi(\langle X(a),Y(a')\rangle)=1$ for every $a\in A_1$.
\begin{enumerate}
\item Let $2\rho+2\leq m$. For a step from $B_\rho$ to $J_\rho$, the pieces cut
  out by the grid are admissible at $x_0$ and split $\Gamma_{2\rho+1}$ at
  $(x_0,y_0)$ (Definition~\ref{th-step}).
\item Let $2\sigma+4\leq m$. For a step from $J_{\sigma+1}$ to $B_\sigma$, the
  pieces cut out by the grid are admissible at $x_1$ and split
  $\Gamma_{2\sigma+2}$ at $(x_1,0)$.
\end{enumerate}
\end{lemma}
\leanref{Development/AbelianVarietyTheory/Definitions/TwistedInnerProducts/IsAuxiliaryVectorWittDec.lean}{GlobalTheory.AscFin.piecesOf}
\leanref{Development/AbelianVarietyTheory/Definitions/TwistedInnerProducts/IsAuxiliaryVectorWittDec.lean}{GlobalTheory.AscFin.pcBF}
\leanref{Development/AbelianVarietyTheory/Definitions/TwistedInnerProducts/IsAuxiliaryVectorWittDec.lean}{GlobalTheory.AscFin.pcFB}
\leanref{Development/AbelianVarietyTheory/Proofs/TwistedOperators/BesselToFourierJacobiMixedPiecesAdmissibleStatement.lean}{GlobalTheory.T.besselToFourierJacobi_mixedPieces_admissible}
\leanref{Development/AbelianVarietyTheory/Proofs/TwistedOperators/FourierJacobiStepPiecesAdmissibleSplitsStatement.lean}{GlobalTheory.T.fourierJacobi_step_pieces_admissible_splits}
\leanref{Development/AbelianVarietyTheory/Proofs/TwistedOperators/GridLatticesCompactSelfDualSplittingStatement.lean}{GlobalTheory.T.grid_lattices_compact_selfDual_splitting}
\leanref{Development/AbelianVarietyTheory/Proofs/TwistedOperators/FiniteLatticeStandardParabolicStatement.lean}{GlobalTheory.T.finite_lattice_standard_parabolic}
\leanref{Development/AbelianVarietyTheory/Proofs/TwistedOperators/BesselToFourierJacobiAnnihilatorPieceCompactOpenStatement.lean}{GlobalTheory.T.besselToFourierJacobi_annihilatorPiece_compactOpen}
\leanref{Development/AbelianVarietyTheory/Proofs/TwistedOperators/BesselStepPiecesBFiniteAdelicLatticesStatement.lean}{GlobalTheory.T.bessel_step_pieces_B_finite_adelic_lattices}
\leanref{Development/AbelianVarietyTheory/Proofs/TwistedOperators/BesselStepPiecesBMutualAnnihilatorsStatement.lean}{GlobalTheory.T.bessel_step_pieces_B_mutual_annihilators}
\leanref{Development/AbelianVarietyTheory/Proofs/TwistedOperators/BidualityDepthLatticePieceBesselStepStatement.lean}{GlobalTheory.T.biduality_depth_lattice_piece_bessel_step}
\leanref{Development/AbelianVarietyTheory/Proofs/TwistedOperators/GenericCharacterTrivialBesselStepUnitriangularStatement.lean}{GlobalTheory.T.genericCharacter_trivial_besselStep_unitriangular}
\leanref{Development/AbelianVarietyTheory/Proofs/TwistedOperators/BesselStepCorePieceFiniteNonzeroVacuumStatement.lean}{GlobalTheory.T.bessel_step_core_piece_finite_nonzero_vacuum}
\leanref{Development/AbelianVarietyTheory/Proofs/TwistedOperators/BesselToFourierJacobiDepthLatticeSplitsStatement.lean}{GlobalTheory.T.besselToFourierJacobi_depthLattice_splits}
\leanref{Development/AbelianVarietyTheory/Proofs/TwistedOperators/JacobiToBesselPiecesFiniteAdelicLatticesStatement.lean}{GlobalTheory.T.jacobi_to_bessel_pieces_finite_adelic_lattices}
\leanref{Development/AbelianVarietyTheory/Proofs/TwistedOperators/FourierJacobiToBesselPiecesMutualAnnihilatorsStatement.lean}{GlobalTheory.T.fourierJacobiToBessel_pieces_mutual_annihilators}
\leanref{Development/AbelianVarietyTheory/Proofs/TwistedOperators/JacobiToBesselPiecesMutualAnnihilatorsStatement.lean}{GlobalTheory.T.jacobi_to_bessel_pieces_mutual_annihilators}
\leanref{Development/AbelianVarietyTheory/Proofs/TwistedOperators/GenericCharacterTrivialFourierJacobiStepStatement.lean}{GlobalTheory.T.genericCharacter_trivial_fourierJacobi_step}
\leanref{Development/AbelianVarietyTheory/Proofs/TwistedOperators/JacobiToBesselCorePieceNonzeroVacuumStatement.lean}{GlobalTheory.T.jacobi_to_bessel_core_piece_nonzero_vacuum}
\leanref{Development/AbelianVarietyTheory/Proofs/TwistedOperators/FourierJacobiToBesselLatticeSplitsStatement.lean}{GlobalTheory.T.fourierJacobiToBessel_lattice_splits}
\leanref{Development/AbelianVarietyTheory/Proofs/TwistedOperators/ExistsNeZeroLatticeVacuumSelfDualStatement.lean}{GlobalTheory.T.exists_ne_zero_latticeVacuum_selfDual}
\begin{proof}
We prove (1); (2) is the same set of computations for the step from
$J_{\sigma+1}$ to $B_\sigma$. Write $\Gamma=\Gamma_{2\rho+1}$ and
$\Gamma'=\Gamma'_{2\rho+1}$. Each item below is a computation with the lattices
of the grid in the bases of the tower decomposition, with the result stated.
\begin{enumerate}
\item[(a)] $\Gamma$, $\Gamma'$ and the depth lattices of the core pairs are compact
  and open in the group of finite elements. The ideal
  $\mathfrak m\bar{\mathfrak m}^{-1}$ has norm $1$, so $\mathfrak m\bar{\mathfrak m}^{-1}$
  times the lattice of $V_{2\rho+1,0}$ is self-dual, and $\Gamma'$ is the lattice
  of a pair of self-dual lattices; by Lemma~\ref{th-lattice-signs}(1) a finite
  element $u$ lies in $\Gamma'$ if and only if $\psi(\langle u,u'\rangle)=1$ for
  every $u'\in\Gamma'$. The lattice of $V'_\rho$ (the preimage of that of
  $V_{2\rho+2,0}$) is self-dual, and at every finite place the lattices of
  $V_{2\rho+2,0}$ and $V_{2\rho+1,0}$ are the direct sums of their intersections
  with the lines of the $e_i$, of the $e'_i$, and of the lattices of $V'_\rho$ and
  $P_{2\rho+1}$.
\item[(b)] Hence $A_1$, $B_1$ and $B_2$, the intersections of $\Gamma$ and
  $\Gamma'$ with closed subgroups of coordinates, are compact and open among the
  finite matrices, and so is $A_2$. By the self-duality of $\Gamma'$ and the
  direct sum decomposition in (a), $B_1$ and $B_2$ are each other's annihilators.
\item[(c)] $A_2$ annihilates $A_1$ by construction. Conversely, a finite $a$ that
  annihilates $A_2$ lies in $A_1$: at the chosen places over the distinguished
  places, the components of the elements of $\Gamma$ lie in $\Gamma'$, and those
  of the elements of $\Gamma'$ lie in $\mathfrak m$ times the lattice of
  homomorphisms, which is the condition defining the depth lattice there.
\item[(d)] If $u\in U_\rho(\mathbf A_E)$, $u-1$ is finite and $x_0u-x_0\in A_1$,
  then $\psi(\operatorname{Tr}u_{i,i+1})=1$ for every $i<\rho$. This is checked
  place by place on the coordinates of $x_0(u-1)$: away from the distinguished
  places it uses the relation between the ideal $\mathfrak A$ of the host and the
  conductor of $\psi$, and at the distinguished places the modulus. So
  $\psi_U(u)=1$.
\item[(e)] $\Gamma_0$ lies in the lattice of the pair for the self-dual lattices of
  $V'_\rho$ and $P_{2\rho+1}$, by (a). So $\Gamma_0$ consists of finite elements.
  The lattice of a pair of self-dual lattices has a nonzero lattice vacuum, by a
  computation in the model of Definition~\ref{th-weil-datum}, and a lattice
  vacuum for a subgroup is one for each of its subgroups; so $\Gamma_0$ has a
  nonzero lattice vacuum.
\item[(f)] Let $u\in\Gamma$. Its components along the $e_i$ of $V_{2\rho+2,0}$,
  along the $e'_i$, and along $V'_\rho$, the last decomposed further along the
  $e_j$, the $e'_j$ and $P_{2\rho+1}$ in $V_{2\rho+1,0}$, lie in $A_1$, $A_2$,
  $B_1$, $B_2$ and $\Gamma_0$, by (a) and (c). The characters
  $\psi(\langle X(x_0),Y(a')\rangle)$ and $\psi(\langle y_0^*,X_T(\xi)\rangle)$
  are trivial on the components $a'$ and $\xi$: the first is
  $\psi(\operatorname{Tr}(\xi_0c))$ for a sum $c$ of diagonal coordinates of $u$,
  the second $\psi(\operatorname{Tr}c')$ for a coordinate $c'$ of $u$ against
  $\zeta_{\rho+1}$, and these coordinates lie in ideals on which the characters
  are trivial, by the choice of the ideal $\mathfrak A$ of the host and of the
  modulus. So the pieces split $\Gamma$.
\end{enumerate}
Items (a) to (e) are the admissibility at $x_0$, and (f) the splitting.
\end{proof}

\begin{lemma}\label{th-standard-one-coset}
Let $D$, the level system, the theta grid, the grid Weil datum and the tower
decomposition be as in Lemma~\ref{th-standard-tower}, with the groups $K_X$ of
Lemma~\ref{th-depth-fixed}.
\begin{enumerate}
\item For $0\leq i\leq m$, $K_{V_{i,0}}$ is compact and open in the group of
  elements of $U(V_{i,0})(\mathbf A)$ with archimedean component $1$.
\item Let $2\rho+2\leq m$, take the pieces cut out by the grid for a step from
  $B_\rho$ to $J_\rho$ (Lemma~\ref{th-standard-pieces-grid}(1)), and let
  $h\in U(V_{2\rho+1,0})(\mathbf A)$ be such that $hk$ satisfies their unfolding
  conditions at $(x_0,y_0)$ at all finite places for every
  $k\in K_{V_{2\rho+1,0}}$ (Definition~\ref{th-standard-pieces}(4)). Then
  $h=rak$ with $r\in R_{x_0}$, $a$ of finite component $1$ and
  $k\in K_{V_{2\rho+1,0}}$.
\item Let $2\sigma+4\leq m$, take the pieces cut out by the grid for a step from
  $J_{\sigma+1}$ to $B_\sigma$ (Lemma~\ref{th-standard-pieces-grid}(2)), and let
  $h\in U(V_{2\sigma+2,0})(\mathbf A)$ satisfy their unfolding conditions at
  $(x_1,0)$ at all finite places. Then $h=rak$ with $r\in R_{x_1}$, $a$ of
  finite component $1$ and $k\in K_{V_{2\sigma+2,0}}$.
\end{enumerate}
\end{lemma}
\leanref{Development/AbelianVarietyTheory/Proofs/TwistedOperators/DepthCongruenceSubgroupCompactOpenPlacesToAdelesStatement.lean}{GlobalTheory.T.depthCongruenceSubgroup_compactOpen_places_to_adeles}
\leanref{Development/AbelianVarietyTheory/Proofs/TwistedOperators/LatticeStabilizerAtPlaceMemDepthCongruenceGridNoStatement.lean}{GlobalTheory.T.latticeStabilizerAtPlace_mem_depthCongruence_gridNode}
\leanref{Development/AbelianVarietyTheory/Proofs/TwistedOperators/FlatMapNonvanishingConditionLocaliseStatement.lean}{GlobalTheory.T.flatMapNonvanishingCondition_localise}
\leanref{Development/AbelianVarietyTheory/Proofs/TwistedOperators/GoodPlaceOneCosetOfBesselLevelStatement.lean}{GlobalTheory.T.goodPlace_oneCoset_of_besselLevel}
\leanref{Development/AbelianVarietyTheory/Proofs/TwistedOperators/GoodPlaceOneCosetOfBesselLevel.lean}{GlobalTheory.OneCosetGood.cond_red}
\leanref{Development/AbelianVarietyTheory/Proofs/TwistedOperators/FourierJacobiLevelOneCosetAtGoodPlaceStatement.lean}{GlobalTheory.T.fourierJacobiLevel_oneCoset_at_goodPlace}
\leanref{Development/AbelianVarietyTheory/Proofs/TwistedOperators/GoodPositionOffBadPlacesFromBesselLevelStatement.lean}{GlobalTheory.T.goodPosition_and_incAdj_mem_latO_iff_mem_intAt}
\leanref{Development/AbelianVarietyTheory/Proofs/TwistedOperators/StandardFramesSuccessiveLevelsStatement.lean}{GlobalTheory.T.standard_frames_successive_levels}
\leanref{Development/AbelianVarietyTheory/Proofs/DecoratedLatticeProofs/IwasawaDecompositionSelfDualWittSplitAtPlace.lean}{GlobalTheory.IwaPlace.iwasawa_at_place}
\leanref{Development/AbelianVarietyTheory/Proofs/TwistedOperators/GoodPlacePointConditionTorusInverseDiagIntegralStatement.lean}{GlobalTheory.T.goodPlace_pointCondition_torus_inverseDiag_integral}
\leanref{Development/AbelianVarietyTheory/Proofs/TwistedOperators/GoodPlaceShiftConditionTorusDiagIntegralStatement.lean}{GlobalTheory.T.goodPlace_shiftCondition_torus_diag_integral}
\leanref{Development/AbelianVarietyTheory/Proofs/TwistedOperators/IsotropicConditionLastTorusDiagonalIntegralStatement.lean}{GlobalTheory.T.isotropic_condition_last_torus_diagonal_integral}
\leanref{Development/AbelianVarietyTheory/Proofs/DecoratedLatticeProofs/WittLeviDiagonalMemLatticeStabilizerOfUnits.lean}{GlobalTheory.FinPl.LeviUnitK.inK_levi}
\leanref{Development/AbelianVarietyTheory/Proofs/TwistedOperators/OrbitAny.lean}{GlobalTheory.AscFin.OrbitAny.orbitAny}
\leanref{Development/AbelianVarietyTheory/Proofs/TwistedOperators/BesselLevelOneCosetAtExceptionalPlacesStatement.lean}{GlobalTheory.T.besselLevel_oneCoset_at_exceptionalPlaces}
\leanref{Development/AbelianVarietyTheory/Proofs/TwistedOperators/BesselStepOrbitAtBadPlaceEqCongruenceTranslateStatement.lean}{GlobalTheory.T.besselStep_orbit_at_badPlace_eq_congruenceTranslate}
\leanref{Development/AbelianVarietyTheory/Proofs/TwistedOperators/SplitPlaceTuplesOneOrbitModUniformizerPowerStatement.lean}{GlobalTheory.T.splitPlace_tuples_one_orbit_mod_uniformizer_power}
\leanref{Development/AbelianVarietyTheory/Proofs/TwistedOperators/OneCosetDecompositionStatementStatement.lean}{GlobalTheory.T.one_coset_decomposition_statement}
\leanref{Development/AbelianVarietyTheory/Proofs/TwistedOperators/ExistsLevelElementOffPlaceStatement.lean}{GlobalTheory.T.exists_level_element_off_place}
\leanref{Development/AbelianVarietyTheory/Proofs/TwistedOperators/ExistsUnitriangularStabilizerOfPrincipalMoveStatement.lean}{GlobalTheory.T.exists_unitriangularStabilizer_of_principal_move}
\begin{proof}
Part (1) holds because $K_{V_{i,0}}$ is the congruence subgroup, for the nonzero
modulus $\mathfrak m$, of the stabilizer of a lattice.

For (2) and (3) write $T$ for $V_{2\rho+1,0}$, resp. $V_{2\sigma+2,0}$, with its
hyperbolic decomposition $T=T_0\oplus H^\sigma$ from $J_\rho$, resp. $B_\sigma$;
write $x$ for $x_0$, resp. $x_1$, $R=R_x$, $K=K_T$, and $\Sigma_0$ for the set of
distinguished places of the grid. An element of $U(T)(\mathbf A)$ is
\emph{supported at} a set $P$ of finite places if its archimedean component and
its components outside $P$ are $1$, and \emph{trivial at} $P$ if its components
at $P$ are $1$. Every $g$ is $g_1g_2$ with $g_1$ supported at $P$ and $g_2$
trivial at $P$, and elements supported at disjoint sets commute. For
$v\notin\Sigma_0$ let $K_v$ be the group of elements supported at $\{v\}$ that
stabilize the lattice of $T$ at $v$; it lies in $K$, since $K$ differs from the
stabilizer of the lattice only at the distinguished places. An element of
$\mathbf A_E$ is \emph{integral at} $v$ if its component at $v$ is integral.

\emph{Step 1: localization.} If $g$ satisfies the unfolding conditions at all
finite places and $g=g_1g_2$ with $g_1$ supported at $P$ and $g_2$ trivial at
$P$, then $g_1$ satisfies them at $P$. Indeed $g_2$ and $g_2^{-1}$ fix every
matrix supported at $P$; so $g^{-1}xN=g_1^{-1}xN$ and $gX_T(\xi)=g_1X_T(\xi)$ for
$N$ and $\xi$ supported at $P$, and $g^{-1}xu-g_1^{-1}xu=(g_2^{-1}-1)g_1^{-1}xu$
has component $0$ at the places of $P$.

\emph{Step 2: assembly over the places.} Suppose that (i) for every
$v\notin\Sigma_0$ and every $h=h_1h_2$ with $h_1$ supported at $\{v\}$ and $h_2$
trivial at $\{v\}$, $h_1=r\kappa$ with $r\in R$ supported at $\{v\}$ and
$\kappa\in K_v$; and (ii) for every $h=h_1h_2$ with $h_1$ supported at
$\Sigma_0$ and $h_2$ trivial at $\Sigma_0$, $h_1=rk$ with $r\in R$ and $k\in K$
both supported at $\Sigma_0$. Then $h=rak$ as in the statement. Indeed, $h$ lies
in $K_v$ at all but finitely many $v$, and multiplying the factors of (i) over
the finitely many other places outside $\Sigma_0$ writes $h=r_1g\kappa_1$ with
$r_1\in R$ supported outside $\Sigma_0$, $\kappa_1$ in the product of the $K_v$,
which lies in $K$, and $g$ trivial outside $\Sigma_0$. Write $g=g_1a$ with $g_1$
supported at $\Sigma_0$ and $a$ trivial at $\Sigma_0$, so of finite component
$1$. Then $h=g_1\cdot r_1a\kappa_1$ with the second factor trivial at
$\Sigma_0$, and (ii) gives $g_1=r_0k_0$. Since $k_0$ commutes with $r_1$ and $a$,
$h=(r_0r_1)\,a\,(k_0\kappa_1)$.

\emph{Step 3: (2) outside $\Sigma_0$.} Let $v\notin\Sigma_0$ and $h=h_1h_2$ as in
(i). For $\kappa\in K_v$, $\kappa\in K$ commutes with $h_2$, so
$h\kappa=(h_1\kappa)h_2$ satisfies the unfolding conditions at all finite places,
and $h_1\kappa$ satisfies them at $\{v\}$ by Step 1. The lattices of the grid are
in good position at $v$: the lattices of $S=V_{2\rho+2,0}$, $V'_\rho$, $T$ and
$T_0=P_{2\rho+1}$ are self-dual, and at $v$ those of $S$ and $T$ are the direct
sums of their intersections with the lines of the $e_i$, of the $e'_i$, and of
the lattices of $V'_\rho$ and $T_0$; for consecutive indices $i,i+1$,
multiplication by an $s$ supported at $v$ maps the lattice of the line of
$e_{i+1}$ of $S$ into that of $e_i$ if and only if $\psi(\operatorname{Tr}(so))=1$
for every $o$ integral at $v$; a multiple $ce_i$ of $e_i$ of $S$ lies in the
lattice at $v$ if and only if $cx_i$ lies in $\mathfrak A$ times the lattice of
$T$ at $v$, where $x_i$ is the $i$-th column of $x$ and $\mathfrak A$ the ideal
fixed by the host; at $v$, $A_1$ and $B_1$ are cut out by the lattice of the pair
$(S,T)$; and for $\rho>0$ and $c$ supported at $v$, $(cy_0)^*$ lies in the
lattice of the pair $(V'_\rho,T)$ at $v$ if and only if $c$ is integral at $v$.
These are computations with the lattices in the bases of the tower
decomposition.
\begin{enumerate}
\item[(a)] By the Iwasawa decomposition at $v$,
  $h_1=n\,m(t)\,\iota(\gamma)\,\kappa$ with $n\in N_T$, $t$ diagonal in
  $GL_\rho$, $\gamma\in U(T_0)$, all supported at $v$, and $\kappa\in K_v$; here
  $m(t)$ acts on the $e_j$ of $T$ by $t$, on the $e'_j$ by
  $(\bar t^{\mathsf T})^{-1}$, and trivially on $T_0$. This is the Iwasawa
  decomposition for the stabilizer of a self-dual lattice adapted to the
  hyperbolic decomposition \cite[\S3.5]{Car79}, which Alpöge proves at $v$ from
  the good position.
\item[(b)] $r=n\iota(\gamma)$ lies in $R$, the Jacobi group of $T$
  (Lemma~\ref{th-bessel-unfolding}(1)), and $m(t)$ commutes with $\iota(\gamma)$,
  so $h_1\kappa^{-1}=r\,m(t)$.
\item[(c)] If $r\in R$ and $rg$ satisfies the unfolding conditions at $\{v\}$, so
  does $g$. Write $r^{-1}x=xu'$ and $rx=xu''$ with $u',u''\in U_\rho(\mathbf A_E)$.
  Condition (a) for $rg$ at $u$ is condition (a) for $g$ at $u'u$. In condition
  (b), $N\mapsto u'^{-1}N$ preserves the strictly upper triangular matrices
  supported at $\{v\}$ and their superdiagonal entries, and
  $(rg)^{-1}x(u'^{-1}N)=g^{-1}xN$. In condition (c), $rX_T(\xi')=X_T(u''\xi')$,
  since $X_T(\xi')=x_0\xi'$; and $y_0u''=y_0$ because $y_0$ has only its last
  column nonzero, so the characters of $\xi'$ and $u''\xi'$ agree. Applied to
  $h_1\kappa^{-1}=r\,m(t)$, this shows that $m(t)$ satisfies the unfolding
  conditions at $\{v\}$.
\item[(d)] Since $m(t)^{-1}xu=xt^{-1}u$, condition (a) and the good position give
  that each $t_i^{-1}$ is integral at $v$. Condition (b) for
  $N=t_is\,E_{i,i+1}$, $i<\rho$, with $s$ supported at $v$, has
  $m(t)^{-1}xN=s\,xE_{i,i+1}$; by the good position it gives
  $\psi(\operatorname{Tr}(t_is))=1$ whenever $\psi(\operatorname{Tr}(so))=1$ for
  all $o$ integral at $v$, so $t_i$ is integral at $v$. Condition (c), for $\xi$
  supported at $v$ in the last row, gives
  $\psi(\langle y_0^*,X_T((t_\rho-1)\xi)\rangle)=1$, and by the last fact of the
  good position $t_\rho$ is integral at $v$.
\item[(e)] An element $m(t)$ supported at $v$ with $t_i$ and $t_i^{-1}$ integral
  at $v$ stabilizes the lattice of $T$ at $v$, by the good position; so
  $m(t)\in K_v$ and $h_1=r\cdot m(t)\kappa$ is as in (i).
\end{enumerate}

\emph{Step 4: (2) at $\Sigma_0$.} Let $h=h_1h_2$ as in (ii), and for
$v\in\Sigma_0$ let $h_v$ be the component of $h_1$ at $v$, extended by $1$.
Since $h=h_v\cdot h_v^{-1}h$ with the second factor trivial at $\{v\}$, Step 1
shows that $h_v$ satisfies condition (a) at $\{v\}$: there are $u$ and
$a\in A_1$ with $h_v^{-1}x_0u\equiv x_0+a$ at $v$. The place $v$ splits in $E$,
and at the chosen place over $v$ the component of $a$ lies in $\mathfrak m$ times
the lattice of homomorphisms. In a basis adapted to the lattice, $U(T)$ at $v$
is a general linear group, and the tuples of vectors congruent to $x_0$ modulo
the power of the uniformizer given by $\mathfrak m$ form one orbit of the
principal congruence subgroup of that level. So there is $k_v\in K$ supported at
$\{v\}$ with $k_vx_0=h_v^{-1}x_0u$ at $v$. Then $r_v=h_vk_v$ satisfies
$r_vx_0=x_0u_v$, with $u_v$ the unitriangular matrix equal to $u$ at $v$ and to
$1$ elsewhere, so $r_v\in R$ is supported at $\{v\}$ and $h_v=r_vk_v^{-1}$. The
components at the places of $\Sigma_0$ commute, which gives (ii). Step 2 now
proves (2).

\emph{Step 5: (3) outside $\Sigma_0$.} Let $v\notin\Sigma_0$ and $h=h_1h_2$ as
in (i); by Step 1, $h_1$ satisfies the unfolding conditions at $\{v\}$. Here
$S=V_{2\sigma+3,0}$ with $S_0=P_{2\sigma+3}$,
$T_0=V'_\sigma=P_{2\sigma+2}\oplus(-L_z)$, and $x=x_1$ has columns the
$e_1,\ldots,e_\sigma$ of $T$ followed by $\zeta_{\sigma+1}\in T_0$. The good
position at $v$ holds as in Step 3, and the lattice of $T_0$ is at $v$ the direct
sum of the self-dual lattices of $P_{2\sigma+2}$ and of $-L_z$ in the grid.
\begin{enumerate}
\item[(a)] By the Iwasawa decomposition at $v$, as in Step 3,
  $h_1=n\,m(t)\,\iota(\gamma)\,\kappa$, and $n\in R$, since an element of $N_T$
  fixes the $e_j$ modulo the earlier ones and moves $\zeta_{\sigma+1}$ by a
  combination of the $e_j$.
\item[(b)] Conditions (a) and (b) for
  $h_1=n\cdot m(t)\iota(\gamma)\cdot\kappa$, with the computations of
  Step 3(d) on the columns $e_1,\ldots,e_\sigma$ of $x_1$, show that $t_i$ and
  $t_i^{-1}$ are integral at $v$, so $m(t)\in K_v$. Condition (a) applied to
  the last column shows that $\gamma^{-1}$ maps every multiple of
  $\zeta_{\sigma+1}$ lying in $\mathfrak A$ times the lattice of $T_0$ at $v$
  into $\mathfrak A$ times that lattice.
\item[(c)] Such a $\gamma$ is $\iota_A(a')k$ with $a'\in U(P_{2\sigma+2})$
  supported at $v$, acting on the first summand of $T_0$, and $k\in U(T_0)$
  supported at $v$ stabilizing the lattice of $T_0$ at $v$. This is a computation
  at $v$ with the self-dual lattices of the two summands of $T_0$.
\item[(d)] $\iota(\iota_A(a'))\in R$, since it fixes the $e_j$ and
  $\zeta_{\sigma+1}$, and $\iota(k)$, $m(t)$ and $\kappa$ lie in $K_v$. Since
  $m(t)$ commutes with $\iota(U(T_0))$,
  $h_1=n\,\iota(\iota_A(a'))\cdot\iota(k)\,m(t)\,\kappa$ is as in (i).
\end{enumerate}

\emph{Step 6: (3) at $\Sigma_0$.} This is Step 4 with $x_1$ in place of $x_0$: at
each $v\in\Sigma_0$, condition (a) gives $k_v\in K$ supported at $\{v\}$ with
$k_vx_1=h_v^{-1}x_1u$ at $v$, and $h_v=r_vk_v^{-1}$ with $r_v\in R$ supported at
$\{v\}$. Step 2 now proves (3).
\end{proof}

\begin{lemma}\label{th-standard-lattices}
Let $D$, the level system, the theta grid, the grid Weil datum, the tower
decomposition and the subgroups $\Gamma_L(\rho)$ be as in
Lemma~\ref{th-standard-tower}, and take the pieces cut out by the grid
(Lemma~\ref{th-standard-pieces-grid}).
\begin{enumerate}
\item Let $2\rho+2\leq m$, let $\Xi_1$ be a lattice vacuum for the depth lattice of
  the core pair $(P_{2\rho+2},P_{2\rho+1})$, and let $\Xi_2$ be a theta function
  for $(-L_z,P_{2\rho+1})$ with $\overline{\Xi_2\circ\kappa_\rho}$ a lattice vacuum
  for $\Gamma_L(\rho)$, where $\kappa_\rho(w)=(\tau_\rho w)^*$. Then
  $\Xi_1\boxtimes\Xi_2$ is a lattice vacuum for the piece $\Gamma_0$ of the step
  from $B_\rho$ to $J_\rho$.
\item Let $2\sigma+4\leq m$, let $\Xi_1$ be a lattice vacuum for the depth lattice
  of the core pair $(P_{2\sigma+3},P_{2\sigma+2})$, and let $\Phi_0$ be a lattice
  vacuum for $\Gamma_L(\sigma+1)$. Then $\Xi_1\boxtimes\Phi_0$ is a lattice
  vacuum for the piece $\Gamma_0$ of the step from $J_{\sigma+1}$ to $B_\sigma$.
\item For $2\rho+2\leq m$, $\Gamma_L(\rho)$ has a nonzero lattice vacuum.
\item The depth lattice $\Gamma$ of the row pair $(V_{1,0},V_{0,0})$, the subgroup
  $\Gamma_L(0)$ and the depth lattice $\Gamma_X$ of the core pair $(P_1,P_0)$
  satisfy the hypotheses on the subgroups in
  Theorem~\ref{th-seesaw-period-nonvanishing} for $W=V_{1,0}$, $W_0=P_1$, $T=Y_0$
  and $X=P_0=Y_0\oplus L_z$.
\end{enumerate}
\end{lemma}
\leanref{Development/AbelianVarietyTheory/Proofs/TwistedOperators/BoxDepthDataMulLineVacuaBesselLevelStatement.lean}{GlobalTheory.T.boxDepthData_mul_lineVacua_besselLevel}
\leanref{Development/AbelianVarietyTheory/Proofs/TwistedOperators/FourierJacobiDepthDatumMulLineVacuumIsVacuumStatement.lean}{GlobalTheory.T.fourierJacobi_depthDatum_mul_lineVacuum_isVacuum}
\leanref{Development/AbelianVarietyTheory/Proofs/TwistedOperators/LastSeesawLatticesNonzeroLineVacuaStatement.lean}{GlobalTheory.T.last_seesaw_lattices_nonzero_line_vacua}
\leanref{Development/AbelianVarietyTheory/Proofs/TwistedOperators/BesselStepCoreComponentMemDepthLatticeGridCoreStatement.lean}{GlobalTheory.T.besselStep_coreComponent_mem_depthLatticeGridCore}
\leanref{Development/AbelianVarietyTheory/Proofs/TwistedOperators/BesselStepLineComponentEqKappaLatticePointStatement.lean}{GlobalTheory.T.besselStep_lineComponent_eq_kappa_lattice_point}
\leanref{Development/AbelianVarietyTheory/Proofs/TwistedOperators/LinkedLineHomTransportLatticeVacuaStatement.lean}{GlobalTheory.T.linked_line_hom_transport_lattice_vacua}
\leanref{Development/AbelianVarietyTheory/Proofs/TwistedOperators/JacobiStepCoreComponentMemCoreDepthLatticeStatement.lean}{GlobalTheory.T.jacobi_step_core_component_mem_core_depth_lattice}
\leanref{Development/AbelianVarietyTheory/Proofs/TwistedOperators/FourierJacobiDepthLatticeLineComponentsMemStatement.lean}{GlobalTheory.T.fourierJacobi_depth_lattice_line_components_mem}
\leanref{Development/AbelianVarietyTheory/Proofs/TwistedOperators/ExistsNeZeroVacuumBesselLineLatticeRankZeroStatement.lean}{GlobalTheory.T.exists_ne_zero_vacuum_besselLineLattice_rankZero}
\leanref{Development/AbelianVarietyTheory/Proofs/TwistedOperators/LastSeesawLatticesNonzeroLineVacua.lean}{GlobalTheory.LastSeeSaw.lastLat_rankZero}
\leanref{Development/AbelianVarietyTheory/Proofs/TwistedOperators/LastSeesawLatticesNonzeroLineVacua.lean}{GlobalTheory.LastSeeSaw.latCore_zero_split}
\begin{proof}
For a theta function $f_A\boxtimes f_L$ along an orthogonal decomposition and an
element $u=(u_A,u_L)$, $\rho(u)(f_A\boxtimes f_L)=\rho(u_A)f_A\boxtimes\rho(u_L)f_L$,
the two components being orthogonal; and $u$ lies in the half attached to a split
place if and only if both components do. So $f_A\boxtimes f_L$ is a lattice
vacuum for a subgroup $\Lambda$ if $f_A$ and $f_L$ are theta functions
(Lemma~\ref{th-tensor-rule}) and are fixed by $\rho(2u_A)$ and $\rho(2u_L)$ for
$u\in\Lambda$, and by $\rho(u_A)$ and $\rho(u_L)$ for $u\in\Lambda$ in a split
half.

\emph{Part (1).} Let $u\in\Gamma_0\subset\mathbb W(V'_\rho,P_{2\rho+1})$, with
components $u_A$ and $u_L$ along $V'_\rho=P_{2\rho+2}\oplus(-L_z)$. A computation
with the lattices of the grid shows that $u_A$ lies in the depth lattice of the
core pair $(P_{2\rho+2},P_{2\rho+1})$, and that $u_L=\kappa_\rho(w)$ for some
$w\in\Gamma_L(\rho)$. Since $\overline{\Xi_2\circ\kappa_\rho}$ is a lattice vacuum
for $\Gamma_L(\rho)$, and $\kappa_\rho$ is additive, preserves the halves and
multiplies the pairing by $-1$ (proof of Lemma~\ref{th-line-comparison}), the
translation $\rho(\kappa_\rho w)$ acts on $\Xi_2$ as $\rho(w)$ acts on
$\overline{\Xi_2\circ\kappa_\rho}$, up to complex conjugation. Hence
$\rho(2\kappa_\rho w)\Xi_2=\Xi_2$, and $\rho(\kappa_\rho w)\Xi_2=\Xi_2$ when
$\kappa_\rho w$ lies in a split half. With the first paragraph this proves (1).

\emph{Part (2).} Let $u\in\Gamma_0\subset\mathbb W(P_{2\sigma+3},V'_\sigma)$, with
components $u_A$ and $u_L$ along $V'_\sigma=P_{2\sigma+2}\oplus(-L_z)$. A
computation with the lattices of the grid shows that $u_A$ lies in the depth
lattice of the core pair $(P_{2\sigma+3},P_{2\sigma+2})$ and $u_L$ in
$\Gamma_L(\sigma+1)$. The first paragraph proves (2).

\emph{Part (3).} The lattice $\mathcal O_E\zeta_0$ of $-L_z$ is self-dual, since
$\zeta_0$ maps to the vector $a$ of the grid, with $h(a,a)=1$. For
$\rho=\sigma+1\geq1$, $\Gamma_L(\rho)$ is the lattice of the pair for the lattice
of $P_{2\rho+1}$ in the grid and $I\zeta_0$, with
$I=\mathfrak a^\rho\bar{\mathfrak a}^{-\rho}\mathfrak m\bar{\mathfrak m}^{-1}$;
since $I\bar I=\mathcal O_E$, the lattice $I\zeta_0$ is self-dual. For $\rho=0$,
the lattice of $P_0$ in the grid projects along $Y_0$ into $\mathcal O_E\zeta_0$,
so $\Gamma_L(0)$ lies in the lattice of the pair for the lattice of $P_1$ and
$\mathcal O_E\zeta_0$. In both cases $\Gamma_L(\rho)$ lies in the lattice of a
pair of self-dual lattices, so it has a nonzero lattice vacuum by the computation
in Lemma~\ref{th-standard-pieces-grid}(e).

\emph{Part (4).} Since $J_0$ has depth $0$, $P_1=V_{1,0}$ and $\iota_0$ is the
identity. The group $\Gamma_X$ is the set of $u=(u_1,u_2)$ along
$P_0=Y_0\oplus L_z$ with $u_1\in\Gamma$ and $u_2\in\Gamma_L(0)$. Indeed, a
computation with the lattices of the grid shows that the orthogonal projection
$u\mapsto(u_1,0)$ preserves $\Gamma_X$ and that $(u_1,0)\in\Gamma_X$ if and
only if $u_1\in\Gamma$; and $(0,u_2)\in\Gamma_X$ if and only if
$u_2\in\Gamma_L(0)$, by the definition of $\Gamma_L(0)$. Let
$w=(w_1,w_2)\in\Gamma_X^\perp$. Testing $w$ against the elements $(u_1,0)$
and $(0,u_2)$ of $\Gamma_X$, which are orthogonal to $(0,w_2)$ and $(w_1,0)$,
gives $w_1\in\Gamma^\perp$ and $w_2\in\Gamma_L(0)^\perp$; the pairing on
$\mathbb W(P_1,-L_z)$ is the negative of that on $\mathbb W(P_1,L_z)$, and
$\psi(-t)=\overline{\psi(t)}$. For the vacua, let $f_1$ be the complex conjugate
of the Fock vacuum vector of the row pair (Definition~\ref{th-grid}), a nonzero
lattice vacuum for the lattice of the pair, hence for $\Gamma$, and $f_2$ a
nonzero lattice vacuum for $\Gamma_L(0)$ from (3). By the description of
$\Gamma_X$ and the first paragraph, $f_1\boxtimes\bar f_2$ is a lattice vacuum
for $\Gamma_X$.
\end{proof}

\begin{lemma}\label{th-standard-tower}
Let $D$ be a definite-side datum in which $p_{de}$ and $p'_{de}$ are prime to
$Q_Sq$, with host element $\xi_0$, a level system with base tuple, a theta
grid, a grid Weil datum and measures as in Definition~\ref{th-core-chain}. For
$1\leq c\leq m$ let $a_0,\ldots,a_{c-2}$ be the basis of the orthogonal
complement of $P_c$ in $V_{c,0}$ given by the grid, with antidiagonal Gram
matrix, and let $a$ span the complement of $Y_0$ in $P_0$. Define a tower
decomposition (Definition~\ref{th-tower-decomposition}) as follows:
\begin{enumerate}
\item $L_z$ is a line with a vector $\zeta_0$, and $P_0=Y_0\oplus L_z$ with
  $\zeta_0\mapsto a$;
\item $\zeta_\rho=(-\xi_0)^\rho\zeta_0$, and $\tau_\rho$ is multiplication by
  $\xi_0$;
\item $B_\rho$ has $e_p=(-\xi_0)^{p+1}a_p$ and $e'_p$ the multiple of
  $a_{2\rho-p}$ with $h(e'_p,e_p)=1$, for $p<\rho$, in the complement of
  $P_{2\rho+2}$, and $V'_\rho=P_{2\rho+2}\oplus(-L_z)$ embedded through
  $P_{2\rho+2}\subset V_{2\rho+2,0}$ and $\zeta_0\mapsto a_\rho$;
\item $J_\rho$ has $e_p=(-\xi_0)^{p+1}a_p$ and $e'_p$ the multiple of
  $a_{2\rho-1-p}$ with $h(e'_p,e_p)=1$, for $p<\rho$, in the complement of
  $P_{2\rho+1}$;
\end{enumerate}
with the Haar measures of volume $1$ on the flag groups. Let
$\Gamma_L(0)$ be the set of $w\in\mathbb W(P_1,-L_z)$ whose image in
$\mathbb W(P_1,P_0)$ lies in the depth lattice of the core pair $(P_1,P_0)$,
and for $\sigma\geq0$ let $\Gamma_L(\sigma+1)$ be the lattice of the pair
$(P_{2\sigma+3},-L_z)$ for the lattice of $P_{2\sigma+3}$ in the grid and the
lattice
$\mathfrak a^{\sigma+1}\bar{\mathfrak a}^{-\sigma-1}\mathfrak m\bar{\mathfrak m}^{-1}\zeta_0$
of $L_z$, where $\mathfrak m$ is the modulus of the grid and
$\mathfrak a=\mathfrak A\bar{\mathfrak m}^{-1}$ for the ideal $\mathfrak A$
fixed by the host. Then the tower decomposition fits the ends and is good for
the subgroups $\Gamma_L(\rho)$.
\end{lemma}
\leanref{Development/AbelianVarietyTheory/Proofs/TwistedOperators/ExistsGoodGridFramesOfDepthFixedStatement.lean}{GlobalTheory.T.existsGoodGridFrames_of_depthFixed}
\leanref{Development/AbelianVarietyTheory/Proofs/TwistedOperators/ExistsStandardAuxiliaryFramesAllRowGridsStatement.lean}{GlobalTheory.T.exists_standardAuxiliaryFrames_allRowGrids}
\leanref{Development/AbelianVarietyTheory/Proofs/TwistedOperators/StandardAuxiliaryFramesFitGridEndsStatement.lean}{GlobalTheory.T.standardAuxiliaryFrames_fit_grid_ends}
\leanref{Development/AbelianVarietyTheory/Proofs/TwistedOperators/StandardAuxiliaryFramesGoodRowGridsStatement.lean}{GlobalTheory.T.standardAuxiliaryFrames_good_row_grids}
\leanref{Development/AbelianVarietyTheory/Proofs/TwistedOperators/StandardAuxiliaryFramesFitGridEnds.lean}{GlobalTheory.AscFin.EndsPf.Chr.psiBes_eq_star_besChar}
\begin{proof}
Since $\xi_0\neq0$, the $e_p$ and $e'_p$ in (3) and (4) form hyperbolic
decompositions, with $e_p$ and $e'_p$ in the hyperbolic planes spanned by
$a_p$ and $a_{c-2-p}$, and in (3) the middle vector $a_\rho$ is orthogonal to
them. In (2),
$\tau_\rho\zeta_\rho=\xi_0(-\xi_0)^\rho\zeta_0=-\zeta_{\rho+1}$ and
$h(\zeta_{\rho+1},\zeta_{\rho+1})=\xi_0\bar\xi_0\,h(\zeta_\rho,\zeta_\rho)$, as
Definition~\ref{th-tower-decomposition}(2) requires.

\emph{The bottom end.} The embedding of $U(Y_0)$ in $U(P_0)$ given by (1) is
the embedding $\iota_0$ of the grid, since (1) is the decomposition of the
grid.

\emph{The top end.} Let $2\rho+2=m$, and let $\varphi$ on $U(P_m)(\mathbf A)$
and $f$ on $U(V_{m,0})(\mathbf A)$ be rapidly decreasing.
\begin{enumerate}
\item[(a)] The flag group of $B_\rho$ is the flag group of the top Bessel
  period: scaling the $e_p$ by nonzero constants does not change the flag
  they span.
\item[(b)] The embedding of $U(P_m)$ in $U(V_{m,0})$ through $V'_\rho$ is the
  embedding of the top Bessel period.
\item[(c)] On the flag group, the character $\psi_{\zeta_{\rho+1}}$ of $B_\rho$
  is the complex conjugate of the character of the top Bessel period. Let $n$
  be in the flag group, and write, in the basis $a_0,\ldots,a_{\rho-1}$ and the
  middle vector $a_\rho$,
  $na_q=a_q+\sum_{p<q}t_{pq}a_p$ and $na_\rho=a_\rho+\sum_{p<\rho}b_pa_p$. The character of the top Bessel period is
  $\psi(\operatorname{Tr}(\xi_0s))$ with $s=\sum_pt_{p,p+1}+b_{\rho-1}$. In the
  basis $e_p=\lambda_pa_p$ of $B_\rho$, $\lambda_p=(-\xi_0)^{p+1}$, the
  superdiagonal coordinates of $n$ are
  $\lambda_p^{-1}t_{p,p+1}\lambda_{p+1}=-\xi_0t_{p,p+1}$, and the middle vector
  $\zeta_{\rho+1}$ is $\lambda_\rho a_\rho$ with $\lambda_\rho=(-\xi_0)^{\rho+1}$,
  so its last coordinate is $\lambda_{\rho-1}^{-1}b_{\rho-1}\lambda_\rho=-\xi_0b_{\rho-1}$.
  Hence $\psi_{\zeta_{\rho+1}}(n)=\psi(\operatorname{Tr}(-\xi_0s))$, which is
  the complex conjugate of $\psi(\operatorname{Tr}(\xi_0s))$.
\end{enumerate}
By (a) to (c),
$\mathcal B_{B_\rho}(\bar\varphi,\bar f)$ is the complex conjugate of the top
Bessel period of $(\varphi,f)$ computed with the Haar measures and fundamental
domains of $B_\rho$. Two Bessel periods that differ only in the choice of Haar
measures and of fundamental domains, those for the flag group of volume $1$,
differ by a positive factor. Hence $\mathcal B_{B_\rho}(\bar\varphi,\bar f)\neq0$
implies $\mathcal B^+(\varphi,f)\neq0$.

\emph{Goodness.} Let $2\rho+2\leq m$, take a step from $B_\rho$ to $J_\rho$ with
Weil datum $\omega^{\mathrm{row}}_{2\rho+1}$ on $(V_{2\rho+2,0},V_{2\rho+1,0})$,
and let $\Gamma$ be the depth lattice of this row pair. Take the pieces cut out
by the grid (Lemma~\ref{th-standard-pieces-grid}), with their class
$\mathcal S$ at $(x_0,y_0)$ and their last piece $\Gamma_0$
(Definition~\ref{th-standard-pieces}), and put $K=K_{V_{2\rho+1,0}}$.
\begin{enumerate}
\item[(G1)] The pieces are admissible at $x_0$ and split $\Gamma$ at $(x_0,y_0)$,
  by Lemma~\ref{th-standard-pieces-grid}(1). The matrix $x_0$ has entries in $E$
  and linearly independent columns (Lemma~\ref{th-bessel-unfolding}(1)).
\item[(G2)] $K$ is compact and open in the group of elements with archimedean
  component $1$, by Lemma~\ref{th-standard-one-coset}(1), and
  $\omega^{\mathrm{row}}_{2\rho+1}(1,k)$ fixes the lattice vacua for $\Gamma$ for
  $k\in K$, by Lemmas~\ref{th-grid-conjugation} and~\ref{th-depth-fixed}(1).
\item[(G3)] Every $h$ such that $hk$ satisfies the unfolding conditions at all
  finite places for every $k\in K$ is $h=rak'$ with $r\in R_{x_0}$, $a$ of finite
  component $1$ and $k'\in K$, by Lemma~\ref{th-standard-one-coset}(2).
\item[(G4)] By Lemma~\ref{th-standard-local-data}, $(\mathcal S,\Gamma_0,K)$ are
  local data for the step, and every member of $\mathcal S$ is a lattice vacuum
  for $\Gamma$.
\end{enumerate}
By Lemma~\ref{th-standard-lattices}(1), $\Xi_1\boxtimes\Xi_2$ is a lattice vacuum
for $\Gamma_0$ under the conditions of
Definition~\ref{th-tower-decomposition}(4). This gives condition (4) of that
definition. Condition (5) is proved in the same way for a step from
$J_{\sigma+1}$ to $B_\sigma$ with Weil datum $\omega^{\mathrm{row}}_{2\sigma+2}$,
with the point $x_1$, by Lemmas~\ref{th-standard-pieces-grid}(2),
\ref{th-standard-one-coset}(3), \ref{th-standard-local-data}
and~\ref{th-standard-lattices}(2). Here $x_1$ has entries in $E$ and linearly
independent columns, since the $e_j$ of $B_\sigma$ are linearly independent and
isotropic, $\zeta_{\sigma+1}$ is orthogonal to them, and
$h(\zeta_{\sigma+1},\zeta_{\sigma+1})\neq0$. Conditions (6) and (7) are
Lemma~\ref{th-standard-lattices}(3) and~(4).
% Goodness is GlobalTheory.AscFin.GoodStd.std_good; the packaging is GoodStd.loc_cls, and the
% independence of the columns of x_1 is GoodStd.xLast_rat.
\end{proof}

\begin{lemma}\label{th-tower-rungs}
Let $D$ be a definite-side datum in which $p_{de}$ and $p'_{de}$ are prime to
$Q_Sq$, with a level system with base tuple, a theta grid, a grid Weil datum and
measures as in Definition~\ref{th-core-chain}.
\begin{enumerate}
\item There are a tower decomposition of the grid that fits the ends and
  subgroups $\Gamma_L(\rho)$ for which it is good
  (Definition~\ref{th-tower-decomposition}).
\end{enumerate}
Let a tower decomposition that fits the ends and is good for subgroups
$\Gamma_L(\rho)$ be given, let $(\Psi_0,\ldots,\Psi_m)$ be a column chain of
core forms, and let $u$ be a function on $U(Y_0)(\mathbf A)$ all of whose
partial row lifts $\theta^{(n)}(u)$, $n\leq m$, are rapidly decreasing for all
row data (Definition~\ref{th-core-chain}). For $2k+2\leq m$ say that
\emph{$(J_k)$ holds} if there are a character $\chi'$, a Weil datum $\omega_z$
for $(P_{2k+1},-L_z)$ with characters $(\chi',\eta_L)$, where $\eta_L$ is the
character of (L3) of the level system (Definition~\ref{th-grid}), row data and
a lattice vacuum $\Phi_0$ for $\Gamma_L(k)$ with
\[
\mathcal J_{J_k}\bigl(\bar\Psi_{2k+1},\overline{\theta^{(2k+1)}(u)};\Phi_0\bigr)\neq0,
\]
the period taken with $\zeta_k$ and $\omega_z$; and that \emph{$(B_k)$ holds}
if there are row data with
$\mathcal B_{B_k}\bigl(\bar\Psi_{2k+2},\overline{\theta^{(2k+2)}(u)}\bigr)\neq0$.
\begin{enumerate}\setcounter{enumi}{1}
\item If $\int_{\mathcal F_{Y_0}}\Psi_0(\iota_0h)u(h)\,dh\neq0$, then $(J_0)$
  holds.
\item If $(J_k)$ holds, then $(B_k)$ holds.
\item If $2k+4\leq m$ and $(B_k)$ holds, then $(J_{k+1})$ holds.
\end{enumerate}
\end{lemma}
\leanref{Development/AbelianVarietyTheory/Definitions/TwistedInnerProducts/TowerEx.lean}{GlobalTheory.Asc.TowerEx}
\leanref{Development/AbelianVarietyTheory/Proofs/TwistedOperators/ExistsGoodGridFramesOfDepthFixedStatement.lean}{GlobalTheory.T.existsGoodGridFrames_of_depthFixed}
\leanref{Development/AbelianVarietyTheory/Definitions/TwistedInnerProducts/PartFixed.lean}{GlobalTheory.Asc.NZJ}
\leanref{Development/AbelianVarietyTheory/Definitions/TwistedInnerProducts/PartFixed.lean}{GlobalTheory.Asc.NZB}
\leanref{Development/AbelianVarietyTheory/Proofs/TwistedOperators/NonvanishingThetaTowerBottomRungStatement.lean}{GlobalTheory.T.nonvanishing_thetaTower_bottom_rung}
\leanref{Development/AbelianVarietyTheory/Proofs/TwistedOperators/FourierJacobiToBesselNonvanishingTowerStatement.lean}{GlobalTheory.T.fourierJacobi_to_bessel_nonvanishing_tower}
\leanref{Development/AbelianVarietyTheory/Proofs/TwistedOperators/BesselToFourierJacobiNonvanishingTowerStatement.lean}{GlobalTheory.T.bessel_to_fourierJacobi_nonvanishing_tower}
\leanref{Development/AbelianVarietyTheory/Proofs/TwistedOperators/NonvanishingThetaTowerBottomRung.lean}{GlobalTheory.AscBottom.rung0}
\leanref{Development/AbelianVarietyTheory/Proofs/TwistedOperators/FourierJacobiToBesselNonvanishingTower.lean}{GlobalTheory.AscRungB.rungB}
\leanref{Development/AbelianVarietyTheory/Proofs/TwistedOperators/BesselToFourierJacobiNonvanishingTower.lean}{GlobalTheory.Asc.RungJac.rungJ}
\begin{proof}
\emph{Part (1)} is Lemma~\ref{th-standard-tower}.

In the proofs of (2) to (4), $(\varrho,\Phi)$ are row data, $i$ is the index of
the row pair at which the step takes place, and $(\varrho',\Phi')$ are the row
data equal to $(\varrho,\Phi)$ except that $(\varrho'_i,\Phi'_i)=(P,\kappa)$ for
a polynomial $P$ and a coset vector $\kappa$. By
Lemma~\ref{th-grid-conjugation}(6),
$\theta^{(i)}_{\varrho',\Phi'}(u)=\theta^{(i)}_{\varrho,\Phi}(u)$, and by
Lemma~\ref{th-grid-conjugation}(5),
\[
\overline{\theta^{(i+1)}_{\varrho',\Phi'}(u)}=\theta_{\bar f_{\kappa,P}}\bigl(\overline{\theta^{(i)}_{\varrho,\Phi}(u)}\bigr)\tag{$*$}
\]
with the lift by $\omega^{\mathrm{row}}_i$ and $f_{\kappa,P}$ as in
Lemma~\ref{th-fock-span}.

\emph{Part (3).} Let $2k+2\leq m$, assume $(J_k)$ with $\chi'$, $\omega_z$,
row data $(\varrho,\Phi)$ and $\Phi_0$, and put $i=2k+1$, which is odd. By
Lemma~\ref{th-grid-conjugation}(1) and~(2), $\omega^{\mathrm{row}}_i$ is a Weil
datum for $(V_{2k+2,0},V_{2k+1,0})$ with characters
$(\chi_{m-i}^{-1}\eta_L,\chi_{m-i}^{-1})$, and $\omega^{\mathrm{core}}_i$ one
for $(P_{2k+2},P_{2k+1})$ with characters $(\chi_{m-i}^{-1},\chi_{m-i}^{-1})$.
Choose a step from $B_k$ to $J_k$ with Weil datum $\omega^{\mathrm{row}}_i$ and
the measure on $U(V_{2k+1,0})$ of the grid; a Weil datum for
$(V'_k,P_{2k+1})$ with the same characters, and the measures, exist. Let
$(\mathcal S,\Gamma_0,K_{V_{2k+1,0}})$ be local data for it as in
Definition~\ref{th-tower-decomposition}(4). Put $\varphi=\bar\Psi_{2k+2}$ and
$G=\overline{\theta^{(2k+1)}_{\varrho,\Phi}(u)}$. Both are rapidly decreasing,
and $G$ is invariant under right translation by $K_{V_{2k+1,0}}$ by
Lemma~\ref{th-depth-fixed}(3).

The core form $\Psi_{2k+1}$ is the lift of $\Psi_{2k+2}$ by a theta function
$f'$. By Lemma~\ref{th-grid-conjugation}(3), $\Xi_1=\bar f'$ is a lattice vacuum
for the depth lattice of the core pair, and by
Lemma~\ref{th-grid-conjugation}(4), $\theta_{\Xi_1}(\varphi)=\bar\Psi_{2k+1}$
for the lift by $\omega^{\mathrm{core}}_i$. By
Lemma~\ref{th-line-comparison}(3) for $\tau_k$, $\zeta_{k+1}$ and $\zeta_k$,
$\Phi_0=\overline{\Xi_2\circ\kappa_k}$ for a theta function $\Xi_2$ for
$(-L_z,P_{2k+1})$, and by Definition~\ref{th-tower-decomposition}(4)
$\Xi_1\boxtimes\Xi_2$ is a lattice vacuum for $\Gamma_0$. The hypothesis reads
$\mathcal J_{J_k}(\theta_{\Xi_1}(\varphi),G;\overline{\Xi_2\circ\kappa_k})\neq0$.
Apply Proposition~\ref{th-step-nonvanishing}(1) with $A=P_{2k+2}$, $L=-L_z$,
$\zeta=\zeta_{k+1}$, $L'=-L_z$, $\zeta'=\zeta_k$, $\tau=\tau_k$,
$\chi_A=\chi_{m-i}^{-1}$, $\chi_z=\eta_L$ and
$\omega_A=\omega^{\mathrm{core}}_i$; the character conditions hold by the
characters just listed. It gives $\Phi_1\in\mathcal S$ with
$\mathcal B_{B_k}(\varphi,\theta_{\Phi_1}(G))\neq0$.

The map $\ell(\Phi)=\mathcal B_{B_k}(\varphi,\theta_\Phi(G))$, with the lift by
$\omega^{\mathrm{row}}_i$, is tame by Lemma~\ref{th-step-functionals}(1). The
function $\Phi_1$ is a lattice vacuum for the depth lattice of the row pair, by
Definition~\ref{th-tower-decomposition}(4), and $\ell(\Phi_1)\neq0$. By
Lemma~\ref{th-fock-span} there are $\kappa$ and $P$ with
$\ell(\bar f_{\kappa,P})\neq0$. By $(*)$,
$\ell(\bar f_{\kappa,P})=\mathcal B_{B_k}\bigl(\bar\Psi_{2k+2},\overline{\theta^{(2k+2)}_{\varrho',\Phi'}(u)}\bigr)$,
so $(B_k)$ holds.

\emph{Part (4).} Let $2k+4\leq m$, assume $(B_k)$ with row data
$(\varrho,\Phi)$, and put $i=2k+2$, which is even. By
Lemma~\ref{th-grid-conjugation}(1) and~(2), $\omega^{\mathrm{row}}_i$ is a Weil
datum for $(V_{2k+3,0},V_{2k+2,0})$ with characters
$(\chi_{m-i}^{-1},\chi_{m-i}^{-1}\eta_L)$, and $\omega^{\mathrm{core}}_i$ one
for $(P_{2k+3},P_{2k+2})$ with characters $(\chi_{m-i}^{-1},\chi_{m-i}^{-1})$.
The character $\eta_L$ is a Hecke character, as the quotient of two characters
of Weil data, and its restriction to $\mathbf A_F^\times$ is $\eta$ by (L3) of
Definition~\ref{th-grid}, the restriction required for the line $L_z$. So a
Weil datum $\omega_z$ for $(P_{2k+3},-L_z)$ with characters
$(\chi_{m-i}^{-1},\eta_L)$ exists by the existence of Weil data with given Hecke
characters of the right restrictions to $\mathbf A_F^\times$. Choose a step
from $J_{k+1}$ to $B_k$ with Weil datum $\omega^{\mathrm{row}}_i$ and the measure
on $U(V_{2k+2,0})$ of the grid, and local data
$(\mathcal S,\Gamma_0,K_{V_{2k+2,0}})$ for it as in
Definition~\ref{th-tower-decomposition}(5). Put $\varphi=\bar\Psi_{2k+3}$ and
$G=\overline{\theta^{(2k+2)}_{\varrho,\Phi}(u)}$; both are rapidly decreasing,
and $G$ is right invariant under $K_{V_{2k+2,0}}$ by
Lemma~\ref{th-depth-fixed}(3).

As in (3), $\Psi_{2k+2}$ is the lift of $\Psi_{2k+3}$ by a theta function $f'$,
$\Xi_1=\bar f'$ is a lattice vacuum for the depth lattice of the core pair, and
$\theta_{\Xi_1}(\varphi)=\bar\Psi_{2k+2}$ for the lift by
$\omega^{\mathrm{core}}_i$. By Definition~\ref{th-tower-decomposition}(6) there
is a nonzero lattice vacuum $\Phi_0$ for $\Gamma_L(k+1)$, and by
Definition~\ref{th-tower-decomposition}(5) $\Xi_1\boxtimes\Phi_0$ is a lattice
vacuum for $\Gamma_0$. The hypothesis reads
$\mathcal B_{B_k}(\theta_{\Xi_1}(\varphi),G)\neq0$. Apply
Proposition~\ref{th-step-nonvanishing}(2) with $A=P_{2k+2}$, $L=-L_z$,
$\zeta=\zeta_{k+1}$, $\chi_A=\chi_{m-i}^{-1}$, $\chi_z=\eta_L$ and
$\omega_A=\omega^{\mathrm{core}}_i$. It gives $\Phi'_1\in\mathcal S$ with
$\mathcal J_{J_{k+1}}(\varphi,\theta_{\Phi'_1}(G);\Phi_0)\neq0$.

The map $\ell(\Phi)=\mathcal J_{J_{k+1}}(\varphi,\theta_\Phi(G);\Phi_0)$ is tame
by Lemma~\ref{th-step-functionals}(2), $\Phi'_1$ is a lattice vacuum for the
depth lattice of the row pair, and $\ell(\Phi'_1)\neq0$. By
Lemma~\ref{th-fock-span} there are $\kappa$ and $P$ with
$\ell(\bar f_{\kappa,P})\neq0$, and by $(*)$ this is
$\mathcal J_{J_{k+1}}\bigl(\bar\Psi_{2k+3},\overline{\theta^{(2k+3)}_{\varrho',\Phi'}(u)};\Phi_0\bigr)\neq0$.
So $(J_{k+1})$ holds, with $\chi'=\chi_{m-i}^{-1}$.

\emph{Part (2).} Put $i=0$ and $a=\chi_m$. By Lemma~\ref{th-grid-conjugation}(1)
and~(2), $\omega^{\mathrm{row}}_0$ is a Weil datum for $(V_{1,0},Y_0)$ with
characters $(a^{-1},a^{-1}\eta_L)$, and $\omega^{\mathrm{core}}_0$ one for
$(P_1,P_0)$ with characters $(a^{-1},a^{-1})$. As in (4), a Weil datum
$\omega_z$ for $(P_1,-L_z)$ with characters $(a,\eta_L)$ exists; the
restriction of $\eta_L$ to $\mathbf A_F^\times$ is the one required for a line
because $\operatorname{rk}P_0=\operatorname{rk}Y_0+1$. The core form $\Psi_0$ is
the lift of $\Psi_1$ by a theta function $f'$; $\Xi=\bar f'$ is a lattice vacuum
for the depth lattice $\Gamma_X$ of the core pair $(P_1,P_0)$, and
$\theta_\Xi(\bar\Psi_1)=\bar\Psi_0$ for the lift by $\omega^{\mathrm{core}}_0$
(Lemma~\ref{th-grid-conjugation}(3) and~(4)). Since the tower decomposition fits
the ends, the embedding $h\mapsto(h,1)$ of $U(Y_0)$ in $U(P_0)$ is $\iota_0$, so
\[
\int_{\mathcal F_{Y_0}}\theta_\Xi(\bar\Psi_1)\bigl((h,1)\bigr)\,\bar u(h)\,dh
=\overline{\int_{\mathcal F_{Y_0}}\Psi_0(\iota_0h)\,u(h)\,dh}\neq0 .
\]
Apply Theorem~\ref{th-seesaw-period-nonvanishing} with $W=V_{1,0}$,
$W_0=P_1$, $\iota_0$ the embedding of $J_0$, $X=P_0=Y_0\oplus L_z$, $T=Y_0$,
$L''=L_z$, the subgroups of Definition~\ref{th-tower-decomposition}(7),
$\omega=\omega^{\mathrm{row}}_0$, $\omega_X=\omega^{\mathrm{core}}_0$, $\omega_z$,
$\chi_S=a^{-1}$, $\chi_T=a^{-1}\eta_L$, $\chi_z=\eta_L$, $\varphi=\bar\Psi_1$
and $\bar u$ in place of $u$. It gives lattice vacua $\Phi'_1$ for the depth
lattice of the row pair $(V_{1,0},V_{0,0})$ and $\Phi_0$ for $\Gamma_L(0)$ with
$\mathcal J_{J_0}(\bar\Psi_1,\theta_{\Phi'_1}(\bar u);\Phi_0)\neq0$; for depth
$0$ the Fourier-Jacobi period of Definition~\ref{th-periods}(2) is the period in
the conclusion of Theorem~\ref{th-seesaw-period-nonvanishing}. As in (3), the
map $\Phi\mapsto\mathcal J_{J_0}(\bar\Psi_1,\theta_\Phi(\bar u);\Phi_0)$ is tame,
Lemma~\ref{th-fock-span} gives $\kappa$ and $P$ on which it is nonzero, and
$(*)$ with $\theta^{(0)}(u)=u$ shows that $(J_0)$ holds, with $\chi'=a$.
\end{proof}

\begin{theorem}[Ascent of nonvanishing]\label{th-bessel-ascent}
Let $D$ be a definite-side datum (Definition~\ref{th-definite-datum}) in which
$p_{de}$ and $p'_{de}$ are prime to $Q_Sq$. Let a level system with base tuple,
a theta grid, a grid Weil datum, and Haar measures with fundamental domains for
the row spaces $V_{i,0}$ and the core spaces $P_i$ be given
(Definition~\ref{th-core-chain}). Let $(\Psi_0,\ldots,\Psi_m)$ be a column
chain of core forms, and let $u$ be a function on $U(Y_0)(\mathbf A)$ all of
whose partial row lifts $\theta^{(n)}(u)$, $n\leq m$, are rapidly decreasing
for all row data (Definition~\ref{th-core-chain}). Suppose
\[
\int_{\mathcal F_{Y_0}}\Psi_0(\iota_0h)\,u(h)\,dh\neq0,
\]
where $\iota_0\colon U(Y_0)\to U(P_0)$ is the inclusion. Then there are row
data $(\varrho,\Phi)$ with
$\mathcal B^+\bigl(\Psi_m,\theta^{(m)}_{\varrho,\Phi}(u)\bigr)\neq0$
(Definition~\ref{th-periods}).
\end{theorem}
\leanref{Development/AbelianVarietyTheory/Proofs/TwistedOperators/BesselPeriodNonvanishingAscentOfDepthDataStatement.lean}{GlobalTheory.T.bessel_period_nonvanishing_ascent_of_depth_data}
\leanref{Development/AbelianVarietyTheory/Definitions/TwistedInnerProducts/OffQOfSizeRows.lean}{GlobalTheory.SepRows.Ascent}
\leanref{Development/AbelianVarietyTheory/Proofs/TwistedOperators/BesselPeriodAscentOfTowerRungs.lean}{GlobalTheory.Asc.ascent_of_rungs}
\leanref{Development/AbelianVarietyTheory/Proofs/TwistedOperators/DependenceFixedStatement.lean}{GlobalTheory.T.dependence_fixed}
\begin{proof}
Choose a tower decomposition that fits the ends and is good for subgroups
$\Gamma_L(\rho)$, by Lemma~\ref{th-tower-rungs}(1). We show by induction on $k$
that $(B_k)$ holds for every $k$ with $2k+2\leq m$. For $k=0$, the hypothesis
and Lemma~\ref{th-tower-rungs}(2) give $(J_0)$, and
Lemma~\ref{th-tower-rungs}(3) gives $(B_0)$. If $(B_k)$ holds and
$2k+4\leq m$, Lemma~\ref{th-tower-rungs}(4) gives $(J_{k+1})$, and
Lemma~\ref{th-tower-rungs}(3) gives $(B_{k+1})$. Since $m$ is even, the case
$2k+2=m$ gives row data $(\varrho,\Phi)$ with
$\mathcal B_{B_k}\bigl(\bar\Psi_m,\overline{\theta^{(m)}_{\varrho,\Phi}(u)}\bigr)\neq0$.
The partial row lift $\theta^{(m)}_{\varrho,\Phi}(u)$ is rapidly decreasing by
hypothesis, and so is $\Psi_m$. Since the tower decomposition fits the ends,
$\mathcal B^+(\Psi_m,\theta^{(m)}_{\varrho,\Phi}(u))\neq0$.
\end{proof}

\begin{definition}\label{th-bessel-identity}
Let $D$ be a definite-side datum (Definition~\ref{th-definite-datum}), with
Bessel pairings $B_u$, standard Bessel periods $\mathcal B_u$, division
lattices $N_u\subset\mathcal M_u$ and $R\subset\mathcal E$ with embeddings
$j_N$ and $j_R$, and division pairings $B^{\mathrm{div}}_u$ with values in $L$.
The \emph{Bessel identity} holds for $D$ if:
\begin{enumerate}
\item the Bessel pairings are rational:
  $B_u(j_N\varphi,j_Rx)=\iota\bigl(B^{\mathrm{div}}_u(\varphi,x)\bigr)$ for all
  words $u$, $\varphi\in N_u$ and $x\in R$, where $\iota\colon L\to\mathbf C$
  is the embedding of the coefficient field;
\item there is a unitary map $V$ of $\mathcal E$ such that for every word $u$
  there are $a_u\in\mathbf C$ with $a_u\neq0$ and a unitary map $U_u$ of
  $\mathcal M_u$ with
  \[
  B_u(\varphi,F)=a_u\,\mathcal B_u(U_u\varphi,VF)\qquad(\varphi\in\mathcal M_u,\ F\in\mathcal E).
  \]
\end{enumerate}
\end{definition}
\leanref{Formalism/SharpRestriction/GlobalTheoremInputs.lean}{Formalism.MainNotions.BesselId}
\leanref{Formalism/LocalRingCardinality/RhoDetPrimeEqMuTildeMulChiPow.lean}{Formalism.BasicNotions.D7.BuRational}
\leanref{Formalism/LocalRingCardinality/DataRows.lean}{Formalism.BasicNotions.Rows.BuDef}

When the Bessel pairing is determined by the periods, in the sense that
$\operatorname{per}_{u,\mu_0}(\varphi,F)=B_u(\varphi,\iota_{\mu_0}F)$ for the
top forms $F$ of the base tuple, the Bessel identity says that the period of
the iterated lift equals, up to the nonzero constant $a_u$ and unitary changes
of variables, the standard Bessel period $\mathcal B_u$, which does not depend
on the tuple of Hecke characters. It is one of the conditions of the division
link (Definition~\ref{is-division-link}), and it is proved for the families
used in the argument together with that link
(Theorem~\ref{is-division-link-holds}).

\begin{remark}\label{th-scope}
The four parts of this chapter are proved in Alpöge's proof with different
generality.
\begin{enumerate}
\item \emph{Definite unitary groups.} The spaces, column chains and data of
  Definitions~\ref{th-spaces} to~\ref{th-definite-datum} are definitions.
  The data used in the argument, with all the properties needed later, are
  supplied by Theorem~\ref{asm-definite-side-exists}.
\item \emph{Theta kernels.} The iterated lift of
  Definition~\ref{th-iterated-lift} is defined for arbitrary kernels
  $\vartheta_i(c,c')$. That they are theta kernels of lattice pairs is the
  condition of Remark~\ref{th-lift-kernels}, imposed and used where integrality
  of periods is proved (Chapter~\ref{chapter-period-divisibility}).
\item \emph{The lifts $\pi\mapsto\pi\boxplus\mu_1[t+1]\boxplus\cdots$.} No
  statement of this chapter identifies the global parameter of the iterated
  lift. The block shape enters through the weights and Harish-Chandra
  parameters of Definition~\ref{th-definite-datum}(7). The theta
  correspondence is used through Theorems~\ref{th-rallis-inner-product}
  and~\ref{th-first-lift-nonzero}, which concern the lifts of the trivial
  representation from totally definite groups, and their hypotheses include
  Howe duality at the archimedean places only for the trivial representation
  (Theorem~\ref{lit-howe-duality-real}).
\item \emph{Seesaw and Bessel periods.} The tensor rule
  (Lemma~\ref{th-tensor-rule}), the seesaw identity
  (Theorem~\ref{th-seesaw-identity}), the irreducibility lemma
  (Lemma~\ref{th-vacuum-irreducibility}), the nonvanishing step
  (Theorem~\ref{th-seesaw-period-nonvanishing}) and the steps of the tower
  (Propositions~\ref{th-vanishing-descends} and~\ref{th-step-nonvanishing})
  are proved for arbitrary Hermitian spaces and Weil data over the host field. The ascent
  (Theorem~\ref{th-bessel-ascent}) is proved for every theta grid of every
  datum in which $p_{de}$ and $p'_{de}$ are prime to $Q_Sq$. The identification of the period of the iterated lift with the
  standard Bessel period (Definition~\ref{th-bessel-identity}) is a condition
  on the datum, verified only for the data used, as part of the division link
  of Chapter~\ref{chapter-integral-structure}.
\end{enumerate}
\end{remark}
